% !TeX program = xelatex
\documentclass[UTF8,11pt,a4paper,oneside,openany,fontset=fandol]{ctexbook}
\usepackage[left=24mm,right=24mm,top=25mm,bottom=25mm,headheight=16pt]{geometry}
\usepackage{amsmath,amssymb,booktabs,tabularx,array,float,graphicx,enumitem,makecell}
\usepackage{xcolor,tikz,fancyhdr,etoolbox,pgfplots}
\pgfplotsset{compat=1.18}
\usetikzlibrary{calc,positioning,arrows.meta,shapes.arrows,shapes.geometric,decorations.pathreplacing}
\usepackage[most]{tcolorbox}
\usepackage{hyperref}
\hypersetup{unicode=true,colorlinks=true,linkcolor=black,urlcolor=FinanceBlue,bookmarksnumbered=true,pdftitle={金融学基础：十五章中文译稿}}
\definecolor{FinanceBlue}{cmyk}{1,0.65,0,0.35}
\setcounter{secnumdepth}{2}
\setcounter{tocdepth}{1}
\setlength{\parskip}{0.25em}
\setlength{\emergencystretch}{3em}
\linespread{1.16}
\ctexset{chapter={format=\Large\bfseries,beforeskip=8pt,afterskip=22pt},section={format=\large\bfseries},subsection={format=\normalsize\bfseries}}
\setlist{itemsep=0.35em,topsep=0.5em}
\newcolumntype{Y}{>{\raggedright\arraybackslash}X}
\NewDocumentEnvironment{openingcase}{m}{%
\begin{tcolorbox}[enhanced,breakable,colback=FinanceBlue!3,colframe=FinanceBlue,boxrule=0pt,borderline west={2pt}{0pt}{FinanceBlue},sharp corners,left=10pt,right=10pt,top=10pt,bottom=10pt,before skip=1em,after skip=1em]
{\small\bfseries\color{FinanceBlue}章首案例\par}\smallskip
{\large\bfseries #1\par}\medskip
}{\end{tcolorbox}}
% The source chapters use two mandatory example arguments (title and ID).
\NewDocumentEnvironment{example}{m m}{%
\begin{tcolorbox}[enhanced,breakable,colback=gray!5,colframe=FinanceBlue!75,title={例：#1},fonttitle=\bfseries,boxrule=0.5pt]
\label{ex:#2}
}{\end{tcolorbox}}
\newcommand{\TailAnchor}[1]{\phantomsection\hypertarget{sec:\thechapter:#1}{}\label{sec:\thechapter:#1}}
\newif\iffinancereadings
\newcommand{\TailModule}[2]{\TailAnchor{#1}\section*{#2}\ifstrequal{#1}{furtherread}{\begingroup\financereadingstrue\small\setlist[itemize]{itemsep=0.2em}}{}}
\newcommand{\TailLink}[3]{\hyperlink{sec:#1:#2}{#3}\,(\pageref*{sec:#1:#2})}
\ExplSyntaxOn
\cs_new:Npn \finance_tail_title:n #1 {
\str_case:nnF {#1} {
{summary}{本章小结} {mindmap}{思维导图} {keyterms}{关键术语}
{questions}{讨论题} {appprobs}{应用题} {aiprompt}{金融~AI~实验室}
{dataapp}{数据应用} {caseanalysis}{案例分析} {furtherread}{延伸阅读}
}{#1}}
\bool_new:N \l__finance_tail_first_bool
\cs_new_protected:Npn \finance_tail_list:nn #1#2 {
\bool_set_true:N \l__finance_tail_first_bool
\clist_map_inline:nn {#2} {
\bool_if:NF \l__finance_tail_first_bool {\quad}
\TailLink{#1}{##1}{\finance_tail_title:n{##1}}
\bool_set_false:N \l__finance_tail_first_bool
}}
\NewDocumentCommand{\HyperTailList}{mm}{\finance_tail_list:nn{#1}{#2}}
\ExplSyntaxOff
% A compact guide at each chapter's end avoids overcrowding the main contents.
\newcommand{\ChapterTail}[1]{\par\medskip\noindent{\footnotesize\HyperTailList{\thechapter}{#1}}\par\iffinancereadings\endgroup\fi}
\pagestyle{fancy}
\fancyhf{}
\fancyhead[L]{\small 金融学基础}
\fancyhead[R]{\small\nouppercase{\leftmark}}
\fancyfoot[C]{\thepage}
\renewcommand{\headrulewidth}{0.3pt}
\renewcommand{\chaptermark}[1]{\markboth{第\thechapter 章\quad #1}{}}
\fancypagestyle{plain}{\fancyhf{}\fancyfoot[C]{\thepage}\renewcommand{\headrulewidth}{0pt}}
\begin{document}
\frontmatter
\begin{titlepage}
\centering
\vspace*{38mm}
{\Huge\bfseries 金融学基础\par}
\vspace{12mm}
{\Large 十五章中文译稿\par}
\vspace{8mm}
{\large 金融体系、信息契约、估值、金融市场与金融机构\par}
\vfill
{\normalsize 根据提供的十五章英文 LaTeX 原稿翻译\par}
\vspace{4mm}
{\small 2026 年 11 月\par}
\end{titlepage}
\chapter*{译稿说明}
本稿依据所提供的十五章英文原稿完整翻译，采用中国大陆金融学常用术语。原文的案例背景、货币单位、公式、数值、章节层级及习题均予保留。参考文献保留作者姓名、原文题名和出版信息，并附题名中译，便于检索。文中的 AI 实验、数据练习和案例研究要求均为教材内容，未代为作答。

本稿中的金融概念以原文语境为准，例如“清算”与“结算”、“风险分担”与“风险转换”，以及“预期回报”与“必要报酬率”分别使用。正文及图示中的说明文字已译为中文，公式变量和交叉引用标识保持一致。各章关键术语附英文对照。
\tableofcontents
\mainmatter
\part{金融体系基础}
\chapter{金融的本质：摩擦、功能与金融体系}
\label{chap:The_Essence_of_Finance_Frictions_Functions_and_the_Financial_System}

\begin{openingcase}{项目有价值，为何融不到资金？}

苏州工程师李伟开发了一套用于工业设备的节能控制系统。要建立生产线并将产品推向市场，他现在需要人民币 1,000 万元，但只能投入自己的 100 万元储蓄。

另一边，王成女士有一笔近期不打算使用的储蓄。在没有摩擦的世界里，双方似乎很容易匹配：一方有资金，另一方能够将资金用于生产。然而，交易并不会自行达成。

王成对这项技术的了解不如李伟。她无法观察项目的真实质量，无法核实李伟的所有预测，也无法预先知道他获得资金后会如何行动。项目本身也存在不确定性：需求可能不及预期，成本可能上升，也可能出现竞争性技术。双方偏好的期限也不同。李伟需要能够长期投入的资金，而王成可能希望更早动用自己的储蓄。

即使王成愿意提供剩余的 900 万元，仍有许多问题需要解决：风险如何分配？李伟应披露哪些信息？王成应获得哪些权利？双方如何确信未来的义务会得到履行？如果项目预期在多年内产生现金流，这些不确定的未来现金流在今天究竟值多少钱？

仅有储蓄和投资机会，并不能保证资金从储蓄者流向投资者。需要金融安排来促进跨期交易，并应对不确定性、信息不对称、激励与承诺带来的挑战。

\end{openingcase}

\begin{tcolorbox}[
  enhanced,
  colback=FinanceBlue!6,
  colframe=FinanceBlue,
  boxrule=0.6pt,
  sharp corners,
  left=8pt,
  right=8pt,
  top=7pt,
  bottom=7pt,
  before skip=1em,
  after skip=1.6em,
  title=\textbf{核心问题},
  coltitle=white,
  colbacktitle=FinanceBlue,
  fonttitle=\bfseries,
  attach boxed title to top left={xshift=6pt,yshift=-2mm},
  boxed title style={
    boxrule=0pt,
    sharp corners,
    left=5pt,
    right=5pt,
    top=2pt,
    bottom=2pt
  }
]
\begin{enumerate}[
  leftmargin=*,
  label=\textcolor{FinanceBlue}{\textbf{\arabic*.}},
  itemsep=0.55em,
  topsep=0.2em
]
\item 哪些经济摩擦使金融成为必要？
\item 要实现互利交换，金融体系必须履行哪些功能？
\item 为什么货币、市场、机构和金融工具会采取不同形式？这些形式又为什么会随时间变化？
\end{enumerate}
\end{tcolorbox}

\section{金融为何存在：从经济机会到金融摩擦}

今天拥有资金的人，未必是能够最有效地使用这些资金的人。金融通过以当期资源交换未来支付的安排，将双方联系起来。由于这些支付取决于未来结果，不确定性不可避免。

未来支付无法事先确定，因此不确定性使信息变得重要。参与者必须判断可能发生什么，以及承诺是否可能得到履行。然而，信息往往不完整，且分布不均，因此金融安排还必须明确各方享有哪些权利、承担哪些义务。这些承诺本身也必须可信，并能够得到执行。

\subsection{时间：资源与机会的分离}

经济活动中的现金流入与流出发生在不同时间。个人可能在某一时点获得收入，却希望在另一时点支出；企业也往往先承担成本，后通过销售获得现金。因此，资金可供使用的时间，与其潜在用途所需的时间不一定一致。

这种分离使跨期交换（intertemporal exchange）能够带来收益。储蓄把购买力从现在转移到未来，借款则进行相反的转移。在一个简单的两期基准模型中，若个人能够按相同利率 $r$ 借入或贷出资金，其约束为

\begin{equation}
c_0+\frac{c_1}{1+r}=y_0+\frac{y_1}{1+r},
\label{eq:ch1_intertemporal_budget}
\end{equation}

其中，$y_0$ 和 $y_1$ 分别表示第 0 期和第 1 期的收入，$c_0$ 和 $c_1$ 表示相应时期的消费，$r$ 为利率。该约束要求消费现值等于收入现值。借贷使消费可以在不改变其现值的情况下跨期转移。今天每增加一单位消费，就必须放弃下一期的 $1+r$ 单位消费。

\begin{figure}[H]
\centering
\begin{tikzpicture}[
  >=Latex,
  axis/.style={-Latex, line width=0.8pt},
  budget/.style={line width=1.0pt},
  guide/.style={dashed, gray, line width=0.6pt},
  move/.style={-Latex, line width=0.9pt, FinanceBlue}
]
\draw[axis] (0,0) -- (7.1,0) node[right] {$c_0$};
\draw[axis] (0,0) -- (0,5.8) node[above] {$c_1$};
\draw[budget] (0.7,5.15) -- (6.45,0.65);
\coordinate (E) at (3.55,2.92);
\fill[FinanceBlue] (E) circle (2.2pt);
\node[above right] at (E) {$E=(y_0,y_1)$};
\draw[guide] (E) -- (3.55,0) node[below] {$y_0$};
\draw[guide] (E) -- (0,2.92) node[left] {$y_1$};
\draw[move] (3.43,3.02) -- (2.25,3.94);
\node[above left, align=center] at (4.88,3.95) {储蓄（贷出）};
\draw[move] (3.67,2.82) -- (4.88,1.88);
\node[below right] at (5.38,1.88) {借款};
\node[above right] at (4.72,1.95) {斜率 $=-(1+r)$};
\end{tikzpicture}
\caption{储蓄、借款与跨期预算约束}
\label{fig:ch1-intertemporal-budget}
\end{figure}

在禀赋点 $E$，个人每一期的消费都等于当期收入。沿预算线向左上方移动，意味着当前消费低于收入，并将购买力转移到未来；向右下方移动，则意味着以未来收入为基础借款。预算线说明，收入的时间分布并不必然决定消费的时间分布。

对企业而言，有投资机会并不意味着拥有实施投资所需的资金。企业可能有一个有价值的项目，但内部资金不足，而外部投资者拥有可供投资的资源。如图~\ref{fig:ch1-internal-external-finance} 所示，外部融资通过将外部资源引入企业，并以未来支付请求权作为交换，弥合了这种分离。
\begin{figure}[H]
\centering
\begin{tikzpicture}[
  >=Latex,
  box/.style={
    draw,
    rounded corners,
    minimum width=3.0cm,
    minimum height=1.15cm,
    align=center,
    line width=0.8pt
  },
  flow/.style={
    -Latex,
    line width=0.9pt,
    FinanceBlue
  },
  return/.style={
    -Latex,
    dashed,
    line width=0.8pt,
    gray
  },
  firmbox/.style={
    draw,
    rounded corners,
    dashed,
    line width=0.8pt
  }
]

%------------------------------------------------
% External investors
%------------------------------------------------
\node[box] (outside) at (0.3,1.3)
{外部投资者\\[2pt]
\small 拥有可用资金};

%------------------------------------------------
% Firm
%------------------------------------------------
\node[box] (internal) at (6.2,1.3)
{内部资金\\[2pt]
\small 单靠内部资金不足};

\node[box] (project) at (6.2,3.35)
{投资机会\\[2pt]
\small 需要资金};

% Firm boundary
\draw[firmbox]
  (3.8,0.25) rectangle (8.6,4.45);

\node[
  fill=white,
  inner sep=2pt,
  font=\bfseries
]
at (6.2,4.45) {企业};

%------------------------------------------------
% Internal financing and future return
%------------------------------------------------

% Current financing: slightly left
\coordinate (intfinstart)
  at ($(internal.north)+(-0.18,0)$);

\coordinate (intfinend)
  at ($(project.south)+(-0.18,0)$);

\draw[flow]
  (intfinstart)
  -- node[left, font=\small]
  {内部融资}
  (intfinend);

% Future return: slightly right
\coordinate (intreturnstart)
  at ($(project.south)+(0.18,0)$);

\coordinate (intreturnend)
  at ($(internal.north)+(0.18,0)$);

\draw[return]
  (intreturnstart)
  -- node[right, font=\small]
  {未来回报}
  (intreturnend);

%------------------------------------------------
% External financing and future return
%------------------------------------------------

% External finance: upper path
\coordinate (extfinstart)
  at ($(outside.north)+(-0.15,0)$);

\coordinate (extfinend)
  at ($(project.west)+(0,0.18)$);

\draw[flow]
  (extfinstart)
  |- node[pos=0.72, above, font=\small]
  {外部融资}
  (extfinend);

% Future return: lower parallel path
\coordinate (extreturnstart)
  at ($(project.west)+(0,-0.18)$);

\coordinate (extreturnend)
  at ($(outside.north)+(0.15,0)$);

\draw[return]
  (extreturnstart)
  -| node[pos=0.28, below, font=\small]
  {未来回报}
  (extreturnend);

\end{tikzpicture}

\caption{资源与投资机会分离时的内部融资和外部融资}
\label{fig:ch1-internal-external-finance}
\end{figure}


\subsection{不确定性与风险分配}

未来的不确定性增加了跨期交换的难度。奈特根据人们对未来结果的认知程度区分了风险（risk）与不确定性（uncertainty）。风险是指可能的未来状态已知，而且每种状态的概率可以确定。不确定性则包括两种情况：第一，可能的未来状态已知，但其概率未知；第二，甚至无法完全识别可能出现的未来状态，因此这些状态的概率也未知。

金融安排可以决定，不同状态带来的经济后果如何在参与者之间分配。它们无法消除究竟会出现何种状态的根本不确定性，但可以改变由谁承担随之产生的收益和损失。

如果参与者面临的风险敞口不同，承担风险的能力或意愿也不同，风险分配就可能创造经济价值。金融安排能够将特定风险转移给更有能力或更愿意承担它的人。风险分担（risk sharing）是风险分配的一种形式：一个人难以承受的损失，可以分散给许多参与者。

设状态 $s$ 下可用的总资源为 $Y(s)$，参与者 $i$ 获得的资源为 $X_i(s)$，则任何可行分配都必须满足

\begin{equation}
\sum_i X_i(s)=Y(s)
\qquad \text{对每一种状态 }s.
\label{eq:ch1_risk_sharing}
\end{equation}

这一约束意味着，风险分担可以改变某一状态下资源的分配，却不能改变该状态下可用资源的总量。图~\ref{fig:ch1-risk-sharing} 展示了这一区别。

\begin{figure}[H]
\centering
\begin{tikzpicture}[
  >=Latex,
  title/.style={font=\small\bfseries},
  label/.style={font=\small},
  flow/.style={-Latex, line width=0.9pt, FinanceBlue},
  total/.style={draw, line width=0.8pt, minimum height=0.65cm},
  abar/.style={fill=FinanceBlue!25},
  bbar/.style={fill=gray!20}
]

%------------------------------------------------
% Column headings
%------------------------------------------------
\node[title] at (3.0,5.2) {状态 1};
\node[title] at (8.0,5.2) {状态 2};

%------------------------------------------------
% Before risk sharing
%------------------------------------------------
\node[label, anchor=east] at (0.3,4.15)
{风险分担前};

% State 1: A has more, B has less
\draw[total] (1.4,3.8) rectangle (4.6,4.45);
\fill[abar] (1.4,3.8) rectangle (3.9,4.45);
\fill[bbar] (3.9,3.8) rectangle (4.6,4.45);

\node[label] at (2.65,4.125) {$A$};
\node[label] at (4.25,4.125) {$B$};

% State 2: A has less, B has more
\draw[total] (6.4,3.8) rectangle (9.6,4.45);
\fill[abar] (6.4,3.8) rectangle (7.1,4.45);
\fill[bbar] (7.1,3.8) rectangle (9.6,4.45);

\node[label] at (6.75,4.125) {$A$};
\node[label] at (8.35,4.125) {$B$};

%------------------------------------------------
% Risk-sharing arrows
%------------------------------------------------
\draw[flow] (3.0,3.55) -- (3.0,2.75);
\draw[flow] (8.0,3.55) -- (8.0,2.75);

\node[label, fill=white, inner sep=2pt]
  at (5.5,3.15) {风险分担};

%------------------------------------------------
% After risk sharing
%------------------------------------------------
\node[label, anchor=east] at (0.3,2.0)
{风险分担后};

% State 1: more balanced allocation
\draw[total] (1.4,1.65) rectangle (4.6,2.30);
\fill[abar] (1.4,1.65) rectangle (3.35,2.30);
\fill[bbar] (3.35,1.65) rectangle (4.6,2.30);

\node[label] at (2.38,1.975) {$A$};
\node[label] at (3.98,1.975) {$B$};

% State 2: more balanced allocation
\draw[total] (6.4,1.65) rectangle (9.6,2.30);
\fill[abar] (6.4,1.65) rectangle (7.65,2.30);
\fill[bbar] (7.65,1.65) rectangle (9.6,2.30);

\node[label] at (7.03,1.975) {$A$};
\node[label] at (8.63,1.975) {$B$};

%------------------------------------------------
% Aggregate resource labels
%------------------------------------------------
\draw[decorate,
  decoration={brace, amplitude=4pt, mirror},
  line width=0.7pt]
  (1.4,1.40) -- (4.6,1.40)
  node[midway, below=5pt, label] {$Y(1)$ 不变};

\draw[decorate,
  decoration={brace, amplitude=4pt, mirror},
  line width=0.7pt]
  (6.4,1.40) -- (9.6,1.40)
  node[midway, below=5pt, label] {$Y(2)$ 不变};

\end{tikzpicture}

\caption{风险分担改变分配，不改变总资源}
\label{fig:ch1-risk-sharing}
\end{figure}

在各状态可用资源给定时，金融能够重新分配风险后果，却不能消除总量损失。保险将损失分摊给投保人，股权将经营风险分配给投资者，衍生工具重新配置特定风险敞口，分散化则降低对个别结果的依赖。因此，风险分配关注的是：在不同的世界状态下，由谁承担经济后果。

\subsection{信息与激励摩擦}

金融决策要求在结果实现之前，对项目、交易对手和未来现金流作出评估。这些评估所依据的信息，往往不完整，且在参与者之间分布不均。

这种信息差异构成信息不对称（information asymmetry）。企业管理者对项目质量的了解可能多于外部投资者，而借款人可能掌握贷款人无法直接观察的行为信息。融资之前，如果融资条件影响哪些借款人寻求或接受资金，不可观察的项目质量差异就可能引发逆向选择（adverse selection）。融资之后，如果借款人有动力采取影响贷款人收益、但无法被完全观察的行动，就可能产生道德风险（moral hazard）。

仅靠信息披露无法解决这些问题。信息生产可能成本高昂，信息可能难以核实，而且即使能够获得也可能不完整。金融体系将资源用于筛选、审计、认证与监督。这些活动能够改善金融决策的信息基础，但无法消除未来结果的不确定性。

\begin{tcolorbox}[
  enhanced,
  breakable,
  colback=white,
  colframe=FinanceBlue,
  boxrule=0.7pt,
  sharp corners,
  borderline west={2.2pt}{0pt}{FinanceBlue},
  title=\textbf{金融视角：金融契约为何不只是转移资金},
  coltitle=white,
  colbacktitle=FinanceBlue,
  fonttitle=\bfseries,
  before skip=1.2em,
  after skip=1.2em
]
提供资金之前，贷款人和投资者可以评估借款人、发行人或项目，核实财务信息，并参考信用记录或声誉。贷款人也可能要求抵押品。融资之后，契约可以要求定期披露信息、限制某些行为，并赋予资金提供者监督绩效或在违反契约条件时介入的权利。

这些机制并不能消除信息不对称。筛选和核实主要用于融资之前，此时相关特征可能难以观察；披露、监督和契约限制则在融资之后变得重要，此时行为可能难以观察。
\end{tcolorbox}

信息既影响融资能否实现，也影响融资条件。但信息本身并不界定各方的权利和义务。这些权利和义务必须得到明确规定，而未来承诺必须可信且能够执行。


\subsection{承诺与执行}

许多金融安排依赖于规定未来承诺、权利和义务的契约。契约明确支付、或有情形，以及未遵守约定条款时的救济措施。例如，债务契约规定何时付款，以及借款人违约时债权人享有哪些权利。

只有当各方预期未来义务会得到履行时，契约承诺才可信。如果一方在签订契约时无法可信地约束自己的未来行为，就会出现承诺问题（commitment problem）。融资之后，一方可能有动力采取使自己受益、却降低另一方请求权价值的行动。抵押品、限制性条款、声誉与控制权，可以通过限制某些行为或改变不履约的后果来支持承诺。

执行（enforcement）使契约权利具有实际效力。法院、债权人权利、破产程序和抵押安排，决定了义务被违反时可采取哪些救济措施。对执行情况的预期会影响是否提供融资及其条件。然而，执行既不完备，也并非没有成本：契约无法规定所有未来情形，法律程序也可能耗时且昂贵。因此，金融安排除了依靠正式契约权利，还依靠重新谈判、声誉和酌情处理。

这四类经济问题往往同时存在。例如，银行贷款在跨期转移购买力的同时，分配与不确定偿付有关的风险，依赖借款人的信息，并建立由执行机制支持的契约权利和义务。因此，金融安排通常同时应对多个经济问题。

\section{金融体系履行哪些功能}

上述经济问题构成了金融体系核心功能的依据。尽管各国及不同时期的金融体系各不相同，它们都在不同时间和用途之间配置资源、分担和转换风险、提供流动性并支持支付、生产信息并约束行为，以及促进价格发现。即使履行这些功能的机构和技术发生变化，这些功能依然重要。

\subsection{在不同时间和用途之间配置资源}

金融在不同时点之间重新配置购买力。借款使个人能够提前使用未来购买力，而储蓄和投资把当前购买力转移到未来。在生命周期中，这些安排可以帮助个人在成年初期为教育或住房融资，在工作期间积累资源，并支持退休后的消费。对企业而言，融资同样使投资可以发生在项目现金流实现之前。

\begin{figure}[H]
\centering
\begin{tikzpicture}[
  >=Latex,
  axis/.style={-Latex, line width=0.8pt},
  income/.style={line width=1.1pt, FinanceBlue},
  consumption/.style={line width=1.0pt, gray!75},
  guide/.style={dashed, gray, line width=0.6pt},
  label/.style={font=\small}
]

%------------------------------------------------
% Axes
%------------------------------------------------
\draw[axis] (0,0) -- (9.4,0) node[right] {年龄};
\draw[axis] (0,0) -- (0,5.4)
  node[above, align=center] {收入与\\消费};

%------------------------------------------------
% Income over the life cycle
%------------------------------------------------
\draw[income]
  (0.6,1.25)
  .. controls (1.6,1.55) and (2.3,3.3) ..
  (3.5,4.15)
  .. controls (4.5,4.65) and (5.7,4.55) ..
  (6.5,3.85)
  .. controls (7.0,3.3) and (7.25,1.55) ..
  (8.7,1.05);

%------------------------------------------------
% Consumption over the life cycle
%------------------------------------------------
\draw[consumption]
  (0.6,2.05)
  .. controls (2.1,2.35) and (3.7,2.65) ..
  (5.0,2.70)
  .. controls (6.2,2.70) and (7.5,2.45) ..
  (8.7,2.10);

%------------------------------------------------
% Curve labels
%------------------------------------------------
\node[label, FinanceBlue, anchor=west]
  at (4.15,4.65) {收入（$I$）};

\node[label, anchor=west]
  at (3.75,2.38) {消费（$C$）};

%------------------------------------------------
% Stage separators
%------------------------------------------------
\draw[guide] (2.45,0) -- (2.45,4.9);
\draw[guide] (6.85,0) -- (6.85,4.9);

%------------------------------------------------
% Stage labels
%------------------------------------------------
\node[label, align=center] at (1.35,-0.55)
  {成年初期};

\node[label, align=center] at (4.65,-0.55)
  {工作期间};

\node[label, align=center] at (7.95,-0.55)
  {退休后};

%------------------------------------------------
% Financial allocation labels
%------------------------------------------------
\node[label, align=center] at (1.35,3.25)
  {借款\\[1pt]\footnotesize $C > I$};


\node[label, align=center] at (4.65,1.15)
  {储蓄与投资\\[1pt]\footnotesize $I > C$};


\node[label, align=center] at (7.95,3.25)
  {使用积累的储蓄\\[1pt]\footnotesize $C > I$};


\end{tikzpicture}

\caption{生命周期中的金融配置}
\label{fig:ch1-life-cycle-allocation}
\end{figure}

图~\ref{fig:ch1-life-cycle-allocation} 展示了金融在生命周期中的这一作用。收入和支出需求可能沿着不同路径变化，因此借款和储蓄可以在不同人生阶段之间转移购买力。

金融还在相互竞争的用途之间配置稀缺资源。利率、证券价格、贷款决策与投资授权，会影响哪些项目获得资金以及融资条件。若两家企业争取有限资金，相对于其风险和融资成本而言，预期回报更有吸引力的项目可能获得更多融资或更优惠的条件，而另一项目可能获得较少融资，甚至无法融资。

\subsection{分担与转换风险}

参与者面临的风险及其承担能力或意愿并不相同。这些差异创造了重新分配风险的空间。相对于某一参与者资源而言很大的损失，另一参与者可能更容易吸收；将损失分散给多人，或与其他风险敞口组合，也可以减轻负担。金融安排能够在参与者之间分担风险，或转换风险被承担的形式。
\begin{figure}[H]
\centering
\begin{tikzpicture}[
  >=Latex,
  box/.style={
    draw,
    rounded corners,
    minimum width=3.8cm,
    minimum height=1.05cm,
    align=center,
    line width=0.8pt
  },
  topbox/.style={
    draw,
    rounded corners,
    minimum width=5.4cm,
    minimum height=1.05cm,
    align=center,
    line width=0.8pt
  },
  flow/.style={
    -Latex,
    line width=0.9pt,
    FinanceBlue
  },
  note/.style={
    font=\small,
    align=center
  }
]

%------------------------------------------------
% Source of risk reallocation
%------------------------------------------------
\node[topbox] (heterogeneity) at (5.5,4.6)
{风险敞口与\\风险承担能力的差异};


%------------------------------------------------
% Risk sharing
%------------------------------------------------
\node[box] (sharing) at (2.7,2.45)
{\textbf{风险分担}\\[2pt]
\small 在参与者之间\\重新配置风险敞口};


\node[box] (sharingout) at (2.7,0.55)
{改变风险由谁承担};

%------------------------------------------------
% Risk transformation
%------------------------------------------------
\node[box] (transform) at (8.3,2.45)
{\textbf{风险转换}\\[2pt]
\small 通过状态依存支付\\重塑风险敞口};

\node[box] (transformout) at (8.3,0.55)
{改变风险如何被承担};

%------------------------------------------------
% Branching arrows from common source
%------------------------------------------------
\coordinate (split)
  at ($(heterogeneity.south)+(0,-0.35)$);

% Common stem
\draw[line width=0.9pt, FinanceBlue]
  (heterogeneity.south) -- (split);

% Left branch: Risk sharing
\draw[flow]
  (split)
  -| (sharing.north);

% Right branch: Risk transformation
\draw[flow]
  (split)
  -| (transform.north);

\draw[flow]
  (sharing.south)
  -- (sharingout.north);

\draw[flow]
  (transform.south)
  -- (transformout.north);

\end{tikzpicture}

\caption{风险分担与风险转换}
\label{fig:ch1-risk-sharing-transformation}
\end{figure}

风险分担改变参与者之间风险敞口的分布。例如，保险将特定损失分摊给众多投保人。风险转换（risk transformation）则通过在未来不同状态下具有不同支付的金融请求权来重塑风险敞口。例如，债务和股权把企业现金流划分为优先请求权和剩余请求权，而衍生工具的支付则与价格、利率或其他标的变量的变化相联系。


\subsection{提供流动性并支持支付}

流动性（liquidity）涉及按期履行支付义务的能力，以及资产在不发生明显延误或价值损失的情况下转化为购买力的难易程度。当现金流入与支付义务无法完全同步时，就会产生流动性需求。金融安排通过在需要时提供资金，弥合这一时间缺口。

例如，一家具有清偿能力的企业可能持有价值较高的资产，却只有少量现金。如果必须大幅降价才能迅速出售这些资产，市场流动性（market liquidity）就有限；如果企业无法以这些资产为担保借款，也无法获得短期融资，其融资流动性（funding liquidity）也有限。

现金和其他流动资产有助于履行支付义务，而支付与结算系统使资金转移和义务清偿成为可能。

流动性供给可能涉及期限转换和流动性转换。例如，银行以可随时提取的存款，为期限较长、流动性较低的贷款融资。这种结构为存款人提供流动性，但也造成流动负债与低流动性资产之间的错配；当提款或融资需求上升时，错配可能成为脆弱性的来源。

\subsection{生产信息并约束行为}

金融决策依赖于有关预期现金流、风险以及各方激励的信息。在投入资源之前，必须生产、解释和比较信息。筛选、财务分析、审计、信用评估与市场交易，都对这一过程有所贡献。在市场中，价格还汇集了众多参与者决策所揭示的信息。

融资之后，信息仍然重要。监督提供有关绩效和资金用途的证据，而治理安排与融资条件影响后续决策及激励。金融体系不只是根据现有信息配置资源，还会生产信息，并在信息不完整时建立约束行为的机制。

\subsection{发现价格}

金融体系通过交易形成价格。估值根据预期现金流、时间、风险及或有情形进行估计，而市场价格通过买卖双方的相互作用形成。参与者依据不同的信息、信念、约束和目标采取行动。他们的交易决策影响供求，进而引起价格调整。随着新信息和新交易进入市场，价格不断变化，因此价格发现（price discovery）是一个持续过程。

这些价格又进入后续金融决策。利率影响储蓄、借款和融资成本；债券收益率和股票价格影响发行与投资；汇率改变外币现金流的本币价值；资产价格影响抵押品价值与借款能力。因此，金融价格将交易中的信息传递到后续经济决策中。

\section{金融体系如何组织}

金融体系通过若干相互联系的组成部分履行这些功能，包括货币与支付系统、金融市场、金融机构及金融工具。虽然可以在分析上加以区分，这些部分在实践中相互作用。

\subsection{货币、支付与结算}

金融交易产生最终必须清偿的义务。清偿需要被接受的支付手段。货币提供表达金融义务的共同单位，也提供清偿这些义务的手段。

货币通过支付系统转移时，可能涉及清算与结算。清算（clearing）确定交易产生的义务，并可通过轧差减少这些义务；结算（settlement）通过转移被接受的结算资产完成支付。货币、支付安排、清算与结算共同构成清偿金融义务的制度基础。

\subsection{金融市场}

金融市场组织金融请求权的发行与交换，并提供价格形成机制。新证券的发行使企业、政府和其他借款人能够从投资者处获得资金；二级市场交易则重新配置既有请求权和风险敞口。价格产生于具有不同信息、信念和约束的参与者之间的交易，并为后续融资和投资决策提供依据。

然而，市场交换并不消除评估交易对手、生产和核实信息、设计契约及监督绩效的成本。直接交易也未必总能满足参与者对委托投资、风险汇集以及流动性和期限转换的需求。

\subsection{金融机构}

金融机构通过专门化组织履行金融功能，而不只依靠直接的市场交换。专业化可以降低信息生产、交易执行、风险管理及金融关系监督的成本。银行、资产管理机构、保险公司和其他中介，通过不同形式的汇集、转换与委托，将这些功能结合起来。

通过这些活动，金融中介（financial intermediation）影响金融关系中信息、控制权和风险的配置。替他人作决策的委托关系可能产生代理问题，而期限、流动性和风险转换可能产生需要机构管理的风险敞口。因此，治理、契约安排与监管，既影响金融功能的履行，也影响中介活动的相关风险。

\subsection{金融工具}

金融工具规定未来现金流、风险和权利如何在参与者之间分配。债务明确承诺支付和债权人的优先受偿权；股权赋予剩余请求权，通常还包括投票权。保险和衍生工具的支付，则以特定损失、价格、利率或其他未来事件为条件。

一种工具形成的分配取决于其契约条款。期限、偿付顺序、抵押品、限制性条款、保证金要求和嵌入式期权，决定何时支付、如何吸收损失，以及控制权何时可能转移。金融工具界定了跨时间、跨未来状态配置现金流、风险和控制权的权利与义务。

\par\medskip
\noindent
表~\ref{tab:ch1-architecture} 概括了这些组织形式如何支持上一节讨论的金融功能。

\begin{table}[H]
\centering

\caption{金融体系的组织形式}
\label{tab:ch1-architecture}

\begin{tabularx}{\textwidth}{@{}p{3.0cm}X X@{}}
\toprule
\textbf{形式}
&
\textbf{组织的内容}
&
\textbf{主要功能}
\\
\midrule

货币、支付与结算
&
义务的清偿与转移
&
支付与流动性
\\

金融市场
&
请求权与风险的交换
&
配置、流动性、价格发现
\\

金融机构
&
专业化中介与委托
&
信息、流动性、风险分担
\\

金融工具
&
现金流、权利与风险敞口
&
配置、风险转换、激励
\\

\bottomrule
\end{tabularx}

\end{table}

\section{金融与实体经济}

金融影响家庭、企业和政府能够作出哪些经济选择。它影响哪些项目得以实施、消费如何跨期分配、冲击如何被吸收，以及扰动如何通过资产负债表和市场传播。

\subsection{资本形成与经济发展}

金融将可用资源引向生产性投资，从而支持资本形成。这一过程涉及识别可行的投资机会、安排融资、监督资金使用，并分配相关风险。

在留存收益尚不足以提供资金时，企业可能已经发现可行的投资机会。通过发行对未来现金流的金融请求权，企业可以获得今天的购买力，并用来购置生产性资产。

从总量看，金融对经济发展的贡献不仅取决于融资规模，还取决于金融体系如何配置资本、设计激励和分配风险。融资增加并不必然意味着生产性投资增加。

\subsection{金融创新与金融服务可得性}

金融创新通过改变交易成本、信息要求、契约结构或请求权的可转让性，改变金融功能的履行方式。支付、核实、交易或风险转移成本的下降，可以使过去因成本过高或难以安排而无法实现的交易成为可能。

这些变化可能使面临高固定成本、抵押品不足或与传统中介联系薄弱的家庭和企业，更容易获得金融服务。金融服务可得性提高，有助于放松支付、储蓄、借款、投资和风险管理的约束，使更多经济活动能够获得融资或得以开展。

金融创新对经济发展的贡献，取决于它缓解了哪些金融摩擦，以及由此带来的可得性变化如何影响经济活动。例如，交易成本或信息成本下降，可以使原本安排成本过高的生产性投资或风险分担成为可能。然而，如果创新削弱筛选、便利过度杠杆，或引导资金流向生产率较低的用途，仅仅提高可得性并不能保证生产性投资增加。



\begin{figure}[H]
\centering
\begin{tikzpicture}[
  >=Latex,
  box/.style={
    draw,
    rounded corners,
    text width=5.2cm,
    minimum height=0.95cm,
    align=center,
    inner sep=6pt,
    line width=0.8pt
  },
  sidebox/.style={
    draw,
    rounded corners,
    text width=3.8cm,
    minimum height=0.95cm,
    align=center,
    inner sep=6pt,
    line width=0.8pt
  },
  flow/.style={
    -Latex,
    line width=0.9pt,
    FinanceBlue
  },
  sideflow/.style={
    -Latex,
    dashed,
    line width=0.8pt,
    gray
  }
]

\node[box] (innovation) at (0,4.2)
  {\textbf{金融创新}};

\node[box] (frictions) at (0,2.0)
  {\textbf{金融摩擦的变化}\\[3pt]
   \small 交易、信息与契约成本};

\node[box] (access) at (0,-0.4)
  {\textbf{金融服务可得性的变化}};

\node[box] (activity) at (0,-2.8)
  {\textbf{实体经济活动的变化}\\[3pt]
   \small 储蓄、借款、投资与风险管理};

\node[sidebox] (risks) at (6.2,2.0)
  {\textbf{新的风险与扭曲}\\[3pt]
   \small 杠杆、筛选薄弱与集中度};

\draw[flow] (innovation.south) -- (frictions.north);
\draw[flow] (frictions.south) -- (access.north);
\draw[flow] (access.south) -- (activity.north);

\draw[sideflow] (frictions.east) -- (risks.west);
\draw[sideflow]
  (risks.south) |- (activity.east);

\end{tikzpicture}

\caption{金融创新、金融服务可得性与实体经济活动}
\label{fig:ch1-innovation-access}
\end{figure}

\subsection{脆弱性与金融放大}

使金融发挥作用的安排，也可能造成脆弱性。杠杆放大收益与损失，流动性转换造成资产负债错配；当许多参与者作出相似反应时，市场调整也可能破坏稳定。风险转移可以分散敞口，却也可能使风险的最终承担者更难识别。

金融脆弱性（financial fragility）往往源于反馈效应，而非某一次损失。资产价格下跌会降低抵押品价值，收紧借款约束，迫使资产出售。这些出售又可能进一步压低价格，再次削弱抵押品价值和信贷条件。图~\ref{fig:ch1-fragility-loop} 展示了这一放大机制。

随着金融条件收紧，这些效应传播到实体经济。抵押品价值下降、信贷可得性降低的企业可能削减投资；面临借款条件收紧或财富损失的家庭可能减少消费。金融放大能够把最初的扰动转化为更广泛的支出和生产收缩。

\begin{figure}[H]
\centering
\begin{tikzpicture}[
  >=Latex,
  box/.style={
    draw,
    rounded corners,
    text width=3.8cm,
    minimum height=0.9cm,
    align=center,
    inner sep=5pt,
    line width=0.8pt
  },
  realbox/.style={
    draw,
    rounded corners,
    text width=4.0cm,
    minimum height=0.9cm,
    align=center,
    inner sep=5pt,
    line width=0.8pt
  },
  flow/.style={
    -Latex,
    line width=0.9pt,
    FinanceBlue
  }
]

% Financial amplification loop
\node[box] (prices) at (0,3.6)
  {\textbf{资产价格下跌}};

\node[box] (collateral) at (0,2.0)
  {\textbf{抵押品价值下降}};

\node[box] (credit) at (0,0.4)
  {\textbf{信贷收紧}};

\node[box] (sales) at (-4.8,0.4)
  {\textbf{被迫出售资产增加}};

% Real-economy effect
\node[realbox] (real) at (4.8,0.4)
  {\textbf{实体经济活动减弱}\\[2pt]
   \small 投资与消费下降};

% Main financial loop
\draw[flow] (prices.south) -- (collateral.north);
\draw[flow] (collateral.south) -- (credit.north);
\draw[flow] (credit.west) -- (sales.east);
\draw[flow] (sales.north) |- (prices.west);

% Transmission to the real economy
\draw[flow] (credit.east) -- (real.west);

\end{tikzpicture}

\caption{金融放大及其向实体经济的传导}
\label{fig:ch1-fragility-loop}
\end{figure}

\section{从功能视角理解变化中的金融体系}

技术变化可以改变由谁履行金融功能、如何组合这些功能，以及通过何种组织形式履行。活动可能在银行、市场、数字平台和技术基础设施之间转移；信息处理任务可能越来越多地由算法支持或执行；金融请求权也可能以新方式发行、交易或结算。这些变化能够显著改变金融的组织形式，却不会消除配置资源、生产信息、管理风险、提供流动性，以及支持可信金融义务的根本需求。

人工智能说明了这一区别。AI 可以改变筛选、监督、交易及金融决策的执行方式，而共同数据、模型或决策规则的使用，也可能产生新的集中和相互依赖。数字平台、可编程契约、代币化资产与新型支付安排，也会带来类似的组织形式变化。它们可能改变市场、中介与基础设施的边界，或改变金融请求权的创设、转移和结算方式。

\TailModule{summary}{本章小结}

金融交换往往将今天可用的资源与其他时点的支出和投资机会分开。即使未来结果已知，这种分离也使储蓄、借款和融资具有作用。当未来结果不确定时，决策还依赖于可能不完整或分布不均的信息；契约与制度则有助于界定权利和义务，使未来承诺可信且能够执行。

金融通过支持投资、风险分担和金融服务的获取，扩大经济活动的可能性，但同样的安排也会产生脆弱性。杠杆放大损失，流动性转换造成融资错配，个体理性的反应可能通过价格和资产负债表相互作用，放大整个金融体系的冲击。人工智能、数字平台及其他技术，可以改变履行金融功能的成本和组织形式，却不会消除产生这些功能的经济问题。

\TailModule{keyterms}{关键术语}

跨期交换（Intertemporal exchange）；风险分担（risk sharing）；逆向选择（adverse selection）；道德风险（moral hazard）；承诺（commitment）；执行（enforcement）；资本配置（capital allocation）；货币流动性（monetary liquidity）；市场流动性（market liquidity）；融资流动性（funding liquidity）；信息生产（information production）；价格发现（price discovery）；金融中介（financial intermediation）；金融脆弱性（financial fragility）。




\TailModule{questions}{讨论题}

\begin{enumerate}
\item 即使未来结果完全确定，为什么跨期交换仍然使金融具有作用？
\item 不确定性如何使信息具有经济价值？为什么仅有信息仍不足以支持金融交换？
\item 为什么一种金融安排能够同时解决多个经济问题？请以银行贷款、保险契约或交易中的证券为例说明。
\item 为什么提供流动性既能提高经济效率，又会造成金融脆弱性？
\end{enumerate}

\TailModule{appprobs}{应用题}

\begin{enumerate}
\item 某人今天收入为 40，下一期收入为 66，借入与贷出利率均为 10\%。若当前消费为 50，利用式~\eqref{eq:ch1_intertemporal_budget} 求下一期消费。此人今天是在借款还是储蓄？金额是多少？

\item 一家金融中介主要用可随时兑付的负债为五年期贷款融资。解释这种结构在正常时期的经济收益，以及它在压力时期可能变得不稳定的机制。指出其提供的金融功能和脆弱性来源。
\end{enumerate}

\TailModule{aiprompt}{金融 AI 实验室}

\textbf{实验：追踪金融功能}

选择一种金融创新或数字金融服务，请 AI 系统解释它“新”在何处。随后，围绕本章的分析框架重新组织回答：其背后的经济问题、涉及的金融功能、履行功能的组织形式、现金流和风险的分配，以及这种安排可能失效的条件。识别需要外部核实的事实性主张，以及 AI 回答中将技术新颖性误认为经济功能变化的部分。








\TailModule{furtherread}{延伸阅读}

\begin{itemize}

\item Merton, Robert C. (1995),
“金融中介的功能视角”（``A Functional Perspective of Financial Intermediation''），
\emph{Financial Management}, 24(2), 23--41.

\item Merton, Robert C., and Zvi Bodie (1995),
“分析金融环境的概念框架”（``A Conceptual Framework for Analyzing the Financial Environment''），
载 Dwight B. Crane 等，\emph{The Global Financial System: A Functional Perspective}（《全球金融体系：功能视角》），Harvard Business School Press，第~3--31~页。


\item Levine, Ross (2005),
“金融与增长：理论与证据”（``Finance and Growth: Theory and Evidence''），
载 Philippe Aghion 与 Steven N. Durlauf（编），\emph{Handbook of Economic Growth}（《经济增长手册》），第~1A~卷，Elsevier，第~865--934~页。



\end{itemize}

\ChapterTail{summary,keyterms,questions,appprobs,aiprompt,furtherread}

\chapter{货币、支付与流动性架构}
\label{chap:Money_Payments_and_Liquidity_Architecture}

\begin{openingcase}{支付应用并不是结算资产}

到 2010 年代末，移动支付和网络支付已深深融入中国的日常商业活动。对用户而言，交易往往十分简单：打开应用、扫码、确认金额，支付就完成了。然而，在这个界面背后，支付的制度路径正在变化。

此前，非银行支付机构直接连接多家商业银行，处理涉及银行账户的网络支付。中国人民银行要求支付机构切断这些直连，并通过集中清算基础设施处理相关交易。截至 2019 年 1 月 14 日，非银行支付机构与商业银行共同开展的支付业务已通过网联清算或中国银联处理。\footnote{中国人民银行，《中国人民银行年报 2018》，有关网联平台及非银行支付机构从与商业银行直连模式迁移的讨论。}

对大多数客户而言，支付体验基本没有变化，尽管其背后的制度架构已发生重大改变。

\end{openingcase}

%========================================================
% Core Questions
%========================================================

\begin{tcolorbox}[
  enhanced,
  colback=FinanceBlue!6,
  colframe=FinanceBlue,
  boxrule=0.6pt,
  sharp corners,
  left=8pt,
  right=8pt,
  top=7pt,
  bottom=7pt,
  before skip=1em,
  after skip=1.6em,
  title=\textbf{核心问题},
  coltitle=white,
  colbacktitle=FinanceBlue,
  fonttitle=\bfseries,
  attach boxed title to top left={xshift=6pt,yshift=-2mm},
  boxed title style={
    boxrule=0pt,
    sharp corners,
    left=5pt,
    right=5pt,
    top=2pt,
    bottom=2pt
  }
]
\begin{enumerate}[
  leftmargin=*,
  label=\textcolor{FinanceBlue}{\textbf{\arabic*.}},
  itemsep=0.55em,
  topsep=0.2em
]
\item 为什么有些金融请求权能够充当货币，而另一些不能？
\item 支付指令如何经过清算和结算处理？支付何时具有最终性？
\item 为什么现代货币体系同时依赖中央银行货币和商业银行货币？什么机制支持二者按面值兑换？
\end{enumerate}
\end{tcolorbox}

\section{作为金融约定的货币}

为什么一个人愿意接受他人提供的某种物品或账户余额，以交换商品、劳动或金融请求权？答案不仅取决于货币的形式，还取决于使货币被广泛接受、并能用于清偿义务的制度。

\subsection{作为社会与制度安排的货币}

货币的有用性，依赖于人们共同预期他人也会接受它作为支付。如果每个卖方在接受支付之前，都必须检查发行人的资产支持、法律地位和资产的可接受性，日常交换就会变得昂贵。货币通过建立普遍接受的计量单位和结算惯例，减轻了这种负担。

这一特性具有一定的自我强化作用。接受某种资产的人越多，它用于支付的价值就越高；接受范围越广，其他人持有和使用它的理由也越充分。这种反馈形成货币交换中的网络效应（network effects）。但仅有网络效应还不够，参与者还必须相信货币单位将持续可用，以其进行的支付会被认可并得到结算。
现代货币体系通过法律、支付基础设施、中央银行安排、监管和共同预期来支持这种信心。

货币并不只是一种社会惯例。其使用既依赖共同接受，也依赖支持货币请求权的制度安排。例如，银行存款是存款人的资产，也是发行银行的负债。因此，货币的普遍可接受性既建立在共同预期上，也建立在支持支付和结算的安排上。

\subsection{货币的核心职能}

一项货币请求权在得到广泛接受和制度支持后，就能够履行三项核心货币职能。

\emph{记账单位}（unit of account）为价格和金融义务提供共同语言。贷款协议可以规定支付人民币 100 万元，而不必约定一组不断变化的商品。债券可以用与记录发行人收入、投资者财富相同的单位，承诺支付利息和本金。缺乏共同单位，会增加比较和缔约的成本。

\emph{交换媒介}（medium of exchange）职能使货币能够促进交易并清偿义务。现金可以直接结清面对面的交易，银行存款则可以通过支付系统转移。


\begin{table}[H]
\centering
\caption{货币的核心职能}
\label{tab:ch2-money-functions}
\renewcommand{\arraystretch}{1.4}
\begin{tabularx}{\textwidth}{
    @{}
    >{\raggedright\arraybackslash}p{3.0cm}
    >{\raggedright\arraybackslash}X
    >{\raggedright\arraybackslash}X
    @{}
}
\toprule
\textbf{职能}
& \textbf{经济作用}
& \textbf{金融意义} \\
\midrule

记账单位
& 为价格与义务提供共同计量标准
& 使契约、会计核算和估值具有可比性 \\

交换媒介
& 便利交易并清偿义务
& 支持商品和服务的高效交换 \\

价值储藏
& 将购买力跨期保存
& 使财富以具有流动性的名义形式持有 \\

\bottomrule
\end{tabularx}
\end{table}

\emph{价值储藏}（store of value）职能使购买力能够跨期保存。货币并非完美的价值储藏手段，因为通货膨胀会削弱其实际购买力。即便如此，它通常仍比许多其他资产更稳定、流动性更高。

这些职能为区分货币、其他金融资产与支付工具提供了有用依据。例如，国债可以是良好的价值储藏手段，却不充当交换媒介。支付卡便利交易，但本身并不是货币：借记卡使持卡人能够使用银行存款，信用卡则使其能够动用授信额度。一种外币在某一司法辖区可能履行货币职能，在另一辖区则主要作为金融资产持有。


\subsection{货币的资产负债表结构}

现代经济中的大部分货币，表现为记录在机构资产负债表上的金融请求权。最重要的形式是商业银行存款、中央银行准备金和通货，尽管它们在持有资格和用途上有所不同。

商业银行存款是银行的负债、存款人的资产。中央银行准备金是中央银行的负债、合资格账户持有人的资产，这些持有人主要是银行和其他指定机构。通货也记录为中央银行负债，尽管现代法定货币通常不能按固定比率兑换为黄金或另一种指定资产。

\begin{table}[H]
\centering
\caption{银行存款同时是资产和负债}
\label{tab:ch2-deposit-taccount}

\renewcommand{\arraystretch}{1.25}

\begin{tabularx}{0.9\textwidth}{
    @{}
    >{\raggedright\arraybackslash}X
    >{\raggedright\arraybackslash}X
    >{\raggedright\arraybackslash}X
    @{}
}
\toprule
\multicolumn{1}{c}{\textbf{家庭}}
&
\multicolumn{2}{c}{\textbf{商业银行}}
\\
\cmidrule(r){1-1}
\cmidrule(l){2-3}

\textbf{资产}
&
\textbf{资产}
&
\textbf{负债}
\\

银行存款：人民币 100 元
&
--
&
银行存款：人民币 100 元
\\

\bottomrule
\end{tabularx}
\end{table}

表~\ref{tab:ch2-deposit-taccount} 展示了同一笔人民币 100 元存款如何出现在两张资产负债表上：它是家庭的资产，同时是银行的负债。

资产负债表视角说明，不同形式的货币在经济上并不完全相同。
人民币 100 元银行存款和 100 元纸币采用相同记账单位，通常按面值被接受，但它们的发行人和资产负债表关系不同。存款是商业银行的负债、存款人的资产；纸币则是中央银行的负债、持有人的资产。

\begin{table}[H]
\centering
\caption{部分货币形式及其资产负债表地位}
\label{tab:ch2-money-claims}

\renewcommand{\arraystretch}{1.15}
\setlength{\tabcolsep}{7pt}

\begin{tabularx}{\textwidth}{
    @{}
    >{\raggedright\arraybackslash}p{3.35cm}
    >{\raggedright\arraybackslash}p{2.85cm}
    >{\raggedright\arraybackslash}X
    @{}
}
\toprule
\textbf{货币形式}
& \textbf{发行人}
& \textbf{资产负债表关系} \\
\midrule

通货\newline
{\small （纸币与硬币）}
& 中央银行
& 中央银行负债；公众持有的资产 \\
\addlinespace[0.3em]

中央银行准备金
& 中央银行
& 中央银行负债；合资格账户持有人的资产 \\
\addlinespace[0.3em]

商业银行存款
& 商业银行
& 银行负债；存款人的资产 \\

\bottomrule
\end{tabularx}
\end{table}

同样的资产负债表逻辑，也有助于说明商业银行存款如何创造。假设银行新发放人民币 100 元贷款，并在借款人的存款账户中记入同额资金。在银行的资产负债表上，

\begin{equation}
\Delta \text{贷款}=+100,
\qquad
\Delta \text{存款}=+100.
\label{eq:ch2-loan-deposit}
\end{equation}

因此，贷款在发放时创造了一笔新存款。借款人同时获得存款资产和相应的贷款负债，其净财富并未等额增加。存款创造仍然受到银行整体资产负债表状况和监管约束的限制。

\subsection{信任、可信度与货币稳定}

货币降低信息成本，因为人们不必在每笔交易中重新评估每项货币请求权背后的发行人。银行转账的收款人通常依赖存款的可转移性，而不必评估付款人的开户银行；通货持有人也不必检查中央银行的资产负债表。货币交换的效率，建立在无需反复核实便能将被接受的货币请求权视为可靠的基础上。

这种信任来自支持货币体系的制度和安排的可信度。货币单位必须足够稳定，才能支持定价和缔约；货币请求权必须能够按承诺转移或兑换。可靠的支付系统，以及明确所有权和结算最终性的法律规则，也支持这种可信度。

货币信任下降会通过多种方式影响行为。如果存款人担心银行无法满足提款需求，可能将存款换成通货，或将资金转至其他银行。如果对私人发行货币请求权的信心减弱，该请求权可能折价交易，或不再被接受。如果通胀削弱人们对记账单位的信心，家庭和企业可能缩短契约期限、对价格实行指数化，或转向其他货币与资产。

货币稳定还取决于这样的预期：以同一单位计价的不同货币形式，将保持可互换性。这一预期支持整个经济中的共同定价、缔约和支付。

\section{支付系统与金融网络}

货币只有能够在用户之间转移，才能充当支付手段。支付系统（payment system）是用于转移资金、清偿货币义务的一整套工具、规则、机构、技术与结算安排。

\subsection{支付工具与传统支付基础设施}

支付工具用于发起和传递支付指令。现金将支付工具与结算资产合为一体：交付通货即可直接完成交易。大多数非现金支付，则将支付指令与用于结算的资产分开。

非现金支付工具在由谁发起支付、如何传递指令方面存在差异。支票和贷记转账由付款人发起，直接借记则由收款人在事先授权下发起。银行卡和移动支付依靠支付网络及数字界面来认证用户、授权交易和传送支付指令。

\begin{figure}[H]
\centering

\resizebox{0.94\textwidth}{!}{%
\begin{tikzpicture}[
  >=Latex,
  box/.style={
  draw,
  rounded corners,
  align=center,
  minimum height=0.85cm,
  text width=1.5cm,
  inner xsep=2pt,
  inner ysep=3pt,
  line width=0.8pt
  },
  flow/.style={
    -Latex,
    line width=0.85pt,
    FinanceBlue
  },
  note/.style={
    font=\scriptsize,
    align=center,
    text=FinanceBlue,
    fill=white,
    inner sep=1pt
  }
]

\node[box] (payer) at (0,0) {付款人};
\node[box] (bankA) at (3.6,0) {付款人银行};
\node[box, text width=2.2cm] (infra) at (6.3,0)
{清算与\\结算系统};

\node[box] (bankB) at (9.0,0) {收款人银行};
\node[box] (payee) at (11.8,0) {收款人};

\draw[flow] (payer) -- 
  node[midway, above=6pt, note]{支付\\指令}
  (bankA);

\draw[flow] (bankA) -- (infra);

\draw[flow] (infra) -- (bankB);

\draw[flow] (bankB) -- 
  node[midway, above=5pt, note]{贷记}
  (payee);

\end{tikzpicture}%
}

\caption{简化的跨行支付链条}
\label{fig:ch2-payment-chain}
\end{figure}

图~\ref{fig:ch2-payment-chain} 展示了一条简化的跨行支付链条。如果付款人和收款人在同一家银行开户，交易可以通过减少一个存款账户、增加另一个账户的余额在行内结算；若双方使用不同银行，支付就会产生银行间义务。参与银行随后必须相互结算，通常使用双方均接受的结算资产。

因此，支付网络不只是一套通信系统。其设计影响义务何时产生和清偿、参与者需要多少流动性，以及结算最终完成之前面临的信用和操作风险敞口。

\subsection{清算、轧差与支付义务}

清算（clearing）是传送、核对和确认支付指令，并将其转化为待结算义务的过程。在适用的情况下，这些义务可在结算前进行轧差。

考虑一组分别以 $i$ 和 $j$ 标识的银行。令 $p_{ij}$ 表示银行 $i$ 必须向银行 $j$ 支付的金额，则银行 $i$ 的付款总额为

\begin{equation}
G_i^{\mathrm{out}}
=\sum_{j\neq i} p_{ij},
\label{eq:ch2-gross-outgoing}
\end{equation}
其净结算头寸可表示为

\begin{equation}
N_i=\sum_{j\neq i} p_{ij}-\sum_{j\neq i} p_{ji}.
\label{eq:ch2-net-position}
\end{equation}

$N_i$ 为正表示净付款义务，为负表示净收款。在封闭的清算系统中，

\begin{equation}
\sum_i N_i=0.
\label{eq:ch2-net-zero}
\end{equation}

轧差（netting）可以显著降低结算所需流动性。假设银行 A 欠银行 B 100，银行 B 欠银行 C 90，银行 C 欠银行 A 80。表~\ref{tab:ch2-netting-example} 列示了相应的总额流量和净结算头寸。

\begin{table}[H]
\centering
\caption{总额支付与净结算头寸}
\label{tab:ch2-netting-example}
\begin{tabular}{@{}lrrr@{}}
\toprule
\textbf{银行} & \textbf{付款} & \textbf{收款} & \textbf{净头寸} \\
\midrule
A & 100 & 80 & 支付 20 \\
B & 90 & 100 & 收到 10 \\
C & 80 & 90 & 收到 10 \\
\bottomrule
\end{tabular}
\end{table}

三笔总额支付合计为 270，但经过多边轧差后，系统只需结算合计 20 的净借方义务。银行 A 支付 20，银行 B 和 C 各收到 10。轧差降低了结算所需流动性。

在延迟净额结算系统中，节省流动性也有代价。支付义务可能直到清算周期结束才结算，使参与者面临其他参与者违约的风险。支付系统设计需要在流动性效率、结算最终性和风险控制之间权衡。

\subsection{结算、支付最终性与结算风险}

结算（settlement）是通过结算机构账簿上的转账来清偿支付义务。支付最终性（payment finality）是指根据适用的系统规则和法律，支付被视为完成且不可撤销的时点。

一笔简单的跨行转账可以说明清算与结算的区别。假设银行 A 的客户向银行 B 的客户支付人民币 100 元。如果支付用中央银行准备金结算，主要的资产负债表变化如表~\ref{tab:ch2-cross-bank-payment} 所示。

\begin{table}[H]
\centering
\caption{人民币 100 元跨行支付的资产负债表影响}
\label{tab:ch2-cross-bank-payment}

\renewcommand{\arraystretch}{1.2}

\begin{tabularx}{\textwidth}{
    @{}
    >{\raggedright\arraybackslash}p{3.0cm}
    >{\raggedright\arraybackslash}X
    >{\raggedright\arraybackslash}X
    @{}
}
\toprule
\textbf{机构}
& \textbf{资产}
& \textbf{负债} \\
\midrule

银行 A
& 准备金 $-100$
& 付款人存款 $-100$ \\

银行 B
& 准备金 $+100$
& 收款人存款 $+100$ \\

中央银行
& ---
& 对银行 A 的准备金负债 $-100$；对银行 B 的准备金负债 $+100$ \\

\bottomrule
\end{tabularx}
\end{table}

准备金从银行 A 转至银行 B，清偿了银行间义务。客户的存款并非从一家银行“搬运”到另一家银行，而是银行 A 减少对付款人的负债，银行 B 增加对收款人的负债，同时中央银行准备金在两家银行之间转移。当这一结算根据系统规则成为不可撤销、无条件的转移时，就达到了结算最终性。

结算资产应使义务得以清偿，而不在结算时引入显著的额外信用风险。中央银行货币被广泛用于具有系统重要性的国内支付系统结算。

在达到结算最终性之前，参与者仍面临结算风险：预期结算可能延迟，或无法完成。这类失败可能引发几种相关风险，如表~\ref{tab:ch2-payment-risk} 所示。

\begin{table}[H]
\centering
\caption{支付系统风险的主要形式}
\label{tab:ch2-payment-risk}

\renewcommand{\arraystretch}{1.2}

\begin{tabularx}{\textwidth}{
    @{}
    >{\raggedright\arraybackslash}p{3.0cm}
    >{\raggedright\arraybackslash}X
    @{}
}
\toprule
\textbf{风险}
& \textbf{经济含义} \\
\midrule


信用风险
& 参与者可能在到期时或此后任何时点，都无法全额履行义务。 \\

流动性风险
& 参与者可能无法在到期时履行义务，但可以在较晚时间全额结算。 \\

操作风险
& 技术故障、网络事件、人为差错或基础设施中断，可能妨碍处理或结算。 \\

法律风险
& 规则、所有权、轧差或最终性方面的不确定性，可能削弱可执行性。 \\

\bottomrule
\end{tabularx}
\end{table}

这些风险的重要性，部分取决于结算的组织方式。实时全额结算（real-time gross settlement）系统逐笔、连续地结算支付，缩短银行间风险敞口的持续时间，却增加日间流动性需求。延迟净额结算（deferred net settlement）系统则先累积义务、进行轧差，再结算，因而降低流动性需求，却使风险敞口持续更久。

结算设计需要在流动性效率、实现最终性的速度和风险控制之间权衡。

\subsection{数字支付网络与平台}

数字化改变支付的发起、认证和路由方式，但没有消除区分货币请求权、清算与结算的必要。移动钱包或支付应用可能让用户感觉支付即时完成，但其底层价值仍可能是银行存款、受监管的支付请求权或另一种货币资产。关键问题不仅是如何支付，还包括转移了什么请求权、如何完成最终结算。

支付网络连接用户、商户、金融机构、支付服务提供商，以及清算或结算基础设施。其价值不仅依赖技术，还依赖参与范围、互操作性、规则和信任。更大的网络可以使支付更方便、更便宜，但也可能造成对少数平台或基础设施的依赖。

\begin{figure}[H]
\centering

\begin{tikzpicture}[
  >=Latex,
  layer/.style={
    draw,
    rounded corners,
    align=center,
    minimum height=0.95cm,
    text width=0.82\textwidth,
    inner xsep=6pt,
    inner ysep=5pt,
    line width=0.7pt,
    font=\footnotesize
  },
  arrow/.style={
    -Latex,
    line width=0.75pt,
    FinanceBlue
  }
]

\node[layer] (interface) at (0,2.4)
{\textbf{用户界面层}\\
移动钱包、支付应用、二维码、银行卡或在线收银台};

\node[layer] (network) at (0,0.7)
{\textbf{支付网络与平台层}\\
用户、商户、金融机构、支付服务提供商、规则、互操作性与信任};

\node[layer] (claim) at (0,-1.0)
{\textbf{货币请求权层}\\
银行存款、受监管的支付请求权或其他货币资产};

\node[layer] (settlement) at (0,-2.7)
{\textbf{清算与结算层}\\
计算义务、在适用时轧差，并以结算资产完成最终结算};

\draw[arrow] (interface) -- (network);
\draw[arrow] (network) -- (claim);
\draw[arrow] (claim) -- (settlement);

\end{tikzpicture}

\caption{数字支付的界面、网络、请求权和结算层次}
\label{fig:ch2-digital-payment-layers}
\end{figure}

图~\ref{fig:ch2-digital-payment-layers} 表明，数字支付涉及的并不只是用户界面的变化。同样的可见支付行为，可能基于不同的底层货币请求权、网络安排和结算资产。

在这种分层架构中，平台层协调用户、商户、银行和支付服务提供商之间的互动。平台通过接入管理、身份认证、交易路由和数据处理，降低交易成本与处理延迟。这些安排也产生网络效应：支付服务对参与者的价值可能随着用户和商户数量增加而提高，在创造规模经济的同时，也可能强化市场集中。支付数据能够改善欺诈识别和服务设计，但其收集与使用也引发隐私、市场势力、互操作性和接入方面的问题。

同样的逻辑适用于较新的支付技术。
可编程转账和代币化请求权，可以通过在交易逻辑中嵌入条件，使交易、清算和结算更加紧密地结合。人工智能能够改善欺诈识别、交易监测、流动性预测和支付路由。然而，更快的执行和更好的预测，本身并不能确立法律最终性、保证按面值赎回或消除流动性需求。这一经济问题既涉及计算，也涉及制度。

\section{货币体系中的流动性}

只有在义务到期时，货币请求权能够被使用、转移或兑换，货币才能充当支付手段。因此，支付系统不仅需要货币资产，还需要流动性，即在需要时获得所需结算资产的能力。

\subsection{流动性、可兑换性与经济交换}

流动性不同于资产价值，它描述资产能够多容易地被用于履行到期支付或结算义务。通货可以直接清偿许多零售支付。政府证券可能价值很高、交易活跃，但通常必须先通过出售或抵押借款，将其价值转化为可支用资金，才能支持普通支付。区别在于持有价值与完成支付之间的距离。

可兑换性（convertibility）能够缩短这一距离。这里，可兑换性指将一种货币请求权兑换为另一种的能力，通常按面值且没有显著延迟。银行存款是商业银行的私人负债，但当它能够按面值兑换为通货，并通过支付网络转移时，就能在日常交换中发挥货币作用。

当按面值赎回的信心减弱时，可兑换性的重要性就会显现。存款人可能提取现金，或将余额转至其他机构。更一般地说，不能可靠地按其标示价值赎回的支付请求权，可能不再作为货币的紧密替代物流通。

不过，可兑换性并不消除时间的重要性。假设一家银行预计下午 4 点有大笔资金流入，却必须在上午 10 点付款。即使日终头寸可以平衡，它仍需要上午 10 点的流动性。这种时间错配产生了对准备金、日间信贷，以及其他能在需要时提供结算资产的安排的需求。


\begin{tcolorbox}[
  enhanced,
  breakable,
  colback=white,
  colframe=FinanceBlue,
  boxrule=0.7pt,
  sharp corners,
  borderline west={2.2pt}{0pt}{FinanceBlue},
  title=\textbf{金融视角：可兑换性与流动性},
  coltitle=white,
  colbacktitle=FinanceBlue,
  fonttitle=\bfseries,
  before skip=1.2em,
  after skip=1.2em
]
可兑换性关注金融请求权能否兑换为可用于支付或结算的资产，通常按面值兑换。流动性还要求相关资产在需要的时间、以需要的数量可供使用。因此，一项请求权原则上可兑换，却未必能在义务到期时提供足够流动性。
\end{tcolorbox}

在系统层面，同样的时间错配会引发支付梗阻。机构可能等待收款后才对外付款。当这种行为普遍出现时，即使全日总流动性足够，结算仍可能延迟。因此，流动性安排不仅必须保证结算资产数量充足，还必须保证资产在参与者之间及时可用。

维持流动性具有机会成本。准备金和其他高流动性资产提供即时结算能力，但收益可能低于流动性较低的替代资产。因此，金融机构需要在流动性与收益之间权衡：一方面持有足以履行义务的流动资产，另一方面将其他资金配置到收益更高的用途。

这种权衡反映了流动性和可兑换性在经济交换中的更广泛作用。它们减少支付前变现其他资产的需要，缩短义务清偿的延迟，从而降低交易摩擦，支持货币请求权在整个经济中流通。

\subsection{安全资产与流动性需求}

流动性需求有助于解释安全资产在货币体系中的特殊作用。安全且具有流动性的资产，能够被调动以满足支付和结算需求，同时限制信用损失敞口。

安全性与流动性在概念上不同。这里，安全性主要指信用质量：安全资产承诺的支付无法按期足额实现的风险较低。流动性则关注资产能否在支付或结算义务到期时被便捷地使用或调动。

仅有安全性并不决定资产的流动性。流动性还取决于能否以有利条件迅速出售或质押资产，尤其是在市场承压时。

流动性和抵押服务也会反映在资产收益率中。在其他条件可比的资产中，投资者可能愿意接受较低收益率，以换取更有价值的流动性或抵押服务。这些服务所带来的价值，通常称为便利收益（convenience yield）。

\begin{equation}
y
=
y^{*}-\lambda,
\qquad \lambda>0,
\label{eq:ch2-convenience-yield}
\end{equation}
其中，$y^{*}$ 是其他方面可比、但不提供同等便利服务的资产收益率，$\lambda$ 表示与流动性、抵押及相关服务有关的便利收益。该表达式是一种经济分解，而非定价模型；$\lambda$ 可能随市场条件变化。

在不确定性上升或金融压力时期，对安全和流动资产的需求通常增加。市场参与者可能减少风险较高请求权的持有，转向现金、短期政府证券，或对合资格机构而言，转向中央银行准备金。这种“逃向安全”（flight to safety）的行为能够降低个体风险敞口，却可能收紧其他领域的融资条件和市场流动性。

在机构层面，保全流动性的努力可能放大压力。如果机构同时减少贷款、保留结算资产或出售流动性较差的证券，融资条件可能收紧，资产市场流动性可能恶化。对单个机构审慎的行动，也可能促成全系统流动性压力。

当这种行为广泛出现时，流动性需求就成为整个系统的协调问题。有韧性的货币体系需要在条件变化时支持流动性调动和分配的安排。中央银行流动性工具、抵押品框架和支付系统安排，都是这一制度架构的组成部分。


\section{公共货币与私人货币的架构}
现代货币体系由相互联系的公共和私人资产负债表所发行的货币负债组成。中央银行货币提供共同的结算基础，而商业银行存款构成日常经济活动中所用货币的很大一部分。

\subsection{公共货币与中央银行}

现代货币体系中的公共货币，主要由中央银行负债构成：流通中的通货，以及可在中央银行开户的机构持有的准备金余额。两者都以共同记账单位计价，但在支付与结算系统中履行不同功能。

中央银行发行这些负债，并管理用于银行间结算的准备金账户。一家参与机构向另一家付款时，准备金余额在双方的中央银行账户之间转移。这一转移在机构之间重新分配中央银行负债，并清偿相应的银行间义务。因此，以中央银行货币结算，无须再持有对私人结算中介的额外请求权。

\subsection{私人货币与商业银行}

现代货币体系中的私人货币主要由商业银行存款构成。每笔存款都是存款人的资产和发行银行的负债。作为货币使用的存款以共同记账单位计价，并可通过支付系统转移。尽管不同银行发行的存款采用同一记账单位，每笔存款仍是对特定机构的请求权。

商业银行在扩张资产负债表时创造存款。银行发放贷款时，将贷款记为资产，同时向借款人存款账户记入同额资金，从而创造新的存款负债。存款创造并不要求事先提供等量的中央银行货币。

但存款创造仍受金融和制度约束。随着资产负债表扩张，银行必须维持充足资本和流动性、管理信用风险，并保持履行支付与结算义务的能力。

现代货币体系的一项关键特征是，不同银行发行的存款在法律上属于不同负债，却通常以相同名义价值被接受。这种等价性依赖于连接商业银行存款与中央银行货币、并支持按面值兑换的安排。

\subsection{按面值兑换与货币层级}

不同货币请求权即使是不同发行人的负债，也可能以同一记账单位计价。一家商业银行的人民币 100 元存款，在法律上不同于另一家银行的 100 元存款，也不同于 100 元通货。按面值兑换（par convertibility）意味着这些货币形式通常可以等面值交换，而不按市场决定的折价或溢价交换。

按面值兑换维持名义等价性，却不消除发行人差异。商业银行存款仍是对单家银行的请求权，通货和准备金余额则是中央银行的负债。

这些差异反映为货币层级（hierarchy of money）。商业银行存款可用于家庭和企业的大多数支付，而银行间义务通常以中央银行货币结算。因此，货币层级反映的是发行人和结算地位的差异，而非记账单位的差异。

只要按面值兑换得以维持，这一层级就与货币统一性相容。当对可兑换性的信心减弱时，发行人差异就会变得具有经济意义，货币体系的统一性也可能开始瓦解。

\subsection{维持货币体系的统一性}

按面值兑换使法律上不同的货币负债能够以共同名义价值使用。维持这种货币统一性，取决于私人发行人履行负债的能力，以及结算流动性的可得性。任何一个条件薄弱，都可能破坏货币请求权之间的平价关系。

第一个条件涉及发行银行的财务稳健性。审慎监管限制风险承担并要求充足资本，降低银行无法履行存款负债的可能性。如果银行倒闭，存款保险保护受保障存款，处置安排则为处理倒闭机构提供有序程序。这些机制有助于维持商业银行存款能够按面值兑付的信心。

第二个条件涉及结算流动性。跨行支付产生必须及时清偿的银行间义务。在中央银行持有的准备金余额充当结算资产，流动性工具则可在暂时短缺时提供额外准备金。共同支付规则、互操作性和结算最终性支持可靠的跨行支付，有助于维持按面值兑换。

这些安排共同使法律上不同的货币负债，能够在共同货币体系中作为等价货币发挥作用。

%========================================================
% Summary
%========================================================

\TailModule{summary}{本章小结}

货币不只是被广泛接受的资产。它提供共同记账单位、支持交换，并跨期保存购买力。在现代货币体系中，大多数货币形式也是金融请求权。通货和中央银行准备金是中央银行负债，商业银行存款则是各家银行的负债。它们的有用性依赖于信任、可转移性，以及支持支付和结算的制度安排。

支付不只是传递指令。清算确定参与者之间产生的义务，轧差可以减少待结算金额，结算则通过被接受的结算资产清偿义务。结算的时间和最终性同时影响流动性需求与金融风险。数字支付技术可以改变支付的发起、认证和路由方式，但不会消除这些基本经济功能。

流动性关注按期履行支付与结算义务的能力。可兑换性使货币请求权能够按面值交换，而安全且具有流动性的资产提供额外的流动性和抵押服务。这些服务能够影响资产收益率，并在金融压力时期变得尤为重要。

现代货币体系结合公共货币与私人货币。商业银行通过资产负债表扩张创造存款货币，而中央银行货币提供银行间结算的共同基础。按面值兑换使法律上不同的货币请求权能够以共同名义价值发挥作用。维持这种货币统一性，需要财务稳健的发行人、充足的结算流动性、可靠的支付基础设施，以及能够处理银行困境的制度。

%========================================================
% Key Terms
%========================================================

\TailModule{keyterms}{关键术语}

货币（Money）；记账单位（unit of account）；交换媒介（medium of exchange）；支付系统（payment system）；清算（clearing）；轧差（netting）；结算（settlement）；结算最终性（settlement finality）；流动性（liquidity）；可兑换性（convertibility）；安全资产（safe asset）；便利收益（convenience yield）；中央银行货币（central-bank money）；商业银行货币（commercial-bank money）；按面值兑换（par convertibility）；货币层级。

%========================================================
% Discussion Questions
%========================================================（hierarchy of money

%========================================================
% Discussion Questions
%========================================================）。

\TailModule{questions}{讨论题}

\begin{enumerate}

\item 为什么移动支付应用不一定等同于结算支付所用的资产？数字支付的底层可能有哪些货币请求权？

\item 在延迟净额结算系统中，为什么轧差可以降低流动性要求，却仍使参与者面临结算风险？

\item 商业银行存款虽然是私人机构的负债，为什么仍能充当货币？哪些制度安排支持它按面值被接受？

\item 在金融压力时期，家庭或金融机构为什么可能从银行存款转向通货、其他银行的存款或其他流动资产？这种行为揭示了货币请求权之间的哪些差异？

\end{enumerate}

%========================================================
% Applied Problems
%========================================================

\TailModule{appprobs}{应用题}

\begin{enumerate}

\item 四家银行在清算周期末有以下支付义务：A 欠 B 120；B 欠 C 100；C 欠 D 70；D 欠 A 60；A 还欠 C 30。计算每家银行的付款总额、收款总额和净结算头寸。如果所有义务逐笔全额结算，支付总额是多少？经过完全多边轧差后，哪些银行是净债务人、哪些是净债权人？为完成结算，净债务人总共需要转移多少资金？

\item 银行 X 的一名客户向银行 Y 的客户转账人民币 50,000 元。当银行间义务以准备金余额结算时，列示银行 X、银行 Y 和中央银行的资产负债表变化。再解释：如果两位客户均在银行 X 开户，资产负债表影响有何不同？

\item 两只其他方面可比的一年期证券具有相同预期现金流和相近信用风险。证券 A 的收益率为 3.2\%，可作为中央银行日间信贷的抵押品；证券 B 的收益率为 3.8\%，不具备这一资格。利用式~\eqref{eq:ch2-convenience-yield} 解释 0.6 个百分点的收益率差。要将该差异完全归因于流动性和抵押服务，需要哪些假设？

\end{enumerate}

%========================================================
% Data Application
%========================================================

\TailModule{dataapp}{数据应用}

\textbf{数据中的支付架构：中国，2024 年}

中国人民银行报告显示，2024 年，大额实时支付系统处理业务 3.92 亿笔，金额 8,824.18 万亿元；网联清算平台处理业务约 1.03 万亿笔，金额 520.53 万亿元。\footnote{中国人民银行，《2024 年支付体系运行总体情况》，第三部分“支付系统”。}

\begin{enumerate}

\item 计算两个系统各自的平均每笔交易金额。

\item 比较交易笔数与交易金额。这两个指标揭示了两个系统在经济作用上的哪些差异？


\item 将这些差异与本章讨论的支付系统机制联系起来。为什么大额实时支付系统尤其重视结算流动性和最终性，而处理海量相对小额交易的平台更强调处理能力和较低的单位成本？





\item 使用中国人民银行较新一期支付体系报告重复比较。两个系统的笔数、金额和平均每笔规模如何变化？得出结论之前，检查统计定义或覆盖范围是否发生变化。




\end{enumerate}

%========================================================
% Financial AI Lab
%========================================================

\TailModule{aiprompt}{金融 AI 实验室}

\textbf{实验：追踪支付直至最终结算}

选择一种可以通过可靠来源加以记录的支付方式，例如银行转账、银行卡、二维码支付或数字钱包。请 AI 系统描述从付款人指令到最终结算的完整交易路径。将生成的回答视为待核实的假说，而不是权威描述。

通过五项检查重建交易：识别付款人持有的货币请求权；识别发行该请求权的机构；识别清算安排；识别最终结算使用的资产；识别决定结算最终性的规则或机构。每一步都要求可核实的机构来源，并标注无法确认的陈述。

将核实后的交易路径与 AI 回答比较。识别 AI 在何处将支付界面与底层货币请求权、清算过程、资产负债表转移或结算资产相混淆。最后解释支付过程中哪些部分属于技术，哪些部分依赖金融和制度安排。

%========================================================
% Further Reading
%========================================================

\TailModule{furtherread}{延伸阅读}
\begin{itemize}

\item 中国人民银行（2025），《中国支付体系发展报告（2024）》（\emph{Payment System Report (2024)}）。

\item McLeay, Michael, Amar Radia, and Ryland Thomas (2014),
“现代经济中的货币：导论”（``Money in the Modern Economy: An Introduction''），
\emph{Bank of England Quarterly Bulletin}, 54(1), 4--13.

\item 支付结算体系委员会与国际证券委员会组织技术委员会（2012），《金融市场基础设施原则》（\emph{Principles for Financial Market Infrastructures}），国际清算银行。




\item 国际清算银行（2025），
“下一代货币与金融体系”（``The Next-Generation Monetary and Financial System''），
载《国际清算银行 2025 年度经济报告》（\emph{BIS Annual Economic Report 2025}）第三章。

\end{itemize}

%========================================================
% Tail row in Full Contents
%========================================================

\ChapterTail{summary,keyterms,questions,appprobs,dataapp,aiprompt,furtherread}

\part{信息、契约与估值}
\chapter{信息经济学与金融契约}
\label{chap:Information_Economics_and_Financial_Contracts}

\begin{openingcase}{为什么提高利率反而可能减少放贷？}

假设华信商业银行正在审查荣达机械和鑫海科技两家制造企业的贷款申请。从书面资料看，两家企业很相似：营业收入、杠杆、抵押品和近期财务表现大致相当，但它们的投资项目却截然不同。

荣达计划扩建一条成熟生产线，预期回报不高，但项目成功概率较大。鑫海则考虑一项投机性更强的技术投资：可能获得高得多的收益，但失败风险也明显更大。两家企业的管理者都远比银行更了解这些差异。

应对贷款风险的自然办法似乎是提高利率。贷款风险越高，银行就应要求更多补偿。然而，当借款人的差异无法被银行完全观察时，这一逻辑可能失效。

随着华信提高贷款利率，荣达这样的企业可能认为借款已不再划算。项目回报相对有限，难以承担高融资成本；鑫海这样的企业可能仍愿借款。如果项目成功，上行收益很大；如果失败，部分损失则由贷款人承担。

因此，利率影响的不只是贷款成功偿还时的承诺支付。它还可能影响\emph{谁愿意借款}，以及贷款发放后\emph{借款人愿意承担何种风险}。

当贷款利率足够高时，这些效应可能降低华信的预期回报。银行因而可能限制信贷供给，而不再继续加息，使一些愿意借款的企业无法获得贷款。

\end{openingcase}

\begin{tcolorbox}[
  enhanced,
  colback=FinanceBlue!6,
  colframe=FinanceBlue,
  boxrule=0.6pt,
  sharp corners,
  left=8pt,
  right=8pt,
  top=7pt,
  bottom=7pt,
  before skip=1em,
  after skip=1.6em,
  title=\textbf{核心问题},
  coltitle=white,
  colbacktitle=FinanceBlue,
  fonttitle=\bfseries,
  attach boxed title to top left={xshift=6pt,yshift=-2mm},
  boxed title style={
    boxrule=0pt,
    sharp corners,
    left=5pt,
    right=5pt,
    top=2pt,
    bottom=2pt
  }
]
\begin{enumerate}[
  leftmargin=*,
  label=\textcolor{FinanceBlue}{\textbf{\arabic*.}},
  itemsep=0.55em,
  topsep=0.2em
]
\item 签订契约之前，交易成本和信息不对称如何影响金融交换？

\item 提供融资之后，隐藏行动如何影响激励？

\item 契约、监督、抵押品、治理和执行如何缓解这些问题？为什么无法将其彻底消除？

\end{enumerate}
\end{tcolorbox}

金融契约往往在信息不对称条件下设计：借款人对自身前景和风险类型的了解，可能多于贷款人，而获得资金后的行为可能难以观察。签约前，私人信息影响借款人的选择；签约后，不可观察的行动可能改变风险承担激励。本章考察这些问题如何影响筛选、激励设计及金融风险分配。

\section{交易成本与信息摩擦}

金融交换具有成本。获取信息、评估交易对手、设计契约和执行契约，都需要资源。这些成本影响哪些交易能够达成、金融契约采用何种条款，以及金融关系如何组织。

\subsection{交易成本的组成}

考虑一家寻求外部融资的小型制造企业——明远精密零部件公司。提供资金之前，企业需要联系潜在贷款人或投资者，资金提供者则需要足够信息来评估机会。为此，可能要编制并审计财务报表、评估资产、起草法律文件和确定抵押品。各方还可能规定出现财务困境或违约时的权利、或有情形及救济措施。

上述活动都消耗资源。这些成本降低融资的净收益，并可能决定一笔交易能否产生正剩余。

假设明远考虑建设一条新生产线，今天需投资 $I$，下一期产生不确定支付 $X$。在无摩擦情形下，如果 $X$ 的预期现值超过初始投资，项目在财务上就值得实施。考虑交易成本 $K$ 后，相关剩余变为
\begin{equation}
S = V_{0}-I-K,
\label{eq:ch3-transaction-surplus}
\end{equation}
其中，$V_{0}$ 为项目预期未来支付的现值，$K$ 为安排和维持金融关系的交易成本现值。

\(K\) 的各组成部分发生在金融关系的不同阶段。寻找交易对手、评估相关信息时会产生搜寻和核实成本；契约成本来自对权利、义务及或有情形的规定。融资提供后，观察绩效与行为会产生监督成本；当契约义务未履行时，执行成本就变得重要。

许多活动涉及可以分摊到多笔交易上的固定成本，从而构成专业化的经济依据。家庭评估一笔企业贷款时，难以把信用分析和法律审查成本分摊到多笔交易上；相比之下，银行能够为众多借款人履行这些功能，将相关成本分摊到整个贷款业务中。审计师、信用评级机构和交易所，同样在核实、信息生产和标准化方面提供专业服务。

专业化不能消除交易成本，但可以降低每笔交易履行这些功能的平均成本。

\begin{figure}[H]
\centering
\begin{tikzpicture}[
    >=Latex,
    box/.style={
        draw,
        rounded corners,
        minimum width=2.7cm,
        minimum height=0.85cm,
        align=center
    },
    central/.style={
        draw,
        rounded corners,
        minimum width=3.2cm,
        minimum height=1.0cm,
        align=center,
        line width=0.9pt
    },
    flow/.style={
        -Latex,
        line width=0.8pt,
        FinanceBlue
    }
]

\node[central] (exchange) at (0,0) {金融交换};

\node[box] (search)   at (-4.3,1.6) {搜寻};
\node[box] (verify)   at (-4.3,0)   {核实};
\node[box] (contract) at (-4.3,-1.6) {缔约};

\node[box] (monitor) at (4.3,0.8) {监督};
\node[box] (enforce) at (4.3,-0.8) {执行};

\draw[flow] (search.east) -- (exchange.west);
\draw[flow] (verify.east) -- (exchange.west);
\draw[flow] (contract.east) -- (exchange.west);

\draw[flow] (exchange.east) -- (monitor.west);
\draw[flow] (exchange.east) -- (enforce.west);

\node[align=center] at (-4.3,2.45) {\small 签约之前及过程中};
\node[align=center] at (4.3,1.65) {\small 融资之后};

\end{tikzpicture}

\caption{金融关系中的交易成本}
\label{fig:ch3-transaction-costs}
\end{figure}

图~\ref{fig:ch3-transaction-costs} 是示意图，并不表示严格的时间顺序。签约后仍可能继续核实信息，监督可能产生新信息，而对执行的预期即使在违约发生前，也可能影响行为。

\subsection{信息不对称与有成本的信息获取}

当金融交易的一方掌握另一方不具备的相关信息时，就存在信息不对称。其中一些信息可能能够获得，但获取和核实需要成本。

这种不对称可能涉及特征。借款人对项目质量的了解可能多于贷款人，管理者对预期现金流的了解可能多于外部投资者。

它也可能涉及行动。贷款人可能无法完全观察借款人如何管理获融资项目；资本提供后，股东也可能观察不到重要的管理决策。

信息不对称不同于不确定性。即使所有人都能获得相同信息，未来结果仍可能不确定。区别在于信息在各方之间的分布，而非不确定性本身是否存在。

额外获取信息可以缩小部分信息差距，但需要消耗资源。令 \(V(q)\) 表示精度为 \(q\) 的信号所提供信息的经济价值，\(C(q)\) 表示生产该信号的成本。
假设决策者选择信息精度 \(q\)，以最大化净信息价值 \(V(q)-C(q)\)。在可微的内点最优解处，提高精度的边际价值等于边际成本：

\begin{equation}
V'(q)=C'(q).
\label{eq:ch3-information-acquisition}
\end{equation}

完全信息通常并非目标。关键是在给定成本下，获取多少信息值得。银行、投资者和保险公司因此会有选择地获取信息，在信息的预期价值与获取、处理所需资源之间权衡。

技术进步能够降低信息收集、处理和评估的成本，从而改变信息获取的经济性。数字记录、卫星数据、电子发票、机器学习和自动化文档分析，可以使过去获取成本高昂的信息更易获得，并更容易纳入金融决策。

然而，信息成本降低并不必然提高信息质量或决策价值。数据可能有偏、过时、被策略性操纵，或与关心的经济变量关联很弱。高效处理无法弥补信号经济相关性不足的问题。

\section{逆向选择：签约前的隐藏信息}

当私人掌握的特征影响参与或契约选择时，就会产生逆向选择（adverse selection）。不同类型对给定条款作出不同反应，而不知情一方在签约前无法完全观察每位参与者的类型。

\subsection{隐藏信息下的混同}

假设申请者中存在 $H$ 和 $L$ 两类借款人，$H$ 型项目质量高于 $L$ 型。贷款人知道申请者中 $H$ 型占比为 $\pi$，却无法在签约前观察任何单个申请人的类型。如果没有可以区分两类人的信息，就只能对申请者进行混同评估（pooling）。

若贷款人对两类人的估值分别为 $V_H$ 和 $V_L$，且 $V_H>V_L$，则类型未知的申请者的混同期望价值为

\begin{equation}
\bar V=\pi V_H+(1-\pi)V_L.
\label{eq:ch3-pooling-value}
\end{equation}

混同条件下，融资条款无法针对借款人的私人类型作出调整。如果参与决策取决于给定条款，高质量借款人可能更容易退出，使借款人池的平均质量下降。最初的混同估值便可能不再反映留下的借款人构成。这正是\emph{柠檬问题}（lemons problem）的核心机制。

这一机制出现在多种金融市场。如果贷款人无法区分借款人质量，低风险借款人可能面临部分反映高风险借款人风险的利率，从而减少外部融资。在保险市场，混同保费可能抑制低风险客户参与，却有利于高风险客户参与。在证券市场，发行人质量信息不足，可能导致高质量企业被低估。

\subsection{知情方的信号传递}

知情方可以通过可观察信号传递私人信息。只有当信号的成本或支付后果在不同类型之间存在系统性差异时，信号才有信息含量。

假设企业选择信号水平 $s$。类型 $j$ 的信号成本为 $c_j(s)$，其中 $j\in\{H,L\}$。单交叉条件可写为

\begin{equation}
c_H'(s)<c_L'(s),
\label{eq:ch3-single-crossing}
\end{equation}
即 $H$ 型传递信号的边际成本低于 $L$ 型。这一差异可以支持分离结果：高质量类型选择一种信号，而低质量类型认为模仿它并不值得。

金融信号可以包括保留更高股权比例、提供额外披露、接受更紧密地与绩效挂钩的报酬，或投入资产和个人资本。信号的信息含量，取决于不同类型在成本或承诺方面的差异。

能够以微不足道成本模仿的陈述，几乎不传递信息。如果低质量类型模仿的吸引力更小，付出成本的承诺就可能更可信。

信号传递（signaling）能够降低信息扭曲，但本身也有成本。用于区分类型的资源可能产生信息价值，却不会直接增加项目产出。

\subsection{不知情方的筛选}

筛选（screening）是指不知情方设计契约条款，诱使掌握私人信息的一方通过选择揭示信息。当不同类型偏好不同契约选项、其选择能够揭示私人信息时，这一机制就能发挥作用。

假设贷款人提供两种契约：契约 $A$ 利率较低，但抵押品要求较高；契约 $B$ 抵押品要求较低，但利率较高。不同类型借款人可能以不同方式评价这一权衡。因此，契约选择可以揭示贷款人无法直接观察的类型信息。

分离型契约菜单要求，每种类型都更偏好为其设计的契约。令 $C_H$ 和 $C_L$ 分别表示针对 $H$ 型和 $L$ 型的契约。若 $U_j(C_k)$ 表示类型 $j$ 选择契约 $C_k$ 的效用，其中 $j,k\in\{H,L\}$，则激励相容要求

\begin{equation}
\begin{aligned}
U_H(C_H) &\geq U_H(C_L),\\
U_L(C_L) &\geq U_L(C_H).
\end{aligned}
\label{eq:ch3-screening-ic}
\end{equation}

这些条件使自我选择与类型一致：$H$ 型至少弱偏好 $C_H$，$L$ 型至少弱偏好 $C_L$。当不同类型作出不同选择时，贷款人就能够从所选契约推断借款人类型。

基于契约的筛选在不同金融市场中形式各异。贷款人可以改变抵押品要求或贷款期限，保险公司则可以改变免赔额或保险限额。

如果不同类型偏好相似条款，其选择揭示的信息就较少。此时，可观察数据可以补充契约筛选，算法信用评分能够降低信息处理成本。但预测关系可能不稳定，也可能受到策略性操纵。

\subsection{市场后果：错误定价、信贷配给与市场瓦解}

在混同条件下，如果融资条款没有充分反映借款人质量差异，就可能出现错误定价。与完全信息相比，高质量借款人的条件可能更差，而低质量借款人的条件可能更好。

如果这些条款影响参与，一些高质量借款人可能退出。留下的借款人平均质量下降，进一步强化最初的定价问题。严重时，互利交易可能消失，市场部分或完全瓦解。

信贷配给（credit rationing）是贷款市场中尤其重要的后果。在标准竞争市场中，超额需求通常推动价格上涨。但在借款人质量隐藏的信贷市场，提高利率还可能改变愿意按该利率借款的借款人构成。

将贷款本金标准化为 1，假设非负项目支付可用于偿还贷款人，且不存在破产成本。令 \(\mathbb{E}_i\) 表示对愿意接受利率 \(i\) 的借款人求期望。银行每单位贷款的预期偿还额为

\begin{equation}
\rho(i)=\mathbb{E}_i[\min\{1+i,X\}]
\label{eq:ch3-bank-return}
\end{equation}
其中，$X$ 为每单位借款对应的项目支付。当 $X\geq 1+i$ 时，银行收到 $1+i$，否则收到 $X$。贷款利率变化同时影响契约偿还额和借款人类型的分布。


预期偿还额 $\rho(i)$ 未必随贷款利率单调增加。图~\ref{fig:ch3-credit-rationing} 展示了一种可能的非单调关系，其具体形式取决于借款人构成和项目支付如何随贷款利率变化。

\begin{figure}[H]
\centering

\begin{tikzpicture}[x=1cm,y=1cm]

% Axes
\draw[->, line width=0.8pt]
    (0,0) -- (7.2,0)
    node[right] {贷款利率 $i$};

\draw[->, line width=0.8pt]
    (0,0) -- (0,5.2)
    node[above] {预期偿还额 $\rho(i)$};

% Expected-return curve
\draw[line width=1.1pt, smooth]
    plot coordinates {
        (0.6,0.9)
        (1.4,1.8)
        (2.3,3.0)
        (3.2,4.0)
        (4.0,4.4)
        (4.8,4.2)
        (5.6,3.6)
        (6.5,2.6)
    };

% Maximum point
\fill (4.0,4.4) circle (2pt);

% Guide line
\draw[dashed, gray]
    (4.0,4.4) -- (4.0,0);

% Labels
\node[below] at (4.0,0) {$i^*$};
\node[above right] at (4.0,4.4)
    {\small 最大预期偿还额};

\end{tikzpicture}

\caption{贷款利率与预期偿还额之间的非单调关系示意}

\label{fig:ch3-credit-rationing}
\end{figure}


在图~\ref{fig:ch3-credit-rationing} 的示例中，令 $i^*$ 表示使预期偿还额最大化的贷款利率。若该利率下的贷款需求超过银行愿意提供的数量，银行可能选择配给信贷，而不再提高利率，因为进一步加息会降低预期偿还额。


\section{道德风险：签约后的隐藏行动}

与逆向选择不同，道德风险涉及签约后的隐藏行动。当这些行动影响另一方的支付，却无法被完全观察或写入可执行契约时，就可能出现激励问题。

\subsection{隐藏行动与代理问题}

隐藏行动是代理人在签约后选择的、影响委托人支付，却无法被完全观察或纳入契约的行动。当代理人的激励与委托人目标不完全一致时，这种信息限制就会产生代理问题（principal--agent problem）。

考虑提供资源的委托人和努力程度影响项目结果的代理人。项目成功时产生支付 $R$，失败时为零。代理人选择努力程度 $e$，它影响成功概率 $p(e)$，并产生私人成本 $c(e)$。

为简化分析，假设双方均为风险中性，因而货币支付线性进入其目标函数。若代理人在成功后获得 $w_S$，失败后获得 $w_F$，则预期效用为

\begin{equation}
U(e)=p(e)w_S+[1-p(e)]w_F-c(e).
\label{eq:ch3-agent-utility}
\end{equation}

委托人重视努力对项目回报的贡献，而代理人考虑契约报酬与努力的私人成本。如果努力不可观察，代理人只能获得其创造收益的一部分，便可能选择低于项目总剩余最大化水平的努力。

\subsection{努力、偷懒与激励问题}

有效率的努力水平 $e^*$ 使预期总剩余最大化：


\begin{equation}
e^*\in\arg\max_e\left\{p(e)R-c(e)\right\}.
\label{eq:ch3-efficient-effort}
\end{equation}

若 $e^*$ 是内点最优解，则满足

\begin{equation}
p'(e^*)R=c'(e^*).
\label{eq:ch3-efficient-effort-foc}
\end{equation}

代理人选择努力，以最大化契约提供的效用。将 $U(e)$ 对努力求导，内点最优解 $\hat e$ 满足

\begin{equation}
p'(\hat e)(w_S-w_F)=c'(\hat e).
\label{eq:ch3-agent-effort-foc}
\end{equation}

比较两个条件可识别激励问题的来源。项目从努力获得的边际收益为 $p'(e)R$，代理人的私人边际收益为 $p'(e)(w_S-w_F)$。若 $w_S-w_F<R$，代理人只能获得更高努力所创造增量价值的一部分。在关于 $p(e)$ 和 $c(e)$ 的标准假设下，私人选择的努力可能低于有效率水平。在此情形下，这一不足表现为偷懒（shirking）。

\begin{example}{高努力与低努力}{high-low-effort}

假设代理人在高努力 $e_H$ 与低努力 $e_L$ 之间选择。高努力提高成功概率，但私人成本也更高：

\begin{equation*}
p(e_H)>p(e_L),
\qquad
c(e_H)>c(e_L).
\end{equation*}

若代理人成功后获得 $w_S$、失败后获得 $w_F$，则在以下条件下更偏好高努力：

\begin{equation*}
p(e_H)w_S+[1-p(e_H)]w_F-c(e_H)
\geq
p(e_L)w_S+[1-p(e_L)]w_F-c(e_L).
\end{equation*}

整理得

\begin{equation*}
[p(e_H)-p(e_L)](w_S-w_F)
\geq
c(e_H)-c(e_L).
\label{eq:ch3-effort-spread}
\end{equation*}

左边是成功概率提高带来的额外预期报酬，右边是高努力的额外私人成本。成功与失败之间的报酬差越大，高努力相对于低努力就越有吸引力。

\end{example}

放松风险中性假设会引入另一项考虑。如果代理人厌恶风险，提高报酬对结果的敏感性，既强化激励，也使其在结果部分取决于不可控因素时面临更多收入风险。因此，契约设计可能需要在激励与风险分担之间权衡。

\subsection{风险转移与资产替代}

签约后的激励问题不仅来自努力选择，也来自项目选择。债务发行后，即使高风险项目降低企业总价值，股东仍可能更偏好它。这种激励称为风险转移（risk shifting）。当更安全的项目或资产被风险更高者替代时，通常称为资产替代（asset substitution）。

有限责任下，股东在满足承诺的债务偿还后获得剩余支付，其支付下限为零。如果企业支付为 $X$，承诺债务偿还额为 $D$，股东支付为

\begin{equation}
V_E(X)=\max\{X-D,0\}.
\label{eq:ch3-equity-payoff}
\end{equation}

这种不对称支付，可能使股东认为风险更高的项目更有吸引力，即使它降低企业总价值。

\begin{example}{债务发行后的风险转移}{risk-shifting-debt-issuance}

假设企业债务面值为 \(D=80\)。债务发行后，股东可在两个项目中选择。安全项目确定支付 110，股东获得 \(110-80=30\)。



风险项目以 \(0.5\) 的概率支付 150，以 \(0.5\) 的概率支付 50，其预期总支付为


\begin{equation*}
0.5(150)+0.5(50)=100,
\end{equation*}

低于安全项目的支付 110。在有限责任下，股东在高支付状态获得 70，低支付状态获得零。因此，其预期支付为



\begin{equation*}
0.5(70)+0.5(0)=35.
\label{eq:ch3-risk-shifting-example}
\end{equation*}

因此，尽管风险项目使企业总价值从 110 降至 100，股东仍偏好风险项目。债权人的预期偿还额则由安全项目下的 80 降为风险项目下的



\begin{equation*}
0.5(80)+0.5(50)=65
\end{equation*}

。股东价值增加的同时，债权人遭受更大损失，企业总价值下降。


\medskip

\begin{center}
\small
\textbf{两种项目选择下的支付}

\medskip

\begin{tabularx}{\linewidth}{
  @{}
  >{\raggedright\arraybackslash}X
  >{\centering\arraybackslash}p{2.8cm}
  >{\centering\arraybackslash}p{3.0cm}
  @{}
}
\toprule
 & \textbf{安全项目} & \textbf{风险项目} \\
\midrule
预期总支付
    & 110 & 100 \\
债权人预期偿还额
    & 80 & 65 \\
股东预期支付
    & 30 & 35 \\
企业价值排序 & 较高 & 较低 \\
股东支付排序 & 较低 & 较高 \\
\bottomrule
\end{tabularx}
\end{center}

\end{example}


融资后风险转移的可能性，为限制投资政策或经营风险变化的契约条款提供了依据。债权人关心的不仅是放贷时的条件，还有提供融资之后的决策。

\subsection{私人收益与代理成本}

当决策者获得的收益不按比例与外部投资者分享时，也可能出现签约后的激励冲突。管理者可能通过在职消费、建立企业帝国，或实施扩大自身控制却不增加企业价值的项目，获取私人收益。控股股东也可能以牺牲中小投资者利益为代价获益。

令决策 $a$ 产生企业价值 $V(a)$，并给决策者带来私人收益 $B(a)$。若决策者持有企业比例为 $\alpha$ 的股权，其私人支付为

\begin{equation}
V_{\text{私人}}(a)=\alpha V(a)+B(a).
\label{eq:ch3-private-benefits}
\end{equation}

当 $\alpha<1$ 时，决策者只承担企业价值下降的一部分，却保留全部私人收益。因此，使私人支付最大化的行动，可能不同于使企业总价值最大化的行动。

因此，私人收益可能使决策者私人支付最大化的行动，与企业价值最大化的行动之间产生偏离。这种偏离形成代理成本（agency costs），既包括用于限制激励冲突的资源，也包括冲突未能完全消除时留下的剩余损失。

\section{金融契约：应对信息与激励问题}

上述信息与激励问题，赋予契约设计以作用。金融契约不只是规定支付，还可以影响参与、揭示私人信息、协调激励，并在情形变化时分配权利。

披露、治理、审计、声誉和法律执行，往往支持这些功能。因此，金融契约是在更广泛的制度框架中运作的。

\subsection{抵押品与承诺}

契约设计能够将偿还义务与借款人未履约时仍可用的资产相联系。抵押品通过赋予债权人对指定资产的请求权，建立这种联系。借款人违约时，质押或抵押资产成为债权人的回收来源；违约之前，可能失去这些资产的后果也会影响借款人激励。

若欠款额为 $D$，抵押品可回收价值为 $C$，损失的简化度量为

\begin{equation}
L=\max\{D-C,0\}.
\label{eq:ch3-collateral-loss}
\end{equation}

抵押品还影响违约前的借款人激励。借款人面临失去抵押资产的可能，增加了提高违约风险行为的私人成本。因此，抵押品既是回收机制，也是承诺工具。

如果提供抵押资产的成本在借款人类型之间不同，抵押品还可能传递信息。自愿提供抵押可作为信号，抵押要求则可诱使借款人在契约之间自我选择。因此，其信息含量取决于不同类型提供抵押品的成本是否存在实质差异。

抵押品的有效性取决于抵押资产的价值和权利的可执行性。抵押品价值可能下降，扣押和变现可能代价高昂。抵押要求也可能限制有价值项目但可抵押资产不足的借款人融资。抵押品价值普遍下跌，会收紧借款约束并放大金融压力。

\subsection{限制性条款与控制权}

抵押品将偿还与特定资产相联系，限制性条款（covenants）则约束融资后的借款人行动。此类条款限制或要求特定行动，也可以规定影响债权人权利的条件。

消极性条款可能限制新增借款、资产出售或向股东分配；积极性条款可能要求披露财务信息、保全抵押品，或遵守特定财务比率要求。

限制性条款聚焦于有限的一组可观察或可核实的行动和财务条件。违反条款时，契约可能触发重新谈判，或向债权人转移更多控制权，例如要求补充抵押品、限制分配或加速偿还。

控制权往往取决于状态。满足契约条件时，管理者保留更多自主权；违反条件时，债权人获得更大影响力。

\begin{figure}[H]
\centering
\begin{tikzpicture}[
  >=Latex,
  box/.style={
    draw,
    rounded corners,
    minimum width=3.5cm,
    minimum height=0.85cm,
    align=center
  },
  flow/.style={
    -Latex,
    line width=0.85pt,
    FinanceBlue
  }
]

\node[box] (satisfied) at (-1,2.2)
  {满足契约条件};

\node[box] (manager) at (5.2,2.2)
  {管理者自主权较大};

\node[box] (breach) at (-1,0)
  {违反契约条件};

\node[box] (creditor) at (5.2,0)
  {债权人权利扩大};

\draw[flow] (satisfied) -- (manager);
\draw[flow] (breach) -- (creditor);

\draw[flow] (satisfied) --
  node[left, align=center]{\small 违反\\\small 限制性条款}
  (breach);

\end{tikzpicture}

\caption{通过限制性条款实现状态依存的控制权配置}
\label{fig:ch3-covenant-control}
\end{figure}

财务限制性条款需要在债权人及时介入与不必要的重新谈判之间权衡。门槛过严，可能因未反映财务状况实质恶化的暂时波动而触发违约；门槛过松，则可能在大量价值损失后才允许介入。

\subsection{支付设计与激励相容}

契约设计可以将支付与可观察结果联系起来，或提供鼓励自我选择的条款，以影响私人选择。激励相容（incentive compatibility）要求：筛选中每种类型都偏好针对自己的契约；行动隐藏时，所期望的行动对个人而言也是最优选择。

股权保留、绩效报酬和保险免赔额，能够通过提高个人对自身决策后果的承担来强化激励。然而，如果结果也取决于代理人无法控制的因素，更高的支付敏感性可能削弱风险分担。因此，契约设计需要在激励供给与风险分担之间权衡。

\subsection{监督与信息生产}

监督生产并核实融资后关于借款人绩效、合规和风险的信息。它可以揭示潜在违约，并支持契约权利的行使。行为可能被发现并受到执行的预期，也能抑制违反契约义务的行动。

当金融请求权高度分散时，监督可能不足。单个请求权人承担监督成本，却与其他请求权人共享收益，形成搭便车问题。金融中介可以代表多个资金提供者开展监督，从而缓解这一问题。

技术可以通过自动分析财务报告和交易数据来降低监督成本。AI 辅助系统能够识别潜在异常，但其结果仍需要核实和经济解释。

\subsection{治理与执行}

金融契约在更广泛的治理和执行安排中运行。治理关注管理者、董事会、股东、债权人和其他各方之间决策权与控制权的分配和行使。这些权利可以包括批准重大交易、更换管理者，或在特定条件满足时介入。

仅有权利分配并不保证其有效。执行决定契约和法律权利在实践中能否行使。法律规则和制度影响抵押品回收、加速偿还、争议解决，以及破产时债权优先顺序的落实。管理者或控股股东违反法定义务时，它们也提供救济。

执行不限于诉讼程序。声誉能够影响未来交易机会；承诺遭到违反时，贷款人和投资者可能撤回未来融资。交易所和监管机构也可以通过持续性要求和制裁来强化合规。

预期的执行有效性也影响契约设计。如果回收缓慢或不确定，贷款人可能要求更多抵押品、更短期限、更严条款或更高信用利差。更强的治理和更可靠的执行，可以在投资者较少直接介入的情况下，使外部融资成为可能。因此，同一份书面契约在不同制度环境中可能具有不同经济价值。

\begin{tcolorbox}[
  enhanced,
  breakable,
  colback=white,
  colframe=FinanceBlue,
  boxrule=0.7pt,
  sharp corners,
  borderline west={2.2pt}{0pt}{FinanceBlue},
  title=\textbf{金融视角：更好的信息并不能解决全部问题},
  coltitle=white,
  colbacktitle=FinanceBlue,
  fonttitle=\bfseries,
  before skip=1.2em,
  after skip=1.2em
]

更好的信息可以改善金融决策，却不能消除激励或执行问题。贷款人可能准确评估违约风险，却仍无法控制融资后的借款人行动。详细披露可以改善股东信息，却未必赋予外部投资者有效控制管理决策的权力。AI 模型可以识别信用质量恶化，但该信息的价值仍取决于可据此行使的权利和救济。因此，信息、激励和可执行权利是互补品，而非替代品。

\end{tcolorbox}

%========================================================
% Summary
%========================================================

\TailModule{summary}{本章小结}

金融交易需要有关交易对手、项目和未来行为的信息。生产和核实信息具有成本，而且相关信息往往在各方之间不对称分布。

逆向选择来自签约前的隐藏信息。质量无法直接观察时，混同定价可能抑制高质量参与者进入或留在市场，降低剩余参与者的平均质量，严重时促成市场瓦解。当不同类型模仿信号的成本存在差异时，信号传递使知情方能够揭示类型；筛选则利用契约选择诱导自我选择。在信贷市场，提高利率可能恶化借款人构成，因此银行预期偿还额未必随利率单调增加。即使有借款人愿意接受更高利率，也可能出现信贷配给。

道德风险来自签约后的隐藏行动。偷懒、风险转移和追求私人收益，会使私人最优行动偏离贷款人、投资者或保险公司所偏好的行动。金融契约通过改变支付、承诺和控制权予以应对。抵押品、限制性条款、股权保留、免赔额、监督及状态依存控制可以改善激励，而治理和执行决定权利在实践中如何行使。

这些机制都涉及权衡。强化激励可能削弱风险分担，更严的契约限制可能压缩有价值的灵活性，抵押要求可能限制可抵押资产较少的借款人，而监督本身也有成本。因此，契约设计旨在以合理经济成本减少信息和激励问题，而非彻底消除它们。

%========================================================
% Mind Map
%========================================================

\TailModule{mindmap}{思维导图}
\begin{figure}[H]
\centering
\begin{tikzpicture}[
  >=Latex,
  main/.style={
    draw,
    rounded corners,
    minimum width=3.4cm,
    minimum height=0.9cm,
    align=center,
    fill=gray!10
  },
  sub/.style={
    draw,
    rounded corners,
    minimum width=3.2cm,
    minimum height=0.8cm,
    align=center
  },
  flow/.style={-Latex, line width=0.75pt}
]

\node[main] (friction) at (0,4.2)
  {信息摩擦};

\node[sub] (hiddeninfo) at (-4.0,2.3)
  {签约前的\\隐藏信息};

\node[sub] (hiddenaction) at (4.0,2.3)
  {签约后的\\隐藏行动};

\node[sub] (adverse) at (-4.0,0.5)
  {逆向选择};

\node[sub] (moral) at (4.0,0.5)
  {道德风险};

\node[sub] (responses1) at (-4.0,-1.5)
  {信号传递与筛选};

\node[sub] (responses2) at (4.0,-1.5)
  {支付设计与监督};

\node[main] (responses) at (0,-3.4)
  {契约与\\制度应对};

\draw[flow] (friction) -- (hiddeninfo);
\draw[flow] (friction) -- (hiddenaction);
\draw[flow] (hiddeninfo) -- (adverse);
\draw[flow] (hiddenaction) -- (moral);
\draw[flow] (adverse) -- (responses1);
\draw[flow] (moral) -- (responses2);
\draw[flow] (responses1) -- (responses);
\draw[flow] (responses2) -- (responses);

\end{tikzpicture}

\caption*{签约前后的信息问题}
\end{figure}

%========================================================
% Key Terms
%========================================================

\TailModule{keyterms}{关键术语}

交易成本（Transaction costs）；信息不对称（information asymmetry）；有成本的信息获取（costly information acquisition）；隐藏信息（hidden information）；逆向选择（adverse selection）；混同（pooling）；信号传递（signaling）；筛选（screening）；激励相容（incentive compatibility）；信贷配给（credit rationing）；隐藏行动（hidden action）；道德风险（moral hazard）；委托代理问题（principal--agent problem）；风险转移（risk shifting）；资产替代（asset substitution）；私人收益（private benefits）；代理成本（agency costs）；抵押品（collateral）；限制性条款（covenant）；控制权（control rights）；监督（monitoring）；治理（governance）；执行。


%========================================================
% Financial AI Lab
%========================================================（enforcement


%========================================================
% Financial AI Lab
%========================================================）。

\TailModule{aiprompt}{金融 AI 实验室}

\textbf{实验：设计并审查 AI 信贷筛选框架}

请 AI 系统为一家小企业贷款机构提出信贷筛选框架，要求识别可观察变量，并解释每个变量可能如何关联偿还风险。区分放贷前可获得的信息，与融资后才能获得的信息，并解释后者为什么属于监督而非初始筛选。

随后审查框架：哪些变量包含借款人内在质量的信息，哪些只是与之相关，哪些会在申请人理解筛选规则后被操纵？指出一种潜在的选择偏差来源、一项公平性问题，以及一个使用历史借款人训练的模型可能在信贷条件变化后表现不佳的原因。

最后，假设模型能够高度准确地预测违约风险。解释哪些问题仍需通过契约条款、抵押品、监督、执行或人工审查来解决。为什么准确预测本身并不能消除缔约问题？

%========================================================
% Discussion Questions
%========================================================

\TailModule{questions}{讨论题}

\begin{enumerate}

\item 为什么不确定性不等同于信息不对称？分别举出一个不确定性很高但信息对称的金融例子，以及一个信息不对称的例子。

\item 为什么提高贷款利率，可能在承诺偿还额上升的同时，降低银行每单位贷款的预期偿还额？


\item 比较信号传递与筛选。两种情况下分别由谁主动采取措施？什么机制使不同类型能够被区分？

\item 为什么融资后对借款人行为的信息改善，未必能消除道德风险？解释监督、契约激励、控制权和执行的作用。



\end{enumerate}

%========================================================
% Applied Problems
%========================================================

\TailModule{appprobs}{应用题}

\begin{enumerate}

\item 一位投资者考虑购买对企业未来项目现金流的请求权。他认为其中 70\% 为高质量，每项价值 120；30\% 为低质量，每项价值 60。他无法观察单项请求权的质量。
    \begin{enumerate}

    \item 计算混同期望价值。
    
    \item 假设高质量企业只愿以至少 110 的价格出售，低质量企业只要求至少 55。如果投资者最初按混同期望价值报价，哪些类型会接受？出售请求权的构成如何变化？

    \item 从经济角度解释逆向选择机制。

    \end{enumerate}

\item 一名代理人在高、低努力之间选择。高努力使成功概率为 $0.8$，私人成本为 12；低努力使成功概率为 $0.5$，私人成本为 3。代理人成功后获得 $w_S$，失败后获得零。

  \begin{enumerate}

  \item 推导选择高努力的激励相容条件，并确定 $w_S$ 的门槛值。

  \item 解释为什么提高成功时的支付会影响代理人的努力选择。

  \item 如果代理人厌恶风险，会出现什么额外权衡？

  \end{enumerate}

\item 某企业债务面值为 80。股东可以选择项目 A，确定支付 112；或项目 B，以 $0.5$ 的概率支付 160，以 $0.5$ 的概率支付 48。

  \begin{enumerate}

  \item 计算每个项目的预期总支付。

  \item 计算每个项目下股东的预期支付。

  \item 股东偏好的项目是否使企业总价值最大化？解释冲突的来源。

  \end{enumerate}

\end{enumerate}

%========================================================
% Case Analysis
%========================================================

\TailModule{caseanalysis}{案例分析}

\textbf{瑞幸咖啡：当披露信息不能照单全收}

2020 年 12 月，美国证券交易委员会（SEC）指控瑞幸咖啡公司通过严重错报收入、费用和净营业亏损欺诈投资者。根据 SEC 诉状，至少从 2019 年 4 月至 2020 年 1 月，瑞幸故意虚构了超过 3 亿美元的零售销售额。公司在不承认也不否认指控的情况下，同意一项有待法院批准的和解，其中包括支付 1.8 亿美元民事罚款。

以 SEC 于 2020 年 12 月 16 日发布的执法公告及相关公司披露为一手来源，运用本章框架分析这一事件。外部投资者依赖哪些信息？哪些信息生产和核实机制原本用于减少内部人与外部投资者的信息不对称？此事件反映了外部审计作为核实机制的哪些价值和局限？披露信息受到质疑后，哪些监督、治理、内部控制和执行机制变得重要？最后，区分案例中的\emph{信息失灵}与\emph{激励或治理失灵}。

\emph{一手来源：美国证券交易委员会，“瑞幸咖啡同意支付 1.8 亿美元罚款以和解会计欺诈指控”（``Luckin Coffee Agrees to Pay \$180 Million Penalty to Settle Accounting Fraud Charges''），新闻稿 2020-319，2020 年 12 月 16 日。}
%========================================================
% Further Reading
%========================================================

\TailModule{furtherread}{延伸阅读}

\begin{itemize}

\item Akerlof, George A.（1970），“柠檬市场：质量不确定性与市场机制”（``The Market for `Lemons': Quality Uncertainty and the Market Mechanism''），\emph{The Quarterly Journal of Economics}，84(3)，488--500。

\item Spence, Michael（1973），“劳动力市场信号传递”（``Job Market Signaling''），\emph{The Quarterly Journal of Economics}，87(3)，355--374。

\item Jensen, Michael C. 与 William H. Meckling（1976），“企业理论：管理者行为、代理成本与所有权结构”（``Theory of the Firm: Managerial Behavior, Agency Costs and Ownership Structure''），\emph{Journal of Financial Economics}，3(4)，305--360。

\item Stiglitz, Joseph E. 与 Andrew Weiss（1981），“信息不完全市场中的信贷配给”（``Credit Rationing in Markets with Imperfect Information''），\emph{American Economic Review}，71(3)，393--410。

\item Diamond, Douglas W.（1984），“金融中介与委托监督”（``Financial Intermediation and Delegated Monitoring''），\emph{The Review of Economic Studies}，51(3)，393--414。

\end{itemize}
\ChapterTail{summary,mindmap,keyterms,aiprompt,questions,appprobs,caseanalysis,furtherread}

\chapter{金融中介与金融结构}
\label{chap:Financial_Intermediation_and_Financial_Structure}

\begin{openingcase}{为什么大企业发债，小企业却向银行借款？}

两家制造企业都需要外部融资。一家大型上市企业计划发行债券融资人民币 20 亿元，一家需要 500 万元的小型民营企业则转向银行。这种差异似乎令人困惑：发债可以直接向投资者融资，而无须通过银行借款。为什么小企业要依赖中介？

上市企业拥有经审计的财务报表，融资需求规模也足以分摊发行成本。对小企业而言，外部投资者很难核实有关客户、设备和管理层的信息。与其已有业务关系的银行，可能更适合评估和监督借款人。

金融中介在融资过程中增加了一家机构，却可能降低获得和维持融资的总体成本。

\end{openingcase}


\begin{tcolorbox}[
  enhanced,
  colback=FinanceBlue!6,
  colframe=FinanceBlue,
  boxrule=0.6pt,
  sharp corners,
  left=8pt,
  right=8pt,
  top=7pt,
  bottom=7pt,
  before skip=1em,
  after skip=1.6em,
  title=\textbf{核心问题},
  coltitle=white,
  colbacktitle=FinanceBlue,
  fonttitle=\bfseries,
  attach boxed title to top left={xshift=6pt,yshift=-2mm},
  boxed title style={
    boxrule=0pt,
    sharp corners,
    left=5pt,
    right=5pt,
    top=2pt,
    bottom=2pt
  }
]
\begin{enumerate}[
  leftmargin=*,
  label=\textcolor{FinanceBlue}{\textbf{\arabic*.}},
  itemsep=0.55em,
  topsep=0.2em
]
\item 什么因素决定融资通过证券市场还是金融中介组织？

\item 为什么委托信息生产、监督和资产负债表转换能够创造价值？

\item 为什么市场和中介会共存？什么决定它们在金融体系中的相对重要性？

\end{enumerate}
\end{tcolorbox}

\section{融资渠道：直接融资与间接融资}

融资渠道由连接最终资金提供者与最终使用者的请求权结构来界定。在直接融资（direct finance）中，资金提供者持有最终借款人发行的请求权；在间接融资（indirect finance）中，中介取得借款人的请求权，再通过发行自身请求权为其融资。其他机构的参与并不改变这一区分；关键是中介是否将自己的资产负债表置于双方之间。

\begin{figure}[htbp]
\centering

\begin{tikzpicture}[
    box/.style={
        draw,
        rounded corners,
        align=center,
        minimum width=2.7cm,
        minimum height=0.85cm,
        inner sep=3pt
    },
    fund/.style={
        -{Latex[length=2.2mm]},
        thick
    },
    claim/.style={
        -{Latex[length=2.2mm]},
        dashed,
        thick
    },
    every node/.style={font=\footnotesize}
]

% ------------------------------------------------
% (a) Direct finance
% ------------------------------------------------

\node[font=\small\bfseries] at (0,1.05)
    {（a）直接融资};

\node[box] (providersA) at (-2.7,0)
    {最终资金\\提供者};

\node[box] (borrowerA) at (2.7,0)
    {最终资金\\使用者};

\draw[fund]
    ([yshift=2.2mm]providersA.east)
    --
    node[above=1pt]{资金}
    ([yshift=2.2mm]borrowerA.west);

\draw[claim]
    ([yshift=-2.2mm]borrowerA.west)
    --
    node[below=1pt]{证券}
    ([yshift=-2.2mm]providersA.east);


% ------------------------------------------------
% (b) Indirect finance
% ------------------------------------------------

\node[font=\small\bfseries] at (0,-1.55)
    {（b）间接融资};

\node[box] (providersB) at (-4.0,-2.65)
    {最终资金\\提供者};

\node[box] (intermediary) at (0,-2.65)
    {金融\\中介};

\node[box] (borrowerB) at (4.0,-2.65)
    {最终资金\\使用者};

\draw[fund]
    ([yshift=2.2mm]providersB.east)
    --
    node[above=1pt]{资金}
    ([yshift=2.2mm]intermediary.west);

\draw[claim]
    ([yshift=-2.2mm]intermediary.west)
    --
    node[below=2pt, align=center, font=\scriptsize]
    {中介发行的\\[-1pt]请求权}
    ([yshift=-2.2mm]providersB.east);

\draw[fund]
    ([yshift=2.2mm]intermediary.east)
    --
    node[above=1pt]{资金}
    ([yshift=2.2mm]borrowerB.west);

\draw[claim]
    ([yshift=-2.2mm]borrowerB.west)
    --
    node[below=2pt, align=center, font=\scriptsize]
    {对资金使用\\[-1pt]者的请求权}
    ([yshift=-2.2mm]intermediary.east);

\end{tikzpicture}

\caption{直接与间接融资渠道}
\label{fig:ch4-direct-indirect-finance}
\end{figure}

\subsection{直接融资与证券市场}

直接融资和间接融资的区别，在于连接最终资金提供者与资金使用者的金融请求权结构。直接融资中，资金提供者持有资金使用者发行的请求权；间接融资中，中介持有对使用者的请求权，并发行自己的金融请求权为其融资。关键在于中介是否将自身资产负债表置于双方之间。

直接融资并不意味着没有金融机构。投资银行可以安排或承销发行，评级机构可以评估信用风险，交易所或交易商可以便利交易。这些机构支持融资过程，但并不将自身资产负债表置于发行人与投资者之间。

要向众多投资者分销证券，请求权必须足够标准化，相关信息也必须以广泛投资者群体能够使用的形式传递和核实。支付权利、期限、抵押品、优先顺序和投票权，通过共同契约条款规定；经审计财务报表、披露要求、招股说明书或募集说明书以及评级，则提供广泛可得的信息。这些安排减少了投资者分别调查发行人、协商条款的需要。

证券发行具有固定成本，融资金额越大，可分摊的基数越大。假设规模为 $Q$ 的发行，其总交易成本为

\begin{equation}
C_M(Q)=F_M+c_MQ,
\label{eq:ch4-market-financing-cost}
\end{equation}
其中，$F_M$ 是与披露、法律工作、尽职调查、注册及发行组织有关的固定成本，$c_M$ 是每单位融资的可变交易成本。平均交易成本为

\begin{equation}
AC_M(Q)=\frac{F_M}{Q}+c_M,
\label{eq:ch4-market-average-cost}
\end{equation}
因此
\begin{equation}
\frac{dAC_M(Q)}{dQ}
=-\frac{F_M}{Q^2}
<0.
\label{eq:ch4-market-scale-effect}
\end{equation}

随着发行规模增加，固定成本被分摊到更大融资基数上，平均交易成本下降，市场融资变得更经济。

信息成本也影响融资渠道选择。如果发行人的财务状况和现金流能够可信地披露并核实，分散投资者便能以较低成本评价其证券，使市场融资更有吸引力。

二级市场使投资者可以出售证券，而无须要求发行人偿还或重新谈判融资。因此，可交易性将发行人的融资期限与单个投资者的持有期分离。这一收益依赖市场流动性，而流动性又需要活跃交易、可靠信息和配套市场基础设施，且可能在压力时期急剧恶化。

\subsection{间接融资与金融中介}

间接融资中，金融中介向借款人提供资金，持有由此形成的请求权，并通过向资金提供者发行自身负债，为这些资产融资。

商业银行体现了这一结构。存款人持有对银行的请求权，银行则持有对借款企业和家庭的请求权。简化资产负债表可以清晰展示这种双边结构。
 
\begin{table}[htbp]

\centering

\caption{商业银行的简化资产负债表}

\label{tab:ch4-bank-balance-sheet}

\begin{tabularx}{0.78\textwidth}{@{}X X@{}}

\toprule

\textbf{资产} & \textbf{负债与权益} \\

\midrule

贷款 & 存款 \\
证券 & 批发融资 \\
准备金与流动资产 & 股本 \\

\bottomrule

\end{tabularx}

\end{table}

因此，中介将自身资产负债表置于资金提供者与借款人之间。它持有对借款人的请求权，通常表现为贷款；资金提供者则持有中介发行的请求权。中介评估信用状况、协商贷款条款并监督借款人绩效。

这种安排将信息生产集中于中介内部。通过重复贷款和与借款人直接接触，银行可以积累分散投资者分别生产时成本高昂的信息。其中很多信息可以保持私有，使中介无须完全依赖标准化公开披露来评估借款人。

同样的安排也将大量决策权委托给中介，形成新的代理关系。资金提供者通常无法观察每一项贷款或风险管理决策。因此，中介履行义务的能力，部分取决于贷款决策质量、流动性管理和风险控制。资本、治理、声誉、监督与监管，有助于约束这种受托自主权，并维持对中介负债的信心。

当专业化信息生产与监督的收益超过委托决策的代理成本时，中介可以提高融资效率，尤其是在借款人信息难以标准化或公开核实的情况下。

\subsection{融资链条与混合型中介}

融资链条通过一连串金融请求权，连接最终资金提供者与最终使用者。当证券市场与金融中介在链条不同阶段履行不同功能时，就形成混合型中介。

\begin{figure}[htbp]
\centering

\begin{tikzpicture}[
    box/.style={
        draw,
        rounded corners,
        align=center,
        minimum width=2.45cm,
        minimum height=0.78cm,
        inner sep=3pt,
        font=\footnotesize
    },
    fund/.style={
        -{Latex[length=2.0mm]},
        line width=0.75pt
    },
    claim/.style={
        -{Latex[length=2.0mm]},
        dashed,
        line width=0.75pt
    },
    labeltext/.style={
        font=\scriptsize,
        align=center
    },
    panel/.style={
        font=\small\bfseries
    }
]

% ==================================================
% (a) Mutual fund
% ==================================================

\node[panel] at (0,1.15)
    {（a）共同基金};

\node[box] (hh) at (-4.0,0)
    {家庭};

\node[box] (mf) at (0,0)
    {共同基金};

\node[box] (firm) at (4.0,0)
    {企业};

% funds
\draw[fund]
    ([yshift=2mm]hh.east)
    --
    node[above, labeltext]{资金}
    ([yshift=2mm]mf.west);

\draw[fund]
    ([yshift=2mm]mf.east)
    --
    node[above, labeltext]{资金}
    ([yshift=2mm]firm.west);

% claims
\draw[claim]
    ([yshift=-2mm]mf.west)
    --
    node[below, labeltext]{基金份额}
    ([yshift=-2mm]hh.east);

\draw[claim]
    ([yshift=-2mm]firm.west)
    --
    node[below, labeltext]{公司债券}
    ([yshift=-2mm]mf.east);


% ==================================================
% (b) Securitization
% ==================================================

\node[panel] at (0,-2.0)
    {（b）证券化};

\node[box] (investors) at (-5.7,-3.25)
    {投资者};

\node[box] (spv) at (-1.9,-3.25)
    {特殊目的\\载体};

\node[box] (originator) at (1.9,-3.25)
    {银行或\\财务公司};

\node[box] (borrowers) at (5.7,-3.25)
    {借款人};

% funds
\draw[fund]
    ([yshift=2mm]investors.east)
    --
    node[above, labeltext]{资金}
    ([yshift=2mm]spv.west);

\draw[fund]
    ([yshift=2mm]spv.east)
    --
    node[above, labeltext]{购买\\贷款}
    ([yshift=2mm]originator.west);

\draw[fund]
    ([yshift=2mm]originator.east)
    --
    node[above, labeltext]{贷款\\资金}
    ([yshift=2mm]borrowers.west);

% claims
\draw[claim]
    ([yshift=-2mm]borrowers.west)
    --
    node[below, labeltext]{贷款\\请求权}
    ([yshift=-2mm]originator.east);

\draw[claim]
    ([yshift=-2mm]originator.west)
    --
    node[below, labeltext]{贷款\\资产池}
    ([yshift=-2mm]spv.east);

\draw[claim]
    ([yshift=-2mm]spv.west)
    --
    node[below, labeltext]{资产支持\\证券}
    ([yshift=-2mm]investors.east);


% ==================================================
% (c) Repo funding
% ==================================================
\node[panel] at (0,-5.55)
    {（c）回购融资};
\node[box] (issuer) at (-4.0,-6.75)
    {证券\\发行人};

\node[box] (fi) at (0,-6.75)
    {金融\\中介};

\node[box] (cash) at (4.0,-6.75)
    {资金提供者};

% Intermediary purchases security
\draw[fund]
    ([yshift=2mm]fi.west)
    --
    node[above, labeltext]{资金}
    ([yshift=2mm]issuer.east);

\draw[claim]
    ([yshift=-2mm]issuer.east)
    --
    node[below, labeltext]{可交易\\证券}
    ([yshift=-2mm]fi.west);

% Repo funding
\draw[fund]
    ([yshift=2mm]cash.west)
    --
    node[above, labeltext]{回购融资}
    ([yshift=2mm]fi.east);

\draw[claim]
    ([yshift=-2mm]fi.east)
    --
    node[below, labeltext]{证券作为\\抵押品}
    ([yshift=-2mm]cash.west);

\end{tikzpicture}

\caption{混合融资链条的例子}
\label{fig:ch4-hybrid-financing}
\end{figure}

共同基金持有公司债券就是一个简单例子。公司通过发行可交易证券融资，家庭投资者持有基金份额，而非直接持有公司债券。发行人使用市场融资，家庭储蓄则通过中介进入市场。

证券化通过将贷款发起与筛选，同后续资金和风险分配分开，延伸融资链条。银行或财务公司发起并筛选贷款，将其汇集后转移给特殊目的载体，后者发行以贷款现金流为支持的证券。信息生产仍集中在发起环节附近，而资金供给与风险承担可以通过证券市场分散。

可交易证券还可以支持中介融资。在回购交易中，金融机构以证券获得短期融资，并约定日后回购证券。因此，中介资产负债表上持有的证券，可以成为货币市场融资的抵押品，将证券市场与中介融资联系起来。

因此，市场和中介往往互补。有用的区分不是整条融资链究竟属于市场融资还是中介融资，而是信息在哪里生产、资金如何筹集、流动性在哪里提供，以及最终由谁承担风险。

\section{金融中介为何存在？}

金融中介为融资增加了一个组织层次，但可以通过专业化和规模效应，降低信息生产、缔约、监督及交易成本。这些收益必须与中介成本加以权衡。

\subsection{交易成本、规模与专业化}

重复交易使金融中介能够将信息系统、标准化流程及监督基础设施的成本，分摊到众多借款人上。类似交易的经验，也能够提高其评估和管理风险的能力。

这一优势的一个来源是，有些成本并不随交易量同比例增加。假设某项活动的固定成本为 $F$，每笔交易的可变成本为 $c$，则 $n$ 笔交易的平均成本为

\begin{equation}
AC(n)=\frac{F}{n}+c.
\label{eq:ch4-scale-cost}
\end{equation}

随着 $n$ 增加，每笔交易分摊的固定成本下降。规模也支持专业化：信用分析、缔约、风险管理和交易处理方面的专长，可以服务众多交易，而无须由每个投资者重复建立。类似任务反复开展时，标准化流程还可进一步节约成本。

交易成本下降可以使原本不经济的融资变得可行。假设为一家小企业融资，在扣除信息与缔约成本之前产生预期剩余 $S$。若个人投资者的成本为 $k>S$，直接融资就不经济。若中介提供这些服务的总成本低于 $S$，它就可能使交易成为可行。

规模收益并非无限。中介扩大时，新增组织层级可能削弱内部激励、增加程序僵化，并降低本地或关系专属信息的有效利用。因此，一些交易成本下降的同时，组织成本可能上升。

\subsection{信息生产与委托监督}

金融中介通过收集、核实和解释借款人特定数据来生产信息。融资之前，这些信息支持筛选和融资条款制定；融资之后，监督更新中介对借款人的评估，并识别行为或信用质量变化。当融资来自众多分散投资者时，将这些活动集中在中介内部，可以避免重复生产信息。

考虑融资后的监督。假设 $N$ 位投资者向同一借款人融资，如果分别监督，每位投资者的监督成本为 $m$，则总监督成本为

\begin{equation}
C_D=Nm.
\label{eq:ch4-direct-monitoring}
\end{equation}

也可以将监督委托给金融中介。令 $m_I$ 表示中介的监督成本，$g$ 表示委托产生的代理成本，则委托监督下的总成本为

\begin{equation}
C_I=m_I+g.
\label{eq:ch4-delegated-monitoring}
\end{equation}

当以下条件成立时，委托监督（delegated monitoring）降低总成本：

\begin{equation}
m_I+g<Nm.
\label{eq:ch4-delegation-condition}
\end{equation}

$g$ 反映委托形成的新代理问题。资金提供者无法完全观察中介的监督努力或风险承担决策。因此，中介可能不如资金提供者所希望的那样认真监督，也可能承担更多风险。

$g$ 的大小取决于中介面对的激励和约束。股本使股东承担不良贷款决策或过度风险承担造成的损失。经审计财务报告改善债权人、投资者和监管者对中介资产质量、资本状况及风险敞口的信息。契约限制和审慎监管，可以进一步约束将潜在损失转给资金提供者的风险承担。这些机制能够减少、但不能消除委托代理成本。

\subsection{私人信息与关系型贷款}

重复互动使贷款人能够从偿还行为、账户活动、抵押品状况以及与管理层的直接接触中积累私人信息。其中一些信息难以量化、编码或传递给外部投资者，通常称为\emph{软信息}（soft information），例如对管理能力或财务困难持续性的判断。

关系型贷款（relationship lending）利用持续借贷关系中积累的信息，评估那些难以通过标准化披露判断信用状况的借款人。这些信息可以帮助贷款人区分在外部投资者看来相似的企业，改善公开信息有限的借款人的融资可得性。

提供融资之后，信息优势同样重要。借款人遭遇暂时挫折时，关系型贷款人可能更有能力区分短期现金流冲击与内在信用质量恶化，并判断是否值得继续放贷或重新谈判。

私人信息也会形成借款人依赖。一家贷款人积累的信息可能难以转移给竞争对手，增加转换成本，并强化原有贷款人的议价能力。因此，关系型贷款需要在改善信用评估的收益，与借款人依赖和贷款人市场势力的潜在成本之间权衡。

\subsection{流动性、期限与风险转换}

除了生产信息，金融中介还转换金融请求权的特征。如果资金提供者重视随时或短期内取得资金，而借款人需要为长期、低流动性资产融资，就会产生流动性转换。银行通过发行具有流动性的存款、持有流动性较低的贷款，弥合这种错配。

汇集众多存款人的个体流动性需求，能够降低总提款额的不确定性。单个存款人的流动性需求时间不确定，但只要这些需求不完全相关，总提款额就会更可预测。下面的简单模型说明这一机制。

假设有 $N$ 位存款人，每人以概率 $\lambda$ 在本期需要一单位流动性。如果流动性需求相互独立，需要资金的存款人人数期望值为

\begin{equation}
\mathbb{E}[L]=N\lambda.
\label{eq:ch4-liquidity-expected}
\end{equation}

其方差为

\begin{equation}
\mathrm{Var}(L)=N\lambda(1-\lambda).
\label{eq:ch4-liquidity-variance}
\end{equation}

标准差随 $\sqrt{N}$ 增长，而预期流动性需求随 $N$ 增长，因此相对不确定性随 $N$ 增加而下降：

\begin{equation}
\frac{\sqrt{\mathrm{Var}(L)}}{\mathbb{E}[L]}
=
\sqrt{\frac{1-\lambda}{N\lambda}}.
\label{eq:ch4-liquidity-relative-risk}
\end{equation}

总提款额的相对不确定性下降，使中介能够持有更多长期、低流动性资产，同时仍向存款人提供具有流动性的请求权。

如果提款高度相关，流动性错配就可能使中介面临严重融资压力。中介可能需要额外融资，或出售非流动资产，以满足提款需求。

流动性转换通常伴随期限转换（maturity transformation）。银行可以用较短期限负债为长期贷款融资，使借款人获得长期资金，而资金提供者持有较短期请求权。由此形成的期限错配使银行面临再融资风险：资产尚未产生足够现金流，到期负债就需要续借。

除了流动性和期限转换，金融中介还通过分散化与损失分配进行风险转换。银行组合信用风险不完全相关的贷款，可以降低借款人特定风险敞口。资本结构还决定损失如何在金融请求权之间分配，其中股权先于债权吸收损失。

\begin{table}[htbp]
\centering
\caption{金融中介的转换功能}
\label{tab:ch4-transformation}

\renewcommand{\arraystretch}{1.15}

\begin{tabularx}{\textwidth}{
@{}
>{\raggedright\arraybackslash}p{3.2cm}
>{\raggedright\arraybackslash}X
>{\raggedright\arraybackslash}X
@{}
}

\toprule
\textbf{转换类型}
& \textbf{资产或风险敞口}
& \textbf{向资金提供者发行的请求权} \\
\midrule

流动性转换
& 流动性相对较低的资产
& 更具流动性或易于兑付的请求权 \\

期限转换
& 期限较长的资产
& 期限较短的负债 \\

风险转换
& 分散化贷款组合
& 具有不同损失敞口的股权与债权 \\

\bottomrule
\end{tabularx}

\end{table}



\section{金融结构：市场、银行与混合体系}

金融结构（financial structure）描述融资如何通过金融市场、金融中介及两者的组合来组织。

\subsection{市场主导型与银行主导型金融体系}

不同金融体系通过证券市场和银行向资金使用者输送资金的相对重要性不同。这是一种程度上的差异，而非两类体系之间的严格分界。

市场主导型融资更依赖标准化披露、可转让金融请求权和市场价格。当相关信息可以低成本评估时，这些安排便利分散投资者提供融资。契约权利和治理机制则提供额外的投资者监督手段。

银行主导型融资更依赖私人信息与持续借款关系。如果借款人信息难以公开核实，或融资需要持续调整，集中监督和可灵活调整的贷款契约就可能更有优势。

市场融资与银行融资的相对成本，取决于企业特征和金融体系条件。更大融资金额能够更广泛地分摊固定发行成本，可信公开披露则降低分散投资者的信息成本。如果评估借款人需要私人信息、持续监督或重新谈判，银行融资可能更经济。证券市场深度和流动性，以及银行的融资与监管约束，也影响两条渠道的相对成本。

在体系层面，企业构成影响市场和银行融资的相对重要性。大型、透明企业较多的经济体，可能更依赖证券市场；借款人特定信息和持续监督较重要的经济体，银行融资可能更重要。法律制度、披露标准、市场基础设施和金融监管，也影响两类渠道的发展。

\begin{table}[htbp]
\centering
\caption{市场主导型与银行主导型融资}
\label{tab:ch4-market-bank}

\renewcommand{\arraystretch}{1.15}

\begin{tabularx}{\textwidth}{
@{}
>{\raggedright\arraybackslash}p{3.0cm}
>{\raggedright\arraybackslash}X
>{\raggedright\arraybackslash}X
@{}
}

\toprule
& \textbf{市场融资}
& \textbf{银行融资} \\
\midrule

信息
& 更依赖标准化公开信息
& 更多使用私人生产的借款人信息 \\

监督
& 通过披露、治理和委托监督实施投资者监督
& 集中于专业中介内部 \\

缔约
& 请求权更标准化、更可转让
& 更有空间定制契约 \\

资金提供者的流动性
& 通过二级市场交易
& 通过存款等具有流动性的银行负债 \\

重新谈判
& 请求权高度分散时更困难
& 持续借贷关系内通常更容易 \\

有利于该渠道的条件
& 融资需求大、公开透明度高
& 不透明度较高，更需监督或重新谈判 \\

\bottomrule
\end{tabularx}
\end{table}

\subsection{混合型与平台型金融体系}

混合型金融体系在相互联系的融资链条中，结合市场机制与中介机制。家庭储蓄可以通过投资基金进入证券市场，银行发起的贷款也可以通过证券化转移给投资者。

数字平台能够拆分传统上由单个中介结合提供的金融功能。借贷平台可以收集借款人数据、核实申请、评估信用、标准化契约并管理还款，而形成的请求权可以由平台保留、转移给银行或出售给机构投资者。如果请求权被转移，而服务仍由平台提供，信息处理和贷款服务就与资金供给及风险承担分离。

\begin{tcolorbox}[
  enhanced,
  breakable,
  colback=white,
  colframe=FinanceBlue,
  boxrule=0.7pt,
  sharp corners,
  borderline west={2.2pt}{0pt}{FinanceBlue},
  title=\textbf{金融视角：透过机构名称看经济功能},
  coltitle=white,
  colbacktitle=FinanceBlue,
  fonttitle=\bfseries,
  before skip=1.2em,
  after skip=1.2em
]

机构名称未必能够识别其履行的经济功能。安排和管理贷款、但转移相关请求权的平台，与保留贷款的平台，面临不同的风险敞口。关键区别在于谁提供资金、持有金融请求权、行使控制权，以及承担损失。








\end{tcolorbox}

\begin{example}{中国的混合型金融结构}
{china-hybrid-financial-structure}

中国金融体系结合了规模可观的银行中介与债券、股票市场。中国人民银行的社会融资规模（Aggregate Financing to the Real Economy，AFRE）框架，记录银行贷款、企业债券、股票发行及其他渠道的融资。信贷资产证券化允许银行发起的贷款转移到向投资者发行证券的结构中，进一步连接银行体系与证券市场。









\end{example}

\section{什么塑造金融结构？}

\subsection{法律制度与投资者保护}

法律制度通过金融请求权的保护和执行影响金融结构。对债权人而言，破产程序、优先顺序规则与抵押权影响预期回收；对外部股东而言，法律和治理保护可以限制企业资源被转移，并保障其在公司决策中的权利。

这些保护在证券市场尤其重要，因为分散投资者单独协商和监督的空间有限。标准化披露、投资者权利与有效执行，能够支持广泛投资者群体提供融资。

银行融资更依赖借款人特定契约和集中监督，但同样依赖可执行的债权人权利。执行薄弱或不确定，可能降低预期回收并限制信贷供给。

投资者保护也影响所有权和控制权结构。更强的外部股东保护可以支持分散所有权；保护较弱则可能鼓励集中所有权和大股东监督。然而，控制权集中也可能使控股股东以牺牲中小投资者利益为代价，攫取私人收益。

\subsection{经济发展与金融深化}

经济发展通过增加金融服务需求、使更多金融安排在经济上可行，影响金融结构。收入、财富和经济活动增长，可以扩大金融活动规模。更大的交易量，又能将信息生产、金融基础设施和证券发行的固定成本，分摊到更广基数上。

金融深化（financial deepening）是指相对于经济规模而言，金融中介和市场活动的规模与范围扩大。它可以通过银行信贷、证券市场、机构投资者与金融基础设施扩张实现。这些组成部分能够相互强化：机构投资者增长增加对可交易证券的需求，发达的政府债券市场则可以为其他证券定价提供基准收益率，并为金融交易提供抵押品。

\begin{figure}[htbp]
\centering

\begin{tikzpicture}[
    >=Latex,
    box/.style={
        draw=FinanceBlue,
        rounded corners=2pt,
        align=center,
        minimum height=1.05cm,
        text width=2.8cm,
        inner xsep=1pt,
        inner ysep=4pt
    },
    arrow/.style={
        -Latex,
        line width=0.8pt
    }
]

% Nodes
\node[box] (development) at (0,0)
{经济发展};

\node[box] (demand) at (3.8,1.15)
{金融服务\\需求增加};

\node[box] (volume) at (3.8,-1.15)
{交易量\\扩大};

\node[box] (cost) at (7.1,-1.15)
{平均固定\\成本下降};

\node[box] (deepening) at (10.6,0)
{金融深化};

% Arrows
\draw[arrow] (development.east) -- ++(0.45,0) |- (demand.west);
\draw[arrow] (development.east) -- ++(0.45,0) |- (volume.west);

\draw[arrow] (demand.east) -| ([xshift=-0.45cm]deepening.west) -- (deepening.west);
\draw[arrow] (volume.east) -- (cost.west);
\draw[arrow] (cost.east) -| ([xshift=-0.45cm]deepening.west) -- (deepening.west);

\end{tikzpicture}

\caption{经济发展与金融深化}
\label{fig:ch4-financial-deepening}
\end{figure}

经济发展能够同时支持金融中介和证券市场。随着金融体系深化，两者可能通过融资、投资和流动性关系，变得更加相互联系。

然而，金融深度并不等同于金融效率。信贷、市值或金融资产增加，本身并不意味着资源配置更有效。信贷可能流向预期回报低的项目，杠杆可能增加金融脆弱性，证券市场也可能在披露或执行未相应改善时扩张。金融深化的经济价值，取决于金融扩张改善资本配置、信息生产、风险分担、流动性和服务可得性的收益，相对于它引入的成本与脆弱性如何。

\subsection{技术与金融基础设施}

技术进步可以通过降低信息处理和金融交易成本来影响金融结构。数字信息系统能够降低筛选与核实成本，电子交易系统则可降低证券交易成本。

这些影响对市场和中介未必对称。如果技术降低借款人信息标准化、公开核实的成本，原本分散投资者难以评估的借款人，就可能更容易采用市场融资。反之，如果专有交易数据改善筛选和监督，拥有大量数据的银行或平台就可能获得相对信息优势。

金融基础设施是技术影响金融结构的另一渠道。市场和中介都依赖记录请求权、转移资金、交换信息和结算交易的系统。这些系统的可靠性与互操作性，会影响不同融资渠道的效率。


\subsection{监管与制度变化}

监管可以通过改变不同金融安排的成本和约束来影响金融结构。即使主要目标是审慎监管或投资者保护，资本、流动性、披露和投资规则，仍会改变银行、证券发行人及机构投资者的经营条件。

监管待遇差异影响金融活动的组织方式。如果监管提高银行在资产负债表上持有某类信用敞口的成本，一部分活动可能转向证券市场或非银行中介。即使机构形式改变，提供信贷的底层功能可能仍相似。这些转移可能只是重新配置而非消除风险，因为市场型和非银行安排可能涉及不同的杠杆、融资结构与流动性风险组合。

监管也与金融创新相互作用。新融资安排可能引发监管回应，而回应又影响其可行性与组织形式。

%========================================================
% Summary
%========================================================

\TailModule{summary}{本章小结}

金融结构关注融资、信息生产、监督、流动性供给和风险分配，如何在市场与中介之间组织。直接融资中，投资者持有最终资金使用者发行的请求权；间接融资中，中介持有对借款人的请求权，并向资金提供者发行自身请求权。现代金融体系通常在相互联系的融资链条中结合两种形式。

如果信息能够标准化并由公众评价，融资规模又足以分摊固定发行成本，市场融资会更有效。中介通过专业化、规模经济、私人信息生产和委托监督创造价值。委托可以减少分散投资者的重复监督，但也引入资金提供者与中介之间的代理问题。

中介还转换金融请求权的特征。汇集流动性需求，可以支持以流动性较低资产为基础的流动负债；期限转换使较长期资产能够用较短期负债融资；风险转换则通过汇集风险和在不同金融请求权之间分配损失实现。这些转换在提供金融服务的同时，也形成潜在脆弱性来源。

市场与中介的相对作用，既取决于融资特征，也取决于更广泛的经济与制度环境。法律保护与执行、经济发展、技术、金融基础设施和监管，会改变不同金融安排的成本和可行性。随着市场、中介和混合安排适应变化中的经济与制度条件，金融结构不断演化。

%========================================================
% Mind Map
%========================================================

\TailModule{mindmap}{思维导图}

\begin{figure}[H]
\centering

\begin{tikzpicture}[
  >=Latex,
  main/.style={
    draw,
    rounded corners,
    minimum height=0.9cm,
    align=center,
    fill=gray!10,
    inner xsep=5pt,
    inner ysep=4pt
  },
  sub/.style={
    draw,
    rounded corners,
    minimum height=0.85cm,
    align=center,
    inner xsep=4pt,
    inner ysep=4pt
  },
  flow/.style={
    -Latex,
    line width=0.75pt
  }
]

% Top
\node[main, text width=3.4cm] (functions) at (0,4.2)
{金融功能};

% Organizational forms
\node[sub, text width=3.2cm] (markets) at (-5.0,2.2)
{市场与\\直接融资};

\node[sub, text width=3.2cm] (inter) at (0,2.2)
{中介与\\间接融资};

\node[sub, text width=3.2cm] (hybrid) at (5.0,2.2)
{混合型\\融资安排};

% Mechanisms
\node[sub, text width=3.5cm] (marketmech) at (-5.0,0)
{公开信息\\与可交易性};

\node[sub, text width=3.8cm] (intermech) at (0,0)
{信息、监督\\与转换};

\node[sub, text width=3.8cm] (hybridmech) at (5.0,0)
{市场与中介\\共同履行功能};

% Financial structure
\node[main, text width=3.4cm] (structure) at (0,-2.2)
{金融结构};

% Environment
\node[sub, text width=5.2cm] (environment) at (0,-4.2)
{法律、发展、技术、\\基础设施与监管};

% Arrows
\draw[flow] (functions) -- (markets);
\draw[flow] (functions) -- (inter);
\draw[flow] (functions) -- (hybrid);

\draw[flow] (markets) -- (marketmech);
\draw[flow] (inter) -- (intermech);
\draw[flow] (hybrid) -- (hybridmech);

\draw[flow] (marketmech) -- (structure);
\draw[flow] (intermech) -- (structure);
\draw[flow] (hybridmech) -- (structure);

\draw[flow] (environment) -- (structure);

\end{tikzpicture}

\caption*{从金融功能到金融结构}
\end{figure}

%========================================================
% Key Terms
%========================================================

\TailModule{keyterms}{关键术语}

直接融资（Direct finance）；间接融资（indirect finance）；金融中介（financial intermediary）；融资链条（financing chain）；混合型金融体系（hybrid financial system）；规模经济（economies of scale）；委托监督（delegated monitoring）；关系型贷款（relationship lending）；流动性汇集（liquidity pooling）；流动性转换（liquidity transformation）；期限转换（maturity transformation）；风险转换（risk transformation）；金融结构（financial structure）；市场主导型金融体系（market-based financial system）；银行主导型金融体系（bank-based financial system）；金融深化（financial deepening）；投资者保护（investor protection）；金融基础设施（financial infrastructure）；平台型金融（platform-based finance）；非银行金融中介。

%========================================================
% Financial AI Lab
%========================================================（non-bank financial intermediation

%========================================================
% Financial AI Lab
%========================================================）。

\TailModule{aiprompt}{金融 AI 实验室}

\textbf{实验：重建融资链条}

选择一种当代融资安排，例如数字借贷平台、私募信贷基金、证券化载体、众筹平台或在线理财产品。请 AI 系统将其划为直接或间接融资，并将初步分类视为暂定结论。

重建完整融资链条，识别最终资金提供者和使用者、每个阶段创设的金融请求权，以及每项请求权的持有人。再确定谁开展筛选和监督、谁提供流动性，以及最终由谁承担信用损失。

将你的重建结果与 AI 分类比较。如果同一融资链条结合市场和中介功能，解释简单的直接／间接二分为什么可能不够。利用相关金融机构、监管机构、交易所或平台的一手文件，核实制度结构和风险分配。

%========================================================
% Discussion Questions
%========================================================

\TailModule{questions}{讨论题}

\begin{enumerate}

\item 为什么投资银行参与公司债券发行，并不必然使交易成为间接融资？

\item 解释金融中介的委托监督理论。委托如何在减少重复监督的同时，产生新的代理问题？

\item 当个人流动性需求并非完全相关时，存款人人数增加如何降低总流动性需求的相对不确定性？

\item 为什么应将市场主导型与银行主导型金融体系的区别理解为程度差异，而非二元分类？

\item 同一技术变化如何对证券市场和金融中介产生不同影响？

\end{enumerate}

%========================================================
% Data Application
%========================================================

\TailModule{dataapp}{数据应用}

\textbf{中国金融结构及国际比较}

使用 IMF 金融发展指数数据库，在共同样本期内比较中国和至少三个其他经济体的金融机构指数与金融市场指数。选择金融结构具有明显差异的经济体，考察金融发展主要通过机构、市场还是两者共同实现。应将这些指数理解为金融发展不同维度的度量，而非金融体系优劣排名。

针对中国，使用中国人民银行社会融资规模统计考察融资构成。比较人民币贷款与企业债券、非金融企业境内股票融资等选定的直接融资类别。解释为什么社会融资规模度量的融资流量，与 IMF 指数反映的是金融结构的不同维度。

利用证据评价本章的一个命题，例如银行中介持续重要、市场型融资发展，或两类渠道共存。报告指标定义、样本期、所选经济体和一手数据来源。

%========================================================
% Further Reading
%========================================================

\TailModule{furtherread}{延伸阅读}

\begin{itemize}

\item Diamond, Douglas W.（1984），“金融中介与委托监督”（``Financial Intermediation and Delegated Monitoring''），
\emph{The Review of Economic Studies}, 51(3), 393--414.

\item Allen, Franklin, and Douglas Gale (2000), 
\emph{Comparing Financial Systems}（《金融体系比较》），MIT Press。

\item Levine, Ross（2002），“银行主导型还是市场主导型金融体系：哪一种更好？”（``Bank-Based or Market-Based Financial Systems: Which Is Better?''），
\emph{Journal of Financial Intermediation}, 11(4), 398--428.

\item Beck, Thorsten, Asli Demirg\"{u}\c{c}-Kunt, and Ross Levine (2003), 
“法律与金融：为什么法律起源重要？”（``Law and Finance: Why Does Legal Origin Matter?''），
\emph{Journal of Comparative Economics}, 31(4), 653--675.

\item Svirydzenka, Katsiaryna (2016), 
“介绍一种新的综合金融发展指数”（``Introducing a New Broad-based Index of Financial Development''），
IMF 工作论文，第~16/5~号。

\end{itemize}

\ChapterTail{summary,mindmap,keyterms,aiprompt,questions,dataapp,furtherread}

\chapter{金融估值的逻辑}
\label{chap:The_Logic_of_Financial_Valuation}

\begin{openingcase}{预期现金流相同，价值也相同吗？}

一位投资者正在比较两项一年期金融请求权，两者的预期支付均为人民币 100 元。请求权 A 在经济强劲时支付 120 元，在经济疲弱时支付 80 元；请求权 B 则相反，在经济强劲时支付 80 元，在经济疲弱时支付 120 元。假设两种状态出现的概率相同。

两者的预期支付相同：

\[
0.5(120)+0.5(80)=100.
\]

然而，两项请求权的当前价值未必相同。请求权 B 在经济疲弱时支付更多，而此时额外一单位现金可能特别有价值。因此，即使预期支付相同，支付在经济状态之间的分布仍会影响估值。

\end{openingcase}

\begin{tcolorbox}[
  enhanced,
  colback=FinanceBlue!6,
  colframe=FinanceBlue,
  boxrule=0.6pt,
  sharp corners,
  left=8pt,
  right=8pt,
  top=7pt,
  bottom=7pt,
  before skip=1em,
  after skip=1.6em,
  title=\textbf{核心问题},
  coltitle=white,
  colbacktitle=FinanceBlue,
  fonttitle=\bfseries,
  attach boxed title to top left={xshift=6pt,yshift=-2mm},
  boxed title style={
    boxrule=0pt,
    sharp corners,
    left=5pt,
    right=5pt,
    top=2pt,
    bottom=2pt
  }
]
\begin{enumerate}[
  leftmargin=*,
  label=\textcolor{FinanceBlue}{\textbf{\arabic*.}},
  itemsep=0.55em,
  topsep=0.2em
]

\item 为什么仅凭预期现金流不足以确定金融请求权的价值？

\item 未来现金流的时间、风险和状态依存性，如何影响其现值？

\item 信息、契约权利与市场价格如何相互作用，共同影响金融请求权的价值？

\end{enumerate}
\end{tcolorbox}


\section{金融价值的基础}

金融请求权的价值来自它赋予持有人的未来经济利益。使用估值公式之前，需要回答三个问题：


\begin{itemize}[
  leftmargin=2.2em,
  itemsep=0.45em,
  topsep=0.55em
]

\item 该请求权赋予持有人获取哪些现金流的权利？

\item 这些现金流何时收到？

\item 这些现金流在未来不同世界状态下如何变化？

\end{itemize}

\subsection{金融资产：对未来经济资源的请求权}

金融资产代表对未来经济资源的权利，通常通过未来现金流表达。
债券赋予持有人获取约定支付的契约权利，普通股则赋予持有人在优先请求权获满足之后对公司的剩余请求权。衍生工具创造取决于标的资产、利率、指数或状态变量的或有现金流。

因此，估值对象是由这些权利界定的\emph{请求权}（financial claim），而不只是发行人、项目或标的资产。假设企业经营资产在时点 $t$ 产生现金流 $Y_t$，这一现金流可能分配给多个请求权人：

\begin{equation}
Y_t=C_{D,t}+C_{E,t}+C_{O,t},
\label{eq:ch5-cash-flow-allocation}
\end{equation}
其中，$C_{D,t}$ 表示债权人现金流，$C_{E,t}$ 表示股东现金流，$C_{O,t}$ 表示其他契约请求权人的现金流。

金融请求权的价值不仅取决于底层资产产生的现金流，还取决于决定现金流如何分配的契约规则。

优先顺序提供了一个简单例子。考虑仅由债务和股权融资的企业，假设下一期产生现金流 $Y$，债务承诺偿还额为 $D$。不考虑其他契约条款时，债权人获得

\begin{equation}
C_D=\min\{Y,D\},
\label{eq:ch5-debt-payoff}
\end{equation}
而股东获得

\begin{equation}
C_E=\max\{Y-D,0\}.
\label{eq:ch5-equity-payoff}
\end{equation}

因此，

\begin{equation*}
C_D+C_E=Y.
\label{eq:ch5-claim-additivity}
\end{equation*}

契约结构不改变企业总现金流，而是决定在 $Y$ 的每一种实现值下，如何在请求权人之间分配现金流。图~\ref{fig:ch5-debt-equity-payoff} 展示了这一分配。

\begin{figure}[H]

\centering

\begin{tikzpicture}[
=Latex,
axis/.style={-Latex, line width=0.8pt},
debt/.style={line width=1.0pt, FinanceBlue},
equity/.style={line width=1.0pt, dashed}
 ]

\draw[axis] (0,0) -- (7.2,0) node[right] {企业现金流 $Y$};
\draw[axis] (0,0) -- (0,5.4) node[above] {请求权现金流};

\draw[debt] (0,0) -- (3.6,3.6) -- (6.6,3.6);
\draw[equity] (0,0) -- (3.6,0) -- (6.6,3.0);

\draw[dashed, gray] (3.6,0) -- (3.6,4.6);

\node[below] at (3.6,0) {$D$};
\node[FinanceBlue, right] at (6.0,3.85) {债务};
\node[right] at (6.2,2.55) {股权};

\end{tikzpicture}
\caption{企业现金流在债务与股权之间的分配}
\label{fig:ch5-debt-equity-payoff}

\end{figure}

\subsection{时间与未来现金流的现值}

未来收到的现金流，与今天收到的同额现金流，通常并不等值。个人可能要求因推迟当前消费而获得补偿，今天可用的资金也可以投资并产生未来回报。这些因素构成货币时间价值、利息和贴现的基本经济基础。

将未来现金流按适当利率贴现，可以把它表示为今天的等值金额。

令 \(V_0\) 表示一期后确定支付 \(CF_1\) 的请求权现值。如果相关的一期无风险利率为 \(r_f\)，无套利等价关系要求：

\begin{equation*}
V_0(1+r_{f})=CF_1.
\label{eq:ch5-one-period-equivalence}
\end{equation*}
由此可得

\begin{equation*}
V_0=\frac{CF_1}{1+r_{f}}.
\label{eq:ch5-one-period-pv}
\end{equation*}

因子 $1/(1+r_f)$ 是一期贴现因子（discount factor），表示一期后收到一单位现金流的现值。



假设每期贴现率恒定为 $r$，则 $t$ 期后收到现金流 $CF_t$ 的现值为



\[
V_0=\frac{CF_t}{(1+r)^t}.
\]

一系列未来现金流的现值为


\begin{equation}
V_0=\sum_{t=1}^{T}
\frac{CF_t}{(1+r)^t}.
\label{eq:ch5-pv-cash-flow-stream}
\end{equation}

贴现将不同日期的现金流转换为共同日期的价值。转换到共同基础后，现金流就能够比较和相加。

\begin{figure}[H]
\centering
\begin{tikzpicture}[
  >=Latex,
  axis/.style={-Latex, line width=0.8pt},
  cash/.style={-Latex, line width=1.0pt, FinanceBlue},
  discount/.style={-Latex, dashed, line width=0.8pt}
]

\draw[axis] (0,0) -- (10.2,0) node[right]{时间};

\foreach \x/\t in {0/0,2/1,4/2,6/3,8/4}
{
  \draw (\x,0.12) -- (\x,-0.12);
  \node[below=4pt] at (\x,0) {$t=\t$};
}

\draw[cash] (2,0) -- (2,1.2) node[above] {$CF_1$};
\draw[cash] (4,0) -- (4,1.55) node[above] {$CF_2$};
\draw[cash] (6,0) -- (6,1.05) node[above] {$CF_3$};
\draw[cash] (8,0) -- (8,1.45) node[above] {$CF_4$};

\draw[discount] (2,2.25) -- (0.35,2.25);
\draw[discount] (4,2.25) -- (0.35,2.25);
\draw[discount] (6,2.25) -- (0.35,2.25);
\draw[discount] (8,2.25) -- (0.35,2.25);

\node[above] at (4.1,2.25) {贴现至 $t=0$};

\end{tikzpicture}

\caption{将未来现金流贴现至当前}
\label{fig:ch5-discount-timeline}
\end{figure}

贴现直接带来以下几点含义。

第一，即使名义现金流总额不变，时间仍然重要。若两项请求权支付总额相同，但一项更早支付，那么在贴现率为正时，较早支付通常具有更高现值。

第二，现值对贴现率的敏感性随支付时间推远而提高。因此，贴现率变化对更远期现金流的影响更大。

第三，现金流时间必须与贴现率惯例一致。年利率不能未经适当转换就直接用于季度现金流，还必须区分期初与期末发生的现金流。

\begin{example}{仅时间差异就会改变价值}{timing-alone-changes-value}

假设请求权 A 在一年后支付人民币 1,000 元，请求权 B 在五年后支付 1,000 元。若相关年利率为 5\%，

\[
V_A=\frac{1,000}{1.05}=952.38,
\]

而

\[
V_B=\frac{1,000}{1.05^5}=783.53.
\]

两者未来名义支付相同，现值差异完全来自支付日期不同。

\end{example}

\subsection{风险与必要报酬率}

价值并非仅由时间决定。许多金融请求权的未来现金流具有风险。假设下一期有两种可能状态，请求权分别以概率 $p$ 和 $1-p$ 支付 $CF_H$ 和 $CF_L$，其预期现金流为

\begin{equation}
\mathrm{E}[X]=pCF_H+(1-p)CF_L.
\label{eq:ch5-expected-cash-flow}
\end{equation}

通常，仅有预期现金流不足以估值，请求权价值还取决于现金流风险。如果现金流在整体经济状况较差时往往较低，就使投资者面临更多风险。这类请求权通常需要更高预期回报。因此，预期支付相同的两项请求权，当前价值可能不同。

将风险纳入估值的一种常见方式是使用必要报酬率（required return）。令 $r_f$ 为无风险利率，$\pi$ 为适用于该请求权的风险溢价，则必要报酬率为

\begin{equation}
r=r_f+\pi.
\label{eq:ch5-required-return}
\end{equation}

以预期现金流为基础的简化一期估值为

\begin{equation}
V_0=\frac{\mathrm{E}[CF]}{1+r}.
\label{eq:ch5-risk-adjusted-discount}
\end{equation}

风险溢价（risk premium）补偿投资者承担与估值有关的风险。可分散风险可能只要求很少溢价，甚至没有溢价；在整体经济条件较差时往往造成损失的风险，则可能要求更高溢价。

更一般地，可以通过风险调整现金流、贴现率或状态依存估值来纳入风险。风险处理必须与所估值请求权保持一致。

\medskip

下面概括金融估值的主要维度及相应分析输入。

\begin{table}[H]

\centering

\caption{估值中的三个问题}
\label{tab:ch5-three-valuation-questions}

\begin{tabularx}{\textwidth}{
  @{}
  >{\raggedright\arraybackslash\hyphenpenalty=10000\exhyphenpenalty=10000}p{3.3cm}
  Y
  Y
  @{}
}

\toprule

\textbf{问题}
& \textbf{经济问题}
& \textbf{典型估值输入} \\

\midrule

收到什么？
& 对未来经济资源的权利
& 现金流定义与契约支付 \\[20pt]

何时收到？
& 跨期估值
& 贴现因子或利率 \\[20pt]

在哪些状态下收到？
& 风险与状态依存性
& 风险溢价或状态依存估值 \\

\bottomrule

\end{tabularx}

\end{table}


\section{估值框架}

适当的估值方法取决于金融请求权的性质及可用信息。

现金流贴现（discounted cash-flow，DCF）估值根据归属于某项请求权的预期现金流推导价值；相对估值使用可比资产的观察价格；或有请求权估值则明确考虑状态依存支付，尤其适用于非线性或取决于未来事件的支付。

尽管这些方法的输入和实施方式不同，它们都提供价值估计。令 $V$ 表示金融请求权的估计价值，$P$ 表示观察到的市场价格。在相对估值中，可比请求权的市场价格作为输入，用于估计待估值请求权的价值；随后可将估值与该请求权的市场价格比较。

\subsection{现金流贴现估值}

现金流贴现估值，通过贴现归属于某项请求权的现金流，估计其价值。

若 $r_j$ 表示时点 $j-1$ 到 $j$ 之间适用的一期必要报酬率，一系列预期现金流的价值可写为

\begin{equation}
V_0 = \sum_{t=1}^{T}
\frac{\mathrm{E}[CF_t]}
{\prod_{j=1}^{t}(1+r_j)}.
\label{eq:ch5-general-dcf}
\end{equation}

若必要报酬率恒定为 $r$，则

\begin{equation}
V_0=\sum_{t=1}^{T}
\frac{\mathrm{E}[CF_t]}{(1+r)^t}.
\label{eq:ch5-constant-dcf}
\end{equation}

公式本身简单，但应用取决于如何定义现金流、风险和长期假设。

\paragraph{先定义请求权，再预测现金流。}

会计利润并不直接度量可供某项金融请求权使用的现金流。权责发生制、非现金费用、营运资本变化和投资支出，都可能使报告利润不同于现金流。因此，相关现金流指标取决于估值对象。

对项目而言，相关现金流是实施项目产生的增量现金流。对股权而言，是满足经营需要、投资、税收及优先请求权人支付后，可供股东取得的现金流。对债务而言，是归属于债权的支付，并需考虑违约与回收的可能性。







\paragraph{区分预测与估值。}

预测是在给定假设下估计未来现金流，估值则将这些现金流转化为现值。区分两项任务，有助于识别错误究竟来自预测本身，还是对预测的估值。复杂贴现模型无法弥补不现实的增长假设、被低估的投资需求，或现金流预测中的不一致。

假设项目需要初始投资 $I_0$，产生预期现金流 $\mathrm{E}[CF_t]$，其净现值（net present value，NPV）为

\begin{equation}
NPV=\sum_{t=1}^{T}
\frac{\mathrm{E}[CF_t]}{(1+r)^t}-I_0.
\label{eq:ch5-npv}
\end{equation}

NPV 为正，意味着在既定假设下，预期未来现金流的现值超过今天所需资源。它表示按相关必要报酬率估值时，项目具有正经济剩余；并不意味着在某些实际实现的状态下不会亏损。

\begin{example}{从现金流到净现值}{from-cash-flows-to-npv}


某项目初始投资为人民币 1,000 万元，预计第 1 年末产生 300 万元、第 2 年末产生 400 万元、第 3 年末产生 550 万元。假设必要报酬率为 8\%。

未来现金流的现值约为

$$
PV=\frac{3.0}{1.08}+\frac{4.0}{1.08^2}+\frac{5.5}{1.08^3}=10.57
$$
百万人民币。因此，

$$
NPV=10.57-10.00=0.57
$$

百万人民币。

在给定现金流和必要报酬率假设下，该项目净现值为正。

\end{example}

\paragraph{终值将预测延伸到显式预测期之后。}

许多资产在显式预测期结束后仍继续产生现金流。因此，现金流贴现估值通常将价值分为两部分：截至时点 $T$ 逐期预测的现金流，以及时点 $T$ 的终值，后者代表该时点之后预期现金流的价值。

\begin{equation}
V_0=\sum_{t=1}^{T}
\frac{\mathrm{E}[CF_t]}{(1+r)^t}+\frac{TV_T}{(1+r)^T}.
\label{eq:ch5-dcf-terminal}
\end{equation}

如果预期第 $T$ 年后的现金流将永远按恒定增长率 $g$ 增长，则时点 $T$ 的终值为

\begin{equation}
TV_T=\frac{CF_{T+1}}{r-g},
\qquad r>g.
\label{eq:ch5-growing-perpetuity}
\end{equation}

终值（terminal value）对 $r$ 和 $g$ 都敏感。显式预测期较短时，终值可能占总估值很大比例，此时估值高度依赖这些参数背后的长期假设。

要使增长型永续年金公式产生有限价值，必须满足 $r>g$。所假设的长期增长率，还应与相关市场或经济的可持续增长一致。

\paragraph{敏感性分析识别最重要的假设。}

估值以其假设为条件。敏感性分析（sensitivity analysis）考察关键输入变化时，估计价值如何变化。

敏感性分析度量估值结果对选定假设的依赖，而不是不同结果发生的概率。更远期现金流对贴现率尤其敏感，终值则往往对长期增长假设敏感。

现金流贴现估值的输入，建立在经济、会计、契约和市场分析之上。销售额、利润率、再投资及其他现金流预测，反映对未来经营条件的假设；必要报酬率取决于请求权风险及相关市场证据。DCF 将这些输入纳入一致的估值框架。

\subsection{相对估值：可比对象与市场倍数}

相对估值（relative valuation）依据经济上可比资产的观察价格，估计资产价值。

假设可比企业市盈率约为 15 倍。如果待估值企业每股收益为人民币 2 元，隐含价格为 $15\times 2=30$ 元。

估值倍数表示市场价值相对于利润、账面价值、销售额、现金流或其他经济指标的比例。常见例子包括

\begin{equation*}
\frac{P}{E},
\qquad
\frac{P}{B},
\qquad
\frac{EV}{EBITDA},
\qquad
\frac{EV}{Sales}.
\label{eq:ch5-common-multiples}
\end{equation*}

估值倍数的分子和分母，应对应同一组请求权人。股权市场价值 $P$ 应与归属股东的利润或股权账面价值等指标匹配；企业价值 $EV$ 应与可供所有资本提供者分享的经营指标匹配，例如息税折旧摊销前利润（EBITDA）或销售额。

企业内在经济特征不同，估值倍数也会不同。考虑一只股利预计永久按恒定增长率 $g$ 增长的股票。若下一期每股股利为 $D_1$，股权必要报酬率为 $r$，其价值为

\begin{equation}
P_0=\frac{D_1}{r-g}.
\label{eq:ch5-dividend-growth}
\end{equation}

这一估值关系称为戈登增长模型（Gordon growth model）。
在股利支付率恒定为 $\alpha$ 时，下一期每股股利满足 $D_1=\alpha E_1$，其中 $E_1$ 是下一期每股收益。将其代入式~\eqref{eq:ch5-dividend-growth} 得：

\begin{equation}
\frac{P_0}{E_1}=\frac{\alpha}{r-g}.
\label{eq:ch5-pe-drivers}
\end{equation}

即使在这一简化框架中，市盈率也取决于股利支付率、必要报酬率和预期增长。更一般地，增长取决于再投资规模及再投资回报，盈利能力则影响企业资产基础产生的现金流。因此，估值倍数反映有关增长、风险、分配和盈利能力的内在假设。

相对估值的可靠性，取决于资产在经济上的可比性，以及倍数变量的计量。行业分类是有用起点，但同行业企业也可能在增长、杠杆、经营风险、会计惯例和地域敞口上不同。相对估值纳入当前市场价格，可作为评价估值的有用基准；但全市场错误定价、会计差异、周期性或负收益，以及财务杠杆差异，可能降低观察倍数的有用性。


\subsection{或有请求权估值：状态依存支付}

一些金融请求权的支付随未来状态变化。因此，估值不仅取决于预期现金流，还取决于支付在状态间如何分布。

期权是一个明显例子，但状态依存支付的存在更为广泛。股权体现有限责任，信用工具暴露于违约状态，保险契约取决于约定损失事件。担保、可转换证券和结构化请求权也具有类似特征。

假设时点 $T$ 有 $S$ 种可能状态，$CF_T(s)$ 表示状态 $s$ 发生时请求权在时点 $T$ 的支付。令 $q_s$ 表示一种证券在时点 0 的价格，该证券在时点 $T$ 仅当状态 $s$ 发生时支付一单位货币，否则支付零。$q_s$ 称为状态 $s$ 的状态价格（state price）。

状态价格反映在不同状态下获得资源的价值。资源相对稀缺状态下的一单位支付，可能比有利状态下的同额支付更有价值。因此，状态依存估值根据支付发生的状态赋予不同价值，而不是把每单位预期支付都视为经济上等价。

在状态价格估值下，请求权在时点 0 的价值为

\begin{equation}
V_0=\sum_{s=1}^{S}q_s CF_T(s).
\end{equation}

\begin{example}{状态依存估值}{state-dependent-valuation}

考虑两种等概率状态 $s=1,2$。请求权 A 在状态 1 支付 120、状态 2 支付 80；请求权 B 在状态 1 支付 80、状态 2 支付 120。两者预期支付均为 100。

假设相应状态价格为

$$
q_{1}=0.40,
\qquad
q_{2}=0.55.
$$

两项请求权在时点 0 的价值为

\begin{align}
V_A
&=
0.40(120)+0.55(80)=92,\\
V_B
&=0.40(80)+0.55(120)=98.
\end{align}

请求权 B 在时点 0 的价值更高，因为它在状态价格更高的状态 2 支付更多。

\end{example}

状态价格还与无风险贴现有关。在每种状态下均支付一单位的请求权是无风险的，其在时点 0 的价值为

\begin{equation}
\sum_{s=1}^{S}q_s=\frac{1}{1+r_f},
\label{eq:ch5-state-price-riskfree}
\end{equation}
其中，$r_f$ 是一期无风险利率。因此，状态价格之和等于一期无风险贴现因子。更一般地，状态价格估值将状态依存请求权在时点 0 的价值，表示为各状态支付与相应状态价格的加权和。

状态价格估值与无套利原则紧密相关。令 $CF_T^A(s)$ 和 $CF_T^B(s)$ 分别表示组合 A、B 在时点 $T$、状态 $s$ 的支付，$P_A$ 和 $P_B$ 为时点 0 的价格。根据无套利原则，如果每种可能状态下支付都相同，时点 0 的价格就必须相同：

\begin{equation}
CF_T^A(s)=CF_T^B(s)
\quad \text{对每个 } s
\quad \Longrightarrow \quad
P_A=P_B.
\label{eq:ch5-law-one-price}
\end{equation}

如果价格不同，投资者就可以买入低价组合、卖出高价组合，同时未来支付相互抵消，从而形成套利机会。

状态价格同时反映货币时间价值，以及不同状态支付的相对价值。请求权的估值，是将每种状态依存支付按相应状态价格加权。

状态价格表示法也澄清了估值中风险的性质。与低状态价格状态下的等额支付相比，高状态价格状态下的支付对时点 0 价值贡献更大。因此，仅凭支付方差不能决定价值。即使预期支付相同，支付波动更大也未必意味着价值更低；其影响取决于支付如何分布于具有不同状态价格的状态。

\section{信息、契约与价值形成}

估值受有关未来现金流和风险的可用信息影响，也受待估值请求权所附契约权利影响。

\subsection{信息与现金流预期}

估值依赖对未来现金流和风险的预期，而预期根据可用信息形成。营业收入、违约损失、投资需求和长期增长都无法事先确定，而是根据估值日可用信息估计。

令 $\mathcal{I}_0$ 表示时点 0 的可用信息，请求权价值可写为

\begin{equation}
V_0=\sum_{t=1}^{T}
\frac{\mathrm{E}[CF_t\mid \mathcal{I}_0]}
{\prod_{j=1}^{t}(1+r_j)},
\label{eq:ch5-information-conditional-value}
\end{equation}
其中，预期现金流以可用信息为条件，$r_j$ 为第 $j$ 期必要报酬率。

新信息可以改变预期现金流、必要报酬率或两者，从而改变价值。令 $\mathcal{I}_0'$ 表示时点 0 可用的更新后信息集，则




\begin{equation*}
\mathrm{E}[CF_t\mid \mathcal{I}_0']
\neq
\mathrm{E}[CF_t\mid \mathcal{I}_0]
\end{equation*}

新信息通过两条渠道影响估值：一是改变预期未来现金流，二是改变对请求权风险的判断，进而改变必要报酬率。

新增信息对价值的影响并非单向。它可能提高或降低预期现金流，也可能改变对内在风险的评估。信息质量改善能够减少估计不确定性，却未必降低底层经济风险。估值既依赖信息数量，也依赖相关性和可靠性。








\subsection{契约与现金流权利分配}

契约界定金融请求权的现金流权利，包括支付条件，以及相对于其他请求权的优先顺序。这些契约特征与估值直接相关。

对仅以债务和股权融资的企业，令 $V_A$ 表示资产价值，相应请求权的价值满足

\begin{equation}
V_A=V_D+V_E,
\label{eq:ch5-value-allocation}
\end{equation}
其中，$V_D$ 和 $V_E$ 分别为债务与股权价值。即使底层经营资产价值不变，契约条款也能影响价值在不同请求权之间的分配。

优先顺序决定支付次序，抵押品影响违约回收，担保则在各方之间重新分配损失。

\begin{example}{抵押品与请求权价值}{collateral-claim-value}


假设两项一年期债权均承诺偿还 100。第一项有担保，抵押品违约时可变现价值为 60；第二项无担保。如果违约发生，且企业剩余资产不足以全额偿还两项债权，有担保债权的回收更高。因此，即使承诺支付和期限相同，两项请求权也可能具有不同价值。

\end{example}

契约也通过激励影响价值。限制性条款可以限制新增借款、股利支付，或其他增加债权人风险的行动。这些限制能够保护债权，却可能降低经营灵活性。

其他契约条款创造状态依存权利。例如，可转换债券赋予持有人在特定条件下将债券换成股权的权利。转换权、赎回权等嵌入式期权，改变未来不同状态下的支付。

契约条款通过决定支付顺序、影响激励和界定状态依存支付，影响请求权价值。因此，对复杂金融工具，估值取决于其所附契约权利，而不只是工具名称。

\subsection{市场价格与估值}

\paragraph{市场价格与价格发现。}
市场价格是请求权成交的价格，估值则根据特定信息、假设和定价关系估计其价值。市场价格与估计价值相关，但概念不同。

在具有流动性的市场，价格反映参与者信息、信念和交易指令的相互作用，新信息可能由此纳入价格。然而，市场价格并非不依赖任何假设的基本价值度量。投资者可能在信息、信念、约束、税收状况、融资成本和风险敞口上不同。价格也可能受到流动性需求、杠杆、做市能力、被迫交易和暂时性订单失衡影响。

\paragraph{作为估值基准的市场价格。}
市场价格为估值提供外部基准。估值与观察市场价格存在显著差异，可能反映现金流预期、风险判断、贴现率假设、遗漏因素或市场错误定价的差别。这种差异本身不能证明估值或市场价格哪一个错误，而是说明需要进一步审查底层假设和分歧来源。

无套利对相对价格施加更强约束。如果两项交易资产在同一未来日期、每一种可能状态下的支付完全相同，且不存在相关市场摩擦，它们在时点 0 的价格就必须相同，否则便存在套利机会。这一原则称为一价定律（law of one price），连接具有等价支付的证券价格，是固定收益和衍生工具估值的基础。

\paragraph{复制受限时的估值。}
如果某项请求权无法用交易中的证券紧密复制，估值就更依赖现金流预测和必要报酬率假设。例如，非上市企业和长期投资项目可能没有接近的市场对应物。尽管如此，可比企业与交易请求权的价格，以及利率、信用利差和期权价格，仍提供有用估值输入。

\begin{figure}[H]
\centering

\begin{tikzpicture}[
  >=Latex,
  box/.style={
    draw,
    rounded corners=2pt,
    align=center,
    minimum height=0.9cm,
    inner xsep=12pt,
    inner ysep=5pt
  },
  flow/.style={
    -Latex,
    line width=0.8pt,
    FinanceBlue
  }
]

% Top nodes
\node[box] (info) at (-2.0,2.2)
{信息与\\预期};

\node[box] (contract) at (2.0,2.2)
{契约规定的\\现金流权利};

% Central sequence
\node[box] (valuation) at (0,0.3)
{价值评估};

\node[box] (trading) at (0,-1.6)
{交易与价格发现};

\node[box] (price) at (0,-3.5)
{观察到的市场价格};

% Main flows: 90-degree turns
\coordinate (val-left) at
  ($(valuation.north west)!0.09!(valuation.north east)$);
\draw[flow]
  (info.south) -- (val-left);

\coordinate (val-right) at
  ($(valuation.north west)!0.91!(valuation.north east)$);

\draw[flow]
  (contract.south) -- (val-right);


% Vertical flows
\draw[flow]
  (valuation.south) -- (trading.north);

\draw[flow]
  (trading.south) -- (price.north);

% Feedback from observed market price to valuation
\draw[flow]
  (price.west)
  -- ++(-2.0,0)
  -- node[left, midway, align=center, xshift=-3pt]
  {市场\\证据}
     ($(valuation.west)+(-2.0,0)$)
  -- (valuation.west);

\end{tikzpicture}

\caption{信息、契约权利、估值与市场价格}
\label{fig:ch5-value-formation}

\end{figure}

图~\ref{fig:ch5-value-formation} 概括了这些互动。信息影响对未来现金流和风险的预期，契约界定每项请求权所附的现金流权利。这些因素进入估值，估值又为交易决策提供依据。交易和价格发现形成观察到的市场价格，为评估和修订估值提供市场证据。

图~\ref{fig:ch5-value-formation} 总结了上述相互关系。信息塑造未来现金流与风险预期，契约条款界定每项请求权所附的现金流权利。估值结合这些输入，并为交易决策提供依据。交易和价格发现产生市场价格，后者又为评估和修订估值提供证据。

\begin{tcolorbox}[
  enhanced,
  breakable,
  colback=white,
  colframe=FinanceBlue,
  boxrule=0.7pt,
  sharp corners,
  borderline west={2.2pt}{0pt}{FinanceBlue},
  title=\textbf{金融视角：估值以假设为条件},
  coltitle=white,
  colbacktitle=FinanceBlue,
  fonttitle=\bfseries,
  before skip=1.2em,
  after skip=1.2em
]

估值以分析采用的契约现金流权利、可用信息、时间假设和风险处理为条件。这些输入变化时，估值可能改变。稳健估值明确说明假设并一致应用，使结果能够被评估，并随信息变化而修订。

\end{tcolorbox}

%========================================================
% Summary
%========================================================

\TailModule{summary}{本章小结}

金融估值将未来现金流转换为当前价值。估值对象是金融请求权，因此首先要识别其现金流和契约权利。支付优先顺序、抵押品、转换权及其他条款，会使基于同一底层资产的请求权具有不同价值。

现金流时间很重要，因为不同日期的现金流必须转换到共同日期的价值。风险也影响价值：预期支付相同的两项请求权，如果支付在未来状态间分布不同，价值也可能不同。必要报酬率是一种纳入风险的方式，而状态价格估值则明确呈现价值对未来状态的依赖。

现金流贴现估值将请求权特定现金流预测与适当贴现率结合。可靠性取决于现金流与必要报酬率的一致处理、现实预测，以及经济上合理的终值假设。相对估值使用可比资产的观察价格，但估值倍数反映增长、风险、盈利能力、分配及其他基本面的差异。或有请求权估值则利用状态价格与无套利关系，对状态依存支付估值。

信息、契约和估值紧密相连。信息影响未来现金流和风险预期，契约界定每项请求权所附权利。交易与价格发现形成可观察市场价格，为评估和修订估值提供证据。估值以分析采用的信息、契约权利、时间假设和风险处理为条件，其可靠性取决于内部一致性、经济合理性，以及与可用市场证据的关系。

%========================================================
% Key Terms
%========================================================

\TailModule{keyterms}{关键术语}

金融估值（Financial valuation）；金融请求权（financial claim）；现金流权利（cash-flow right）；现值（present value）；贴现因子（discount factor）；必要报酬率（required return）；无风险利率（risk-free rate）；风险溢价（risk premium）；预期现金流（expected cash flow）；现金流贴现估值（discounted cash-flow valuation）；净现值（net present value）；终值（terminal value）；增长型永续年金（growing perpetuity）；敏感性分析（sensitivity analysis）；相对估值（relative valuation）；可比企业（comparable firm）；市场倍数（market multiple）；市盈率（price--earnings ratio）；企业价值（enterprise value）；或有请求权（contingent claim）；状态依存支付（state-contingent payoff）；状态价格（state price）；一价定律（law of one price）；无套利（no-arbitrage）；价格发现（price discovery）；市场价格。

%========================================================
% Financial AI Lab
%========================================================（market price

%========================================================
% Financial AI Lab
%========================================================）。

\TailModule{aiprompt}{金融 AI 实验室}

\textbf{实验：审查 AI 生成的估值}

选择一家上市公司，请 AI 系统使用公开信息生成简化的现金流贴现估值。要求明确估值对象、预测期、现金流指标、增长假设、终值方法和贴现率。

分五个阶段审查估值。

\begin{enumerate}

\item \textbf{识别请求权。}
确定估值针对股权还是整个企业，并检查现金流和贴现率是否与该请求权保持一致。

\item \textbf{核实计算。}
独立复算现值，检查单位、现金流时间、复利及终值处理。

\item \textbf{评价经济假设。}
考察现金流预测是否具有经济合理性、长期增长能否合理持续、再投资是否与增长一致，以及贴现率是否适合请求权风险。

\item \textbf{进行敏感性分析。}
识别估计价值特别敏感的两项关键假设，报告这些假设变化如何影响估值。

\item \textbf{与市场证据比较。}
将结果与观察到的市场价格及一个简单相对估值指标比较，确保比较对象是同一请求权。不要认为 AI 估值或市场价格必然正确，而要解释能够说明差异的假设或市场因素。

\end{enumerate}

最终提交应识别 AI 估值中至少一个仅看最终估值数字无法发现的问题，可以涉及计算、经济假设、内部一致性或市场证据核实。

%========================================================
% Applied Problems
%========================================================

\TailModule{appprobs}{应用题}

\begin{enumerate}

\item \textbf{时间与现值。}

某项请求权三年后支付人民币 5,000 元，相关年贴现率为 6\%。

\begin{enumerate}

\item 计算该请求权的现值。

\item 如果改为六年后支付，重新计算现值。

\item 解释为什么贴现率为正时，推迟支付会降低现值。

\end{enumerate}


\item \textbf{净现值与必要报酬率。}

某项目初始投资为人民币 1,200 万元，预期第 1 至第 4 年末现金流依次为 300 万元、400 万元、500 万元和 400 万元。

\begin{enumerate}

\item 按 7\% 必要报酬率计算项目净现值。

\item 按 11\% 必要报酬率重新计算净现值。

\item 解释为什么预期现金流不变，项目价值却变化。

\end{enumerate}


\item \textbf{终值敏感性。}

某企业预计第 6 年，即显式预测期后的第一年，产生人民币 2,000 万元现金流。从第 7 年起，现金流预计永久按恒定年增长率 \(g\) 增长。假设所有情景中第 6 年现金流保持不变。

\begin{enumerate}

\item 若 $r=9\%$、$g=3\%$，计算第 5 年末的终值。

\item 保持 $r$ 不变，若 $g=4\%$，重新计算终值。

\item 若 $r=8\%$、$g=3\%$，重新计算终值。

\item 比较第（b）和（c）问结果。解释为什么在给定假设下，$g$ 上升一个百分点，与 $r$ 下降一个百分点，会产生相同终值。

\end{enumerate}


\item \textbf{相对估值与一致性。}

企业 A 的企业价值为人民币 80 亿元，EBITDA 为 8 亿元。企业 B 的 EBITDA 为 5 亿元。

\begin{enumerate}

\item 计算企业 A 的 EV/EBITDA 倍数。

\item 将这一倍数应用于企业 B，估计其隐含企业价值。

\item 假设企业 B 净债务为人民币 12 亿元，计算其隐含股权价值。

\item 解释为什么企业价值倍数不应直接用于仅归属于股东的利润指标。

\end{enumerate}


\item \textbf{状态价格估值。}

未来存在两种等概率状态。请求权 A 在强劲状态支付 140、疲弱状态支付 60；请求权 B 在强劲状态支付 80、疲弱状态支付 120。

\begin{enumerate}

\item 计算两项请求权各自的预期支付。

\item 假设状态价格为 $q_S=0.38$、$q_W=0.56$，计算两项请求权在时点 0 的价值。

\item 比较两项请求权的预期支付与时点 0 价值。

\item 解释为什么预期支付相同，时点 0 的价值却未必相同。

\end{enumerate}


\item \textbf{价值的契约分配。}

某企业下一期现金流为 150 或 70，两种状态等概率。其债务承诺偿还 100。

\begin{enumerate}

\item 计算每种状态下债务和股权的支付。

\item 验证每种状态下债务与股权支付之和等于企业总现金流。

\item 使用状态价格 $q_H=0.40$、$q_L=0.55$，计算债务、股权及企业整体的当前价值。

\item 验证债务和股权价值之和等于企业价值。

\end{enumerate}

\end{enumerate}

%========================================================
% Further Reading
%========================================================

\TailModule{furtherread}{延伸阅读}

\begin{itemize}

\item Fisher, Irving (1930), 
\emph{The Theory of Interest: As Determined by Impatience to Spend Income and Opportunity to Invest It}（《利息理论：由支出收入的时间偏好与投资机会共同决定》），
Macmillan.

\item Williams, John Burr (1938), 
\emph{The Theory of Investment Value}（《投资价值理论》），
Harvard University Press.

\item Modigliani, Franco, and Merton H. Miller (1958), 
“资本成本、公司融资与投资理论”（``The Cost of Capital, Corporation Finance and the Theory of Investment''），
\emph{The American Economic Review}, 48(3), 261--297.

\item Sharpe, William F. (1964), 
“资本资产价格：风险条件下的市场均衡理论”（``Capital Asset Prices: A Theory of Market Equilibrium under Conditions of Risk''），
\emph{The Journal of Finance}, 19(3), 425--442.

\item Damodaran, Aswath (2012), 
\emph{Investment Valuation: Tools and Techniques for Determining the Value of Any Asset}（《投资估值：评估任何资产价值的工具与方法》），
第 3 版，Wiley。

\end{itemize}

\ChapterTail{summary,keyterms,aiprompt,appprobs,furtherread}

\part{金融市场}
\chapter{利率与资本配置}
\label{chap:Interest_Rates_and_Capital_Allocation}

\begin{openingcase}{哪一种利率才是相关利率？}

一家企业正在考虑一项为期一年的投资，今天需要投入人民币~1000~万元，预计一年后产生~1060~万元收入。初看之下，决策似乎很简单：将项目~6\%~的预期回报率与“利率”进行比较。

但应当使用哪一种利率？

银行可能报出一种借款利率，而同期限政府证券提供另一种利率。利率还可能随期限、通胀预期、信用风险和流动性而变化。因而，并不存在适用于所有金融决策的单一市场利率。

同一个项目，在一种必要报酬率下可能具有吸引力，在另一种必要报酬率下却可能不值得投资。理解市场利率为何不同，以及这些差异如何影响融资和投资决策，是理解金融市场如何配置资本的关键。

\end{openingcase}

\begin{tcolorbox}[
  enhanced,
  colback=FinanceBlue!6,
  colframe=FinanceBlue,
  boxrule=0.6pt,
  sharp corners,
  left=8pt,
  right=8pt,
  top=7pt,
  bottom=7pt,
  before skip=1em,
  after skip=1.6em,
  title=\textbf{核心问题},
  coltitle=white,
  colbacktitle=FinanceBlue,
  fonttitle=\bfseries,
  attach boxed title to top left={xshift=6pt,yshift=-2mm},
  boxed title style={
    boxrule=0pt,
    sharp corners,
    left=5pt,
    right=5pt,
    top=2pt,
    bottom=2pt
  }
]
\begin{enumerate}[
  leftmargin=*,
  label=\textcolor{FinanceBlue}{\textbf{\arabic*.}},
  itemsep=0.55em,
  topsep=0.2em
]

\item 为什么利率是跨期交换的价格，而不仅仅是借款成本？

\item 为什么不同期限、借款人和金融工具的利率存在差异？

\item 利率如何影响储蓄、借款、投资和资本配置？

\end{enumerate}
\end{tcolorbox}

\section{利率作为跨期交换的价格}

利率是当期购买力与未来购买力交换的价格。它决定推迟使用资金所获得的回报，以及在未来收入实现之前取得资金的成本。

\subsection{跨期选择与资金机会成本}

对于一期投资，今天的一单位购买力在下一期增长为~\(1+r\)。等价地，下一期收到一单位购买力的现值为

\[
PV_0=\frac{1}{1+r}.
\]

这一关系从未来回报和现值两个角度表达了同一种跨期价格。

储蓄、借款和投资决策都需要比较资金的不同用途或来源。储蓄者考虑推迟消费所能获得的回报，借款人评估不同融资安排的成本，投资者则将项目预期回报与其他可行投资的回报进行比较。相关机会成本取决于可获得的最佳替代选择，同时需要考虑风险、期限和流动性。

假设投资者购买一种一年期证券可获得~4\%~的回报率，其风险与拟投资项目相当。如果该项目的预期回报率仅为~3\%，那么在风险和期限相当的情况下，替代证券提供了更高的回报。

这一机会成本原则将市场利率与资本配置直接联系起来。设某项目今天需要投资~$I_0$，一期后产生确定的支付~$CF_1$。如果资金机会成本为~$r$，则项目在满足下式时创造价值：

\begin{equation*}
\frac{CF_1}{1+r}>I_0.
\label{eq:ch6-one-period-investment-condition}
\end{equation*}

等价地，

\begin{equation*}
\frac{CF_1-I_0}{I_0}>r.
\label{eq:ch6-return-hurdle-simple}
\end{equation*}

相关机会成本是风险和期限相近的可比替代投资的回报率。风险或期限的差异相应意味着必要报酬率的差异。

在实际市场中，借入与贷出资金的利率可能不同，不同借款人面临的信用利差也可能不同。这些差异影响市场参与者面临的机会成本。

\begin{example}{机会成本与一年期投资}{opportunity-cost-one-year-investment}

假设一家企业可以将人民币~1000~万元投入一年期项目，确定获得~1050~万元，项目回报率为~5\%。

如果相关机会成本为~3\%，支付的现值为

$$
PV=\frac{10.5}{1.03}=10.19
$$
百万元，项目的净现值为正。如果机会成本升至~6\%，

$$
PV=\frac{10.5}{1.06}=9.91
$$
百万元，项目的净现值为负。

现金流没有变化；机会成本提高改变了投资决策。

\end{example}

\subsection{利率报价、复利与有效利率}

利率描述购买力的跨期交换，但给定期限内积累的金额取决于利息的计算方式。金融市场的报价规定利率适用的期间及相应的复利计息惯例。

设~$r$~表示每期利率。将~$P_0$~投资一期，其终值为

\begin{equation}
P_1=P_0(1+r).
\label{eq:ch6-one-period-growth}
\end{equation}

如果投资持续~$T$~期，每期复利计息一次，则

\begin{equation}
P_T=P_0(1+r)^T.
\label{eq:ch6-compound-growth}
\end{equation}

在复利计息下，每期利息计入本金，并在以后各期继续产生利息。相比之下，单利仅以初始本金计算利息：

\begin{equation}
P_T=P_0(1+rT).
\label{eq:ch6-simple-interest}
\end{equation}

当利率为正时，单利与复利之间的差距随投资期限延长而扩大。例如，将人民币~100~元按~5\%~的利率投资三年，单利计息下增长为

$$
100(1+0.05\times 3)=115
$$
（单利计息）；按年复利计息则增长为

$$
100(1.05)^3=115.76
$$
因此，复利计息惯例影响利率报价所代表的有效回报率。

\paragraph{名义年利率与复利计息频率。}

许多市场利率以年率报价，即使实际复利计息频率高于每年一次。设~$j$~为每年复利计息~$m$~次的名义年利率，则每期利率为

\begin{equation}
r=\frac{j}{m}.
\end{equation}

相应的有效年利率记为~$EAR$，其表达式为

\begin{equation}
EAR=
\left(1+\frac{j}{m}\right)^m-1.
\label{eq:ch6-ear}
\end{equation}

有效年利率衡量名义年利率及其复利计息频率所隐含的一年期回报率。在名义年利率给定时，复利计息越频繁，有效年利率越高，如表~\ref{tab:ch6-compounding-frequency}~所示。

\begin{table}[H]
\centering
\caption{名义年利率为~6\%~时不同复利计息频率的比较}
\label{tab:ch6-compounding-frequency}

\begin{tabular}{@{}l c c@{}}
\toprule
\textbf{复利计息方式} & \textbf{每年计息期数} & \textbf{有效年利率} \\
\midrule
每年 & 1 & 6.000\% \\
每半年 & 2 & 6.090\% \\
每季度 & 4 & 6.136\% \\
每月 & 12 & 6.168\% \\
\bottomrule
\end{tabular}
\end{table}

有效年利率为比较不同复利计息频率下的报价提供了统一基础。
\paragraph{连续复利。}

随着复利计息频率增加，离散复利逐渐趋近连续复利。设~$r_c$~为连续复利年利率，在~$T$~年内，

\begin{equation}
P_T=P_0 e^{r_c T}.
\label{eq:ch6-continuous-compounding}
\end{equation}

连续复利利率与有效年利率的关系为

\begin{equation}
EAR=e^{r_c}-1,
\label{eq:ch6-continuous-to-ear}
\end{equation}
等价地，

\begin{equation}
r_c=\ln(1+EAR).
\label{eq:ch6-ear-to-continuous}
\end{equation}

连续复利便于比较不同期限的利率，广泛用于金融建模。

\paragraph{不足一年的期间利率。}

利率也可能需要表示为不足一年的期间利率。如果~$EAR$~为有效年利率，则一年中~$\tau$~这一时间比例对应的复利利率为

\begin{equation}
r_{\tau}=(1+EAR)^{\tau}-1.
\label{eq:ch6-fractional-period-rate}
\end{equation}

这一换算假定不足一年的期间也采用复利计息。实践中，一些市场采用不同的报价与日计数惯例，货币市场工具尤其如此。

\begin{tcolorbox}[
enhanced,
breakable,
colback=white,
colframe=FinanceBlue,
boxrule=0.7pt,
sharp corners,
borderline west={2.2pt}{0pt}{FinanceBlue},
title=\textbf{金融视角：在可比基础上比较},
coltitle=white,
colbacktitle=FinanceBlue,
fonttitle=\bfseries,
before skip=1.2em,
after skip=1.2em
]

比较利率之前，应将其换算到统一基础。期限、复利计息频率、日计数惯例，以及名义报价与有效报价的差异，都可能使相同的报价百分比对应不同的有效回报率。

\end{tcolorbox}

\subsection{名义利率与实际利率}

名义利率以货币金额衡量资金回报，实际利率则以购买力衡量回报。当价格水平变化时，两者便出现区别：部分名义回报可能仅反映价格上涨，而非未来金额所能购买的商品和服务数量增加。

假设一单位货币在一年后增长为~$1+i$，其中~$i$~为名义利率。如果同期价格水平上涨~$\pi$，则相对于今天的价格水平，未来金额的购买力增长为

\begin{equation}
1+r
=
\frac{1+i}{1+\pi},
\label{eq:ch6-fisher-exact}
\end{equation}
其中，$r$~为已实现实际利率。

整理这一关系，可得精确的费雪方程：

\begin{equation}
i=r+\pi+r\pi.
\label{eq:ch6-fisher-expanded}
\end{equation}

当~$r$~和~$\pi$~不大时，交叉乘积项~$r\pi$~相对于一次项较小，从而得到近似式

\begin{equation}
i\approx r+\pi.
\label{eq:ch6-fisher-approx}
\end{equation}

在通胀尚未实现时作决策，相关的是预期通胀，而非已实现通胀。如果~$\pi^e$~表示预期通胀率，$r^e$~表示预期实际利率，则近似的费雪关系为

\begin{equation}
i\approx r^e+\pi^e.
\label{eq:ch6-fisher-expected}
\end{equation}

预期通胀与已实现通胀的差异，会使固定名义请求权的已实现实际回报偏离预期实际回报。假设贷款人发放一年期固定利率贷款，名义利率为~6\%，当时预期通胀率为~2\%，则预期实际回报率约为~4\%。如果最终实现的通胀率为~5\%，已实现实际回报率就降至约~1\%。

在固定名义契约下，非预期通胀由此在借款人与贷款人之间重新分配购买力。通胀高于预期会降低贷款人的已实现实际回报，并使借款人相对于原有条款受益；通胀低于预期则产生相反效果。

名义利率与实际利率服务于不同的分析目的。金融契约通常采用名义金额表述：借款人支付名义金额，债券承诺名义票息，政策利率也以名义利率报价。实际利率则以购买力衡量回报与融资成本，尤其适用于从实际量角度评估储蓄、投资及资本配置决策。名义利率与预期通胀的关系，将名义金融契约与其实际经济影响联系起来。

\section{利率的决定因素}

利率是受多种相互作用力量影响的市场价格。这些力量包括资金供求、通胀预期、货币与流动性条件，以及风险和期限差异。它们的相对重要性因市场、工具和期限而异。

\subsection{储蓄、投资与资金需求}

一个简单的分析起点，是将利率视为平衡资金供给与资金需求的价格。储蓄者通过推迟当期支出提供资金，借款人则寻求资金，用于投资和其他支出。

设资金供给为

\begin{equation}
S=S(r),
\qquad
S'(r)>0,
\label{eq:ch6-saving-supply}
\end{equation}
资金需求为

\begin{equation}
D=D(r),
\qquad
D'(r)<0.
\label{eq:ch6-investment-demand}
\end{equation}

均衡利率~$r^{*}$~满足

\begin{equation}
S(r^{*})=D(r^{*}).
\label{eq:ch6-funds-equilibrium}
\end{equation}

向上倾斜的供给曲线反映：利率提高通常使推迟当期支出更具吸引力。向下倾斜的需求曲线反映：利率提高通常会减少预期回报超过资金成本的投资机会。

\begin{figure}[H]
\centering
\begin{tikzpicture}[
  >=Latex,
  axis/.style={-Latex, line width=0.8pt},
  curve/.style={line width=1.0pt},
  guide/.style={dashed, gray, line width=0.6pt}
]

\draw[axis] (0,0) -- (7.2,0) node[right] {资金数量};
\draw[axis] (0,0) -- (0,5.5) node[above] {利率};

\draw[curve, FinanceBlue] (1.0,1.0) -- (6.2,4.5);
\draw[curve] (1.0,4.6) -- (6.2,1.1);

\node[FinanceBlue, right] at (6.2,4.5) {$S(r)$};
\node[right] at (6.2,1.1) {$D(r)$};

\coordinate (E) at (3.674,2.80);
\fill (E) circle (2pt);
\node[above, yshift=3pt] at (E) {$E$};

\draw[guide] (E) -- (0,2.80) node[left] {$r^{*}$};
\draw[guide] (E) -- (3.674,0) node[below] {$F^{*}$};

\end{tikzpicture}
\caption{简化的资金市场}
\label{fig:ch6-funds-market}
\end{figure}

该模型为比较静态分析提供了简单框架。预期生产率提高或其他有利可图的投资机会增加，会使资金需求曲线向右移动。在其他条件不变时，均衡利率与资金数量均上升。各利率水平下储蓄意愿增强，会使资金供给曲线向右移动，降低均衡利率，并增加均衡资金数量。

这些结果描述的是简化模型中的比较静态关系，并非市场利率的完整理论。实际金融体系中的融资条件还反映货币政策、银行资产负债表、资本流动和制度约束。资金供给不应被理解为可供放贷的、事先形成的固定储蓄存量。这一框架仅以简单模型说明资金供求变动如何影响利率与融资规模。

\subsection{通胀预期与货币条件}

预期通胀通过改变维持购买力所需的名义回报，影响名义利率。费雪关系表明，在预期实际利率给定时，预期通胀越高，对应的名义利率越高。

这种变化未必一比一发生。预期实际利率也可能变化，风险溢价可能调整，货币政策也可能对通胀预期作出反应。费雪关系识别的是名义利率的一项组成因素，并不能完整解释观察到的利率变动。

货币条件对短期名义利率尤其重要。中央银行通常通过政策利率，以及准备金和短期融资的供给条件，影响短期市场利率。政策立场变化还可能通过对未来短期利率、通胀、经济活动和风险溢价的预期，影响较长期利率。

政策利率降低可能减少短期融资成本。如果市场预期较低的政策利率将持续，较长期收益率也可能下降。不过，如果通胀预期上升，长期名义收益率的降幅可能较小，甚至可能上升。即使政策利率下降，风险溢价或流动性溢价扩大，也可能提高部分私人部门借款人的借款成本。

\subsection{风险结构与利差}

在其他方面可比、期限相同的工具之间，利率差异构成利率的风险结构。信用风险与市场流动性是这些差异的主要决定因素，契约、税收和制度因素也会影响必要报酬率。

\paragraph{信用风险。}

信用风险来自承诺支付可能无法足额或按时履行。预期信用损失越大，或者偿还的不确定性越高，投资者要求的回报通常就越高。在债务市场中，相对于可比低风险基准的利率差额，通常称为信用利差。

\begin{example}{信用风险与必要报酬率} {credit-risk-required-return}

假设一种一年期低风险政府证券的回报率为~3\%，另一种其他特征可比的一年期企业债务工具，其必要报酬率为~5.5\%。两者之差为

\[
5.5\%-3.0\%=2.5\%
\]
该差额代表企业债务工具相对于基准的信用利差。

较高的必要报酬率反映企业承诺支付可能无法足额或按时履行。利差未必等于预期信用损失，因为投资者还可能要求对偿还不确定性给予补偿。

\end{example}

风险补偿不仅取决于预期损失，还取决于损失发生的状态。两种一年期工具可能具有相同的预期支付，但损失集中的状态不同。一种工具若主要在衰退期提供较低支付，那么在投资者收入及其他资产价值也可能较低时，它提供的保护就较少。相较于主要在经济景气时支付较低的工具，投资者可能要求更高的补偿。

其他风险只要改变预期现金流、偿还风险，或投资者对不利状态的风险暴露，也可能影响必要报酬率。

\begin{table}[htbp]
\centering
\caption{影响可比工具必要报酬率的因素}
\label{tab:ch6-rate-differences}

\begin{tabularx}{\textwidth}{@{}p{3.2cm}Y Y@{}}
\toprule
\textbf{因素} &
\textbf{经济原因} &
\textbf{对必要报酬率的影响} \\
\midrule

信用风险 & 承诺支付可能延迟、减少或无法足额履行 & 预期信用损失及承担违约风险所需的补偿可能提高必要报酬率 \\

不利状态下的风险暴露 & 当投资者收入及其他资产价值较低时，支付也可能较低 & 此类风险暴露在不利条件下提供的保护较少，可能提高必要补偿 \\

市场流动性 & 若不作较大价格让步，工具可能难以迅速出售 & 较低的流动性可能提高必要报酬率 \\

契约特征 & 嵌入式期权及其他条款改变未来现金流的时点与分布 & 对必要报酬率的影响取决于这些特征如何改变各方支付与风险 \\

税收与制度待遇 & 税收规则、抵押品资格和监管约束可能影响持有工具的净收益 & 税后回报和制度性收益的差异可能导致必要报酬率不同 \\

\bottomrule
\end{tabularx}
\end{table}

\paragraph{市场流动性。}

市场流动性影响金融工具在不作较大价格让步的情况下出售的难易程度。投资者通常重视较高的流动性，而流动性较低的工具可能需要提供高于其他方面可比工具的回报。

金融压力时期，流动性可能恶化。交易商可能减少库存，融资条件可能收紧，买方可能要求更大的价格让步。即使预期信用损失变化不大，流动性恶化也可能扩大利差。

\paragraph{基准利率与利差。}

市场利率常与基准利率比较，后者为期限相近的工具提供参照。某种工具利率与基准利率之差称为其利差。对于期限为~$T$~的工具~$i$，

\begin{equation}
r_{i,T}=r_{b,T}+s_{i,T},
\label{eq:ch6-benchmark-spread}
\end{equation}
其中，$r_{b,T}$~为期限~$T$~的基准利率，$s_{i,T}$~为工具~$i$~相对于该基准的利差。该工具利率的变化可写为

\begin{equation}
\Delta r_{i,T}=\Delta r_{b,T}+\Delta s_{i,T}.
\label{eq:ch6-benchmark-spread-change}
\end{equation}

如果许多同期限工具的利率共同变动，基准利率变化可能解释其中相当一部分。货币条件和通胀预期可能促成此类共同变化。如果基准利率变化不大，而特定工具利率上升，利差扩大可能反映信用风险增加、市场流动性下降，或其他工具特有因素。

观察到的利差不能揭示单一的深层原因。信用条件、市场流动性、契约特征、税收待遇、监管和投资者需求，都可能影响相对于基准利率的利差。

图~\ref{fig:ch6-china-10y-yields}~使用中国债券市场的十年期收益率说明这些模式。各序列期限相同，共同变动可能反映整体市场条件的变化，而收益率水平差异可能反映信用风险、流动性及其他工具特有特征。

\begin{figure}[H]
\centering

\IfFileExists{ChinaBond_10Y_Yields_2016_2026.pdf}{\includegraphics[width=0.9\textwidth]{ChinaBond_10Y_Yields_2016_2026.pdf}}{\fbox{\parbox[c][32mm][c]{0.85\textwidth}{\centering 原稿外部图文件尚未提供\\[4pt]\small\texttt{\detokenize{ChinaBond_10Y_Yields_2016_2026.pdf}}}}}

\caption{不同债券类别的中债十年期收益率，2016--2026~年}
\label{fig:ch6-china-10y-yields}

\vspace{0.15em}

\begin{minipage}{0.9\textwidth}
\footnotesize
\noindent\textit{资料来源：中央国债登记结算有限责任公司（中债）。}
\end{minipage}

\end{figure}

跨期限比较引入了利率差异的另一个维度，即利率的期限结构。

\section{利率期限结构}

利率期限结构描述其他方面可比的工具，其利率与期限之间的关系。其图形表示为收益率曲线。

\subsection{即期利率、远期利率与收益率曲线}

即期利率是当前适用于某项零息请求权的年化利率，该请求权仅在特定未来日期支付一次。远期利率是当前隐含的、适用于未来某一期间借入或贷出资金的利率。

设~$s_t$~表示无违约零息请求权的年化~$t$~期即期利率。按年复利计息时，$t$~时点支付一单位的现值为

\begin{equation}
P(0,t)=\frac{1}{(1+s_t)^t}.
\label{eq:ch6-spot-discount-factor}
\end{equation}

例如，假设一年期即期利率为~3\%，两年期即期利率为~4\%，相应的贴现因子为

\[
P(0,1)=\frac{1}{1.03},
\qquad
P(0,2)=\frac{1}{(1.04)^2}.
\]

为推导从今天起一年后开始的一年期远期利率，记该利率为~$f_{1,2}$。在无套利条件下，按两年期即期利率投资两年，与先按一年期即期利率投资一年、再按事先锁定的远期利率投资一年，应产生相同终值。因此，

\begin{equation}
(1+s_2)^2
=
(1+s_1)(1+f_{1,2}),
\label{eq:ch6-forward-from-spot}
\end{equation}
其中，$s_1$~和~$s_2$~分别为一年期和两年期即期利率。因此，

\begin{equation}
f_{1,2}
=
\frac{(1+s_2)^2}{1+s_1}-1.
\label{eq:ch6-forward-rate}
\end{equation}

代入例中的利率，得到隐含远期利率约为~5.01\%。远期利率并非两种即期利率的算术差。

更一般地，设~$s_m$~和~$s_n$~分别表示~$m$~年期和~$n$~年期即期利率，$f_{m,n}$~表示第~$m$~年至第~$n$~年的年化远期利率，其中~$n>m$。按年复利计息时，

\begin{equation}
f_{m,n}
=
\left[
\frac{(1+s_n)^n}{(1+s_m)^m}
\right]^{\frac{1}{n-m}}
-1.
\label{eq:ch6-general-forward-rate}
\end{equation}

将零息即期利率与其期限绘制在图中，可得即期利率曲线，即由零息利率构成的收益率曲线。在更广泛的市场用法中，收益率曲线表示可比债务证券的收益率与期限之间的关系。附息债券的到期收益率未必等于即期利率。

\begin{example}{构造即期利率曲线}
{spot-rate-curve}

假设一年、两年、三年、五年和十年期的年化即期利率如下：

\begin{center}
\begin{tabular}{lccccc}
\toprule
\textbf{期限（年）} & 1 & 2 & 3 & 5 & 10 \\
\midrule
\textbf{即期利率（\%）} & 3.00 & 4.00 & 4.40 & 4.70 & 4.90 \\
\bottomrule
\end{tabular}
\end{center}

将这些利率与期限对应绘图，可得下图所示的即期利率曲线。

\begin{center}
\begin{tikzpicture}

\begin{axis}[
    width=0.88\linewidth,
    height=5.2cm,
    xmin=0,
    xmax=10.5,
    ymin=2.0,
    ymax=5.2,
    xtick={1,2,3,5,10},
    ytick={2.5,3.0,3.5,4.0,4.5,5.0},
    xlabel={剩余期限~$T$（年）},
    ylabel={即期利率~$s_T$（\%）},
    axis lines=box,
    tick align=outside,
    enlargelimits=false,
    grid=none,
]

\addplot[
    FinanceBlue,
    thick,
    smooth,
    mark=*,
    mark size=1.8pt
]
coordinates {
    (1,3.00)
    (2,4.00)
    (3,4.40)
    (5,4.70)
    (10,4.90)
};

\end{axis}

\end{tikzpicture}
\end{center}
\end{example}

\begin{figure}[H]
\centering

\IfFileExists{China_Government_Yield_Curves_Typical.pdf}{\includegraphics[width=1.0\textwidth]{China_Government_Yield_Curves_Typical.pdf}}{\fbox{\parbox[c][32mm][c]{0.85\textwidth}{\centering 原稿外部图文件尚未提供\\[4pt]\small\texttt{\detokenize{China_Government_Yield_Curves_Typical.pdf}}}}}

\caption{中国国债收益率曲线示例}
\label{fig:ch6-china-gov-yield-curves}

\vspace{0.15em}
\begin{minipage}{1.0\textwidth}
\footnotesize
\noindent\textit{资料来源：中央国债登记结算有限责任公司（中债）；作者计算。}
\end{minipage}

\end{figure}

即期利率曲线可能向上倾斜、近似水平、倒挂，或呈现更复杂的形态。市场收益率曲线也有类似形态。图~\ref{fig:ch6-china-gov-yield-curves}~列示了中国国债市场的若干例子。

传统期限结构理论通过未来短期利率预期、不同期限市场之间的分割，以及投资者偏好的期限区间，解释这些形态。现代期限结构分析通常将较长期利率分解为预期短期利率路径与时变期限溢价。

\subsection{预期假说}
纯预期假说认为，在投资者的投资组合中，不同期限债券是完全替代品。投资者依据预期回报选择期限，不要求为持有较长期债券获得额外期限溢价。 

考虑两年投资期限，今天投入一单位资金。投资者可以采用两种策略。第一种是按当前两年期即期利率~$s_2$~投资两年；第二种是按当前一年期即期利率~$s_1$~投资一年，再按届时的一年期短期利率投资第二年。设~$\widetilde{r}_{1,2}$~表示第~1~年至第~2~年之间届时适用的一年期短期利率。在时点~0，该利率尚未确定，其当前预期为~$\mathrm{E}_0[\widetilde{r}_{1,2}]$。

\begin{figure}[H]
\centering

\begin{tikzpicture}[
    >=Latex,
    timeline/.style={line width=0.8pt},
    lab/.style={font=\small}
]

% 策略~A
\node[lab, anchor=east] at (-0.3,2.2) {策略~A};

\draw[timeline,->] (0,2.2) -- (7,2.2);

\foreach \x/\t in {0/0,3.5/1,7/2} {
    \fill (\x,2.2) circle (1.5pt);
    \node[lab, below=3pt] at (\x,2.2) {$t=\t$};
}

\node[lab, above=5pt] at (3.5,2.4)
    {按~$s_2$~投资两年};

\node[lab, above=5pt, anchor=east] at (7,2.25)
    {$(1+s_2)^2$};

% 策略~B
\node[lab, anchor=east] at (-0.3,-0.0) {策略~B};

\draw[timeline,->] (0,-0.0) -- (7,-0.0);

\foreach \x/\t in {0/0,3.5/1,7/2} {
    \fill (\x,-0.0) circle (1.5pt);
    \node[lab, below=3pt] at (\x,-0.0) {$t=\t$};
}

\node[lab, above=5pt] at (1.75,0.2)
    {按~$s_1$~投资};

\node[lab, above=5pt] at (5.25,0.2)
    {按~$\widetilde{r}_{1,2}$~再投资};

\node[lab, above=5pt, anchor=east] at (7,0.15)
    {$(1+s_1)(1+\widetilde{r}_{1,2})$};

\end{tikzpicture}

\caption{两年期限内的两种投资策略}
\label{fig:ch6-expectations-strategies}
\end{figure}

图~\ref{fig:ch6-expectations-strategies}~比较了两年期限内的两种策略，其终值分别为~$(1+s_2)^2$~和~$(1+s_1)(1+\widetilde{r}_{1,2})$。在纯预期假说下，不同期限债券是完全替代品。投资者因而对两种策略无差异，它们在相同投资期限内应提供相等的预期回报：

\begin{equation}
(1+s_2)^2=(1+s_1)
\left(1+\mathrm{E}_0[\widetilde{r}_{1,2}]\right).
\label{eq:ch6-expectations-gross}
\end{equation}

展开式~(\ref{eq:ch6-expectations-gross})，可得

\[
1+2s_2+s_2^2
=
1+s_1+\mathrm{E}_0[\widetilde{r}_{1,2}]
+s_1\mathrm{E}_0[\widetilde{r}_{1,2}].
\]

当利率不大时，二阶项较小。忽略这些项，可得

\begin{equation}
s_2
\approx
\frac{s_1+\mathrm{E}_0[\widetilde{r}_{1,2}]}{2}.
\label{eq:ch6-expectations-two-year}
\end{equation}

两年期即期利率近似等于当前一年期即期利率与第~1~年至第~2~年之间预期一年期短期利率的平均值。

这一假说也赋予远期利率一种解释。将式~(\ref{eq:ch6-expectations-gross})~与前面推导的远期利率关系结合，可得

\begin{equation}
f_{1,2}
=
\mathrm{E}_0[\widetilde{r}_{1,2}].
\label{eq:ch6-forward-expectation}
\end{equation}

在预期假说下，远期利率等于预期未来短期利率。

同样的推理可扩展至更长期限。$n$~年期即期利率与一系列一年期远期利率之间的关系为

\begin{equation}
(1+s_n)^n
=
(1+s_1)
\prod_{j=1}^{n-1}
(1+f_{j,j+1}).
\label{eq:ch6-spot-forward-general}
\end{equation}

利用
\[
f_{j,j+1}
=
\mathrm{E}_0[\widetilde{r}_{j,j+1}],
\]
上述关系变为

\begin{equation}
(1+s_n)^n
=
(1+s_1)
\prod_{j=1}^{n-1}
\left(1+\mathrm{E}_0[r_{j,j+1}]\right).
\label{eq:ch6-expectations-general-gross}
\end{equation}

当利率不大时，一阶近似给出

\begin{equation}
s_n
\approx
\frac{1}{n}
\left[
s_1
+
\sum_{j=1}^{n-1}
\mathrm{E}_0[r_{j,j+1}]
\right].
\label{eq:ch6-expectations-general}
\end{equation}

\begin{example}{在预期假说下解读期限结构} {expectations-term-structure}

假设当前一年期即期利率为~3\%，第~2~年至第~5~年的预期一年期短期利率依次为~3.5\%、4.0\%、4.5\%~和~5.0\%。在预期假说下，

\[
s_5
\approx
\frac{
3.0\%+3.5\%+4.0\%+4.5\%+5.0\%
}{5}
=
4.0\%.
\]

在预期假说下，五年期即期利率~\(s_5\approx4.0\%\)~高于当前一年期即期利率~\(s_1=3.0\%\)，反映预期未来短期利率的平均水平高于当前短期利率。

\end{example}

当不同期限证券并非完全替代品，且各期限区间的供求条件不同时，仅凭预期无法决定期限结构。

\subsection{市场分割与偏好期限}

投资者需求可能集中于特定期限区间。货币市场基金和管理短期流动性需求的企业往往持有短期工具；养老基金与寿险公司通常具有长期负债，因而对较长期资产形成需求。

市场分割理论将不同期限证券视为不完全替代品。各期限区间的供求可能影响相应利率。在供给不变时，长期债券需求增加会提高其价格、降低其收益率。因此，即使未来短期利率预期不变，长期利率也可能相对于短期利率下降。

偏好期限理论承认跨期限替代并不完全。投资者可能偏好特定期限区间，但预期补偿足够时，也愿意持有该区间之外的证券。期限偏好能够影响利率，却并不意味着不同期限市场完全分割。

关键区别在于跨期限替代程度：市场分割假定替代十分有限，而偏好期限理论允许投资者在补偿足够时转向其他期限。

供给条件也可能因期限而异。例如，政府债券发行期限结构的变化，会改变短期与长期证券的相对供给，并可能影响其相对利率。

\subsection{期限溢价}

现代期限结构分析允许较长期利率通过时变期限溢价，偏离预期未来短期利率路径。因此，远期利率未必等于预期未来短期利率。

与其他特征相近的短期限证券相比，长期限固定收益证券通常对市场利率变化更敏感。市场利率上升时，较远期固定现金流的现值下降更明显。这种更高的价格敏感性增加了投资者的利率风险暴露，是期限溢价的重要来源。

考虑期限溢价后，$n$~期即期利率可以近似表示为

\begin{equation}
s_n
\approx
\frac{1}{n}
\left[
s_1
+
\sum_{j=1}^{n-1}
\mathrm{E}_0[r_{j,j+1}]
\right]
+
TP_n,
\label{eq:ch6-term-premium-decomposition}
\end{equation}
其中，$TP_n$~为期限溢价，即较长期利率中无法由预期短期利率路径解释的部分。

期限溢价的符号与大小均可能变化。正的期限溢价表明，相较于反复投资短期限证券，投资者要求持有长期限证券获得额外预期回报。当长期限证券具有分散化或套期保值效益时，期限溢价可能较低，在某些情况下甚至可能为负。

向上倾斜的期限结构可能反映对未来短期利率上升的预期、正的期限溢价，或两者共同作用。仅凭斜率无法识别其相对贡献。

\begin{example}
{同一长期利率的不同分解}
{term-premium-sources}

例~\ref{ex:expectations-term-structure}~表明，五年期即期利率~$s_5\approx4.0\%$，可以与当前一年期即期利率~$s_1=3.0\%$，以及第~2~年至第~5~年依次为~$3.5\%$、$4.0\%$、$4.5\%$~和~$5.0\%$~的预期一年期短期利率相一致。在预期假说下，这一组合意味着~$TP_5\approx0$。

现保持~$s_1=3.0\%$~和~$s_5=4.0\%$~不变，但假设第~2~年至第~5~年的预期一年期短期利率均为~$3.0\%$。预期短期利率路径的平均值便为~3.0\%。式~(\ref{eq:ch6-term-premium-decomposition})~意味着

\[
TP_5
\approx
4.0\%-3.0\%
=
1.0\%.
\]

同样的五年期即期利率，此时反映正的期限溢价，而非对未来短期利率上升的预期。

\end{example}

长期利率变动可能反映预期未来短期利率或期限溢价的变化。由预期未来短期利率上升推动的利率上升，与主要由期限溢价提高推动的利率上升，具有不同的经济含义。期限溢价提高表明，投资者要求在预期未来短期利率之外，获得更多补偿，以持有长期限证券并承担相应利率风险。

\begin{tcolorbox}[
  enhanced,
  breakable,
  colback=white,
  colframe=FinanceBlue,
  boxrule=0.7pt,
  sharp corners,
  borderline west={2.2pt}{0pt}{FinanceBlue},
  title=\textbf{金融视角：解读收益率曲线},
  coltitle=white,
  colbacktitle=FinanceBlue,
  fonttitle=\bfseries,
  before skip=1.2em,
  after skip=1.2em
]

收益率曲线的斜率不应被解释为对未来短期利率的直接预测。长期利率还可能反映时变期限溢价与不同期限的供求条件。同一条观察到的收益率曲线，可以对应不同的预期、期限溢价和各期限供求组合。
\end{tcolorbox}

\section{资本配置与投资决策}

利率通过改变资金机会成本，影响储蓄、借款和投资决策，也影响资本在不同投资用途之间的配置。通过这些决策，利率将金融市场价格与实际资源配置联系起来。

\subsection{利率与储蓄、借款和投资}

利率变化会改变当期消费与未来消费的相对价格。对于储蓄者，利率提高会增加给定当期储蓄所支持的未来消费。这种替代效应通常鼓励储蓄。总体反应还取决于收入效应和财富效应。利息收入增加提高终身资源后，净储蓄者可能增加当期消费，净借款人则面临相反影响。总储蓄对利率变化的反应未必一致。

借款对融资成本的反应更直接。借款利率提高，会增加给定借款金额所要求的未来偿还额，通常降低信贷需求。当较高利率与逆向选择、抵押约束或贷款人资产负债表条件相互作用时，这种反应可能被放大，使信贷需求与信贷可得性同时下降。

投资会对资金机会成本变化作出反应。市场利率上升会提高替代金融请求权可获得的回报，在其他条件不变时，提高投资的必要报酬率。预期回报相对较低的投资机会更难获得融资，而回报足以补偿其风险的投资机会仍较有吸引力。

在总量层面，必要报酬率变化同时影响总投资和资本配置。市场利率的共同变化未必同等影响所有投资机会，因为项目的预期回报和风险不同，借款人面临的融资条件也不同。利率变化由此既影响总体投资水平，也影响资本在不同投资用途之间的分布。

\begin{figure}[H]
\centering
\begin{tikzpicture}[
  >=Latex,
  box/.style={
    draw,
    rounded corners=2pt,
    align=center,
    text width=2.4cm,
    minimum height=0.9cm,
    inner sep=4pt
  },
  box1/.style={
    draw,
    rounded corners=2pt,
    align=center,
    text width=2.0cm,
    minimum height=0.9cm,
    inner sep=4pt
  },
  flow/.style={
    -Latex,
    line width=0.8pt,
    FinanceBlue
  }
]

% Financing conditions
\node[box] (rates) at (-0.6,1.6)
  {基准利率\\与利差};

\node[box] (finance) at (2.5,1.6)
  {融资成本\\与可得性};

% Investment fundamentals
\node[box] (projects) at (2.5,-1.5)
  {预期项目\\现金流与风险};

% Investment choice
\node[box] (ranking) at (5.8,0.05)
  {不同投资用途的\\相对吸引力};

\node[box1] (allocation) at (8.8,0.05)
  {资本\\配置};

% Arrows
\draw[flow] (rates.east) -- (finance.west);

\draw[flow] (finance.east) -| (ranking.north);

\draw[flow] (projects.east) -| (ranking.south);

\draw[flow] (ranking.east) -- (allocation.west);

\end{tikzpicture}

\caption{融资条件与资本配置}
\label{fig:ch6-capital-allocation}
\end{figure}

图~\ref{fig:ch6-capital-allocation}~概括了这些关系。基准利率与利差影响融资成本和外部资金可得性，预期项目现金流与风险则决定不同投资用途的预期回报。由于各借款人的融资条件不同，市场利率的共同变化未必同等影响所有项目。利率变化因而既可能改变总投资，也可能改变资本在借款人与经济活动之间的配置。

\subsection{必要报酬率与资本配置}

市场利率为必要报酬率提供基准。某一借款人或金融请求权的必要报酬率，可以表示为基准利率加上反映借款人或工具特有特征的利差。即使基准利率相同，信用风险、流动性、抵押品、信息质量和契约条款的差异，也可能产生不同的必要报酬率。

基准利率与利差变化，对不同借款人的融资条件可能产生不均衡影响。表~\ref{tab:ch6-borrower-financing-rates}~说明了这一区别。

\begin{table}[H]

\centering

\caption{基准利率、利差与融资利率的变化}

\label{tab:ch6-borrower-financing-rates}

\begin{tabular}{@{}lccc@{}}

\toprule

\textbf{借款人} &
\textbf{基准利率} &
\textbf{利差} &
\textbf{融资利率} \\

\midrule

较低风险借款人
& 3.0\% $\rightarrow$ 4.0\%
& 1.0\% $\rightarrow$ 1.0\%
& 4.0\% $\rightarrow$ 5.0\% \\

较高风险借款人
& 3.0\% $\rightarrow$ 4.0\%
& 3.0\% $\rightarrow$ 4.0\%
& 6.0\% $\rightarrow$ 8.0\% \\

\bottomrule

\end{tabular}

\end{table}

两类借款人面临相同的基准利率增幅，但利差扩大使较高风险借款人的融资利率上升更多。其借款与投资可能对融资条件恶化作出更强反应。

融资成本差异影响资本配置。基准利率上升会提高一般的资金机会成本，利差变化则改变借款人之间的相对融资成本。当利差不均衡扩大时，资本可能转向融资条件较有利的借款人。利率既影响总投资，也影响投资在借款人与经济活动之间的配置。

%========================================================
% Summary
%========================================================

\TailModule{summary}{本章小结}

利率是当期购买力与未来购买力交换的价格。它既是贷款人的回报、借款人的成本，也是投资决策的基准，将不同日期的现金流联系起来。利率报价取决于期限和复利计息惯例。有效年利率为比较不同计息频率或连续复利下的报价提供了统一基础。名义利率与实际利率的区别，将货币金额变化与购买力变化区分开来。

观察到的利率反映多种经济力量。储蓄与投资影响融资供求，预期通胀影响名义利率，货币条件影响短期利率及预期，风险和流动性差异则形成工具之间的利差。仅凭观察到的利率较高，无法识别单一原因。

期限结构描述利率如何随期限变化。即期利率用于贴现特定未来日期的现金流，远期利率则适用于未来期间，可在无套利条件下由当前即期利率推导。预期假说将较长期利率与预期未来短期利率联系起来。市场分割、期限偏好及各期限的特定供求，也可能影响不同期限的利率。现代期限结构分析将较长期利率分解为预期短期利率成分与时变期限溢价。

利率既影响资本数量，也影响资本配置。基准利率变化改变一般的资金机会成本，利差差异使不同借款人面临不同融资条件。必要报酬率提高通常会减少边际投资机会获得的融资，但对各借款人或资本用途的影响未必一致。金融市场通过相对融资成本以及总体利率水平配置资本。

%========================================================
% Key Terms
%========================================================

\TailModule{keyterms}{关键术语}

利率（interest rate）；资金机会成本（opportunity cost of funds）；单利（simple interest）；复利（compound interest）；复利计息频率（compounding frequency）；名义年利率（nominal annual rate）；有效年利率（effective annual rate）；连续复利（continuous compounding）；名义利率（nominal interest rate）；实际利率（real interest rate）；预期通胀（expected inflation）；费雪关系（Fisher relation）；政策利率（policy rate）；基准利率（benchmark rate）；信用利差（credit spread）；风险溢价（risk premium）；利率期限结构（term structure of interest rates）；即期利率（spot rate）；远期利率（forward rate）；收益率曲线（yield curve）；预期假说（expectations hypothesis）；市场分割（market segmentation）；偏好期限（preferred habitat）；期限溢价（term premium）；必要报酬率（required return）；资本配置（capital allocation）。

%========================================================
% Financial AI Lab
%========================================================

\TailModule{aiprompt}{金融~AI~实验室}

\textbf{实验~1：审核利率比较}

请~AI~系统比较三种采用不同复利计息惯例的利率报价，其中至少包含一种连续复利利率。要求~AI~将各利率换算到相同的时间和计息基础，并列明换算使用的全部假设。

独立审核回答。检查期限是否可比，名义利率与有效利率是否正确区分，以及换算后各替代选择的排序是否变化。

引入预期通胀，扩展练习。要求~AI~区分名义利率和实际利率，并运用费雪关系。检查其是否区分基于预期通胀的预期实际利率，以及基于实际通胀的已实现实际利率。

\textbf{实验~2：审核收益率曲线解释}

请~AI~系统解释一条向上倾斜的收益率曲线，并说明其对未来短期利率有何含义。

检查回答是否将远期利率视为确定预测，或将曲线斜率完全归因于短期利率上升预期。结合预期未来短期利率与期限溢价的区别，以及各期限供求条件，修正这一解释。

对水平或倒挂的收益率曲线重复练习。判断哪些结论可由观察到的曲线推出，哪些需要对预期或期限溢价增加假设。

%========================================================
% Discussion Questions
%========================================================

\TailModule{questions}{讨论题}

\begin{enumerate}

\item 即使企业完全用内部资金投资，为什么利率仍是一种机会成本？

\item 为什么两笔年利率报价相同的贷款，可能具有不同的有效借款成本？

\item 为什么中央银行降低政策利率时，高风险企业面临的借款利率仍可能上升？

\item 为什么不应直接将远期利率解释为市场对未来短期利率的无偏预测？

\item 为什么相同的基准利率增幅，可能对不同借款人的投资产生不同影响？

\end{enumerate}

%========================================================
% Applied Problems
%========================================================

\TailModule{appprobs}{应用题}

\begin{enumerate}

\item \textbf{复利与有效利率。}

一家银行报出的存款名义年利率为~5.4\%，按月复利计息。第二种存款按年复利计息，利率为~5.5\%。第三种存款的连续复利年利率报价为~5.35\%。

\begin{enumerate}

\item 计算第一种存款的月利率。

\item 计算第一种存款的有效年利率。

\item 计算连续复利存款的有效年利率。

\item 比较三种存款的有效年度回报率。

\item 将三种存款按有效年度回报率从高到低排序。

\end{enumerate}

\item \textbf{名义回报与实际回报。}

某投资者一年内获得~7.5\%~的名义回报率。

\begin{enumerate}

\item 如果通胀率为~2.5\%，计算精确的实际回报率。

\item 使用费雪近似关系计算近似实际回报率。

\item 如果通胀率为~6\%，重新计算上述两项指标。

\item 解释为什么随着利率和通胀率上升，近似的准确性下降。

\end{enumerate}

\item \textbf{即期利率与远期利率。}

一年期即期利率为~3.0\%，两年期即期利率为~4.2\%，均按年复利计息。

\begin{enumerate}

\item 计算从今天起一年后开始的一年期远期利率。

\item 验证按两年期即期利率投资两年，与先按一年期即期利率投资一年、再按隐含远期利率投资一年，产生相同终值。

\item 解释为什么远期利率未必等于预期未来一年期短期利率。

\end{enumerate}

\item \textbf{预期与期限溢价。}

当前一年期利率为~3\%。第~2~年和第~3~年的预期一年期短期利率分别为~4\%~和~5\%。

\begin{enumerate}

\item 在预期假说下，计算近似的三年期即期利率。

\item 假设三年期期限溢价为~0.6~个百分点，相应的近似三年期即期利率为多少？

\item 解释为什么仅观察三年期即期利率，无法识别预期未来短期利率。

\end{enumerate}

\item \textbf{基准利率、利差与资本配置。}

两类借款人最初面临相同的基准利率~3.0\%。较低风险借款人的信用利差为~1.0\%，较高风险借款人的利差为~3.0\%。

随后，基准利率升至~4.0\%。较低风险借款人的利差不变，较高风险借款人的利差扩大至~4.0\%。

\begin{enumerate}

\item 计算各类借款人最初的融资利率。

\item 计算各类借款人新的融资利率。

\item 哪类借款人的融资成本增幅更大？

\item 解释相对融资条件变化如何影响资本在借款人之间的配置。

\item 解释为什么仅凭基准利率变化，不足以确定整个市场的融资条件如何变化。

\end{enumerate}

\end{enumerate}

%========================================================
% Data Application
%========================================================

\TailModule{dataapp}{数据应用}

\textbf{解读收益率曲线}

从权威数据来源获取所选日期的多种期限政府证券收益率。中国数据可采用中央国债登记结算有限责任公司发布的中债收益率曲线数据。开展国际比较时，可使用中央银行、财政部门、政府债务管理机构或其他权威市场来源的数据。

选取信用质量、工具类型和收益率定义尽可能可比的期限数据，构造收益率曲线。

描述曲线是向上倾斜、近似水平、倒挂，还是非单调。考虑观察到的形态可能来自以下三方面：

\begin{itemize}

\item 对未来短期利率的预期；

\item 期限溢价；

\item 各期限特定的供给、需求或流动性条件。

\end{itemize}

仅凭观察到的曲线，无法识别预期未来政策利率。解释作出这种解读，需要对期限溢价及其他期限特有效应作出哪些假设。

如果能获得两个日期的数据，比较曲线如何变化。判断变化集中在短端、长端，还是遍及整条曲线，并讨论在给出因果解释之前，还需要哪些证据。

报告观察日期、期限、收益率定义和原始数据来源。

%========================================================
% Further Reading
%========================================================

\TailModule{furtherread}{延伸阅读}

\begin{itemize}

\item Fisher, Irving (1930), \emph{The Theory of Interest: As Determined by Impatience to Spend Income and Opportunity to Invest It（利息理论：由消费收入的急切程度与投资收入的机会所决定）}, The Macmillan Company.

\item Hicks, John R. (1939), \emph{Value and Capital: An Inquiry into Some Fundamental Principles of Economic Theory（价值与资本：对经济理论若干基本原理的探讨）}, Clarendon Press.

\item Lutz, Friedrich A. (1940), ``The Structure of Interest Rates（利率结构）,'' \emph{每季度 Journal of Economics}, 55(1), 36--63.

\item Modigliani, Franco, and Richard Sutch (1966), ``Innovations in Interest Rate Policy（利率政策创新）,'' \emph{American Economic Review}, 56(1/2), 178--197.

\item Campbell, John Y., Andrew W. Lo, and A. Craig MacKinlay (1997), \emph{The Econometrics of Financial Markets（金融市场计量经济学）}, Princeton University Press.

\end{itemize}

\ChapterTail{summary,keyterms,aiprompt,questions,appprobs,dataapp,furtherread}

\chapter{货币市场与短期流动性契约}
\label{chap:Money_Markets_and_Short_Term_Liquidity_Contracts}

\begin{openingcase}{当隔夜融资突然变得昂贵}

2013~年~6~月，中国银行间货币市场经历了异常剧烈的流动性紧张。此前通常处于较低个位数水平的短期回购利率突然上升，6~月~20~日，部分基准隔夜和七天期回购利率在日内达到异常高位。银行同业放贷更为谨慎，交易活动也随之收缩。

这次紧张首先出现在期限极短的银行间融资中，并非源于银行资产长期价值的明显恶化。季节性现金需求、税款缴纳、资本流动变化，以及央行流动性条件收紧，与短期利率上升同时出现。随着不确定性增加，银行同业放贷更加谨慎，交易活动减弱。

一个日常用于调剂流动性的市场，为什么会在如此短的时间内陷入严重紧张？

\end{openingcase}

\begin{tcolorbox}[
  enhanced,
  colback=FinanceBlue!6,
  colframe=FinanceBlue,
  boxrule=0.6pt,
  sharp corners,
  left=8pt,
  right=8pt,
  top=7pt,
  bottom=7pt,
  before skip=1em,
  after skip=1.6em,
  title=\textbf{核心问题},
  coltitle=white,
  colbacktitle=FinanceBlue,
  fonttitle=\bfseries,
  attach boxed title to top left={xshift=6pt,yshift=-2mm},
  boxed title style={
    boxrule=0pt,
    sharp corners,
    left=5pt,
    right=5pt,
    top=2pt,
    bottom=2pt
  }
]
\begin{enumerate}[
  leftmargin=*,
  label=\textcolor{FinanceBlue}{\textbf{\arabic*.}},
  itemsep=0.55em,
  topsep=0.2em
]

\item 即使基础经济活动持续较长时间，企业和金融机构为什么仍需要短期融资？

\item 无担保和有抵押货币市场契约如何重新分配流动性，并决定短期融资条款？

\item 短期融资条件如何定价？通常流动性充裕的货币市场为什么会在压力下变得脆弱？

\end{enumerate}
\end{tcolorbox}

\section{短期融资的经济逻辑}

短期融资使企业和金融机构无需持有大量闲置现金，也能满足支付需求。其作用在于提供流动性和灵活性，但较短期限也带来了反复再融资的需要。

\subsection{现金流时点与短期融资需求}

现金收入与现金支出并不总在同一时点发生。具有偿付能力的企业仍可能面临暂时性现金短缺。短期融资弥合这些时间缺口，减少参与者需要持有的闲置现金。

假设一家企业期初现金余额为人民币~20~万元，须在一周内向员工和供应商支付~80~万元，但预计在~45~天后从客户处收到~100~万元。\footnote{译注：原稿此处为~45~天，下文与图示均为第~28~天。译稿保留原数值，待作者核定。}该企业的现金状况变化如下：

\begin{itemize}
\item 期初，企业持有人民币~20~万元现金。

\item 第~7~天付款到期时，企业面临~60~万元现金短缺。

\item 第~28~天，客户支付~100~万元，企业得以偿还短期借款，并保留正的现金余额。
\end{itemize}

因此，即使该期间预期现金流入超过计划现金流出，企业仍需要~60~万元临时融资。这说明，收付款时点差异可能产生无法从期末现金余额看出的融资需求。

金融机构也存在类似的时点问题。银行可能在款项流入前先发生结算资金流出；证券交易商可能需要为持有较长时间的证券库存续借隔夜融资。保证金支付和集中的支付义务也可能产生暂时性融资需求。

短期现金短缺并不必然意味着资不抵债。偿付能力关注资产经济价值是否足以覆盖负债，流动性则关注能否在支付到期时履行义务。因此，一家机构即使具有偿付能力，若无法及时将资产变现或取得新融资以履行支付义务，也可能面临流动性短缺。
\begin{figure}[H]
\centering
\begin{tikzpicture}[
  >=Latex,
  axis/.style={-Latex, line width=0.8pt},
  inflow/.style={-Latex, line width=1.0pt, FinanceBlue},
  outflow/.style={-Latex, line width=1.0pt}
]

\draw[axis] (0,0) -- (10.5,0) node[right]{时间};

% Tick labels
\node[below=4pt] at (0,0) {0};
\node[below=4pt] at (2.2,0) {7};
\node[below=4pt] at (8,0) {28};

% Initial cash
\draw[inflow] (0,0) -- (0,0.5)
  node[above=2pt, align=center]
   {\small 期初现金余额\\\small 人民币~20~万元};
% Cash outflow
\draw[outflow] (2,0) -- (2,-1.5)
  node[below=2pt, align=center]
  {\small 向供应商支付\\\small 及向员工支付\\\small 人民币~80~万元};

% Cash inflow
\draw[inflow] (8,0) -- (8,1.6)
  node[above=2pt, align=center]
  {\small 客户应收账款回款\\\small 人民币~100~万元};

\end{tikzpicture}

\caption{现金流时点缺口与短期融资的作用}
\label{fig:ch7-cash-flow-timing}
\end{figure}

金融体系内部也存在类似的时点问题。一家机构当天收到的现金可能超过其需要，另一家则面临暂时性缺口。货币市场重新分配这些资金余额，使参与者能够节约闲置流动性，并支持支付和证券交易结算。

融资安排的期限影响风险。如果时点缺口确属暂时性，短期借款可以较好地匹配实际需要。但如果短期债务用于融资的资产多年内都缺乏流动性，借款人就必须反复再融资。这种期限错配产生滚动融资风险，即现有融资到期时，替代融资变得更昂贵或无法获得的风险。

\begin{example}{管理多项现金流义务}
{multiple-cash-flow-obligations}

一家制造企业在一个~15~天期间开始时持有人民币~200~万元现金。期间，企业面临多项支付义务，并预计在不同日期收到客户付款。新增融资前的预计现金流如下。

\begin{center}
\begin{tabular}{lrr}
\toprule
\textbf{日期} &
\shortstack[c]{\textbf{现金流}\\\textbf{（人民币百万元）}} &
\shortstack[c]{\textbf{累计}\\\textbf{现金余额}} \\
\midrule
期初 & 2.0 & 2.0 \\
第~2~天：供应商付款 & -5.0 & -3.0 \\
第~4~天：客户回款 & 1.5 & -1.5 \\
第~5~天：工资和税款 & -1.0 & -2.5 \\
第~8~天：客户回款 & 4.0 & 1.5 \\
第~10~天：运费和利息 & -0.5 & 1.0 \\
第~15~天：客户回款 & 2.0 & 3.0 \\
\bottomrule
\end{tabular}
\end{center}

企业在第~2~天出现最大现金缺口，预计现金余额降至负~300~万元。因此，企业至少需要~300~万元临时融资，才能履行支付义务。此后的客户付款使企业能够减少或偿还借款。

这一例子说明，短期融资需求取决于现金流的先后顺序，而不仅是给定期间内总收入与总支出的差额。虽然企业期末现金余额为正，但若无临时性外部融资，就无法按时履行全部义务。

承诺信贷额度或短期贷款可以弥合这一缺口。融资选择取决于可得性、成本、期限及预期客户付款的可靠性。

\end{example}

\subsection{营运资本、流动性管理与财务灵活性}

对非金融企业而言，短期融资与营运资本密切相关。在现金回收之前，经营活动会使资金占用于存货和应收账款，而供应商提供的商业信用则是短期融资来源。

经营性营运资本的一种简化计量为

\begin{equation*}
OWC=\text{存货}+\text{应收账款}-\text{应付账款}.
\label{eq:ch7-operating-working-capital}
\end{equation*}

存货或应收账款增加会占用资金，应付账款增加则提供短期融资。如果必须在客户付款之前投入资金，即使企业盈利，销售快速增长也可能增加短期融资需求。

流动性管理涉及用于满足短期现金需求的资产和融资安排。企业可以持有现金和有价证券、维持银行承诺授信，或取得短期市场融资。这些选择反映流动性保障与融资成本之间的权衡。更大的流动性缓冲可降低暂时性融资中断的影响，但流动资产的回报通常低于生产性投资。承诺信贷额度以一定成本保障资金获取，而依赖极短期市场融资则会增加滚动融资风险暴露。

财务灵活性指在不承担高额成本、不干扰投资与经营的情况下，获得融资或调整融资安排的能力。它不仅取决于当前现金持有量，还取决于未使用的借款能力、进入其他融资市场的渠道、未设权利负担的抵押品、负债期限结构和借款人信用状况。这些因素决定条件变化时能够筹集或重组融资的程度。

对金融机构而言，流动性管理和财务灵活性高度依赖资产构成及融资市场准入。银行可以持有可出售或质押的高质量流动性资产，在银行间市场融资，并以合格抵押品取得央行流动性。证券交易商通常通过回购协议融资证券库存，其融资能力因而部分取决于抵押品资格和折扣率。货币市场基金持有短期资产，但若这些资产无法在不作重大价格让步的情况下迅速出售，快速赎回仍可能带来流动性压力。

因此，短期融资在提供灵活性的同时，也产生对再融资的依赖。较短期限使融资条件能够迅速调整，但也要求借款人更频繁地回到贷款人或市场寻求资金。期限不仅决定偿还时点，也决定融资重新定价、续借或撤回的频率。

\begin{tcolorbox}[
    enhanced,
    breakable,
    colback=white,
    colframe=FinanceBlue,
    boxrule=0.7pt,
    sharp corners,
    borderline west={2.2pt}{0pt}{FinanceBlue},
    title=\textbf{金融视角：短期限与滚动融资依赖},
    coltitle=white,
    colbacktitle=FinanceBlue,
    fonttitle=\bfseries,
    before skip=1.2em,
    after skip=1.2em
]

短期借款允许融资条款频繁调整，从而提供灵活性，但同一特征也增加了再融资依赖。一笔一天期贷款的贷款人只承诺将资金提供至到期日，可在下次续借时重新定价或拒绝续借。如果基础融资需求仍然存在，借款人就继续面临市场条件变化。短期限因而提高了续借决策的重要性，并使融资风险集中于再融资时点。

\end{tcolorbox}

\section{货币市场契约与工具}

货币市场工具和契约是用于满足临时融资需求、管理现金余额及融资证券头寸的短期请求权与融资安排。它们在发行人或交易对手、抵押品、信用风险暴露、可交易性和报价惯例方面存在差异，进而影响定价、流动性和风险。

\subsection{国库券与短期主权请求权}

国库券是中央政府发行的短期债务证券，通常不定期支付票息，而在到期时支付固定面值，回报来自购买价格与到期价值之差。

若以~\(P\)~购买面值为~\(F\)~的国库券，持有~\(d\)~天至到期，其持有期回报率为

\begin{equation}
HPR=\frac{F-P}{P}.
\label{eq:ch7-tbill-hpr}
\end{equation}

为比较不同期限国库券，可将持有期回报率按单利年化。所得收益率取决于市场报价惯例。以购买价格为基础的货币市场收益率可写为

\begin{equation}
y_m=\frac{F-P}{P}\frac{D}{d},
\label{eq:ch7-money-market-yield}
\end{equation}
其中，$D$~为市场惯例规定的年度计息天数，通常为~360~或~365。

一些国库券市场还使用以面值为基础的贴现收益率：

\begin{equation}
y_d=\frac{F-P}{F}\frac{D}{d}.
\label{eq:ch7-discount-yield}
\end{equation}

\begin{example}{报价惯例改变报告收益率}{tbill-quotation-convention}

一张期限为~91~天、面值为~100~的国库券，售价为~98.80。采用每年~360~天的计息基础，持有期回报率为

\[
HPR=\frac{100-98.80}{98.80}=1.2146\%.
\]

年化货币市场收益率为

\[
y_m=\frac{1.20}{98.80}\frac{360}{91}=4.805\%,
\]
贴现收益率则为

\[
y_d=\frac{1.20}{100}\frac{360}{91}=4.747\%.
\]

两种收益率对应相同的基础现金流，数值差异完全来自报价惯例。

\end{example}

国库券是短期主权债务最常见的形式，但具体工具、期限和发行惯例因国家而异。

短期主权证券的作用不限于政府融资。当信用风险较低、二级市场流动性较高时，它们可以作为临时现金投资工具、有抵押融资的抵押品，以及其他短期请求权定价的基准。

信用风险较低，并不意味着所有持有期内市场价格都无风险。如果到期前出售国库券，其价格可能随市场利率变化。结算、市场流动性及主权信用条件也可能产生影响。与其他特征相近的较长期固定利率证券相比，短期限通常使其利率敏感性较低。

\subsection{商业票据与企业短期融资}

商业票据（commercial paper，CP）是企业和金融机构为满足短期融资需求而发行的、通常无担保的短期债务工具。投资者主要依赖发行人的信用状况和流动性状况。商业票据可采用贴现或付息方式发行。贴现票据低于面值发行，到期按面值偿还；付息票据则在本金之外支付约定利息。

两种结构的区别在于投资者获得回报的方式。贴现票据的回报来自购买价格与面值之差，付息票据的回报来自约定利息支付。

\begin{example}{贴现商业票据与付息商业票据}
{commercial-paper-payment-structures}

考虑两笔~90~天期商业票据，其本金或面值均为人民币~1000~万元。假设计息年度为~360~天。

\textbf{贴现商业票据。}

假设票据以~985~万元发行，90~天后按面值~1000~万元偿还，投资者获得~15~万元。年化贴现率为

\[
d=\frac{10-9.85}{10}\times\frac{360}{90}
=6\%.
\]

以购买价格为分母计算的相应单利年化投资收益率为

\[
y=\frac{10-9.85}{9.85}\times\frac{360}{90}
\approx 6.09\%.
\]

\textbf{付息商业票据。}

假设票据按面值发行，年利率为~6\%，到期应付利息为

\[
I=10\times 6\%\times\frac{90}{360}
=0.15\text{ 百万元}.
\]

投资者到期收到本息合计人民币~1015~万元，单利年化收益率为~6\%。

这一例子说明，贴现率与投资收益率不同。贴现率以面值为基础计算，投资收益率则以购买价格为基础计算。按面值发行的付息票据通过约定利率表达其回报。

虽然两种报价均为~6\%，但贴现票据的贴现率以面值而非购买价格计算，因而提供略高的投资收益率。

\end{example}

投资者要求对货币时间价值，以及与发行人相关的信用和流动性风险给予补偿。因此，必要收益率反映可比期限的基准利率，以及与发行人特定信用和流动性条件相关的利差。

对信用质量较高的发行人而言，商业票据提供直接获取短期市场融资的渠道，其成本可能低于银行借款。不过，依赖反复发行会产生滚动融资风险。依赖~30~天期票据的企业，必须在每次到期时以可用现金偿还，或获得新融资。如果市场对其信用质量的评价恶化，或投资者需求减弱，利差可能扩大，新发行渠道也可能受限。因此，企业常以承诺银行信贷额度或流动性储备支持商业票据计划。即使发行人仍有偿付能力，滚动融资风险也可能出现。

相关的短期市场融资形式是资产支持商业票据（ABCP）。与普通企业商业票据不同，ABCP~由特殊目的载体发行，以应收账款或贷款等金融资产池为支持。该计划还可能获得发起人或其他机构提供的流动性便利或信用增级。

\subsection{存单与银行融资}

存单是银行负债，承诺在约定到期日偿还本金，通常还支付利息。可转让存单可以在到期前转让，使初始投资者能够通过出售请求权获得流动性，而无需要求发行银行提前偿还。

从银行角度看，可转让存单是负债管理工具。它使银行能够按市场决定的利率筹资，而不必完全依赖零售客户存款。报价利率反映请求权期限、市场流动性、银行信用质量，以及其他短期投资的供给。

可交易性改变了投资者在到期前获取流动性的方式。不可转让的定期存款通常要求存款人持有至到期，或向银行申请提前支取；可转让存单则可在二级市场出售，但价格可能随市场利率及市场对发行银行信用质量的判断而变化。

\begin{table}[H]
\centering
\caption{定期存款与可转让存单}
\label{tab:ch7-cd-comparison}

\begin{tabularx}{\textwidth}{lXX}
\toprule
 & \textbf{不可转让定期存款} & \textbf{可转让存单} \\
\midrule

可转让性 & 通常不可转让 & 可在二级市场转让 \\

到期前流动性 & 通常需提前支取，或与银行另作安排 & 可通过出售请求权获得 \\

投资者风险暴露 & 主要面向发行银行 & 面向发行银行，并承担到期前的市场价格风险 \\

对银行的融资作用 & 存款融资 & 基于市场的批发融资 \\

\bottomrule
\end{tabularx}
\end{table}

中国银行间同业存单是市场化银行融资的重要例子。根据中国人民银行~2013~年暂行规定推出的此类记账式定期存款凭证，由符合条件的存款类机构在全国银行间市场发行。发行利率由市场决定，公开发行的存单可以在二级市场交易。它们为银行提供额外的批发融资来源，也为投资者提供对发行机构的可转让请求权。

\subsection{同业贷款与隔夜无担保融资}

银行日常代客户收付款，这些资金流动使各银行的流动性头寸暂时不同。同业拆借使资金盈余机构能够向存在短期缺口的机构提供资金。

无担保同业贷款是对借款银行的直接请求权。偿还不以特定抵押品作支持，贷款人依赖借款人的信用状况及双边或全市场信用限额。因此，利率同时反映一般流动性条件和交易对手信用风险。

隔夜期限在流动性快速再分配中具有核心作用。银行当天遇到非预期结算资金流出时，可以借入隔夜资金，并在次日资金流入后偿还。数天或数月期限的融资可在更长期间内锁定流动性，但也延长贷款人对借款人的风险暴露。

压力时期，请求权的无担保性质更加重要。银行若无法确定交易对手的信用状况，即使持有现金也可能拒绝放贷。许多机构采取类似行为时，市场可能在不同交易对手之间分割：流动性仍在体系内，却无法向需要它的机构重新分配。

信用限额与市场利率同样重要。如果贷款人不愿增加无担保风险暴露，即使报价利率大幅提高，也未必带来额外放贷。此时，可用信贷数量及交易对手身份与报价同样重要。

\subsection{回购协议与有抵押融资}

回购协议（repo）是一种短期融资契约，一方以证券换取现金，并约定未来按指定回购价格进行反向交易。从经济功能看，回购属于有抵押短期融资，但其法律结构因市场而异。

从现金借入方看，交易是正回购；从现金贷出方看，同一交易是逆回购。设市场价值为~$V_0$~的抵押品支持一笔回购，如果贷款人采用折扣率~$h$，则提供的现金为

\begin{equation}
C_0=(1-h)V_0.
\label{eq:ch7-repo-cash-haircut}
\end{equation}

折扣为抵押品价值下跌及借款人违约时的清算成本提供缓冲。如果回购期限为一年的~$\tau$，回购利率为~$r_{repo}$，则回购支付额为

\begin{equation}
C_T=C_0(1+r_{repo}\tau),
\label{eq:ch7-repo-repurchase-payment}
\end{equation}
这里对回购期间采用单利计息惯例。

\begin{example}{回购融资与折扣率}{repo-funding-haircut}

一家交易商持有市值人民币~1~亿元的政府证券，并通过一周期回购融资。现金贷款人采用~2\%~的折扣率。

提供的现金为

\[
C_0=(1-0.02)\times100=98
\]

百万元。如果年化回购利率为~2.5\%，市场采用每年~365~天的计息基础，则七天后的回购支付额为

\[
C_T=98\times\left(1+0.025\times\frac{7}{365}\right)
\approx 98.047
\]
百万元。

折扣率与回购利率发挥不同作用。回购利率是现金融资的价格；折扣率决定相对于抵押品价值能够获得多少现金。

\end{example}

抵押品能够减少贷款人的损失暴露。但抵押品价格可能下跌，证券可能难以清算，法律和操作问题也可能延迟贷款人取得抵押品。回购契约通常规定抵押品资格、估值程序、折扣率和保证金安排。

回购存续期内抵押品价值变化时，保证金调整可能要求补充现金或证券，以维持约定的保护水平。这种盯市机制保护贷款人，但抵押品价值下降时，可能给借款人带来流动性压力。

回购使交易商和其他机构能够融资债券及其他合格证券库存，是证券市场中介活动的核心安排。它也将货币市场与债券市场联系起来：证券价值与流动性影响其作为抵押品的适用性，回购融资的可得性与成本则影响中介持有该证券的意愿。

表~\ref{tab:ch7-money-market-instruments}~概括了本节所讨论工具的主要契约差异。

\begin{table}[H]
\centering
\caption{主要货币市场工具及其契约结构}
\label{tab:ch7-money-market-instruments}

\begin{tabularx}{\textwidth}{
@{}
>{\raggedright\arraybackslash
\hyphenpenalty=10000
\exhyphenpenalty=10000}p{2.3cm}
>{\raggedright\arraybackslash
\hyphenpenalty=10000
\exhyphenpenalty=10000}p{2.5cm}
>{\raggedright\arraybackslash
\hyphenpenalty=10000
\exhyphenpenalty=10000}p{2.6cm}
Y
@{}
}

\toprule

\textbf{工具}
&
\makecell[l]{\textbf{典型发行人}\\\textbf{或借款人}}
&
\makecell[l]{\textbf{抵押品}\\\textbf{或信用基础}}
&
\makecell[l]{\textbf{主要经济作用与风险}}
\\

\midrule

国库券 & 中央政府 & 主权支付义务；通常无担保 & 政府短期融资与流动性投资；优质主权市场中信用风险较低，到期前仍有市场价格风险 \\

商业票据 & 企业或金融机构 & 通常无担保；发行人信用 & 灵活的短期市场融资；发行人信用风险及滚动融资风险 \\

可转让存单 & 银行 & 银行存款支付义务；通常无担保 & 银行批发融资及可转让银行请求权；银行信用风险及滚动融资风险 \\

同业贷款 & 银行或其他金融机构 & 通常为无担保交易对手信用 & 短期流动性再分配；交易对手信用风险及滚动融资风险 \\

回购 & 交易商、银行、基金或其他参与者 & 以卖出并回购协议提供抵押 & 证券融资与流动性管理；抵押品估值、折扣率、交易对手及滚动融资风险 \\

资产支持商业票据 & 特殊目的载体 & 资产池，常附流动性或信用支持 & 资产支持的短期融资；基础资产质量、流动性及滚动融资风险 \\

\bottomrule
\end{tabularx}
\end{table}

\section{货币市场架构与定价}

货币市场价格形成于借款人、现金投资者、银行、交易商和投资基金之间的互动，并依托市场交易与结算基础设施。契约条款决定工具现金流，市场架构则影响融资获取、流动性再分配、价格发现与基准形成，以及交易执行与结算。

\subsection{银行、交易商与货币市场基金}

银行在货币市场中兼为借款人与贷款人。它们通过借款管理结算需求和资产负债表短期融资，并在出现临时流动性盈余时放贷。这些活动将短期融资市场与支付体系及更广泛的银行体系联系起来。

证券交易商将短期融资与债券等证券交易联系起来。持有证券库存的交易商可以通过回购融资部分头寸。这种资产负债表承载能力使交易商能够持有库存、为买卖双方提供即时交易服务，从而支持二级市场流动性。回购融资成本和可得性影响交易商的做市意愿。

货币市场基金汇集投资者现金，主要投资于短期限、高质量工具。它们为家庭、企业和机构提供分散化投资货币市场资产的渠道，同时将现金导向政府、银行、企业和交易商。在适用于基金及所在司法辖区的规则约束下，投资组合可包括政府短期证券、回购、商业票据和存单。

货币市场基金的结构使资产流动性与投资者流动性之间形成重要区别。投资者可能预期短时间内即可赎回份额，而部分组合资产在市场压力时期可能难以迅速出售。基金遭遇大量赎回时，可以动用现金和到期证券回款、出售资产，或减少新增购买。这些行为将投资者提款传导至短期借款人使用的市场。

\begin{figure}[H]
\centering
\begin{tikzpicture}[
  >=Latex,
  box/.style={
    draw,
    rounded corners=2pt,
    align=center,
    minimum height=1.0cm,
    text width=2.7cm,
    inner sep=5pt
  },
  widebox/.style={
    draw,
    rounded corners=2pt,
    align=center,
    minimum height=1.0cm,
    text width=3.6cm,
    inner sep=5pt
  },
  fundbox/.style={
    draw,
    rounded corners=2pt,
    align=center,
    minimum height=1.0cm,
    text width=1.5cm,
    inner sep=5pt
  },
  flow/.style={
    -Latex,
    line width=0.9pt,
    FinanceBlue
  },
  claim/.style={
    -Latex,
    dashed,
    line width=0.9pt,
    FinanceBlue
  },
abovelab/.style={
  font=\small,
  text=FinanceBlue,
  fill=white,
  align=center,
  fill=white,
  inner sep=1.5pt
},
  belowlab/.style={
    font=\small,
    text=black,
    fill=white,
    align=center,
    inner sep=1.5pt
  }
]

% =================================================
% Top row: Money market funds
% =================================================

\node[fundbox] (inv) at (-0.5,3.0)
  {现金投资者};

\node[fundbox] (mmf) at (3.8,3.0)
  {货币市场\\基金};

\node[widebox] (iss) at (9.2,3.0)
  {短期请求权发行人\\
   （政府、企业、银行）};

% Investor subscriptions and fund shares
\draw[flow]
  ([yshift=3pt]inv.east) --
  node[abovelab, midway, above=3pt] {现金\\认购}
  ([yshift=3pt]mmf.west);

\draw[claim]
  ([yshift=-3pt]mmf.west) --
  node[belowlab, midway, below=3pt] {基金份额}
  ([yshift=-3pt]inv.east);

% MMF investment in short-term claims
\draw[flow]
  ([yshift=3pt]mmf.east) --
  node[abovelab, midway, above=3pt] {现金\\投资}
  ([yshift=3pt]iss.west);

\draw[claim]
  ([yshift=-3pt]iss.west) --
  node[belowlab, midway, below=3pt]
  {国库券、\\商业票据、存单}
  ([yshift=-3pt]mmf.east);

% =================================================
% Middle row: Interbank lending
% =================================================

\node[box] (bank1) at (0,0)
  {流动性\\盈余银行};

\node[widebox] (bank2) at (9.2,0)
  {有临时融资\\需求的银行};

% Cash and loan claim
\draw[flow]
  ([yshift=3pt]bank1.east) --
  node[abovelab, midway, above=3pt]
  {同业资金}
  ([yshift=3pt]bank2.west);

\draw[claim]
  ([yshift=-3pt]bank2.west) --
  node[belowlab, midway, below=3pt]
  {贷款请求权}
  ([yshift=-3pt]bank1.east);

% =================================================
% Bottom row: Repo financing
% =================================================

\node[widebox] (cashl) at (0,-3.0)
  {现金贷出方\\（银行、货币基金、其他投资者）};

\node[widebox] (dealer) at (9.2,-3.0)
  {持有证券库存\\的交易商};

% 回购资金 and collateral
\draw[flow]
  ([yshift=3pt]cashl.east) --
  node[abovelab, midway, above=3pt]
  {回购资金}
  ([yshift=3pt]dealer.west);

\draw[claim]
  ([yshift=-3pt]dealer.west) --
  node[belowlab, midway, below=3pt]
  {证券抵押品}
  ([yshift=-3pt]cashl.east);

% =================================================
% Legend
% =================================================

\draw[flow] (0.5,-4.5) -- (1.4,-4.5);
\node[anchor=west, draw=none] at (1.5,-4.5)
  {\small 现金流};

\draw[claim] (4.2,-4.5) -- (5.1,-4.5);
\node[anchor=west, draw=none] at (5.2,-4.5)
  {\small 请求权或抵押品转移};

\end{tikzpicture}

\caption{货币市场融资的机构渠道}
\label{fig:ch7-money-market-channels}
\end{figure}

该图为示意图。银行可同时处于市场借贷两侧，货币市场基金与交易商则履行不同的中介功能。图示旨在说明短期资金如何通过多种机构渠道流动。

\subsection{基准利率、利差与价格发现}

中央银行通过结算余额供给、公开市场操作和流动性便利影响货币市场条件。这些操作影响极短期利率及总体融资可得性。但交易对手信用风险、抵押品条件、市场分割和机构特有融资条件，仍可能使各工具利率不同。

货币市场契约常参照基准利率定价。有效基准应代表界定清楚的市场区间，依据可靠信息，并有透明方法和适当治理支持。其经济含义取决于基础交易及利率构造方法。

基准在所衡量的融资类型和构造方式上均存在差异。一项区别是融资是否有担保。有担保利率反映抵押品支持的融资交易，受抵押品质量和可得性影响；无担保利率反映无抵押品借款，更直接体现交易对手信用条件。另一项区别是利率形成方式：交易型基准由合格市场交易推导，报价型基准依赖选定市场参与者提交的利率。

中国货币市场体现了这些差异。上海银行间同业拆放利率（Shibor）是报价型无担保基准，由报价银行团提交利率形成，并发布多个期限品种。相比之下，存款类机构间回购利率（DR）是交易型有担保利率，依据存款类机构之间以利率债为抵押品的质押式回购交易形成。七天期利率~DR007~广泛用作银行体系流动性指标和货币市场定价参考。

国际基准改革增加了交易型隔夜利率的使用。美国~SOFR~基于有担保的国债回购交易，英国~SONIA~和欧元区~€STR~则基于无担保隔夜交易。

表~\ref{tab:ch7-benchmark-comparison}~概括了这些基准的主要特征。

\begin{table}[H]
\centering
\caption{主要货币市场基准及其特征}
\label{tab:ch7-benchmark-comparison}

\begin{tabularx}{\textwidth}{
@{}
>{\raggedright\arraybackslash}p{1.7cm}
>{\raggedright\arraybackslash}p{3.2cm}
>{\raggedright\arraybackslash}p{2.0cm}
Y
@{}
}

\toprule

\textbf{基准}
&
\textbf{基础市场}
&
\textbf{融资基础}
&
\textbf{利率形成方式}
\\

\midrule

Shibor & 人民币同业融资 & 无担保 & 银行团报价 \\

DR007 & 七天期同业质押式回购 & 有担保 & 实际交易加权利率 \\

SOFR & 隔夜美国国债回购 & 有担保 & 基于交易的成交量加权中位数 \\

SONIA & 隔夜英镑存款 & 无担保 & 基于交易的截尾均值 \\

€STR & 隔夜欧元批发融资 & 无担保 & 基于交易计算 \\

\bottomrule
\end{tabularx}
\end{table}

对于短期工具及期限和市场特征可比的基准，利差为

\begin{equation}
s_i=r_i-r_b.
\label{eq:ch7-money-market-spread}
\end{equation}
其中，$r_i$~表示短期工具利率，$r_b$~表示适当的基准利率。

利差扩大可能反映信用风险增加、市场流动性下降、资产负债表约束收紧，或工具需求减弱。因此，同一观察到的利差可能来自不同基础条件。

借款人与贷款人根据融资需求、新信息、货币操作和交易条件调整利率，形成价格发现。货币市场请求权期限短、重新定价频繁，使极短期利率对融资条件变化尤其敏感。纳税或报表日期前后、央行流动性变化后，或机构放贷意愿下降时，隔夜利率可能剧烈波动。

基准利率也用于货币市场之外的契约。浮动利率贷款、债券、互换及衍生品可能参照短期基准确定契约支付。因此，货币市场价格发现影响更广泛金融请求权的定价与估值。

\subsection{抵押品、清算与结算}

除借款人的一般信用状况外，抵押品为贷款人提供额外回收来源。其有效性取决于三个问题：哪些资产合格、如何估值，以及违约后贷款人能多快行使相关权利。

政府证券及其他高质量有价证券常用作货币市场抵押品。其价值较为透明，通常比流动性较低的资产更易出售或用于融资。其他证券也可能合格，但在价格波动、信用风险或清算成本较高时，贷款人可能采用更高折扣率。

初始折扣为抵押品价值下跌提供缓冲，但市场价格变化后，缓冲可能不足。因此，回购契约可能要求每日或日内对抵押品盯市。如果抵押品价值低于协议要求，借款人可能需要补充证券或现金。这类追缴保证金要求限制贷款人风险暴露，但也可能给借款人带来流动性压力，抵押品价格急跌时尤其如此。

清算确定交易产生的义务，并在适用情况下确定待结算金额。结算通过转移现金和证券完成义务。证券交易中，卖方须交付证券，买方须支付相应现金。如果两项转移发生在不同时点，任何一方都可能面临对方不履约的风险。券款对付（DvP）将证券交付与相应现金支付相联系，降低这种本金风险。

中国银行间市场体现了交易服务与证券登记、托管和结算之间的分工。中国外汇交易中心（CFETS），亦称全国银行间同业拆借中心，为同业存单提供发行、交易和信息服务；上海清算所（SHCH）负责登记、托管和结算。这一安排将工具交易与负责维护所有权记录、完成结算的机构职能区分开来。

这些安排无法消除流动性风险。结算延误或操作问题妨碍抵押品动用时，即使有价值的抵押品也可能无法及时提供融资。货币市场流动性因此不仅取决于金融请求权质量，也取决于其转移所依托的基础设施。

\begin{tcolorbox}[
  enhanced,
  breakable,
  colback=white,
  colframe=FinanceBlue,
  boxrule=0.7pt,
  sharp corners,
  borderline west={2.2pt}{0pt}{FinanceBlue},
  title=\textbf{金融视角：抵押品必须能够动用},
  coltitle=white,
  colbacktitle=FinanceBlue,
  fonttitle=\bfseries,
  before skip=1.2em,
  after skip=1.2em
]

抵押品价值不仅取决于信用质量和市场价格，还取决于融资需要出现时，能否转移、被接受和清算。法律、操作与结算摩擦可能限制其作为即时流动性来源的作用。

\end{tcolorbox}

\section{货币市场的流动性压力}

正常条件下，货币市场可能表现出较高流动性。短期限和定期再融资使资金能够在机构间顺畅流动。但贷款人不愿续借时，同样的特征也可能加剧压力。大量债务可能短时间内到期，迫使借款人在长期资产能够出售或到期之前筹集现金。

\subsection{融资挤兑与市场冻结}

融资挤兑指短期债权人撤回融资、拒绝续借到期请求权，或缩短新贷款期限。即使机构依赖短期批发融资而非存款，其动态过程也类似银行挤兑。市场冻结则指参与者不愿或无法按现行条款交易，导致交易或新增融资急剧下降。

假设借款人以每日或每周到期的债务融资长期资产。只要贷款人继续续借，就能维持头寸。对借款人资产质量、信用状况的担忧，或贷款人自身流动性需求，可能降低续借意愿。借款人必须在现有请求权到期时筹集现金。

为履行义务，借款人可动用承诺授信、使用现金储备、追加抵押品，或出售资产。如果许多机构面临相似融资压力，这些反应可能相互强化。资产出售可能压低市场价格，流动性需求上升则进一步挤压融资条件。资产价格下跌还可能削弱资产负债表、降低抵押品价值，使融资压力进一步增大。这一过程可能自我强化：

\begin{figure}[H]
\centering
\begin{tikzpicture}[
  >=Latex,
  box/.style={draw, rounded corners, align=center, text width=4.0cm, minimum height=0.9cm},
  flow/.style={-Latex, line width=0.8pt, FinanceBlue}
]

\node[box] (concern) at (0,3.2) {对信用、流动性\\或抵押品的担忧};
\node[box] (roll) at (6.0,3.2) {贷款人缩短期限\\或拒绝续借};
\node[box] (cash) at (6.0,0.7) {借款人筹集现金\\并出售资产};
\node[box] (price) at (0,0.7) 
{资产价格下降，削弱\\抵押品价值与市场流动性};

\draw[flow] (concern) -- (roll);
\draw[flow] (roll) -- (cash);
\draw[flow] (cash) -- (price);
\draw[flow] (price) -- (concern);

\end{tikzpicture}
\caption{短期融资的反馈循环}
\label{fig:ch7-funding-run-loop}
\end{figure}

在无担保市场中，即使总体流动性仍充裕，对借款人质量的不确定性也可能减少放贷。机构可能更愿保留流动性，而非增加对状况难以评估的交易对手的风险暴露。融资压力同时影响价格与数量：利差可能上升，部分借款人则可能完全丧失融资渠道。

商业票据市场与货币市场基金也可能出现类似压力。投资者可能从商业票据转向政府证券或现金，减少发行人可获得的融资，增加其对承诺银行信贷额度的依赖。如果货币市场基金减少购买商业票据与存单，或出售资产应对提款，快速赎回就会强化这一过程。压力由此从现金投资者传向短期借款人，再传向银行资产负债表。

2013~年~6~月中国银行间市场的情况说明了这些机制。随着融资不确定性增强，机构更不愿放贷，削弱了银行体系内部的流动性再分配。短期利率急升，随后在流动性条件与预期稳定后回落。

\subsection{回购折扣率、抵押品压力与避险转移}

回购融资条件同时取决于回购利率与折扣率。回购利率决定现金融资成本，折扣率决定给定数量抵押品能够借入多少现金。

市场压力时期，这一区别尤其重要。考虑完全由回购借款与交易商自有资金融资的简化交易商头寸。设~$V$~为证券市值，$E$~为交易商投入的自有资金。按标准折扣惯例，回购折扣率~$h$~表示贷款人针对价值~$V$~的抵押品提供~$(1-h)V$~现金。因此，交易商必须以自有资金承担剩余部分：

\begin{equation*}
E=hV.
\label{eq:ch7-haircut-equity}
\end{equation*}

在这一简化融资结构下，资产权益比为

\begin{equation*}
\frac{V}{E}=\frac{1}{h}.
\label{eq:ch7-haircut-leverage}
\end{equation*}

市场压力时期折扣率上升，其融资影响尤为显著。

\begin{example}{折扣率上升产生融资需求}{haircut-funding-need}

一家交易商通过回购融资市值人民币~1~亿元的证券。折扣率为~2\%~时，贷款人提供~9800~万元，交易商投入~200~万元自有资金。

如果抵押品价值仍为~1~亿元，而折扣率升至~5\%，回购融资降至~9500~万元。交易商必须额外提供~300~万元现金或抵押品，才能维持相同头寸。

如果许多交易商面临相同变化，它们可能缩减头寸或出售证券释放现金。保护单个贷款人的折扣率，在全市场同时上升时，也可能推动整个体系去杠杆。

\end{example}

折扣率提高与抵押品价值下降同时发生时，抵押品压力加剧。借款人既面临每单位抵押品借款能力下降，也面临现有抵押品价值缩水。追缴保证金可能迫使其在下跌市场中出售资产，强化融资流动性与市场流动性的相互作用。

抵押品质量还影响流动性流向。压力时期，贷款人往往更愿接受政府证券，而较少愿意接受私人部门或流动性较低的抵押品。投资者同样将现金转向被视为更安全、更有流动性的请求权。这种\emph{避险转移}可能降低高质量短期政府请求权的收益率，同时扩大私人部门货币市场工具的利差。它体现流动性在资产与交易对手之间的重新配置，未必意味着总体流动性下降。

抵押品、折扣率、交易商资产负债表与市场流动性的相互作用，解释了回购市场为何影响其直接参与者之外的市场。如果交易商融资能力下降，其为债券市场提供中介服务的能力与意愿可能减弱。市场流动性下降又可能增加价格波动，降低用作抵押品的证券价值。贷款人可能进一步提高折扣率，从而收紧融资条件。
\begin{tcolorbox}[
  enhanced,
  breakable,
  colback=white,
  colframe=FinanceBlue,
  boxrule=0.7pt,
  sharp corners,
  borderline west={2.2pt}{0pt}{FinanceBlue},
  title=\textbf{金融视角：利率不变时，更高折扣率也能收紧融资},
  coltitle=white,
  colbacktitle=FinanceBlue,
  fonttitle=\bfseries,
  before skip=1.2em,
  after skip=1.2em
]

回购利率不变时，融资条件仍可能显著收紧。更高折扣率减少给定数量抵押品能够筹集的现金。借款人须投入更多自有资金、追加抵押品，或缩减头寸。因此，即使回购报价利率变化不大，折扣率也会影响可获得的融资数量。

\end{tcolorbox}

%========================================================
% Summary
%========================================================

\TailModule{summary}{本章小结}

货币市场重新分配流动性，并为短期现金流错配提供融资。如果现金流入与流出时点不同，即使基础资产和项目具有正的长期价值，借款人也可能需要短期资金。短期融资减少持有大量闲置现金的需要，但反复再融资产生滚动融资风险。

货币市场工具的契约结构不同。国库券是短期主权请求权；商业票据通常是无担保企业或金融债务；可转让存单是可转让银行负债；同业贷款在金融机构间提供无担保融资；回购则以证券为支持提供现金。报价惯例十分重要，贴现工具尤其如此；还应区分回购利率与抵押品折扣率。

货币市场架构影响定价与流动性。银行兼任借贷双方，交易商融资证券库存，货币市场基金则将投资者现金导向短期请求权。基准利率因融资是否有担保、依据交易还是报价而异。抵押品资格与估值，以及清算结算安排，共同影响使用短期请求权和证券获取流动性的能力。

贷款人拒绝续借短期融资，或大量赎回减少对货币市场借款人的资金供给时，可能出现流动性压力。回购市场中，更高折扣率减少抵押品支持的可融资金额。如果折扣率上升与抵押品价格下跌同时发生，追缴保证金和资产出售可能强化融资流动性与市场流动性的相互作用。短期限在正常条件下提供灵活性，但再融资不确定时，也可能增加脆弱性。

%========================================================
% Key Terms
%========================================================

\TailModule{keyterms}{关键术语}

货币市场（money market）；短期融资（short-term finance）；流动性（liquidity）；现金流时点（cash-flow timing）；营运资本（working capital）；财务灵活性（financial flexibility）；滚动融资风险（rollover risk）；国库券（Treasury bill）；持有期回报率（holding-period return）；贴现收益率（discount yield）；商业票据（commercial paper）；可转让存单（negotiable certificate of deposit）；同业贷款（interbank loan）；回购协议（repurchase agreement）；逆回购（reverse repo）；回购利率（repo rate）；抵押品（collateral）；折扣率（haircut）；追缴保证金（margin call）；货币市场基金（money market fund）；基准利率（benchmark rate）；交易型基准（transaction-based benchmark）；报价型基准（quotation-based benchmark）；上海银行间同业拆放利率（Shibor）；DR007；价格发现（price discovery）；券款对付（delivery versus payment）；融资挤兑（funding run）；市场冻结（market freeze）；避险转移（flight to quality）；融资流动性（funding liquidity）；市场流动性（market liquidity）。

%========================================================
% Financial AI Lab
%========================================================

\TailModule{aiprompt}{金融~AI~实验室}

\textbf{实验：重建回购交易并进行压力测试}

请~AI~系统构造一笔一周期回购，使证券交易商融资债券头寸。要求回答列明抵押品市值、折扣率、提供的现金、回购利率、日计数惯例及回购支付额。

独立审核交易。核实折扣率用于抵押品价值并决定提供的现金金额，回购利率则决定融资成本。同时检查~AI~是否正确区分正回购与逆回购视角。

随后加入两项冲击：抵押品价值下降~4\%，折扣率提高~3~个百分点。保持原有回购融资金额不变，请~AI~计算产生的融资缺口，以及恢复原融资头寸所需的额外现金或合格抵押品。用本章契约关系核验结果。

最后，请~AI~解释折扣率上升如何推动全市场去杠杆。检查回答是否错误地将抵押品视为消除了流动性风险，或假设抵押品总能立即按其市场报价出售或质押。

%========================================================
% Discussion Questions
%========================================================

\TailModule{questions}{讨论题}

\begin{enumerate}

\item 为什么盈利且有偿付能力的企业仍可能需要短期融资？区分流动性问题与偿付能力问题。

\item 比较商业票据与回购协议。有无抵押品如何影响贷款人风险暴露、融资条款和借款人的滚动融资风险？

\item 为什么即使银行体系总体流动性看似充足，无担保同业市场仍可能缺乏流动性？

\item 为什么无担保货币市场利率与有担保回购利率之间的利差可能因不同原因扩大？解释这种变化需要哪些信息？

\end{enumerate}

%========================================================
% Applied Problems
%========================================================

\TailModule{appprobs}{应用题}

\begin{enumerate}

\item \textbf{现金流时点与短期融资。}

一家企业月初持有现金人民币~400~万元。该月预计收到现金~1800~万元、支付~2700~万元。

\begin{enumerate}

\item 计算新增融资前的净现金头寸。

\item 计算短期融资需求。

\item 假设原定本月收回的~700~万元应收账款推迟至下月收回，重新计算融资需求。

\item 解释为什么这一变化本身不足以证明企业资不抵债。

\end{enumerate}

\item \textbf{货币市场报价惯例。}

一张~90~天期票据，面值为~1000，售价为~988。采用每年~360~天的计息基础。

\begin{enumerate}

\item 计算持有期回报率。

\item 计算以购买价格为基础的年化货币市场收益率。

\item 计算以面值为基础的贴现收益率。

\item 解释为什么基础现金流相同，两种年化收益率仍不同。

\end{enumerate}

\item \textbf{商业票据利差。}

一家高评级企业以~3.8\%~的年化利率发行~60~天期商业票据，此时可比短期基准利率为~2.9\%。

\begin{enumerate}

\item 计算相对于基准的利差。

\item 假设基准升至~3.4\%，利差扩大至~1.6~个百分点，新的商业票据利率为多少？

\item 将融资利率增幅分解为基准成分和利差成分。

\item 给出利差可能扩大的两个经济原因。

\end{enumerate}

\item \textbf{回购利率与折扣率。}

一家交易商持有价值人民币~2.5~亿元的证券，并以隔夜回购融资。初始折扣率为~3\%，年化回购利率为~2.2\%，按每年~365~天计息。

\begin{enumerate}

\item 计算提供的现金金额。

\item 计算次日的回购支付额。

\item 如果抵押品价值不变、折扣率升至~6\%，交易商需额外投入多少自有资金，才能维持相同资产头寸？

\item 解释为什么折扣率提高与回购利率提高相同百分点，具有不同经济影响。

\end{enumerate}

\end{enumerate}

%========================================================
% Data Application
%========================================================

\TailModule{dataapp}{数据应用}

\textbf{中国货币市场的流动性条件}

从中国货币网或中国人民银行等权威来源，获取至少两种中国短期货币市场利率的日度数据。一组可用的组合为~DR007~与七天期银行间回购利率或七天期~Shibor。如果加入央行七天期公开市场操作利率，应将其与市场决定的利率明确区分。

选择同时包括正常时期和至少一次明显短期波动的样本期间，绘制所选利率。计算市场利率之间的利差，识别利差或利率水平急剧变化的日期。

结合基础市场结构解释变化。明确各利率是否有担保、哪些机构参与、哪些抵押品合格，以及基准依据交易还是报价。检查月末、季末、缴税、节假日或其他流动性管理日期，是否与异常波动重合。

不要仅将利差变化归因于信用风险或货币政策。解释需要哪些额外证据，才能区分总体流动性、交易对手风险、抵押品条件和市场分割的变化。报告样本期间、利率定义、计算惯例及原始数据来源。

%========================================================
% Case Analysis
%========================================================

\TailModule{caseanalysis}{案例分析}

\textbf{案例：2020~年~3~月美国短期融资市场的压力}

2020~年~3~月，新冠疫情迅速蔓延，引发对现金和高流动性资产的异常需求。优质型货币市场基金（prime money market funds），尤其是面向机构投资者的此类基金，遭遇大量赎回；商业票据及其他私人部门短期工具市场承受严重压力。基金应对赎回时，减少持有和购买私人部门短期请求权，使企业和金融机构为到期义务再融资更加困难。

美联储通过多项紧急便利作出应对。2020~年~3~月~17~日宣布的商业票据融资便利（CPFF），通过特殊目的载体购买合格商业票据，提供流动性后盾。3~月~18~日宣布的货币市场共同基金流动性便利（MMLF），向合格金融机构提供贷款，以其从货币市场基金购买的合格资产作担保。这些及其他干预宣布并实施后，短期融资市场压力缓解。

运用本章框架分析该事件，追踪从货币市场基金投资者到商业票据发行人的融资链条。解释即使发行人经营现金流尚未立即恶化，投资者赎回为什么仍可能减少其可获得的融资。区分市场流动性与融资流动性，并识别滚动融资风险在哪些环节进入链条。

随后比较~CPFF~与~MMLF~的作用。各项便利针对什么市场摩擦？即使后续实际使用有限，可信后盾的存在为何仍可能改善市场运行？如果参与者预期未来短期融资压力时期会获得公共流动性支持，可能产生哪些激励问题？

\emph{原始资料来源：美国联邦储备理事会，“商业票据融资便利”，2020~年~3~月；美国联邦储备理事会，“货币市场共同基金流动性便利”，2020~年~3~月。}

%========================================================
% Further Reading
%========================================================

\TailModule{furtherread}{延伸阅读}

\begin{itemize}

\item Stigum, Marcia L., and Anthony Crescenzi (2007),
\emph{Stigum's Money Market（斯蒂格姆货币市场）}, 4th ed., McGraw-Hill.

\item Gorton, Gary, and Andrew Metrick (2012),
``Securitized Banking and the Run on Repo（证券化银行与回购挤兑）,''
\emph{Journal of Financial Economics}, 104(3), 425--451.

\item Covitz, Daniel, Nellie Liang, and Gustavo A. Suarez (2013),
``The Evolution of a Financial Crisis: Collapse of the Asset-Backed Commercial Paper Market（金融危机的演进：资产支持商业票据市场的崩溃）,''
\emph{The Journal of Finance}, 68(3), 815--848.

\item Krishnamurthy, Arvind, Stefan Nagel, and Dmitry Orlov (2014),
``Sizing Up Repo（衡量回购市场规模）,''
\emph{The Journal of Finance}, 69(6), 2381--2417.

\item Committee on the Global Financial System (2017),
\emph{Repo Market Functioning（回购市场运行）},
CGFS Papers No.~59, Bank for International Settlements.

\end{itemize}

\ChapterTail{summary,keyterms,aiprompt,questions,appprobs,dataapp,caseanalysis,furtherread}

\chapter{债券市场与债务契约}
\label{chap:Bond_Markets_and_Debt_Contracts}

\begin{openingcase}{为什么公司债券的收益率更高？}

一家大型制造企业计划发行七年期债券，筹集人民币~30~亿元。发行时，七年期政府债券收益率约为~2.2\%。虽然该企业信用记录良好、现金流稳定，投资者却不愿按相同收益率购买其债券。为完成发行，企业必须提供约~3.0\%~的收益率。

几个月后，债券在二级市场交易。企业已履行全部约定支付，经营前景也变化不大。即便如此，债券市场价格仍然下跌，收益率升至~3.4\%。

\end{openingcase}


\begin{tcolorbox}[
  enhanced,
  colback=FinanceBlue!6,
  colframe=FinanceBlue,
  boxrule=0.6pt,
  sharp corners,
  left=8pt,
  right=8pt,
  top=7pt,
  bottom=7pt,
  before skip=1em,
  after skip=1.6em,
  title=\textbf{核心问题},
  coltitle=white,
  colbacktitle=FinanceBlue,
  fonttitle=\bfseries,
  attach boxed title to top left={xshift=6pt,yshift=-2mm},
  boxed title style={
    boxrule=0pt,
    sharp corners,
    left=5pt,
    right=5pt,
    top=2pt,
    bottom=2pt
  }
]
\begin{enumerate}[
  leftmargin=*,
  label=\textcolor{FinanceBlue}{\textbf{\arabic*.}},
  itemsep=0.55em,
  topsep=0.2em
]

\item 债务的清偿顺位、抵押品、契约限制条款以及违约条款如何决定债务契约中所包含的权利？

\item 初级市场和次级市场是如何将有关发行者、利率以及流动性的信息纳入债券价格和收益率中的呢？

\item 基准利率和信用利差的变化如何影响债券价值？而久期和凸性又如何衡量债券价格对所需收益率变化的敏感性呢？

\end{enumerate}
\end{tcolorbox}

\section{债务的契约结构}

债务契约规定承诺支付，以及请求权附带的权利与限制。这些条款影响预期回收、借款人激励和债权人控制，进而影响投资者要求的回报及债券价值。

\subsection{承诺现金流、期限与偿还}

最基本的债务契约规定借款金额、到期前应支付的款项及到期偿还额。传统固定利率附息债券承诺定期支付票息，并在债券到期时偿还面值。

设 $F$ 表示面值，$c$ 表示年票面利率，$m$ 表示每年票息支付的次数。每个期间支付的票息为

\begin{equation*}
C=\frac{cF}{m}.
\label{eq:ch8-coupon-payment}
\end{equation*}

如果债券还有 $N$ 个票息期，那么其承诺的现金流序列为

\begin{equation*}
(C_t)_{t=1}^N
=
(\underbrace{C,\; C,\; \ldots,\; C}_{N-1\text{ 保留LaTeX命令、公式、交叉引用和数字。}},\; C+F).
\label{eq:ch8-promised-cash-flow-sequence}
\end{equation*}

图~\ref{fig:ch8-coupon-bond-cashflows}展示了这些承诺支付的时间安排。

\begin{figure}[H]
\centering
\begin{tikzpicture}[
  >=Latex,
  axis/.style={-Latex, line width=0.8pt},
  cash/.style={-Latex, line width=1.0pt, FinanceBlue},
  tick/.style={line width=0.6pt}
]

% Time axis with a break
\draw[line width=0.8pt] (0,0) -- (4.8,0);

\draw[axis] (5.8,0) -- (10.3,0)
  node[right]{时间};

% Omitted periods
% \node at (5.3,0) {$\cdots$};

% Tick marks
\foreach \x in {0,1.8,3.6,7.0,8.8}
{
  \draw[tick] (\x,0.12) -- (\x,0);
}

% Time labels
\node[below=4pt] at (0,0) {$t=0$};
\node[below=4pt] at (1.8,0) {$t=1$};
\node[below=4pt] at (3.6,0) {$t=2$};
\node[below=4pt] at (7.0,0) {$t=N-1$};
\node[below=4pt] at (8.8,0) {$t=N$};

% Dashed segment indicating omitted coupon periods
\draw[densely dashed, line width=0.8pt]
  (4.8,0) -- (5.8,0);


% Coupon payments
\draw[cash] (1.8,0) -- (1.8,0.65)
  node[above=3pt] {$C$};

\draw[cash] (3.6,0) -- (3.6,0.65)
  node[above=3pt] {$C$};

\draw[cash] (7.0,0) -- (7.0,0.65)
  node[above=3pt] {$C$};

% Final coupon and principal
\draw[cash] (8.8,0) -- (8.8,1.25)
  node[above=3pt] {$C+F$};

\end{tikzpicture}

\caption{传统附息债券的承诺现金流}
\label{fig:ch8-coupon-bond-cashflows}
\end{figure}

票面利率决定契约票息支付，并不等同于市场收益率。传统固定利率债券的票面利率由契约固定，而市场收益率随债券价格及投资者必要报酬率而变化。

债券契约在票息和本金偿还结构上均存在差异。零息债券不定期支付票息，附息债券则定期付息。本金可以采用到期一次还本或分期还本结构。前者在到期时偿还本金，后者则在到期前逐步还本。如果到期时仍有未偿余额，就属于部分分期还本结构，最后一笔本金支付称为尾期大额还本。

这些偿还结构使发行人的现金支付义务具有不同的时间分布。到期一次还本结构使全部本金在到期前保持未偿，但将偿还集中于到期日；分期还本要求更早支付本金，逐步降低未偿余额及到期剩余本金。

期限决定债务存续的契约时间范围。期限越长，投资者承受发行人信用状况变化的时间越长。对其他特征相近的固定利率债券而言，较长期限通常也增加其对市场利率变化的敏感性。

\begin{example}{到期一次还本与分期还本结构}{bullet-amortizing-structures}

假设有两种面值为人民币~1,000的三年期债券，其票面利率为6\%。

到期一次还本债券的利息在头两年分别为人民币~60元，而在到期时则达到人民币~1,060元。

\[
60,\qquad 60,\qquad 1,060.
\]

这类传统附息债券采用到期一次还本结构，到期前不偿还本金。“到期一次还本”指本金偿还时点，而非票息支付方式。到期时偿还全部面值的零息债券，同样采用这一本金结构。

\medskip

假设本金在每年分三次等额偿还。利息是根据每年年初的存续本金来计算的。然后进行的支付大致为

\[
393.33,\qquad 373.33,\qquad 353.34.
\]

到期一次还本结构使得本金在到期时全部偿还，而分期偿还则逐步减少本金，同时降低后续的利息支付。 

\end{example}


\subsection{清偿顺位与请求权优先次序}

债务请求权具有不同的优先次序。清偿顺位规定发行人无法足额满足全部请求权时，债权人有权获得价值分配的先后顺序。在契约条款与适用破产法约束下，优先请求权排在次级请求权之前。所有承诺支付都能履行时，优先次序的影响可能并不明显；但在困境中，它决定不同请求权如何分担短缺。

设发行人扣除相关成本及更高顺位请求权后，可供优先和次级债权人分配的价值为~$A$。优先请求权为~$D_S$，次级请求权为~$D_J$。在简化的优先顺序瀑布中，优先债权人的回收额~$R_S$~为

\begin{equation}
R_S=\min\{D_S,A\},
\label{eq:ch8-senior-recovery}
\end{equation}
而次级债权人获得的赔偿金额则有限。

\begin{equation}
R_J
=
\min\left\{D_J,\max\{A-R_S,0\}\right\}
=
\min\left\{D_J,\max\{A-D_S,0\}\right\}.
\label{eq:ch8-junior-recovery}
\end{equation}

简化的瀑布模型忽略了抵押品权利、担保、互信协议以及其他法律复杂性。

\begin{example}{简化的优先顺序瀑布}{priority-waterfall}

企业有人民币~1~亿元优先债务和~6000~万元次级债务。如果可供这些债权人分配的价值为~1.3~亿元，优先债权人回收~1~亿元，次级债权人获得剩余~3000~万元，次级债务回收率为~50\%。

如果能够获得8000万元人民币的资金，那么高级债权人将获得8000万元人民币，而次级债权人则一无所获。

\end{example}

优先次序也可能来自公司结构。母公司发行的债务，可能在结构上次于经营性子公司的债务，因为子公司剩余价值分配给母公司之前，子公司债权人优先对其资产享有请求权。担保使债权人对集团内另一实体拥有直接请求权，可以减轻这种结构性次级地位。

到期时间相近的债券，在优先级不同的情况下，其预期回收价值和数值可能会有所不同。

\subsection{抵押品与债权人保护}

抵押品赋予债权人对特定资产的担保物权，借款人未履行契约义务时可依法执行。因此，有担保债权人既拥有对借款人的支付请求权，也拥有对指定抵押品的权利，但须以担保物权有效且可执行为前提。

抵押品的保护程度取决于其可变现价值、流动性以及法律强制执行力。如果债务请求权为$D$，而抵押品的可变现价值为$V_C$，那么简单的抵押品覆盖比率就是

\begin{equation}
CC=\frac{V_C}{D}.
\label{eq:ch8-collateral-coverage}
\end{equation}

较高的抵押品覆盖率可能会提高回收的预期，但并不能确保全额偿还。在违约之前，抵押品的价值可能会下降，相互竞争的请求权可能会降低债权人可获得的价值，而清算或执行过程可能会非常昂贵。

抵押品也会影响借款人的激励因素。抵押给债权人的资产不易被转让或重新抵押，而在违约时这些资产的损失会降低那些增加债权人风险的行为的吸引力。这些保护措施的有效性也取决于证券协议的条款以及适用的法律框架。

不同的资产提供不同程度的保护。可交易的证券通常比专用设备、房地产或应收账款更容易变现，即使它们的当前评估价值相似。

债券持有人也可以从担保中受益。担保使债权人能够对担保人享有请求权，而不是对特定资产的担保物权。抵押品和担保通过不同的法律机制加强债权人的保护。

\begin{tcolorbox}[
  enhanced,
  breakable,
  colback=white,
  colframe=FinanceBlue,
  boxrule=0.7pt,
  sharp corners,
  borderline west={2.2pt}{0pt}{FinanceBlue},
  title=\textbf{财务展望：安全债券并非无风险之选},
  coltitle=white,
  colbacktitle=FinanceBlue,
  fonttitle=\bfseries,
  before skip=1.2em,
  after skip=1.2em
]

有担保债券并非无风险债券。违约发生时，抵押品可能改善预期回收，却不能确保承诺支付得到履行。因此，信用分析既考虑违约概率，也考虑债权人在违约时可能承担的损失。

\end{tcolorbox}

\subsection{债券契约与持有人权利}

债券契约规定了管理发行人的财务政策及企业行为的契约限制条款。一般条款限制了额外的借款、资产处置、向股东的分配，或企业控制的变更。其他条款则要求发行人提供财务信息或维持特定的财务状况。这些条款定义了债券持有人可以获得的合同保护。

契约限制条款的违反并不必然构成立即的违约事件。其后果取决于契约限制条款的条款规定，包括通知和补救条款。如果违反行为符合违约事件的构成要件，债券持有人可以获得诸如加速偿还或执行担保物权等权利（如适用）。此外，违反契约限制条款还可能导致关于放弃或部分修改债券契约的谈判。

在广泛发行的债券中，由于债券持有人的分散性以及集体行动的问题，行使这些权利变得非常复杂。通常，受托人会代表债券持有人进行监控、与发行方沟通，并管理合同规定的救济措施。契约规定了受托人的权力范围，以及进行集体决策所需的投票门槛。

\subsection{违约、重组与回收}

发行人未履行契约义务，或发生债务协议规定的违约事件时，即构成违约。未支付利息或本金是最明确的情形。其他违约事件可能包括特定契约限制条款被违反、交叉违约或破产事件。

抵押品本身并不能决定债务请求权的损失。债权人可以通过持续的支付、抵押品的变现、重组或清算来收回价值。$R$用“恢复的值”来表示这个数值。$D$违约时的合同请求权。回收率为

\begin{equation}
RR=\frac{R}{D}.
\label{eq:ch8-recovery-rate}
\end{equation}
而违约损失率则是

\begin{equation}
LGD=1-RR.
\label{eq:ch8-lgd}
\end{equation}

在指定的时间范围内，一个简化的预期损失衡量方法如下：

\begin{equation}
EL=PD\times LGD\times EAD,
\label{eq:ch8-expected-credit-loss}
\end{equation}
其中，$PD$表示违约概率，$EAD$表示违约时的敞口。这个表达式将违约的可能性与违约条件下的损失区分开来。

债务重组会改变原有的债务契约。它可能包括延长到期时间、减少票息或本金支付、用新证券替换现有债务，或将债务转换为权益。在清算会破坏企业价值的情况下，债务重组可以保护企业的持续经营价值。

重组的结果受到优先级、抵押品、法律规则以及发行方与其债权人之间以及不同债权类别之间的协商影响。分散的债券持有人也可能面临协调问题。通过规定必须由特定多数票批准变更，可以减少这些问题。

回收金额既反映了在清偿过程中保留的价值，也体现了请求权的合同地位。重组或清算会影响债权人可获得的价值，而清偿顺位、抵押品、担保以及法律规则则决定了这些价值如何在请求权之间进行分配。同一家公司发行的债务请求权可能会有不同的回收率。

\begin{example}{违约概率相同，预期损失不同}{default-and-recovery}

假设有两个一年的债务请求权，每个请求权的暴露度为100，违约概率为3\%。请求权A在违约时的预期回收率为70\%，而请求权B的预期回收率为20\%。

使用简化的预期损失度量方法。

\[
EL_A=0.03\times(1-0.70)\times100=0.9,
\]

而

\[
EL_B=0.03\times(1-0.20)\times100=2.4.
\]

这两个请求权具有相同的违约概率，但由于它们的预期恢复金额不同，因此预期损失也有所不同。

\end{example}

\section{债券市场架构与价格发现}

债券契约规定了承诺的现金流和债权，但市场价格是通过投资者之间的交易形成的。发行、交易、清算和结算的组织方式影响着寻找交易对手、执行交易以及转移所有权的成本。这些摩擦因素影响市场流动性，以及新信息在债券价格和收益率中的体现程度。

\subsection{发行人、投资者与市场参与者}

债券发行者在融资目的、偿还能力和法律结构方面存在差异。政府发行债务以资助公共支出和管理公共负债。金融机构和非金融企业则利用债券来获得期限性资金、重新融资现有债务或进行投资融资。不同的偿还来源和资产负债表结构使得这些请求权具有不同的信用特征。

投资者的负债、投资期限、流动性需求以及监管限制也存在差异。银行可能会持有债券用于投资、提高市场流动性或作为抵押品。保险公司和养老基金通常会持有久期较长的资产来应对长期负债，而资产管理人在执行投资组合任务时，可能会更加重视市场流动性。

这些差异导致了不同到期时间、信用等级和债券类型的多样化需求。某些证券更适合特定投资者的资产负债表需求或投资限制。因此，投资者群体的构成会影响发行者能够在市场上发行的债券种类。

专业机构支持发行人与投资者之间的连接。承销商和安排人协助新证券的发行，评级机构评估信用质量，交易商支持二级市场的交易，而保管人、托管人和结算系统则负责证券的持有与转移。它们的主要职责在表~\ref{tab:ch8-market-participants}中有所总结。

\begin{table}[H]
\centering
\caption{债券市场的主要参与者与基础设施}
\label{tab:ch8-market-participants}

\begin{tabularx}{\textwidth}{@{}p{3.1cm}Y Y@{}}
\toprule
\textbf{参与者} &
\textbf{主要角色} &
\textbf{相关的经济问题} \\
\midrule

政府和公共发行机构  
& 增加公共债务融资  
& 财政能力、债务可持续性、期限结构  
\\

企业和金融机构  
& 筹集债务资金  
& 资金需求、杠杆使用、再融资、信用质量 \\

银行、保险公司、养老基金和资产管理者&投资并交易债券&久期、流动性、信用风险、监管\\

承销商和安排人
& 制定结构、承销并发行债券
& 尽职调查、定价、投资者需求、发行风险 \\

信用评级机构  
& 评估发行机构和项目的信用风险  
& 评估评级质量、评级方法、相关信息以及利益冲突情况  
\\

交易商
& 进行市场交易并支持二级市场交易
& 交易能力、库存风险、融资、市场流动性\\

邦德信托人&代表债券持有人并管理合同权利&契约限制条款的遵守、集体行动、执行\\

证券保管处、托管机构和结算系统
& 负责保管证券记录，保护资产安全，并协助进行结算
& 保管功能、记录准确性、结算的确定性，以及运营风险管理 \\

\bottomrule
\end{tabularx}
\end{table}

\subsection{债券一级市场与二级市场}

一级市场是发行新债券并向投资者筹集资金的地方。二级市场则是已发行的存续债券进行交易的地方。这种区分将融资交易与随后产生的请求权分离开来。

政府债券通常通过拍卖方式发行，而企业和金融债券的发行则常采用承销商、联合承销团或簿记建档的方式。不同司法管辖区和金融工具之间的发行机制各不相同，但它们都解决了一个共同的问题：即如何使新发行的债券的数量、到期日、合同条款以及定价与投资者的需求相匹配。

一级市场发行也汇集信息。拍卖、订单簿和投资者意向揭示不同收益率或价格下的需求。这些信息有助于确定发行能够完成的条款，以及债券如何在投资者之间分配。发行价格为新请求权提供初始市场估值。

一旦债券发行完毕，投资者可以在二级市场上进行转让，而发行方的合同义务并不会因此发生变化。这样一来，发行方的融资期限与单个投资者的持有期限便可以分离开来。即使某些投资者只持有十年期的债券，该债券仍然可以提供十年的融资支持。

\begin{figure}[H]
\centering

\begin{tikzpicture}[
  >=Latex,
  node/.style={
    align=center,
    font=\normalsize
  },
  flow/.style={
    Latex-Latex,
    line width=0.8pt,
    draw=FinanceBlue
  },
  info/.style={
    Latex-Latex,
    dashed,
    line width=0.7pt,
    draw=FinanceBlue
  },
  infra/.style={
    draw=FinanceBlue,
    line width=0.7pt,
    rounded corners=2pt,
    align=center,
    minimum width=10.6cm,
    minimum height=1.0cm,
    inner sep=6pt
  }
]

% Primary market
\node[node] (issuer) at (0,2.0) {发行方};
\node[node] (primary) at (3.5,2.0) {一级市场};
\node[node] (investors1) at (7.0,2.0) {投资者};

\draw[flow] (issuer) -- (primary);
\draw[flow] (primary) -- (investors1);

% Market information
\node[node, text=FinanceBlue] (info) at (5.85,1.25)
{市场信息};

% Secondary market
\node[node] (investors2) at (0,0.45) {投资者};
\node[node] (secondary) at (3.5,0.45) {二级市场};
\node[node] (investors3) at (7.0,0.45) {投资者};

\draw[flow] (investors2) -- (secondary);
\draw[flow] (secondary) -- (investors3);

\draw[info] (primary.south) --  (secondary.north);

% Infrastructure
\node[
  align=center,
  minimum width=10.6cm,
  minimum height=1.0cm,
  inner sep=6pt
] (infra) at (4.0,-1.)
{存储库 \quad 保管人 \quad 清算 \quad 结算};

\draw[
  dashed,
  draw=FinanceBlue,
  line width=0.7pt
]
([yshift=0.55cm]infra.west) -- ([yshift=0.55cm]infra.east);

\draw[
  dashed,
  draw=FinanceBlue,
  line width=0.7pt
]
([yshift=-0.55cm]infra.west) -- ([yshift=-0.55cm]infra.east);

\node[
  align=center,
  text=FinanceBlue
] at (4.0,-2.15)
{同时支持两个市场。};

\end{tikzpicture}

\caption{债券一级市场与二级市场}
\label{fig:ch8-primary-secondary-markets}
\end{figure}

二级市场交易为存续的债券产生了可观测的价格和收益。这些价格包含了关于利率、信用状况以及投资者需求的不断变化的信息。它们为现有债券提供了估值参考，并有助于确定后续发行的定价和条款。例如，二级市场收益率的上升可能会提高重新进入市场的发行者的融资成本。

\subsection{交易商、交易与市场流动性}

许多债券市场在很大程度上依赖交易商的中介服务。投资者可能会向交易商请求报价，而不是将每个订单提交给中央订单簿。交易商可以从一位投资者处购买债券，然后在另一位投资者处出售这些债券，在两次交易之间持有这些债券。

交易商的买入价是指交易商愿意购买的价格，而卖出价则是指交易商愿意出售的价格。买卖价差反映了交易成本、库存风险、逆向选择问题，以及资金使用和资产负债表状况等因素的影响。

市场流动性有多个维度。紧密度指的是执行交易的成本，这通常体现在买卖价差中。深度则指的是在不引起显著价格变动的情况下可以交易的金额。韧性则是指在大额交易或市场动荡之后，价格和交易条件恢复的速度。

不同债券的流动性各不相同。近期发行的政府基准债券可能会经常以大额交易方式成交，而较旧或规模较小的公司债券则可能只偶尔有交易机会。因此，由同一借款人发行的债券在流动性方面也会存在显著差异。

交易商的资产负债表影响着流动性的提供。如果交易商在找到其他买家之前就收购了某只债券，那么他们就必须为此承担融资风险，并承受价格下跌带来的损失。较高的融资成本、更大的价格波动性或更严格的资产负债表限制可能会降低交易商愿意持有的库存量，从而使中介活动变得更加昂贵。

电子交易改变了这一过程的组织方式。询价系统可以促进交易商之间的竞争，而其他平台则能够自动化部分报价和执行流程。交易后的透明度提供了关于最近交易的信息，不过更高的透明度也可能影响交易商愿意处理大型或复杂交易的程度。

\begin{figure}[H]
\centering

\begin{tikzpicture}[
  >=Latex,
  box/.style={
    draw=FinanceBlue,
    line width=0.7pt,
    rounded corners=2pt,
    align=center,
    text width=2.8cm,
    minimum height=0.9cm,
    inner sep=5pt
  },
  cash/.style={
    -Latex,
    line width=0.9pt,
    draw=FinanceBlue
  },
  bond/.style={
    -Latex,
    dashed,
    line width=0.9pt,
    draw=FinanceBlue
  },
  label/.style={
    font=\small,
    fill=white,
    inner sep=1.5pt
  }
]

% Participants
\node[box] (seller) at (-4.8,0)
  {投资者\\希望出售};

\node[box] (dealer) at (0,0)
  {交易商};

\node[box] (buyer) at (4.8,0)
  {投资者\\希望购买};

% Seller to dealer: bond transfer
\draw[bond]
  ([yshift=6pt]seller.east) --
  node[label, above=3pt] {债券}
  ([yshift=6pt]dealer.west);

% Dealer to seller: cash payment
\draw[cash]
  ([yshift=-6pt]dealer.west) --
  node[label, below=3pt] {现金}
  ([yshift=-6pt]seller.east);

% Dealer to buyer: bond transfer
\draw[bond]
  ([yshift=6pt]dealer.east) --
  node[label, above=3pt] {债券}
  ([yshift=6pt]buyer.west);

% Buyer to dealer: cash payment
\draw[cash]
  ([yshift=-6pt]buyer.west) --
  node[label, below=3pt] {现金}
  ([yshift=-6pt]dealer.east);

\end{tikzpicture}

\caption{债券二级市场交易中的交易商中介}
\label{fig:ch8-dealer-intermediation}
\end{figure}

作为本金的交易商，会暂时持有债券库存，并承担与库存和资金相关的风险。买卖价格之间的差额构成了交易商的交易收入。

\begin{tcolorbox}[
  enhanced,
  breakable,
  colback=white,
  colframe=FinanceBlue,
  boxrule=0.7pt,
  sharp corners,
  borderline west={2.2pt}{0pt}{FinanceBlue},
  title=\textbf{财务前景：低信用风险并不意味着高流动性。},
  coltitle=white,
  colbacktitle=FinanceBlue,
  fonttitle=\bfseries,
  before skip=1.2em,
  after skip=1.2em
]

债券可能具有较低的违约风险，但交易起来仍然很困难或成本高昂。信用风险涉及因未能履行合同付款而产生的损失，而市场流动性则涉及到在不产生重大成本或价格影响的情况下交易请求权的能力。在压力之下，这两种风险可能会相互作用，但它们仍然是不同维度的风险。

\end{tcolorbox}

\section{债务契约的估值与风险}

债券估值将合同中的现金流与市场贴现率联系起来。对于没有违约风险的无期权债券而言，其价值等于承诺的票息和本金的现值。基准利率的变化会影响这些现金流的贴现方式，而信用风险则会影响预期支付金额以及投资者要求的利差。

\subsection{债券价格、收益率与契约现金流}

考虑一种债券，其在期限$t$结束时承诺的现金流为$CF_t$，其中$t=1,\ldots,T$。如果使用连续复利的收益率$y$来贴现所有承诺的现金流，那么债券的价格就是

\begin{equation}
P(y)
=
\sum_{t=1}^{T}
CF_t e^{-yt}.
\label{eq:ch8-bond-price}
\end{equation}

对于普通的债券来说，$CF_t=C$对于…$t=1,\ldots,T-1$同时保留LaTeX命令、公式表达、交叉引用以及数字格式。$CF_T=C+F$.

到期收益率（YTM）是使债券承诺现金流现值等于观察到的市场价格的单一贴现率。它未必等于投资者的已实现回报。已实现回报取决于承诺支付是否履行、期间票息的再投资利率，以及在到期前出售债券时获得的售价。

对于无违约风险的债券，每个现金流都可以按照其到期日对应的即期利率进行贴现。用$s_t$表示在日期$t$收到的现金流的连续复利即期利率。那么

\begin{equation}
P
=
\sum_{t=1}^{T}
CF_t e^{-s_t t}.
\label{eq:ch8-bond-price-spot-rates}
\end{equation}

对于具有固定正现金流的普通债券来说，价格与收益率呈负相关关系。对方程~\ref{eq:ch8-bond-price}关于$y$求导，可以得到……

\begin{equation}
\frac{dP}{dy}
=
-\sum_{t=1}^{T}
tCF_t e^{-yt}
<0.
\label{eq:ch8-price-yield-negative}
\end{equation}

这种反向关系还决定了传统固定利率债券是以平价交易、溢价交易还是折价交易。当票息和到期收益率采用相同的周期表示时，如果票面利率等于到期收益率，则债券以平价定价；如果票面利率高于到期收益率，则以溢价定价；如果票面利率低于到期收益率，则以折价定价。

\begin{example}{按收益率为附息债券定价}{coupon-bond-price-yield}

一张面值为100的五年期债券的票息为5\%。在连续复利的4\%收益率下，其价格是多少？

\[
P
=
5e^{-0.04\times1}
+
5e^{-0.04\times2}
+
5e^{-0.04\times3}
+
5e^{-0.04\times4}
+
105e^{-0.04\times5}
=
104.08.
\]

该债券的交易价格高于面值。用$y_{\mathrm{par}}$表示以面值定价时该债券的连续复利收益率。它满足

\[
100
=
5e^{-y_{\mathrm{par}}}
+
5e^{-2y_{\mathrm{par}}}
+
5e^{-3y_{\mathrm{par}}}
+
5e^{-4y_{\mathrm{par}}}
+
105e^{-5y_{\mathrm{par}}}.
\]

对于该债券，

\[
y_{\mathrm{par}}
=
\ln(1.05)
\approx
4.879\%.
\]

假设连续复合的收益率为4\%，这一数值低于正常水平，这与价格高于100美元的情况是一致的。如果收益率高于$y_{\mathrm{par}}$水平，那么价格就会低于正常水平。

\end{example}

图~\ref{fig:ch8-price-yield-curve}展示了债券价格与收益率之间的反比关系，其中包含了上述计算得出的两个数值。


\begin{figure}[H]
\centering

\begin{tikzpicture}
\begin{axis}[
  width=0.82\textwidth,
  height=6.0cm,
  xmin=0.5,
  xmax=10.5,
  ymin=75,
  ymax=125,
  xlabel={连续复合收益（\%）},
  ylabel={债券价格},
  axis lines=box,
  grid=none,
  tick align=outside
]

\addplot[
  FinanceBlue,
  thick,
  domain=0.5:10.5,
  samples=200
]
{
  5*exp(-x/100)
  +5*exp(-2*x/100)
  +5*exp(-3*x/100)
  +5*exp(-4*x/100)
  +105*exp(-5*x/100)
};

\addplot[
  only marks,
  mark=*,
  mark size=1.8pt
]
coordinates {(4,104.08) (4.879,100)};

% dashed guide lines for (4, 104.08)
\addplot[
  dashed,
  FinanceBlue,
  line width=0.6pt
]
coordinates {(4,75) (4,104.08)};

\addplot[
  dashed,
  FinanceBlue,
  line width=0.6pt
]
coordinates {(0.5,104.08) (4,104.08)};

% dashed guide lines for (4.879, 100)
\addplot[
  dashed,
  FinanceBlue,
  line width=0.6pt
]
coordinates {(4.879,75) (4.879,100)};

\addplot[
  dashed,
  FinanceBlue,
  line width=0.6pt
]
coordinates {(0.5,100) (4.879,100)};

\node[
  anchor=south west,
  font=\small
] at (axis cs:4,104.08)
{$(4\%,\,104.08)$};

\node[
  anchor=west,
  font=\small
] at (axis cs:4.879,100)
{$(4.879\%,\,100)$};

\end{axis}
\end{tikzpicture}

\caption{连续复利下五年期、票面利率为5\%的债券价格—收益率关系}
\label{fig:ch8-price-yield-curve}
\end{figure}

\subsection{利率风险、久期与凸性}

利率风险是指市场收益率的变化会改变固定收益请求权价值的风险。对于具有固定合同现金流的债券而言，收益率上升会降低现值，而收益率下降则会增加现值。价格反应的大小取决于债券的现金流发生的时间：当较大部分的现值与后期支付相关时，敏感度通常会更高。

久期是衡量这种敏感性的一个重要指标。它概括了债券在当前时刻的现金流的时间安排。对于以收益率 $y$ 定价、且在时间 $t$ 时有现金流 $CF_t$ 的债券来说，其久期就是

\begin{equation}
D=\frac{1}{P(y)}
\sum_{t=1}^{T}
t\,CF_t e^{-yt}.
\label{eq:ch8-macaulay-duration}
\end{equation}

权重~$w_t = CF_t e^{-yt}/P(y)$~表示每笔现金流在债券现值中的占比。因此，久期是以现值为权重的债券现金流收取时间的加权平均。

零息债券全部现金流在到期时取得，因此久期等于期限。附息债券较早支付票息，使部分现值权重分布于到期日之前，因而久期短于期限。在期限和收益率不变时，较高票息通常会降低久期。

久期衡量的是债券价格对连续复合收益率变化的第一阶敏感性。对方程~(\ref{eq:ch8-bond-price})关于$y$求导，可以得到……

\begin{equation}
\frac{dP(y)}{dy}=-\sum_{t=1}^{T}
tCF_t e^{-yt}.
\label{eq:ch8-price-yield-derivative}
\end{equation}

\begin{equation}
D=
\frac{1}{P(y)}
\sum_{t=1}^{T}
tCF_t e^{-yt}
=
-\frac{1}{P(y)}
\frac{dP(y)}{dy}
.
\label{eq:ch8-duration}
\end{equation}

于是， 

\begin{equation}
\frac{\Delta P(y)}{P(y)}
\approx
-D \Delta y.
\label{eq:ch8-duration-price-change}
\end{equation}

\begin{example}{用久期近似债券价格变化}{duration-price-approximation}

再次考虑例~\ref{ex:coupon-bond-price-yield}中的5年期、票息率为5\%的债券。在连续复利的收益率为4\%的情况下，其价格约为104.08，而其久期则约为4.556。

假设收益率增加了50个基点，那么……

\[
\Delta y=0.005.
\]

使用久期进行计算。

\[
\frac{\Delta P}{P}
\approx
-D\Delta y
=
-4.556(0.005)
=
-2.278\%.
\]

新的债券价格约为

\[
P_{\mathrm{new}}
\approx
104.08\times(1-0.02278)
=
101.71.
\]

直接重新定价得出的连续复利收益率为4.5\%，即

\[
P
=
5e^{-0.045\times1}
+
5e^{-0.045\times2}
+
5e^{-0.045\times3}
+
5e^{-0.045\times4}
+
105e^{-0.045\times5}
=
101.74.
\]

基于久期的估计值为101.71，而精确的重新定价价值为101.74。

\end{example}

债券收益率也可能以每年超过一次的方式计算复利。假设$y$是名义年收益率，每年复利$n$次。对于在每年$t$时刻收到的现金流，债券的价格就是

\begin{equation}
P(y)
=
\sum_{t}
\frac{CF_t}{(1+y/n)^{nt}}.
\label{eq:ch8-bond-price-n-compounding}
\end{equation}

相应的修正久期为 
\begin{equation}
D_{\mathrm{mod}}
=
\frac{D}{1+y/n}.
\label{eq:ch8-modified-duration-relation}
\end{equation}

因此，修正久期是指与离散复利的收益率相关的价格敏感性。在连续复利的情况下，之前定义的久期直接衡量了这种敏感性。

久期提供了一种对价格-收益率关系的局部线性近似。当收益率的变化变大时，这种关系的凸性就变得更为重要，久期的近似也就不再那么准确了。凸性能够捕捉到这种二阶敏感性。在连续复利的情况下，

\begin{equation}
\mathrm{C}
=
\frac{1}{P(y)}
\frac{d^2P(y)}{dy^2}
=
\frac{1}{P(y)}
\sum_{t=1}^{T}
t^2 CF_t e^{-yt}.
\label{eq:ch8-convexity}
\end{equation}

考虑到凸性，可以得到二阶近似

\begin{equation}
\frac{\Delta P}{P}
\approx
-D\,\Delta y
+
\frac{1}{2}\mathrm{C}(\Delta y)^2.
\label{eq:ch8-duration-convexity-approx}
\end{equation}

凸性项随着收益率的变化而提高久期近似值的准确性。

对于现金流固定且为正的普通债券，凸性为正。收益率下降引起的价格涨幅，高于仅用久期预测的涨幅；收益率同幅上升引起的价格跌幅，则小于仅用久期预测的跌幅。

\begin{example}{作为风险度量的久期与凸性}{duration-convexity-risk}

考虑例~\ref{ex:coupon-bond-price-yield}中的5年期、票息率为5\%的债券，其连续复利的收益率为4\%。该债券的价格约为104.08，其久期为4.556。 

对于收益率增加50个基点的情况，久期近似法给出的结果如下：

\[
\frac{\Delta P}{P}
\approx
-D\Delta y
=
-4.556\times 0.005
=
-2.278\%.
\]

该键的凸性约为21.911。

包含二阶项后，结果如下：LaTeX命令、公式、交叉引用以及数字都得到了正确的表示。

\[
\frac{\Delta P}{P}
\approx
-2.278\%
+
\frac{1}{2}\times 21.911\times (0.005)^2
=
-2.251\%.
\]

\[
P_{\mathrm{new}}
\approx
104.08\times(1-0.02251)
=
101.73.
\]

在连续复利收益率为4.5\%的情况下，直接久期定价得出的价格约为101.74。而经过凸性调整后的估计值则更接近这个数值，比仅基于久期的估计值更为准确。

\end{example}

\begin{figure}[H]
\centering

\begin{tikzpicture}
\begin{axis}[
  width=0.7\textwidth,
  height=7.0cm,
  xmin=-2.2,
  xmax=2.2,
  ymin=-18,
  ymax=22,
  axis lines=middle,
  xlabel={$\Delta y$},
  ylabel={$\Delta P$},
  xtick=\empty,
  ytick=\empty,
  tick align=outside,
  enlargelimits=false,
  clip=false,
 legend style={
  at={(1.5,0.97)},
  anchor=north east,
  draw=none,
  fill=none,
  font=\small
},
  legend cell align=left
]

% Exact price-yield relation
\addplot[
  FinanceBlue,
  very thick,
  domain=-2:2,
  samples=200
]
{-8*x + 2*x^2 - 0.12*x^3};
\addlegendentry{确切的价格与收益率关系}

% Duration approximation
\addplot[
  black,
  dashed,
  thick,
  domain=-2:2,
  samples=2
]
{-8*x};
\addlegendentry{久期近似}

% Duration-convexity approximation
\addplot[
  black,
  dotted,
  thick,
  domain=-2:2,
  samples=200
]
{-8*x + 2*x^2};
\addlegendentry{久期-凸性近似}

% Origin
\addplot[
  only marks,
  mark=*,
  mark size=1.7pt
]
coordinates {(0,0)};

\end{axis}
\end{tikzpicture}

\caption{久期与凸性近似债券价格变化的示意图}
\label{fig:ch8-duration-convexity-concept}
\end{figure}

图~\ref{fig:ch8-duration-convexity-concept}展示了久期所提供的局部线性近似，以及凸性所实现的曲率调整。

久期也有其经济含义。期限较长的债券会将更多的现值分配到后期的现金流中，因此对于贴现率的变化更为敏感。这与一般的估值原则是一致的，即距离较远的现金流对贴现率的敏感度更高。

久期和凸性并不能完全反映债券风险的所有方面。信用利差可能会独立于基准利率发生变化，收益率曲线不一定呈平行趋势，而嵌入式期权可能会在利率变动时改变现金流。实际上，这些指标通常会辅以其他衡量曲线风险、利差风险和期权风险的指标。

\subsection{债券定价中的信用利差}

信用等级较低的债券，其收益率通常高于同类到期日债券或曲线中的基准债券。当两种债券的收益率以相同的基础进行衡量时，它们之间的差异就被称为收益率差：

\begin{equation}
spread
=
y_{\mathrm{credit}}
-
y_{\mathrm{benchmark}}.
\label{eq:ch8-credit-spread}
\end{equation}

观察到的利差可能反映预期违约损失、在不利状态下承担损失的补偿、市场流动性、不确定性、税收待遇、嵌入式期权，以及其他工具特征。不应直接将其解释为违约概率。

在估值过程中，还可以通过将某种 spread 加入到基准利率曲线中来进行处理。用 $r_t$ 表示期限为 $t$ 的连续复利基准利率，而 $z$ 则表示附加在每个期限上的固定 spread。债券的价格 accordingly 为……

\begin{equation}
P(z)
=
\sum_{t=1}^{T}
CF_t e^{-(r_t+z)t}.
\label{eq:ch8-spread-price}
\end{equation}

使承诺现金流现值等于观察到的债券价格的常数~$z$，通常称为~\emph{Z利差}。在假定各期限利差保持相同的条件下，它概括了相对于基准即期利率曲线所要求的额外收益率。包含嵌入式期权或更复杂或有现金流的债券，需要进一步处理。

即使基准利率保持不变，信用利差的变化仍可能导致债券价格发生变化。假设某企业债券的收益率从5.0上升到6.0\%，这完全是由于其信用利差增加了1个百分点所导致的。在基准利率不变的情况下，较高的信用利差会降低该债券承诺的现金流的现值。

当基准利率上升时，债券的收益率也可能上升，而其信用利差保持不变。虽然由此导致的价格下降看起来类似，但其经济根源却不同。因此，在分析债券回报和风险时，区分基准利率变动与信用利差变动是非常重要的。

契约特征也会影响信用利差。清偿顺位、抵押品和担保措施可以提高预期回收能力，而契约限制条款则能够限制那些削弱债权人保护的行为。到期日的变化会改变承担信用风险的时间范围。因此，由同一家公司发行的债券，当其契约条款不同时，其信用利差也会有所不同。

流动性也可能造成这些差异。与信用特征相近但流动性更好的债券相比，交易稀少或立即出售需作重大价格让步的债券，可能具有更宽利差。市场压力时期，信用与流动性效应可能共同变化，难以区分。

\begin{example}{区分基准利率变动与利差变动}{benchmark-spread-decomposition}

五年期企业债券的收益率最初为5.2\%，而五年期基准收益率则为3.4\%。这两种收益率都是基于相同的标准来计算的。该债券的利差为……

\[
5.2\%-3.4\%=1.8\%,
\]

或者180个基点。

一个月后，该企业的收益率达到了6.1\%，而基准收益率则为3.6\%。新的利差为250个基点。该企业收益率上升了90个基点，其中20个基点是由于基准收益率的上升导致的，而70个基点则是因为利差扩大所致。

因此，企业收益率的变化包含两个经济上不同的因素：基准利率的变动和信用利差的变动。

\end{example}

%========================================================
% Summary
%========================================================

\TailModule{summary}{本章小结}

债券是一种债务契约，它规定了承诺的现金流、到期日、清偿顺位以及其他债权人的权利。票息和本金偿还结构决定了承诺支付的时机，而清偿顺位和抵押品则影响在困境中损失的分配方式。契约限制条款对特定行为进行了约束或要求，而违约条款则规定了债权人可以在何时行使额外权利。回收情况则取决于合同中的清偿顺位、抵押品、担保、重组情况以及法律和制度环境。

债券市场通过一级市场和二级市场将发行人与投资者联系起来。一级市场的发行决定了新债务的出售条款，而二级市场的交易则允许投资者转让现有的请求权，而不改变发行人的合同义务。交易商通过在买卖双方之间起到中介作用来支持交易和流动性，而交易平台则有助于报价和成交。交易商的介入受到库存风险、融资条件、信息以及资产负债表能力的影响。

债券估值从契约现金流出发。到期收益率是使承诺现金流现值等于观察价格的单一贴即期利率，即期利率估值则对各期限采用相应贴即期利率。现金流固定且为正时，价格与收益率反向变化。久期为价格对收益率变化的敏感性提供一阶度量，凸性则刻画价格—收益率关系的曲率，改善较大变化下的近似。

信用利差衡量的是与基准利率相比，信用风险较高的债券的额外收益。不过，信用利差并非纯粹用于衡量违约概率的指标。预期违约损失、对不良状况下损失的补偿、流动性、契约条款以及市场条件等因素都会对信用利差产生影响。因此，债券价格既反映了债务契约中规定的权利，也反映了市场对这些权利的重新估值。

%========================================================
% Key Terms
%========================================================

\TailModule{keyterms}{关键术语}

债券（Bond）；面值（face value）；票面利率（coupon rate）；票息支付（coupon payment）；期限（maturity）；到期一次还本债券（bullet bond）；分期还本债券（amortizing bond）；零息债券（zero-coupon bond）；清偿顺位（seniority）；次级债务（subordinated debt）；抵押品（collateral）；担保（guarantee）；契约限制条款（covenant）；控制权（control right）；违约（default）；重组（restructuring）；回收率（recovery rate）；违约损失率（loss given default）；债券一级市场（primary bond market）；债券二级市场（secondary bond market）；交易商（dealer）；买卖价差（bid--ask spread）；市场流动性（market liquidity）；到期收益率（yield to maturity）；即期利率估值（spot-rate valuation）；溢价债券（premium bond）；折价债券（discount bond）；久期（duration）；修正久期（modified duration）；凸性（convexity）；基准收益率（benchmark yield）；信用利差（credit spread）；Z利差（Z-spread）；流动性溢价（liquidity premium）。

%========================================================
% Financial AI Lab
%========================================================

\TailModule{aiprompt}{金融~AI~实验室}

\textbf{实验室：审核债券估值及风险报告}

选择一个活跃交易、无期权的政府或企业债券，该债券的面值、票面利率、票息频率、到期日期和市场价格可以通过权威渠道获得。请一个AI系统计算该债券的到期收益率、久期和凸性，并说明所使用的收益率和复利计算方式。

将报告分为四个阶段进行审核。

\begin{enumerate}

\item \textbf{确认合同的有效性。}

请检查票面利率、支付频率、到期日、面值，以及该债券是否包含任何嵌入式期权。确认报价市场价格是干净的还是不干净的，并明确用于计算应计利息的惯例。

\item \textbf{重现该估值结果。}

独立构建合同中的现金流，并验证以报告中的收益率进行贴现后，结算价格能够保持一致。同时检查现金流的时间安排、复利计算原则、计量单位、票息频率以及应计利息的处理方式。

\item \textbf{强调屈服点的重要性。}

将收益率增加或减少50个基点。计算重新定价的债券价值，并首先将它们与久期近似值进行比较，然后与久期-凸性近似值进行比较。解释近似误差的方向和大小。

\item \textbf{区分基准测试和市场风险。}

如果选择的是信用风险较高的债券，将其收益率与类似到期时间的适当基准进行比较。请让AI系统解释所观察到的利差。确定其中哪些部分是由数据直接支持的，而哪些部分需要更多关于违约风险、预期回收金额、流动性、合同条款或市场条件的证据来支持。

\end{enumerate}

最终报告中应至少提及一个例子，即尽管计算结果是准确的，但最终得出的结论却具有误导性。

%========================================================
% Discussion Questions
%========================================================

\TailModule{questions}{讨论题}

\begin{enumerate}

\item 为什么同一家公司发行的两种债券在到期日却可能具有不同的收益率？请结合契约限制条款、抵押品以及流动性因素进行解释。

\item 为什么到期收益率并不一定是投资者在债券上实际获得的回报呢？

\item 为什么久期只能提供债券价格对收益率变化反应的局部近似？凸性如何改善这种近似效果？

\item 债券的收益率上升了80个基点，而基准收益率仅上升了20个基点。在将剩余的变化归因于违约风险之前，还需要哪些额外的信息呢？

\end{enumerate}

%========================================================
% Applied Problems
%========================================================

\TailModule{appprobs}{应用题}

\begin{enumerate}

\item \textbf{合同规定的现金流。}

该债券的面值为人民币~1,000元，票面利率为6\%，票息支付方式为每半年一次，到期时间为四年。

\begin{enumerate}

\item 计算每半年的票息支付。

\item 还剩下多少票息支付？

\item 在到期时支付最终的合同款项。

\item 解释在债券发行后市场收益率立即上升的情况下，哪些合同现金流会发生变化。

\end{enumerate}

\item \textbf{优先级与回收。}

一家陷入困境的公司在次级债务和优先债务上分别拥有1.5亿元人民币的资金。优先债务的请求权为1.2亿元人民币，而次级债务的请求权为8000万元人民币。

\begin{enumerate}

\item 计算每个类别在严格优先级下的回收量。

\item 计算每类物品的回收率。

\item 如果债权人手中的资产价值降至9000万人民币以下，请重新计算回收金额。

\item 解释为什么这两类债券在违约前可能具有不同的信用利差，尽管它们都是由同一发行方发行的债务。

\end{enumerate}

\item \textbf{债券价格和收益率。}

一张面值为100的五年期债券的票息为4\%。

\begin{enumerate}

\item 如果连续复利的到期收益率为3\%，计算其价格。

\item 将价格重新计算，假设年收益率为5\%。

\item 请确定在每种情况下，该债券是以溢价还是折价进行交易的。

\item 请确认价格变动的方向与价格与收益率之间的关系是一致的。

\end{enumerate}

\item \textbf{即期定价法。}

三年期债券的面值为100，年票息为5\%。连续复利的1年、2年和3年的即期利率分别为3.0、3.6和4.1\%。

\begin{enumerate}

\item 使用即期汇率曲线来估算该债券的价格。

\item 解释为什么以三年期利率来贴现每一笔现金流是不相同的。

\item 找出一个能够近似再现您价格的连续复利到期收益率。

\item 解释现汇利率和到期收益率之间的经济差异。

\end{enumerate}

\item \textbf{久期与凸性。}

该债券的久期为6.2，凸性为48。其当前价格为102。假设收益率是连续复利的。

\begin{enumerate}

\item 仅通过久期来估算收益率增加40个基点的价格变动百分比。

\item 添加凸性调整。

\item 使用久期-凸性近似法来估算新的价格。

\item 解释为什么在收益出现相等增减时，凸性项具有相同的符号。

\end{enumerate}

\item \textbf{基准利率和信用利差。}

七年期的企业债券收益率为5.8\%，而同类基准债券的收益率则仅为3.7\%。这两种收益率都是基于相同的标准来计算的。

\begin{enumerate}

\item 计算初始的基点差距。

\item 一个月后，该公司的收益率达到了6.4\%，而基准收益率则为3.9\%。请计算出新的利差。

\item 将企业收益率的60个基点增长分解为基准利率变化和利差变化的影响。

\item 给出导致利差发生变化的三个可能原因，并说明需要哪些证据来区分这些原因。

\end{enumerate}

\end{enumerate}

%========================================================
% Data Application
%========================================================

\TailModule{dataapp}{数据应用}

\textbf{中国的政府曲线、信贷曲线以及债券市场定价}

使用中央国债登记结算有限责任公司发布的中债收益率曲线数据。选择一条国债收益率曲线，以及至少一条信用相关曲线，例如高评级企业债或商业银行债券曲线。采用共同期限，并尽可能保持收益率定义与观察口径一致。

首先，绘制出选定曲线的数据，这些数据对应的是某个共同的观测日期。然后，计算在每个共同到期日时，信用曲线与政府债券曲线之间的差距。观察这个差距在各个到期日之间是否保持恒定，或者是否沿着曲线发生变化。

其次，对更早的观测日期进行分析。对于每一种到期时间，将信用收益率的变化分解为政府收益率的变化以及利差的变化。

\[
\Delta y_{\mathrm{credit},t}
=
\Delta y_{\mathrm{government},t}
+
\Delta s_t.
\]

确定信用利差发生最大变化的到期时间，并将这些变化与政府曲线的变动进行比较。

第三，谨慎解读这些结果。信用曲线与政府曲线之间的差距越大，可能反映出预期信用损失、风险溢价、市场流动性、投资者需求、发行机构特征，或是曲线构建方式的差异。观察到的差距变化不应直接被视为违约概率的变化。

报告观测日期、曲线名称、到期日、收益率定义以及数据来源。如果数据源提供了不同类型的曲线，例如收益率至到期曲线和即期利率曲线，请说明使用了哪种曲线，并解释选择的原因。在对其进行经济分析之前，尽可能考察相关证券在流动性、发行方类型、信用质量或合同结构方面的差异。

%========================================================
% Further Reading
%========================================================

\TailModule{furtherread}{延伸阅读}

\begin{itemize}

\item Macaulay, Frederick R. (1938),
\emph{Some Theoretical Problems Suggested by the Movements of Interest Rates, Bond Yields and Stock Prices in the United States since 1856（1856年以来美国利率、债券收益率和股票价格变动所提出的若干理论问题）},
National Bureau of Economic Research.

\item Fisher, Lawrence, and Roman L. Weil (1971),
``Coping with the Risk of Interest-Rate Fluctuations: Returns to Bondholders from Na\"ive and Optimal Strategies（应对利率波动风险：简单策略与最优策略下债券持有人的回报）,''
\emph{The Journal of Business}, 44(4), 408--431.

\item Merton, Robert C. (1974),
``On the Pricing of Corporate Debt: The Risk Structure of Interest Rates（公司债务定价：利率的风险结构）,''
\emph{The Journal of Finance}, 29(2), 449--470.

\item Duffie, Darrell, and Kenneth J. Singleton (1999),
``Modeling Term Structures of Defaultable Bonds（可违约债券期限结构建模）,''
\emph{The Review of Financial Studies}, 12(4), 687--720.

\item Bessembinder, Hendrik, William Maxwell, and Kumar Venkataraman (2006),
``Market Transparency, Liquidity Externalities, and Institutional Trading Costs in Corporate Bonds（公司债券市场透明度、流动性外部性与机构交易成本）,''
\emph{Journal of Financial Economics}, 82(2), 251--288.

\end{itemize}

\ChapterTail{summary,keyterms,aiprompt,questions,appprobs,dataapp,furtherread}

\chapter{股票市场与股权请求权}
\label{chap:Stock_Markets_and_Equity_Claims}

\begin{openingcase}{为什么在没有现金交换的情况下，市场价值仍然会发生变化？}

一家上市公司有~1~亿股已发行在外股份，股价为人民币~10~元，股权市值为~10~亿元。某日早晨，公司宣布一项新产品取得重要技术进展。公司未支付股利、未发行新股，也没有收到现金，但投资者调整了对未来利润的预期，股价升至~12~元。

公司股权市值增加了人民币~2~亿元，却并未获得~2~亿元新增融资。这一变化发生在投资者相互转让现有股份的二级市场中。

\end{openingcase}

\begin{tcolorbox}[
  enhanced,
  colback=FinanceBlue!6,
  colframe=FinanceBlue,
  boxrule=0.6pt,
  sharp corners,
  left=8pt,
  right=8pt,
  top=7pt,
  bottom=7pt,
  before skip=1em,
  after skip=1.6em,
  title=\textbf{核心问题},
  coltitle=white,
  colbacktitle=FinanceBlue,
  fonttitle=\bfseries,
  attach boxed title to top left={xshift=6pt,yshift=-2mm},
  boxed title style={
    boxrule=0pt,
    sharp corners,
    left=5pt,
    right=5pt,
    top=2pt,
    bottom=2pt
  }
]
\begin{enumerate}[
  leftmargin=*,
  label=\textcolor{FinanceBlue}{\textbf{\arabic*.}},
  itemsep=0.55em,
  topsep=0.2em
]

\item 股权请求权提供哪些权利？剩余现金流权、有限责任以及控制权与债务所附带的权利有何不同？

\item 所有权结构、投票权以及公司治理机制如何影响着控制权分配，以及股东与管理者之间的代理冲突呢？

\item 在股票市场中，股权请求权是如何发行和交易的？信息、增长预期以及风险又如何反映在市场价格和必要报酬率中呢？

\end{enumerate}
\end{tcolorbox}

\section{作为剩余所有权请求权的股权}

股权是对公司的剩余请求权。股东并不持有获得固定支付的契约承诺，其权利价值取决于企业履行契约义务后的剩余价值，以及股份附带的治理权利。

\subsection{剩余现金流权与有限责任}

公司是独立法律实体。股份具有由公司法、公司章程及相应股份类别条款界定的财务与治理权利。普通股权对公司现金流和资产享有剩余请求权。清算时，普通股股东排在债权人和其他优先请求权人之后。

对于已缴足股款的股份，有限责任通常意味着股东不因持有股份而必须额外投入个人资产。但个人担保、借款、衍生品或其他契约产生的独立义务，仍可依法执行。

剩余地位并不意味着立即分配。盈利公司可以留存收益并再投资，而不向股东支付。股东对这些资源所产生的价值继续享有剩余请求权，其价值受未来经营结果及约束资源使用的治理安排影响。股利是向股东分配价值的一种方式，并非剩余请求权的定义。

\subsection{股权融资与风险承担}

股权融资提供资金，却不产生固定的契约利息或本金支付。公司发行普通股，以未来剩余现金流请求权及通常附带的投票权换取资本。经营现金流不确定时，这种支付灵活性可能具有价值。

不存在固定支付义务，同时意味着承担剩余风险。经营表现恶化时，普通股股利可能减少或取消，股权价值可能大幅下降；经营表现良好时，优先请求权得到满足后，股东分享上涨收益。因此，在考虑契约请求权后，普通股权承担企业价值变化的很大一部分风险。

这种承担风险的结构能够支持对不确定项目的投资。虽然仍然存在潜在的业务风险，但融资结构会影响负面经营结果如何传递至不同的利益相关者手中。

\begin{example}{商业风险与风险承担}
{business-risk-bearing}

假设一家公司正在开发一项新技术，初始投资为1亿元人民币，其中4000万元人民币用于债务融资，6000万元人民币用于股权融资。假设债务契约要求在债务到期时偿还4400万元人民币。假设项目所产生的收益足以满足这些请求权，且该公司没有其他资产或负债。

如果该项目在到期时产生1.4亿元人民币的收入，那么债务持有人将获得约定的4400万元人民币，而股东则能获得剩余的9600万元人民币。如果项目价值降至6000万元人民币，债务持有人仍然能收到4400万元人民币的回报，这样股东就能获得1600万元人民币的收益。然而，如果项目价值进一步降至3000万元人民币，股东将失去他们的全部投资，而债务持有人则只能收回3000万元人民币。


\end{example}

\subsection{债务与股权：优先次序、控制与风险分配}

债务和股权是对同一企业不同的请求权。它们的区别不仅仅体现在支付方式上。它们在优先权、控制权以及对企业价值变化的敏感度方面都存在差异。

债务赋予债权人契约规定的请求权，以及相对于普通股权的优先地位。普通债务的上行收益通常限于契约支付额，承诺支付未履行或未足额回收时则产生损失。债权人通常不行使一般公司控制权，但契约条件被违反或企业陷入财务困境时，契约限制条款、执行权利及重组安排可能赋予其更大影响力。



\begin{table}[htbp]
\centering

\caption{债务和股权作为金融请求权}
\label{tab:ch9-debt-equity-claims}

\begin{tabularx}{\textwidth}{@{}p{3.0cm}Y Y@{}}

\toprule

\textbf{维度}
& \textbf{债务}
& \textbf{普通股权} \\

\midrule

现金流权利
& 根据合同规定的利息和/或本金支付，在违约情况下执行
& 在满足合同义务后，对企业现金流的剩余请求权 \\

优先级&通常，清偿顺位高于普通股权；债务请求权的优先级取决于清偿顺位、证券以及适用法律的规定&清算中的剩余索赔人\\

到期日  
& 通常会有明确的指定日期，不过也存在永久性的债务安排。  
& 没有固定的偿还期限  
\\

控制项  
&通常缺乏普通的表决权；合同保护及执行权利遵循债务协议中的规定  
& 表决权与治理权，受股份类别及适用法律的约束  
\\

Upside  
& 通常仅限于合同金额的支付  
& 参与剩余收益的实现以及企业价值的提升  
\\

缺点
& 如果合同中的请求权未能完全得到偿还，可能会遭受损失
& 投资可能会全部损失，但股东通常仅因持有股份而无需对公司的债务承担个人责任 \\

\bottomrule

\end{tabularx}

\end{table}

普通股权是剩余请求权。股东没有约定的到期支付，在考虑优先请求权后获得剩余价值。在有限责任下，股份价值可能降至零，而上行收益不受契约偿还金额封顶。正常情况下，投票权还赋予普通股股东治理权利，包括选举董事和批准特定公司事项。

这些差异可能造成债权人与股东的利益冲突。债务存续期间，股东可能偏好上行收益更大的项目，即使额外下行损失部分由债权人承担。债权人通过定价、抵押品、契约限制条款、期限条款及其他契约保护作出应对。财务困境中，议价能力和控制权也可能向债权人转移。


\section{所有权、控制与公司治理}

公司形式将所有权与日常管理分离开来。股东拥有剩余请求权，而董事和经理则对公司的资源行使委托控制权。这种分离可能会在决策者无法承担其决策的全部后果时引发代理问题。公司治理机制负责分配控制权，并提供监督和惩戒管理的机制。

\subsection{投票权与董事会控制}

股东通常不会做出日常运营决策。他们主要通过投票权来行使控制权，包括选举董事和批准特定的重要事项。董事会任命并监督高级管理层，同时行使根据公司法以及公司治理安排所赋予的决策权。因此，董事会是股东投票权与管理层控制之间的重要纽带。

图~\ref{fig:ch9-governance-chain}展示了一个简化的委托控制链。

\begin{figure}[htbp]
\centering

\begin{tikzpicture}[
  >=Latex,
  box/.style={
    draw,
    rounded corners,
    align=center,
    minimum width=4.2cm,
    minimum height=0.85cm,
    inner xsep=6pt,
    inner ysep=4pt,
    font=\normalsize
  },
  flow/.style={
    -Latex,
    line width=0.85pt,
    FinanceBlue
  },
  lab/.style={
    font=\small,
    text=FinanceBlue,
    fill=white,
    inner xsep=3pt,
    inner ysep=1pt
  },
  node distance=0.75cm
]

\node[box] (share) {股东们};

\node[box, below=of share] (board)
  {董事会成员};

\node[box, below=of board] (mgmt)
  {管理};

\node[box, below=of mgmt] (assets)
  {企业资产与运营};

\draw[flow]
  (share) --
  node[lab, right] {选举董事}
  (board);

\draw[flow]
  (board) --
  node[lab, right] {指定并监督实施。}
  (mgmt);

\draw[flow]
  (mgmt) --
  node[lab, right] {保留LaTeX命令、公式、交叉引用和数字。}
  (assets);

\end{tikzpicture}

\caption{简化的委托控制链}
\label{fig:ch9-governance-chain}
\end{figure}

在一股一票结构下，正式投票权与持有的有投票权股份数量成正比。因此，持有~30\%~已发行在外有投票权股份的股东，也持有相应~30\%~的投票权。但实际影响力还取决于所有权集中度、投票参与程度、投票协议、董事选举规则及需股东批准的事项。

在表决权具有经济价值的情况下，股东能够影响公司的决策制定。即使某位股东的持股比例较低，只要公司股权分散，而其他股东参与投票的积极性不高，该股东仍然可以拥有相当大的影响力。因此，表决权的分配并不一定意味着实际控制权的分配也是如此。

\subsection{代理问题与所有权结构}

当决策权被下放时，就会出现代理问题，因为决策者与利益相关者的利益往往并不一致。在上市公司中，股权结构会影响对 corporate 决策的监督力度，同时也会引发由控制权带来的各种冲突。其中两种关系尤为重要。

第一种情况发生在管理者和股东之间。管理者掌控着公司的资源，但他们可能只拥有少量的剩余现金流权利。因此，他们的目标可能与股东的目标有所不同。管理者可能会重视私人利益、组织扩张、更高的工作稳定性，或者更低的个人风险，即使这些选择会降低股东的价值。

分散化的所有权结构会削弱监督激励机制的作用。小股东可能承担大部分对管理层进行调查的成本，而他们获得的公司价值提升却只占很小的一部分。其他股东则可以从中受益，而无需承担任何成本。这种“搭便车”现象会降低个人进行监督的积极性，同时也会增强董事会、机构投资者、信息披露、审计以及其他治理机制的重要性。

第二种关系涉及控制性股东与少数股东权益之间的关系。当股权高度集中时，需要加强对这种关系的监管。大型股东能够从中获得更多的利益，因此他们更有动力来监督管理层的行为。然而，过度集中的股权结构可能会引发控制性股东与少数股东之间的冲突。控制性股东可能会通过影响关联方交易、资产转让、融资决策以及其他企业行为等方式，谋取私利，同时降低少数投资者的价值。


\begin{table}[H]

\centering

\caption{两个重要的股权代理问题}

\label{tab:ch9-agency-problems}

\begin{tabularx}{\textwidth}{@{}p{3.0cm}Y Y@{}}

\toprule

\textbf{代理问题} & \textbf{冲突的来源} & \textbf{潜在的治理措施} \\

\midrule

管理者与股东&管理者拥有决策权，但只能持有部分剩余请求权&董事会监督、薪酬设计、信息披露、监控、投票权、收购压力\\

控制性股权与少数股东
& 控制性股权持有者能够影响决策，而这些决策的影响同样会被少数投资者所承担
& 关联方交易规则、独立董事制度、信息披露要求、投票保护机制以及法律保护 \\

\bottomrule

\end{tabularx}

\end{table}


\subsection{股东权利与股份类别}

一份股份代表着一系列财务和治理权利。特定股份所附带的权利取决于公司法律、公司的管理文件以及相关的股份类别。不同股份类别之间的差异可能会影响财务利益的分配以及投票权的分配。

不同股份类别可以不同方式分配这些权利。例如，优先股可能在股利或清算分配中优先于普通股，但投票权有限或没有投票权；其地位仍次于债权人请求权。转换、赎回或其他契约特征，可使优先股兼具股权和债务的某些特点。

存续的二元类和加权投票结构更直接地将投票权力与经济所有权分离开来。假设类\(j\)有\(N_j\)份存续股份，每股拥有\(v_j\)票。该类的总投票权比例为

\begin{equation}
\omega_j
=
\frac{v_jN_j}{\sum_k v_kN_k}.
\label{eq:ch9-voting-power}
\end{equation}

如果所有班级的每股现金流权益都相同，那么班级 \(j\) 所拥有的现金流份额为

\begin{equation}
\alpha_j
=
\frac{N_j}{\sum_k N_k}.
\label{eq:ch9-cash-flow-share}
\end{equation}

当不同类别的投票权存在差异时，\(\omega_j\)和\(\alpha_j\)不必相等。

这样的结构可以让创始人或其他控股股东在从外部投资者那里筹集股权的同时，仍然保持投票控制权。这可能会支持企业战略的连续性，尤其是在投资项目具有长期前景的情况下。同时，这种分离可能会削弱对外部股东的问责制，因为这些股东的经济影响力可能超过他们的投票影响力。治理方面的影响则取决于公司法律和上市规则所提供的保障措施。

\begin{tcolorbox}[
  enhanced,
  breakable,
  colback=white,
  colframe=FinanceBlue,
  boxrule=0.7pt,
  sharp corners,
  borderline west={2.2pt}{0pt}{FinanceBlue},
  title=\textbf{财务展望：现金流权益与投票权不必一致},
  coltitle=white,
  colbacktitle=FinanceBlue,
  fonttitle=\bfseries,
  before skip=1.2em,
  after skip=1.2em
]

某些类别的股份可能仅占公司总现金流的很小一部分，但它们却拥有较大的表决权。现金流权利决定了参与剩余经济价值分配的权利，而表决权则决定了影响公司决策的能力。

\end{tcolorbox}

\subsection{公司行动与所有权请求权的变化}

公司行为可以改变企业持有的现金数量、存续股票的数量，或者改变股东之间的权利分配。这些行为的直接影响可以通过分析现金是进入还是离开企业、股票是增加还是减少，以及所有权份额或权利是否发生变化来分析。

新的股权发行会增加存续的股数，并可能改变现有的股权份额。假设一家公司有\(N_0\)存续股份和发行股份\(N_1\)新股。持有这些股票的投资者……\(n\)在发行前持有的股份只占一小部分。$\alpha_0=n/N_0$. 如果投资者不参与本次发行，那么其持有股份比例将变为……$\alpha_0=n/(N_0+N_1)$  

投资者的持股比例受到稀释，但比例稀释并不必然意味着经济价值损失。如果新股按公平价格发行，公司收到的现金在发行时使其价值大致增加相应筹资金额，因此较小的持股比例可能对应更大的总股权价值。如果以异常低价向特定投资者发行，价值则可能从未参与发行的股东转移给新投资者。

在LaTeX命令、公式、交叉引用以及数字方面，保持原有的格式即可。  

股票发行允许现有股东购买新股份，如果他们参与这一计划，就可以保持其持股比例。而不参与该计划的股东则可能会面临股份比例减少的情况；不过，可转让权利可以为这些股东提供一定的补偿，以弥补因放弃认购特权而造成的损失。

\begin{example}{百分比稀释与数值转移}{new-equity-dilution}

该公司有8000万股存续股票，每股交易价格为10元人民币左右，因此其市场股权价值约为8亿元人民币。一位投资者持有800万股股票，占该公司总股本的10\%。

该公司发行了2000万股新股票，而该投资者并未参与此次发行。因此，该投资者的持股比例下降到了……

\[
\frac{8}{100}=8\%.
\]

如果新发行的股票价格为人民币10元，那么公司筹集了约2亿元人民币的资金。假设其他价值没有变化，发行后的股权价值将变为10亿元人民币，而股价则保持在10元人民币。投资者的持股比例会下降，但800万股股票的价值仍然为8000万元人民币。

如果新发行的股份以人民币约5元的价格分配给选定的投资者，那么公司仅筹集了约1亿元人民币的资金。在相同的假设下，发行后的股权价值约为9亿元人民币，这意味着每股的股价为

\[
\frac{900}{100}=\text{RMB }9.
\]

该投资者的股份价值约为7200万元人民币。在这种情况下，这种低于市价的发行方式实际上是将价值从现有股东转移到了新投资者手中。

因此，百分比稀释以及经济价值的损失是两种不同的情况。经济效应取决于所有权比例的变动以及发行价格的因素。

\end{example}

股利和股票回购都是从公司向股东分配现金，但它们对所有权请求的影响不同。现金股利是将资源从公司转移到股东手中，而不会改变存续股份的数量。如果一家公司每股支付股利\(D\)，并且没有其他相关信息，那么在股票开始交易时，没有该股利的股东们的股票价格应该会下降。忽略税收和市场摩擦因素的话……

\begin{equation}
P_{\mathrm{ex}}
\approx
P_{\mathrm{cum}}-D,
\label{eq:ch9-ex-dividend-price}
\end{equation}

其中，\(P_{\mathrm{cum}}\)表示包含股利权利的股票价格，而\(P_{\mathrm{ex}}\)则代表除股利权利后的股票价格。股东会获得现金，同时公司的现金持有量也会减少。

股份回购也将公司现金转移给股东，但会减少已发行在外股份。出售股份的股东获得现金，在其他条件不变时，保留股份的股东持有剩余股权的比例提高。回购可能通过支付价格、交易传递的信息、税收或资本结构变化影响价值。仅减少股份数量，并不会创造经济价值。

股票分割改变了股份的数量，但并未转移现金或改变所有权比例。在 \(m\) 对 \(n\) 的分割中，股份数量会发生变化。

\begin{equation}
N'=\frac{m}{n}N,
\label{eq:ch9-stock-split-shares}
\end{equation}
如果没有其他因素的影响，那么每股的价格大约会调整为

\begin{equation}
P'=\frac{n}{m}\times P.
\label{eq:ch9-stock-split-price}
\end{equation}

因此，如果拆分操作并未传递任何新的信息，也不会对市值产生任何影响，那么总市值将保持不变。在输出结果时，请保留所有LaTeX命令、公式表达、交叉引用以及数字等内容。

\begin{table}[H]
\centering

\caption{选定的企业行为及其直接的影响}
\label{tab:ch9-corporate-actions}

\begin{tabularx}{\textwidth}{@{}p{2.8cm}Y Y@{}}

\toprule

\textbf{操作}
& \textbf{立即更改}
& \textbf{关键的经济问题} \\

\midrule

现金股利
& 现金分配给股东；存续股份保持不变。在除息日，股价通常会下降。
& 这种分配方式如何影响公司保留的资源和股东价值？\\

股份回购  
& 现金流出公司，导致存续股份普遍减少。  
& 以什么价格回购股份？这些操作如何影响股份持有比例和股东价值？\\

股票分割
& 存续股份增加，每股价格相应调整；总股权价值和持股比例保持不变，其他条件相同。
& 这种分割是否传递了信息或影响了交易条件？\\

新股权发行  
& 现金进入公司，存续股份增加，可能会稀释现有股东的持股比例。  
& 新股份以什么价格发行？这种发行对现有股东的经济价值有何影响？\\

\bottomrule

\end{tabularx}

\end{table}

\section{股票市场架构与价格发现}

股票市场的运作具有双重功能。一级市场允许企业筹集股权资本，而二级市场则使得现有股份能够在投资者之间转让。在次级交易中，基于投资者的信息和预期的竞争性买卖订单决定了市场价格，这一价格由交易所、经纪商、清算系统和存托机构共同支撑。

\subsection{股权发行与一级市场}

一家公司通过向投资者发行股票在一级市场上筹集新的股权资本。首次公开募股是指首次向公众投资者发行股票，而成熟股权融资则是指在上市之后继续筹集额外的公众股权。权益发行和私人配售则是发行新股票的另一种方式。

股权发行需要确定新请求权的条款和价格。承销商决定提供多少股份以及这些股份附带哪些权利，而投资者则评估公司的预期现金流、风险、治理情况以及资金的使用方式。承销商、簿记人以及财务顾问等中介机构可以协助进行尽职调查、信息披露、簿记建档以及股份发行工作。他们的角色是促进发行过程，并将新股份与投资者的需求相匹配。

公开发行可能同时包含公司新发行股份和现有股东出售的股份。前者是公司新发行的股份，向公司注入新资金；后者由现有股东出售，只转移所有权，出售所得不进入公司。解释发行筹资额时，这一区别十分重要。

\begin{figure}[H]
\centering

\begin{tikzpicture}[
  >=Latex,
  box/.style={
    draw,
    rounded corners=2pt,
    line width=0.7pt,
    align=center,
    text width=1.9cm,
    minimum height=0.9cm,
    inner sep=5pt
  },
  cash/.style={
    -Latex,
    line width=0.9pt,
    FinanceBlue
  },
  claim/.style={
    -Latex,
    dashed,
    line width=0.9pt,
    FinanceBlue
  },
  lab/.style={
    font=\small,
    align=center,
    fill=white,
    inner sep=1.5pt
  }
]

% Participants
\node[box] (firm) at (-1.0,0)
  {公司};

\node[box] (market) at (4.5,0)
  {一级市场};

\node[box] (investors) at (10.0,0)
  {投资者};

% New equity claims: company to investors
\draw[claim]
  ([yshift=5pt]firm.east) --
  node[lab, above=3pt] {新的股权请求权}
  ([yshift=5pt]market.west);

\draw[claim]
  ([yshift=5pt]market.east) --
  node[lab, above=3pt] {新的股权请求权}
  ([yshift=5pt]investors.west);

% Capital: investors to company
\draw[cash]
  ([yshift=-5pt]investors.west) --
  node[lab, below=3pt] {大写字母}
  ([yshift=-5pt]market.east);

\draw[cash]
  ([yshift=-5pt]market.west) --
  node[lab, below=3pt] {大写字母}
  ([yshift=-5pt]firm.east);

\end{tikzpicture}

\caption{一级市场的资本与股权请求权流动}
\label{fig:ch9-primary-equity}
\end{figure}

发行价格会影响所筹集的资本数量，以及现有股东与新股东之间的所有权分配。如果发行价格相对于投资者的估值来说过高，那么本次发行可能会吸引不足的需求。而如果发行价格过低，发行人可能筹集到的资金会少于原本能够获得的资金。初始定价折扣可能反映了估值的不确定性、信息不对称以及发行风险。

\subsection{交易所与其他交易场所}

交易所为证券的上市和交易提供规则与基础设施。上市与交易虽然相关，但却是两个不同的过程。一家公司在某个交易所上市后，其股票的交易也可以通过其他允许的交易平台进行。

在共同的交易规则、准入安排和市场数据系统的框架下，交易商将订单汇集在一起进行交易。现代股权市场还可能包含其他交易系统，例如交叉网络和暗池，以及以交易商为基础的交易平台。

假设有一家公司的股票在纽约证券交易所上市。投资者并不一定要在纽约证券交易所进行这些股票的交易。根据相关的规则和订单处理机制，经纪人可以将这些订单提交给其他交易所或采用替代交易系统来处理，比如暗池交易系统。因此，同样的股票可以在多个交易平台上进行交易，每个平台都有自己独特的订单流和流动性状况。

多个交易场所会增加执行服务的竞争，同时允许不同的交易机制共存。不过，这也可能导致订单流出现混乱，使得市场中的流动性难以被有效监控。因此，交易场所的结构决定了订单之间的相互作用方式、流动性的集中程度以及价格的形成方式。

交易场所在更广泛的市场基础设施中运作。经纪商负责传递订单，清算机构负责建立结算机制，并在适用的情况下处理与交易相关的义务。而证券保管机构则负责维护所有权记录，结算系统则负责完成现金和证券的转移。电子交易改变了这些过程的速度和组织方式，但并未改变其基本经济功能。

\subsection{交易机制、订单与市场流动性}

一种交易机制决定了买卖订单如何相互作用。在订单驱动的市场中，订单被提交到限价单簿中，当兼容的订单相遇时，交易就会发生。而在交易商市场里，交易商会报出他们愿意买入和卖出的价格。许多现代股权市场结合了这两种结构的元素。

市价单优先考虑按可得最优价格成交，但受可用数量限制。限价单规定买方愿意支付的最高价格或卖方愿意接受的最低价格，提供更强的价格控制，却可能无法成交。市价单通常需求现有流动性，而不能立即成交的限价单可为其他交易者供给流动性。

如果最高出价和最低出价都是\(P_{\mathrm{bid}}\)以及那些使用LaTeX命令、公式表达、交叉引用和数字的格式。\(P_{\mathrm{ask}}\)，所引用的买卖价差为

\begin{equation*}
s=P_{\mathrm{ask}}-P_{\mathrm{bid}}.
\label{eq:ch9-bid-ask-spread}
\end{equation*}

中点是

\begin{equation*}
M=\frac{P_{\mathrm{ask}}+P_{\mathrm{bid}}}{2},
\label{eq:ch9-midquote}
\end{equation*}
以及相应的比例区间范围如下：

\begin{equation*}
s_{\mathrm{rel}}
=
\frac{P_{\mathrm{ask}}-P_{\mathrm{bid}}}{M}.
\label{eq:ch9-relative-spread}
\end{equation*}

价差是衡量流动性的一方面。深度指的是当前价格附近可获得的交易量，而价格影响则反映了执行价格对订单大小的响应程度。对于小额交易来说，市场可能会呈现较小的价差，但对于大额订单而言，市场的深度则可能有限。

市场流动性会随着市场状况的变化而波动。当不确定性或逆向选择风险增加时，流动性提供者可能会报出更宽的价差，展示更少的深度，或者撤回订单。因此，与更了解市场的投资者进行交易的可能性，使得信息条件与市场流动性紧密相连。

交易限制也会影响信息和投资者观点如何进入订单流。卖空机制允许投资者借入股票进行出售，之后再回购这些股票以偿还借款。这样一来，即使投资者并不实际持有这些股票，他们仍然可以通过卖出订单来表达负面信息或预期。不过，相关的限制或较高的借款成本可能会限制这种交易方式的效果。因此，卖空行为确实会对订单流、流动性以及价格形成产生影响，但其具体影响程度则取决于市场状况和交易规则。

\subsection{信息、价格与价格发现}

股票价格是股权请求权的实际交易价格。当投资者对未来现金流或必要报酬率的预期发生变动，或者与请求权相关的权利发生变化时，股票价格就会发生变化。

用 \(\mathcal{I}_t\) 表示在时间 \(t\) 时可获得的信息。一种简化的基于股利的估值可以表示为

\begin{equation}
P_t
=
\sum_{j=1}^{\infty}
\frac{\mathrm{E}[D_{t+j}\mid\mathcal{I}_t]}
{(1+r_{e,t})^j},
\label{eq:ch9-information-price}
\end{equation}
其中，\(r_{e,t}\)被视为从时间\(t\)开始所需的报酬率。新的信息可能会改变预期的未来股利、所需报酬率，或者两者都会发生变化，因此也会改变投资者对这支股票的估值。

价格形成是指信息和投资者评估通过交易被反映到市场价格过程中的过程。不同投资者的信息、模型、预期以及流动性需求各不相同，因此他们的订单只能反映出他们部分已知的信息。因此，价格形成虽然能够提供有用信息，但并不一定是即时或无需成本的。图~\ref{fig:ch9-price-discovery}展示了信息及投资者评估如何被反映到观察到的市场价格中，同时也说明了市场摩擦对这一过程的影响。


\begin{figure}[H]

\centering

\begin{tikzpicture}[
  >=Latex,
  box/.style={
    draw,
    rounded corners,
    align=center,
    text width=2.55cm,
    minimum height=0.9cm
  },
  flow/.style={
    -Latex,
    line width=0.85pt,
    FinanceBlue
  }
]

\node[box] (info) at (0,1.4)
  {公称和私称\\信息};

\node[box] (beliefs) at (3.3,1.4)
  {现金流与风险评估\\};

\node[box] (orders) at (6.6,1.4)
  {命令、公式、交叉引用和数字部分保持不变。\\交易};

\node[box] (price) at (9.9,1.4)
  {已观察到\\市场价格};

\draw[flow] (info) -- (beliefs);
\draw[flow] (beliefs) -- (orders);
\draw[flow] (orders) -- (price);

\node[
  draw,
  rounded corners,
  align=center,
  minimum width=5.8cm,
  minimum height=0.75cm
] (frictions) at (6.6,-0.3)
  {流动性、资金供应以及市场中的摩擦因素。};

\draw[flow] (frictions.north) -- (orders.south);

\end{tikzpicture}

\caption{信息、交易与价格发现}
\label{fig:ch9-price-discovery}
\end{figure}

有效市场假说为评估市场价格如何吸收信息提供基准，其含义取决于所考虑的信息集。在该假说下，扣除交易和信息成本后，投资者不应能够系统地利用相应市场有效形式所涵盖的信息，获得异常的风险调整后回报。

根据所考虑的信息类型，假设通常可以用三种方式来表达。弱形式效率关注的是过去的价格和交易活动中包含的信息。半强形式效率则扩展了信息的范围，包括了公开可用的信息，比如财务报告和企业公告等。而强形式效率则同时考虑了公开信息和私有信息，这一标准更可以被视为一种要求较高的基准，而非大多数市场的真实描述。

价格发现仍然受到经济摩擦的影响。获取信息可能成本高昂，投资者之间可能存在分歧，此外，流动性状况、资金限制以及套利的限制也会对观察到的价格产生影响。如果价格完全反映了所有难以获取的信息，那么投资者就没有动力去承担获取这些信息的成本。

\begin{tcolorbox}[
  enhanced,
  breakable,
  colback=white,
  colframe=FinanceBlue,
  boxrule=0.7pt,
  sharp corners,
  borderline west={2.2pt}{0pt}{FinanceBlue},
  title=\textbf{财务展望：价格发现并非完美估值},
  coltitle=white,
  colbacktitle=FinanceBlue,
  fonttitle=\bfseries,
  before skip=1.2em,
  after skip=1.2em
]

观察到的市场价格反映了投资者在当前市场条件下的信息、预期以及评估。这些价格虽然能反映经济价值，但它们并非对客观存在的基本价值的直接体现。

\end{tcolorbox}

\section{股权价值与风险}

股权价值反映了可供股东使用的现金流的现值，同时考虑了这些现金流的风险因素。股权估值有三个核心要素：预期股东现金流、再投资对未来增长的影响，以及与股权风险相一致的必要报酬率。

\subsection{股权现金流与估值输入}

应当区分账面股权和市场股权。账面股权是一种基于已确认的资产和负债得出的会计计量方式，而市场股权则是存续的请求权的总市值。这两者并不一定要一致，因为市值反映了对未来现金流、增长机会、风险以及与股份相关的权利的预期。

股利是指公司直接支付给股东的现金。股票回购则是一种另一种形式的支付方式：公司向出售股票的股东支付现金，同时使存续的股数减少。因此，股利和股票回购以不同的方式分配价值。

留存收益在保留时并不作为对股东的分配。其对股权价值的贡献取决于公司内部保留的资源所产生的未来现金流。只有当再投资能够支持具有足够价值的未来现金流时，保留收益才会创造价值。

对于计划持有股票直到日期\(T\)的投资者来说，当前的价值可以表示为在持有期间预期获得的股利与预期出售价格的现值之和。

\begin{equation}
P_0
=
\sum_{t=1}^{T}
\frac{\mathrm{E}[D_t]}{(1+r_e)^t}
+
\frac{\mathrm{E}[P_T]}{(1+r_e)^T},
\label{eq:ch9-finite-horizon-equity-value}
\end{equation}
其中，\(r_e\) 是在恒定利率简化条件下，股权请求权的必要报酬率。

如果终端价格被视为在日期$T$之后预期股利的现值，那么无限延长估值范围，就可以得到股利折现表示法：

\begin{equation}
P_0
=
\sum_{t=1}^{\infty}
\frac{\mathrm{E}[D_t]}{(1+r_e)^t}.
\label{eq:ch9-dividend-discount-model}
\end{equation}


对于股权估值，分析师可以预测自由现金流至股权，即\(FCFE_t\)，这是扣除运营需求、税收、再投资以及净债务融资后，可供股权持有人使用的现金流。一个简化的估值表达式为

\begin{equation}
V_E
=
\sum_{t=1}^{\infty}
\frac{\mathrm{E}[FCFE_t]}{(1+r_e)^t}.
\label{eq:ch9-fcfe-valuation}
\end{equation}

现金流指标与贴现率应对应同一请求权。可供全部资本提供者获得的现金流，需要采用适合整个企业的贴现率；仅可供股权持有人获得的现金流，则需要采用股权必要报酬率。

可比公司的市场价格提供了额外的估值基准。相对估值则使用诸如市盈率、市净率等指标来比较类似的请求权。当考虑到增长、盈利能力、杠杆率、风险以及会计计量方面的差异时，这种比较更具参考价值。

\subsection{增长预期与市场价值}

增长通过未来现金流及实现增长所需的再投资影响股权价值。预期销售或利润提高，本身并不意味着价值提高。新增股权投入所获得的回报超过股东要求的回报时，再投资才创造价值。

股利支付比率将预期收益与下一期的股利联系起来。假设下一期的每股收益为$E_1$，公司保留一部分收益$k$，并将剩余部分作为股利分配。那么

\begin{equation}
D_1=(1-k)E_1.
\label{eq:ch9-dividend-payout}
\end{equation}

如果增长是通过留存收益来融资的，那么就不会发行新的普通股股权，而期初的股权回报率保持稳定。在这种情况下，一个简化的可持续增长关系为

\begin{equation}
g=k\,ROE.
\label{eq:ch9-sustainable-growth}
\end{equation}

如果股利以恒定速率
$g$ 无限期地增长，那么戈登增长模型，也被称为恒定增长股利折现模型，会表示当前的股价为

\begin{equation}
P_0=\frac{D_1}{r_e-g},
\qquad
r_e>g.
\label{eq:ch9-gordon-growth}
\end{equation}

代入收益和增长关系后，结果如下：

\begin{equation}
P_0
=
\frac{(1-k)E_1}
{r_e-k\,ROE}.
\label{eq:ch9-growth-reinvestment-value}
\end{equation}

较高的留存比率虽然会减少当前的支付额，但能够提升未来的增长潜力。其对价值的影响取决于所保留资源所获得的报酬率与股权必要报酬率之间的关系。

\begin{example}{增长与价值创造}{growth-reinvestment-value}

假设下一期的每股收益为10，留存率为40\%，而股权的必要报酬率为10\%。

如果留存收益能够保持稳定的水平\(ROE\)15中的…\%,

\[
g=0.40\times 0.15=6\%,
\]
与
\[
P_0=\frac{(1-0.40)\times 10}{0.10-0.06}=\frac{6}{0.04}=150.
\]

如果股权回报为8\%，

\[
g=0.40\times 0.08=3.2\%,
\]
与
\[
P_0=\frac{(1-0.40)\times 10}{0.10-0.032}=\frac{6}{0.068}\approx88.24.
\]

在这种简化的设定中，当再投资获得的收益超过必要报酬率时，增长会增加价值；而当再投资的收益低于必要报酬率时，增长则会降低价值。

\end{example}

这一原则不仅适用于恒定增长模型。快速增长可能需要大量投资，而适度的当前收益在保留资源能够产生足够强劲的未来现金流时，仍然可以体现出较高的价值。增长的预期也会影响价值的实现时间。当预期中的更多现金流发生在更远的未来时，股权价值对必要报酬率的变化就变得更加敏感。

\begin{tcolorbox}[
  enhanced,
  breakable,
  colback=white,
  colframe=FinanceBlue,
  boxrule=0.7pt,
  sharp corners,
  borderline west={2.2pt}{0pt}{FinanceBlue},
  title=\textbf{财务展望：增长并不等同于价值创造},
  coltitle=white,
  colbacktitle=FinanceBlue,
  fonttitle=\bfseries,
  before skip=1.2em,
  after skip=1.2em
]

只有当所需再投资产生的回报足够高时，增长才能提升股权价值，而这一回报必须高于股东所要求的回报。

\end{tcolorbox}

\subsection{股权风险溢价与必要报酬率}

投资者在某一时期内持有的股票的实际回报，是所收到的股利与市场价格变化的总和。如果初始价格为 \(P_0\)，期末价格为 \(P_1\)，股利为 \(D_1\)，那么实际持有的期间回报就是

\begin{equation}
r
=
\frac{P_1-P_0+D_1}{P_0}.
\label{eq:ch9-equity-realized-return}
\end{equation}

必要报酬率不同于已实现回报率。它是投资者基于股权请求权风险及其他可得投资回报，要求为持有该请求权获得的回报。能够通过多项投资分散的风险，通常比与广泛经济条件相关的风险需要更少补偿。

令 \(r_M\) 表示广泛市场投资组合的回报，而 \(r_f\) 表示无风险利率。市场股权风险溢价为

\begin{equation}
ERP
=
\mathrm{E}[r_M]-r_f.
\label{eq:ch9-equity-risk-premium}
\end{equation}

股权风险溢价具有前瞻性。历史股票已实现回报超过无风险回报的平均值，可为过去结果提供证据，却未必等于特定日期投资者要求的溢价。

资本资产定价模型是一种被广泛使用的基准方法，用于将股票的必要报酬率与系统性风险联系起来。股票\(i\)的贝塔值为

\begin{equation}
\beta_i
=
\frac{\mathrm{Cov}(r_i,r_M)}
{\mathrm{Var}(r_M)}.
\label{eq:ch9-beta}
\end{equation}

在资本资产定价模型假设下，必要报酬率为

\begin{equation}
r_i
=
r_f
+
\beta_i\left(\mathrm{E}[r_M]-r_f\right).
\label{eq:ch9-capm}
\end{equation}

贝塔系数衡量了股票对市场整体回报变动的敏感度。贝塔系数越高，意味着在资本资产定价模型中的风险溢价要求也就越高。

\begin{example}{A CAPM 必要报酬率}{capm-required-return}

假设无风险利率为3\%，预期市场股权风险溢价为6\%，某股票的贝塔值为1.2。那么CAPM的必要报酬率为

\[
r_i
=
3\%+1.2(6\%)
=
10.2\%.
\]

该估计结果是基于CAPM模型以及所使用的输入数据而得出的。

\end{example}

CAPM~是一个基准，而非对预期回报的完整描述。其更广泛的经济含义是：与全市场不利条件相关的风险，在保护最有价值时难以分散，因而通常需要更高补偿。贝塔通过与市场回报的协方差表示这种风险暴露，其他资产定价模型则允许存在额外的定价风险来源。

%========================================================
% Summary
%========================================================

\TailModule{summary}{本章小结}

普通股是一种剩余所有权请求权。股东在优先 contractual 请求权得到满足之后获得价值，并参与剩余的增值收益；而有限责任通常可以防止完全支付的股份所带来的损失超出该股份所投入的价值。股权融资可以提供资本，而不需要产生固定的 contractual 利息或本金支付，并且新投资者可以获得剩余现金流的权利，通常还有表决权。即使新发行的股份会稀释股东的百分比所有权，但它也会增加公司可用的资源。

公司治理分配对企业资源的控制权。股东通常通过投票和选举董事来间接行使控制权，而董事会则负责任命和监督管理层。分散化的所有权可能会削弱监督的激励作用，而集中的所有权则可以加强监督，同时增加控股股东与少数股东之间的潜在冲突。不同的股票类别可以区分现金流权和表决权。公司的各种行为，如股利分配、股票回购、股票分割和新股发行，都会以不同的方式影响现金、股票数量以及各股东的请求权。

股票市场的运作既具有融资功能，也具备交易功能。首次发行可以为公司引入新的资本，而二次交易则使得现有股份在投资者之间流转。交易所及其他交易平台为订单的传递提供了机制，而市场的流动性则取决于价差、深度、价格影响、信息状况，以及市场参与者提供流动性的意愿。价格的形成反映了市场上的信息、预期、风险评估以及流动性状况，但这并不意味着每一个观察到的价格都等同于某种客观认定的基本价值。

股权估值要求股东现金流与用于折现的必要报酬率之间保持一致性。股利和回购会释放现金，而留存收益则通过未来再投资的现金流为股权创造价值。当预期再投资的报酬率远高于股权必要报酬率时，股权价值会上升。市场股权风险溢价代表了广泛股权投资的预期超额回报，而CAPM则提供了一个基准，将个股的必要报酬率与其系统性风险联系起来。

%========================================================
% Key Terms
%========================================================

\TailModule{keyterms}{关键术语}

普通股权（Common equity）；剩余请求权（residual claim）；有限责任（limited liability）；持股比例稀释（ownership dilution）；投票权（voting rights）；董事会（board of directors）；公司治理（corporate governance）；代理问题（agency problem）；控股股东（controlling shareholder）；中小股东（minority shareholder）；股份类别（share class）；优先股（preferred share）；不同投票权（weighted voting rights）；现金股利（cash dividend）；股份回购（share repurchase）；股票拆分（stock split）；配股（rights issue）；首次公开发行（initial public offering）；上市后增发（seasoned equity offering）；公司新发行股份（primary shares）；现有股东出售股份（secondary shares）；交易所（exchange）；其他交易场所（alternative trading venue）；市价单（market order）；限价单（limit order）；买卖价差（bid--ask spread）；市场流动性（market liquidity）；价格发现（price discovery）；信息有效性（informational efficiency）；股利贴现模型（dividend discount model）；股权自由现金流（free cash flow to equity）；留存比例（retention ratio）；可持续增长（sustainable growth）；股权风险溢价（equity risk premium）；贝塔（beta）；资本资产定价模型（Capital asset pricing model）。

%========================================================
% Financial AI Lab
%========================================================

\TailModule{aiprompt}{金融~AI~实验室}

\textbf{实验：对人工智能的股权估值和治理情况进行审计与比较}

选择两家属于同一大行业的上市公司。让人工智能系统使用简单的价格倍数法，以及基于自由现金流与股权或股利的评估方法来比较它们的股权估值。要求该系统明确指出所评估的股类、相关的存续股份数量、预期的增长假设、必要报酬率，以及每个事实输入的来源。

将结果分为四个阶段进行审核。

\begin{enumerate}

\item 验证市值和每股估值所使用的股份数量是否正确，以及每个估值倍数是否与相应的股权指标相匹配。

\item 研究估值对增长假设和必要报酬率的响应。至少测试一种不同的增长率和一种不同的必要报酬率，或者，在适用的情况下，测试股权风险溢价。

\item 请比较所选股票类别所附带의投票权和现金流权利。确定是采用“一股一票”的原则，还是通过治理结构来分离投票权与经济所有权的关系。

\item 请对照公司的公开信息或官方的消息来源来核实每一个事实性声明。至少找出一个声明，人工智能提供了合理的解释，但缺乏足够的支持证据。

\end{enumerate}

%========================================================
% Discussion Questions
%========================================================

\TailModule{questions}{讨论题}

\begin{enumerate}

\item 为什么有限责任并不能使普通股权成为低风险的请求权呢？需要区分的是，从拥有股票本身所产生的最大损失，以及股票市场价值的变化。

\item 如何使得集中所有权能够减少某个代理问题，同时增加另一个代理问题的解决呢？

\item 为什么在利好消息公布后，即使公司没有从股市获得新的资金注入，股价仍然能够大幅上涨呢？

\item 为什么更快的盈利增长并不一定意味着更高的股权价值？请结合再投资和股权必要报酬率来阐述你的答案。

\item 为什么将历史平均实现的市场回报与无风险回报之间的差额自动视为当前预期的股权风险溢价是不合理的呢？

\end{enumerate}

%========================================================
% Applied Problems
%========================================================

\TailModule{appprobs}{应用题}

\begin{enumerate}

\item \textbf{剩余支付与有限责任。}

一家以债务和股权为融资方式的公司，在下一个时期将产生160或70的现金流，这两种情况出现的概率都是相等的。承诺的债务偿还金额为100。

\begin{enumerate}

\item 计算每个状态下的债务和股权支付情况。

\item 计算股权部分的预期支付。

\item 解释股权支付结构中的哪一部分体现了有限责任。

\item 解释为什么在发行债务之后，股东可能会从公司现金流风险性的转变中获益，而债权人则可能因此受损。

\end{enumerate}

\item \textbf{投票权和现金流控制权。}

有一家公司拥有2000万股A类股票，每股对应的投票数为10票；同时，该公司还有1.8亿股B类股票，每股对应的投票数为一票。假设这两类股票的每股现金流权利是相同的。

\begin{enumerate}

\item 计算该公司总现金流中A类资产的份额。

\item 计算A类股份在总投票权中的份额。

\item 请解释为什么这两个百分比之间存在差异。

\item 明确指出这种结构的一个潜在优势以及一个潜在的治理风险。

\end{enumerate}

\item \textbf{市价订单和限价订单。}

该股票的最佳买入报价为24.96美元，对应的股票数量为6,000股；而最佳卖出报价则为25.04美元，对应的股票数量则是4,000股。

\begin{enumerate}

\item 计算报价差值和中间值。

\item 计算相应的报价价差比例。

\item 如果报价仍然有效，那么以什么价格可以执行一笔购买2000股的市场订单呢？

\item 如果下达了一个买入限价订单，订单数量是2000股，价格目标是25.00元，而没有任何卖方愿意在25.00元或更低的价格出售股票，那么这个订单在限价订单簿中将会如何处理呢？

\end{enumerate}

\item \textbf{增长与再投资。}

一家公司预计下一期的每股收益为8，其股权必要报酬率为11\%。该公司将保留30\%的收益。假设增长是通过留存收益来融资的，没有新的普通股股本被发行，且\(ROE\)保持在16\%的水平不变。

\begin{enumerate}

\item 计算可持续增长率。

\item 计算下一周期的每股股利。

\item 使用恒定增长模型来计算股票的价值。

\item 如果需要进行重新计算，请点击此处。保留LaTeX命令、公式、交叉引用和数字。\(ROE\)变为9\%虽然保留率保持不变。

\item 请说明估值变化的方向。

\end{enumerate}

\item \textbf{股权风险与必要报酬率。}

无风险利率为2.5\%，预期市场股权风险溢价为5.5\%。

\begin{enumerate}

\item 计算贝塔值为0.8的股票的CAPM必要报酬率。

\item 计算贝塔值为1.4的股票的必要报酬率。

\item 解释在资本资产定价模型（CAPM）中，为什么高贝塔值的股票具有更高的必要报酬率。

\item 请给出一个理由，说明为什么分析师可能不愿意仅依赖CAPM估值来进行决策。

\end{enumerate}

\end{enumerate}

%========================================================
% Data Application
%========================================================

\TailModule{dataapp}{数据应用}

\textbf{中国股权市场的信息与价格发现}

请选择一家在上海证券交易所或深圳证券交易所上市的中国A股公司，并找到一份包含经济相关新信息的企业公告，例如收益公告、利润预警、重大投资决策、股利提案、股票回购计划或新的股权发行。以交易所披露平台和该公司的正式公告作为获取这些信息的首要来源。

收集在披露事件前后几天的股票每日收盘价，同时记录同一时期的市场指数。计算每天的股票回报率以及简单的市场调整后的异常回报值。

\begin{equation*}
r_t=r_{i,t}-r_{M,t},
\end{equation*}

其中，\(r_{i,t}\) 表示第 \(t\) 天的股票回报，而 \(r_{M,t}\) 则表示市场指数的回报。

确定信息首次公开披露的具体时间。如果披露信息是在市场关闭后发布的，那么下一个交易日应被视为投资者可以通过交易进行响应的第一天。绘制股票价格的走势图，并报告在披露信息发布后的市场调整收益情况。判断最大的价格波动发生在披露信息后的第一个交易机会之前、期间还是之后。

谨慎解读这一结果。价格反应可能反映了对预期现金流或必要报酬率的调整，或者两者都有。同时发布的全市场新闻以及其他公司特定的信息也可能影响观察到的回报。较短的事件窗口和每日收盘价本身并不能证明因果关系或日内价格发现的速度。

请列出该公司的名称、股票代码、披露标题、披露日期和时间、样本周期、基准指数以及数据来源。同时说明在得出市场对这些信息反应迟缓或过度反应的结论之前，还需要哪些额外的证据。

%========================================================
% Case Analysis
%========================================================

\TailModule{caseanalysis}{案例分析}

\textbf{案例：小米与香港的加权投票权}

小米公司在2018年7月9日开始在香港上市，成为在2018年上市制度改革后首批采用加权投票权结构在香港上市的公司。小米的A类股票在大多数股东决议中拥有每股十票的权利，而B类股票则拥有每股一票的权利。对于某些特定保留事项，每股股票拥有一票表决权。目前，小米公开上市的股票属于B类股票。

这种结构将投票权与经济所有权分离开来。它允许创始人在从公众投资者那里筹集到股权后，仍然对公司的决策拥有相当大的影响力。香港的上市规则允许采用特定的加权投票结构，但前提是必须满足相关资格要求、持续履行义务，并为非加权投票的股东提供相应的保护措施。

从融资和治理的角度分析该结构。为什么创始人可能重视能够筹集到公开股权，而不放弃相应的投票控制权？集中的投票权力如何支持战略上的连续性或长期投资？当投票控制权超过现金流所有权时，可能会出现哪些代理问题？当外部股东的投票权力较弱时，哪些投资者保护机制变得更为重要？

那么考虑一下市场价值的影响。解释为什么小米公开发行B类股票的市场价格不仅反映了预期的剩余现金流，还反映了与投资相关的治理权和控制权。如果预期运营现金流最初没有变化，那么表决权的变化是否一定会使市场价值保持不变？解释控制权的变化、代理成本、未来现金流的分配以及必要报酬率如何影响股权价值。

\emph{清算：香港交易及结算所有限公司，2018年7月4日发布的《HKEX将为投资者提供更多选择，助力小米股票上市》；小米公司，2018年招股说明书及2018年中期报告；香港证券交易所上市规则，关于加权投票权的第8A章。}

%========================================================
% Further Reading
%========================================================

\TailModule{furtherread}{延伸阅读}

\begin{itemize}

\item Gordon, Myron J. (1959),
``Dividends, Earnings, and Stock Prices（股利、收益与股票价格）,''
\emph{The Review of Economics and Statistics}, 41(2, Part 1), 99--105.

\item Fama, Eugene F. (1970),
``Efficient Capital Markets: A Review of Theory and Empirical Work（有效资本市场：理论与实证研究综述）,''
\emph{The Journal of Finance}, 25(2), 383--417.

\item Grossman, Sanford J., and Joseph E. Stiglitz (1980),
``On the Impossibility of Informationally Efficient Markets（论信息有效市场的不可能性）,''
\emph{American Economic Review}, 70(3), 393--408.

\item Kyle, Albert S. (1985),
``Continuous Auctions and Insider Trading（连续拍卖与内幕交易）,''
\emph{Econometrica}, 53(6), 1315--1335.

\item La Porta, Rafael, Florencio Lopez-de-Silanes, Andrei Shleifer, and Robert W. Vishny (2000),
``Investor Protection and Corporate Governance（投资者保护与公司治理）,''
\emph{Journal of Financial Economics}, 58(1--2), 3--27.

\end{itemize}

\ChapterTail{summary,keyterms,aiprompt,questions,appprobs,dataapp,caseanalysis,furtherread}

\chapter{外汇市场与跨境金融}
\label{chap:Foreign_Exchange_Markets_and_Cross_Border_Finance}

\begin{openingcase}{外币应收账款价值多少？}

一家中国制造企业向境外客户出售设备，金额为~1000~万美元，三个月后付款。按当前每美元兑人民币~7.20~元的即期汇率，应收账款价值为人民币~7200~万元。但企业三个月后才能收到美元，届时即期汇率可能不同。

企业可以保留未套期保值的外汇风险暴露，在收款时兑换美元；也可以远期卖出未来美元收入，今天就锁定汇率。销售契约决定美元现金流，外汇市场决定其人民币价值，以及汇率风险如何分配。

\end{openingcase}

\begin{tcolorbox}[
  enhanced,
  colback=FinanceBlue!6,
  colframe=FinanceBlue,
  boxrule=0.6pt,
  sharp corners,
  left=8pt,
  right=8pt,
  top=7pt,
  bottom=7pt,
  before skip=1em,
  after skip=1.6em,
  title=\textbf{核心问题},
  coltitle=white,
  colbacktitle=FinanceBlue,
  fonttitle=\bfseries,
  attach boxed title to top left={xshift=6pt,yshift=-2mm},
  boxed title style={
    boxrule=0pt,
    sharp corners,
    left=5pt,
    right=5pt,
    top=2pt,
    bottom=2pt
  }
]
\begin{enumerate}[
  leftmargin=*,
  label=\textcolor{FinanceBlue}{\textbf{\arabic*.}},
  itemsep=0.55em,
  topsep=0.2em
]

\item 当现金流以不同货币计价时，外汇市场如何与国际贸易、投资和跨境融资相互关联？

\item 通过货币供需关系，现货汇率是如何形成的？资本流动、预期以及市场流动性又如何影响汇率的波动和价格发现呢？

\item 现货汇率和远期汇率是如何通过汇率差异以及覆盖性利息平价来相互关联的？企业和投资者如何利用外汇工具来管理汇率风险呢？

\end{enumerate}

\end{tcolorbox}

\section{外汇市场架构}

汇率是指一种货币相对于另一种货币的价格。外汇市场提供了将跨境支付、投资和资金在不同货币之间进行转换的机制。

\subsection{国际贸易、投资与融资需求}

契约现金流的计价币种与企业或投资者衡量其价值的币种不同时，跨境交易便产生货币风险暴露。计价币种规定契约支付以何种货币表示，估值币种则用于衡量请求权或负债的经济价值。两者不同时，契约金额可能保持固定，其以估值币种衡量的价值却随汇率变化。

设~$S_t$~表示每单位外币对应的本币数量。因此，$S_t$~上升表示本币贬值。

假设一家国内出口商在日期$T$时会收到$F$单位的外币。如果这笔收入按照当时有效的汇率进行兑换，那么其以本国货币表示的价值就是

\begin{equation}
V_{D,T}=S_T V_{F,T}.
\label{eq:ch10-export-receivable-value}
\end{equation}

商业契约固定外币金额，但其本币价值随~$S_T$~变化。较高的~$S_T$~提高外币应收款的本币价值。外币应付款产生相反的净风险暴露，因为较高的~$S_T$~增加履行该义务所需的本币成本。

跨境投资将外国资产的收益与汇率波动结合在一起。假设一位国内投资者以$S_0$的价格将本国货币兑换成外币，然后获得一定的外币收益$r^{*}$，之后又用同样的货币进行反向兑换，此时使用的汇率为$S_1$。那么，本国货币的净收益就是……

\begin{equation}
1+r_d=(1+r^{*})\frac{S_1}{S_0}.
\label{eq:ch10-foreign-investment-return}
\end{equation}

本币回报既取决于外国资产回报，也取决于汇率变化。本币升值可能降低国外投资回报的本币价值，本币贬值则可能提高其本币价值。

跨境借款会使负债方面临货币风险。外币贷款可能具有较低的合同利率，但其本国货币的偿还成本也取决于未来的汇率。因此，在比较国内和外资时，需要同时考虑借款条款和货币风险。

\subsection[市场参与者、交易商与电子交易平台]{市场参与者、交易商以及电子交易平台\ }

外汇市场是一个去中心化的网络，而不是单一的交易所。银行、非银行交易商、企业、资产管理者、对冲基金、中央银行以及其他机构通过交易商网络、电子平台以及双边关系进行交易。

参与者进入市场，是为了执行与贸易、投资和融资交易相关的货币转换。交易商通过报价汇率、提供即时执行服务以及管理由此产生的货币头寸来中介这些交易。

交易商通过中介行为创造了多层次的交易过程。企业客户可以与银行进行交易，而银行则可以将这些头寸与另一客户、另一位交易商或通过电子渠道进行对冲。因此，一笔客户交易可能会引发更多的交易活动，因为货币头寸会在市场中重新分配。

\begin{figure}[H]
\centering

\begin{tikzpicture}[
  >=Latex,
  box/.style={
    draw,
    rounded corners,
    align=center,
    minimum width=3.0cm,
    minimum height=0.9cm,
    inner sep=5pt
  },
  two/.style={
    Latex-Latex,
    line width=0.85pt,
    FinanceBlue
  }
]

\node[box] (corp) at (-4.3,1.7)
  {企业和机构\\跨境借款人};

\node[box] (asset) at (-4.3,-1.7)
  {资产管理者和\\机构投资者};

\node[box] (dealer) at (0,0)
  {银行与\\外汇交易商};

\node[box] (inter) at (4.3,1.7)
  {交易与中介活动\\电子平台};

\node[box] (policy) at (4.3,-1.7)
  {中央银行和\\官方机构};

\draw[two] (corp) -- (dealer);
\draw[two] (asset) -- (dealer);
\draw[two] (dealer) -- (inter);
\draw[two] (dealer) -- (policy);

\end{tikzpicture}

\caption{简化的外汇交易架构}
\label{fig:ch10-fx-market-architecture}
\end{figure}

该图是示意图。一些机构投资者和非银行金融机构也直接访问电子交易平台和其他流动性提供者。中央银行主要通过官方交易参与，包括储备管理和外汇干预。

电子交易通过允许从多个来源获取报价和订单，从而减少了市场许多领域的搜索和执行成本。然而，大额交易、流动性较低的货币对以及定制合约可能更依赖于交易商的关系和双边信贷安排。

外汇交易也产生结算义务。双方交付不同币种，往往通过不同司法辖区的支付系统完成。如果一种货币已交付，而另一种尚未收到，已付款一方就面临结算风险。通过将两种货币的结算相联系，同步交收（PvP）安排可以降低这种风险暴露，但流动性、操作和法律风险仍可能存在。

这些机构安排为汇率的报价以及货币头寸的交易提供了环境。

\subsection{即期报价、市场流动性与价格发现}

即期汇率是指根据市场标准在短期内进行结算时，两种货币之间的兑换价格。不同货币对的结算时间可能有所不同，因此具体的结算日期仍然是交易条款的一部分。

根据这种引用规则，$S_t$ 表示从国内参与者视角直接引用的内容。而相互引用则……

\begin{equation}
\widetilde{S}_t=\frac{1}{S_t},
\label{eq:ch10-reciprocal-quote}
\end{equation}
以单位外币计算单位国内货币。如果汇率是每美元约7.20人民币，那么反向报价大约为$1/7.20=0.1389$以人民币计算的美元价值。

报价惯例决定汇率变动如何解释。采用每单位外币对应本币数量的报价时，$S_t$~上升意味着购买一单位外币需要更多本币。因此，从每美元兑人民币~7.20~元变为~7.40~元，表示人民币对美元贬值。

现货交易的报价通常包含买入价和卖出价。对于以本国货币兑换外币的报价，买入价是交易商购买外币的价格，而卖出价则是交易商出售外币的价格。如果报价是每美元7.1980至7.2020人民币，那么报价的价差就是

\begin{equation*}
\text{保留LaTeX命令、公式、交叉引用和数字。}=7.2020-7.1980=0.0040.
\label{eq:ch10-fx-spread}
\end{equation*}

买卖价差是市场流动性的一个维度。较窄的买卖价差通常意味着交易报价规模的交易 immediate 交易成本较低。流动性还取决于市场深度和价格影响。市场深度指的是当前价格附近可获得的交易量，而价格影响则指的是与较大订单相关的执行价格的变化。这些条件可能因货币对、交易时间、市场状况以及交易规模而有所不同。

交叉汇率可以连接不同货币对之间的价格。例如，如果市场报价显示$S^{\mathrm{CNY/USD}}=7.20$人民币兑换1美元，而$S^{\mathrm{USD/EUR}}=1.10$美元兑换1欧元，那么得出的人民币兑换欧元的交叉汇率就是如此。

\begin{equation*}
S^{\mathrm{CNY/EUR}}
=S^{\mathrm{CNY/USD}}\times
S^{\mathrm{USD/EUR}}
=7.92.
\label{eq:ch10-cross-rate}
\end{equation*}

当直接报价和隐含交叉汇率之间的差距足够大，足以覆盖交易成本和执行限制时，就会产生套利机会。在相关货币对之间进行交易通常可以使这些价格保持紧密一致。

\begin{example}{交叉汇率与三角套利}{fx-triangular-arbitrage}

假设市场的报价为：1美元兑换7.20人民币，1欧元兑换1.10美元。这意味着实际汇率为1欧元兑换7.92人民币。如果直接汇率是1欧元兑换8.00人民币，那么在实际市场中，欧元的价格会比通过美元兑换时更高。

忽略交易成本的话，交易者可以通过美元间接获得欧元，并以较高的直接人民币/欧元汇率出售这些欧元。这样的交易通常可以减少不同货币对之间的价格差异。实际上，这种差异必须足够大，以覆盖买卖价差以及其他执行成本。

\end{example}

汇率决定了交易商给出的报价。当客户和投资者的订单与交易商的报价相互交流时，订单流就能反映出买卖压力的变化，进而影响交易商的库存状况，导致报价的调整和头寸的重新分配。因此，观察到的汇率实际上反映了交易活动与可用流动性以及市场制造能力之间的相互作用。

因此，观察到的汇率变动可能反映了多种力量的共同作用。政策利率的变化、通货膨胀预期、经济状况以及风险感知等因素可能会相互影响，所以汇率的变动取决于这些更广泛的信息如何影响不同货币的预期回报。


\section{汇率度量、制度与动态变化}

货币条件的变化、资本流动、预期和风险都会对货币价值产生压力。这些压力如何体现在市场汇率、官方干预或其他政策调整上，取决于汇率制度。

\subsection{名义汇率、实际汇率与有效汇率}

名义汇率表示一种货币相对于另一种货币的价格，而不考虑国内和外国价格水平的差异。它可以将外币价格转换为本国货币，但本身并不能确定国内商品和外国商品的相对价格。

让 $P_t$ 表示国内价格水平，而 $P_t^*$ 表示国外价格水平。术语 $S_tP_t^*$ 用于将国外价格水平转换为本国货币。将这一金额与国内价格水平进行比较，就可以得到实际汇率：

\begin{equation}
q_t
=
\frac{S_t P_t^{*}}{P_t}.
\label{eq:ch10-real-exchange-rate}
\end{equation}

按这一惯例，$q_t$~上升意味着外国商品相对本国商品变贵，表示本币实际贬值；下降则表示实际升值。

对于比例性的变化，

\begin{equation}
\frac{\Delta q}{q}
\approx
\frac{\Delta S}{S}
+
\pi^{*}
-\pi,
\label{eq:ch10-real-exchange-rate-change}
\end{equation}

其中，$\pi$和$\pi^{*}$分别代表国内和国外的通货膨胀率。因此，实际汇率反映了名义汇率的变动以及通货膨胀的差异。如果国内价格上升的速度远远快于国外价格，那么名义贬值并不一定会导致实际贬值。

\begin{example}{名义汇率与实际汇率变动}{nominal-real-exchange-rate}

假设本国货币汇率从每单位外币7.00单位贬值至7.21单位，贬值幅度约为3\%。在同一时期，国内价格上升了5\%，而国外价格则上升了2\%。

使用近似关系，

$$
\frac{\Delta q}{q}
\approx
3\%+2\%-5\%=0.
$$

名义贬值大致被通货膨胀的差异所抵消，因此实际汇率基本保持不变。

\end{example}

双边汇率反映了两种货币之间的关系。当一个国家与多个合作伙伴进行贸易时，相对于单一货币的变动可能无法反映其货币的整体外部价值。因此，有效的汇率是通过使用一组权重来结合双边汇率得出的，这些权重通常与贸易模式有关。

名义有效汇率概括了加权后的名义货币变动情况。实际有效汇率则加入了相对价格或成本调整因素，这一逻辑同样适用于双边实际汇率，而实际有效汇率则适用于一组贸易伙伴的情况。

因此，当一种货币相对于不同的贸易伙伴发生不同变化时，实际汇率可能会与个别的双边汇率有所不同。然而，这些汇率的解释取决于所使用的权重以及价格或成本指数，因此它们仍然只是对外部竞争力的概括性衡量，而非完整的衡量标准。


\subsection{固定、浮动与管理汇率制度}

不同汇率制度在处理外汇市场压力时，对汇率波动和货币当局政策行动的影响方式有所不同。

在固定或紧密挂钩的汇率制度下，当局试图将汇率维持在一定的平价水平或狭窄的范围内。因此，市场压力更多地通过干预和相关政策调整来加以缓解，不过，如果市场压力持续存在，最终可能需要对汇率安排进行改变。

在浮动汇率制度下，汇率的变动会导致更多的调整。没有固定汇率制度通常意味着货币政策有更大的空间来应对国内状况，尽管仍然可能需要进行干预。

管理型汇率制度结合了这两种调整方式。汇率会根据市场状况进行调节，同时当局仍有权力通过干预或其他操作手段来影响汇率的变动。中国就是一个例子：人民币被正式定义为实行管理型浮动汇率制度，其调整依据的是市场供求情况，同时参考了多种货币篮子。

汇率制度与资本流动性和货币政策独立性之间也存在相互作用。开放经济政策的三重困境意味着，严格固定的汇率、较高的资本流动性以及完全独立的货币政策通常无法同时维持。在汇率固定且资本流动相对开放的情况下，国内货币状况往往会与外部条件更加紧密地联系在一起。

\begin{table}[H]
\centering
\caption{简化的汇率制度与调整机制}
\label{tab:ch10-exchange-rate-regimes}
\begin{tabularx}{\textwidth}{@{}p{2.7cm}X X@{}}
\toprule
\textbf{制度/体制}
& \textbf{主要调整通道}
& \textbf{政策影响} \\
\midrule
已修复或调整了相关设置。
& 干预与政策调整
& 较少的货币政策自主权 \\
已管理好。
& 汇率波动与干预措施
& 中间阶段的政策灵活性 \\
保留LaTeX命令、公式、交叉引用和数字。
& 汇率波动
& 更大的货币政策自主权 \\
\bottomrule
\end{tabularx}
\end{table}

\subsection{资本流动、预期与汇率动态}

投资者会比较国内和外国资产的预期回报，同时考虑利率差异和预期的汇率变动。较高的汇率或资产风险可能会提高投资者要求的回报，而有限的市场流动性则会增加进入或退出市场的成本。资本管制和监管限制可能会进一步限制资金在不同货币之间的流动。

当风险调整后的预期回报相互之间发生变化时，投资者可能会调整他们的投资组合以跨货币进行配置。这些调整会产生资本流动，并改变货币需求，从而影响汇率变动。由于预期汇率变化也会影响预期回报，因此资本流动和汇率是共同决定的。

这些投资组合的调整也受到投资者资金能力和承担风险能力的影响。当杠杆限制收紧或亏损触发保证金要求时，投资者可能会减少外币投资或将资金汇回本国。这种调整可能会加剧汇率波动，并导致短期市场波动，而这种波动并不一定与长期经济前景的变化相匹配。

\begin{tcolorbox}[
  enhanced,
  breakable,
  colback=white,
  colframe=FinanceBlue,
  boxrule=0.7pt,
  sharp corners,
  borderline west={2.2pt}{0pt}{FinanceBlue},
  title=\textbf{金融展望：利率和汇率},
  coltitle=white,
  colbacktitle=FinanceBlue,
  fonttitle=\bfseries,
  before skip=1.2em,
  after skip=1.2em
]

本国利率较高，并不必然导致本币升值。其汇率影响取决于货币政策、通胀与风险的预期路径，以及由此引起的不同币种预期回报变化。

\end{tcolorbox}


\section{平价关系与跨币种估值}

平价关系通过套利和预期，将不同币种的价格和回报联系起来。套利促使不同市场中等价现金流的价值趋于一致，基于预期的关系则取决于对未来价格或汇率的判断。这些关系是在给定假设下的估值基准，而非普遍成立的经验恒等式。

\subsection{一价定律与购买力平价}

一价定律适用于在不同市场出售的相同可交易商品。用$p_{i,t}$表示商品$i$的本国货币价格，用$p_{i,t}^{*}$表示其外币价格。在贸易壁垒和交易成本可以忽略不计的情况下，套利行为通常会在现货汇率转换后使两种价格趋于一致：

\begin{equation}
p_{i,t}=S_t p_{i,t}^{*}.
\label{eq:ch10-law-of-one-price}
\end{equation}

如果国内价格超过转换后的国外价格，且这种差距超过了购买和运输商品的成本，那么交易者就可以在国外购买商品并在国内出售。这种套利行为往往会缩小价格差异。因此，对于那些可交易性强、运输成本低且市场分割程度有限的商品来说，一价定律更有可能成立。

购买力平价将这一逻辑从个别商品扩展到总体价格水平。绝对购买力平价要求可比的商品和服务篮在用共同货币表示时具有相同的成本。考虑到之前定义的实际汇率，绝对购买力平价相当于$q_t=1$。因此，这意味着

\begin{equation}
S_t=\frac{P_t}{P_t^{*}}.
\label{eq:ch10-absolute-ppp}
\end{equation}

实际上，各国的价格指数通常会包含各种可交易与不可交易的商品和服务。消费模式、税收、运输成本以及市场条件的差异都可能导致购买力平价指数无法保持恒定。

相对平价汇率主要关注汇率和价格水平的变化，而非其绝对数值。这一理论基于这样一个假设：实际汇率在一段时间内保持大致不变。正如之前推导出的实际汇率关系所表明的那样，名义汇率会进行调整，以抵消国内与国外之间的通货膨胀差异。根据这里所使用的报价规则，较高的国内通货膨胀意味着本国货币贬值。

汇率并不要求绝对购买力不变。
实际汇率可能会保持在一个不同的水平上，保持稳定。

可以通过一个简单的数值例子来说明这种关系。

\begin{example}{相对购买力平价}
{relative-ppp-example}

假设在未来一年里，国内通货膨胀率预计为5\%，而国外通货膨胀率为2\%。当前的汇率是每单位外币对应7.20个本国货币单位。

在相对购买力平价的基础上，预计本国货币将贬值约3\%。未来汇率的相应购买力平价基准是

\[
S_1
\approx
7.20(1+0.05-0.02)
=
7.416
\approx 7.42.
\]

这个基准反映了预期的通胀差异，而不是对未来汇率的精确预测。

\end{example}

观察到的汇率在较长时间内可能会偏离绝对购买力平价所暗示的水平。运输成本、贸易壁垒、非可交易商品以及市场分割限制了商品市场的套利行为，而价格粘性则减缓了国家价格水平的调整。金融市场的消息也可能使名义汇率的变动速度比商品价格调整的速度快。因此，绝对购买力平价最好被视为汇率水平的长期基准，而相对购买力平价则提供了与通货膨胀差异相关的长期汇率变化的基准。但两者都无法提供可靠的短期汇率预测。


\subsection{抛补利率平价与远期汇率}

抛补利率平价（CIP）是一项无套利条件，使两种投资策略的本币回报相等：直接投资本币资产；或将资金兑换为外币、投资外币资产，并用远期契约对冲汇率风险。投资期限和信用特征可比时，在不存在套利机会的条件下，两种策略必须提供相同的本币回报。

假设有一单位本国货币被投资了一段时间。通过国内投资可以获得$1+r_d$的收益。或者，投资者可以选择按照即期汇率$S_0$将本国货币兑换成外币，然后在国外以$r_f$的利率进行投资，之后再以$F_0$的价格提前出售这些外币收益。这样得到的本国货币支付结果为

\begin{equation*}
\frac{F_0}{S_0}(1+r_f).
\label{eq:ch10-cip-foreign-payoff}
\end{equation*}

如果这两种策略具有相同的到期时间，并且信用和流动性特征相当，那么套利往往会使它们的回报率达到一致：

\begin{equation*}
1+r_d
=
\frac{F_0}{S_0}(1+r_f).
\label{eq:ch10-cip}
\end{equation*}

求解远期汇率后，得到的结果如下：

\begin{equation}
F_0
=S_0\frac{1+r_d}{1+r_f}.
\label{eq:ch10-forward-from-cip}
\end{equation}

对于同一时期内的适度利率水平来说，

\begin{equation}
\frac{F_0-S_0}{S_0}
\approx
r_d-r_f.
\label{eq:ch10-forward-premium-approx}
\end{equation}

根据这种报价规则，当国内利率高于国外利率时，相应的外币远期价格也会高于其现价。

\begin{figure}[H]
\centering
\begin{tikzpicture}[
> =Latex,
box/.style={
draw,
rounded corners,
align=center,
minimum width=3.3cm,
minimum height=0.9cm,
inner sep=5pt
},
flow/.style={-Latex, line width=0.85pt, FinanceBlue}
]
\node[box] (start) at (0,0)
{1个单位的\\国内货币};

\node[box] (dom) at (5.0,1.6)
{国内投资\\在$r_d$处};

\node[box] (for) at (5.0,-1.6)
{转换，投资海外，\\并提前进行套期保值};

\node[box] (domend) at (10.0,1.6)
{$1+r_d$};

\node[box] (forend) at (10.0,-1.6)
{$\dfrac{F_0}{S_0}(1+r_f)$};

\draw[flow] (start) -- (dom);
\draw[flow] (start) -- (for);
\draw[flow] (dom) -- (domend);
\draw[flow] (for) -- (forend);

\end{tikzpicture}
\caption{利率平价中的两条抛补投资路径}
\label{fig:ch10-cip-paths}
\end{figure}

在$T$年的时间内，以连续复合的利率计算，可以用相同的公式来表示这一关系。

\begin{equation}
F_{0,T}
=S_0 e^{(r_d-r_f)T},
\label{eq:ch10-cip-continuous}
\end{equation}
其中，$r_d$和$r_f$分别表示相同到期时间的国内和外国利率。

\begin{example}{由抛补利率平价推导远期汇率}{cip-forward-example}

假设现汇汇率是每单位外币对应7.20个本国货币单位。一年内的本国利率为4\%，而相应的国外利率为2\%。在简化的CIP假设下，

$$
F_0
=
7.20\times\frac{1.04}{1.02}
\approx
7.34.
$$

因此，根据这种报价方式，外币的远期溢价会更高。

\end{example}

实践中，交易成本、抵押品条款、融资条件、交易对手风险、监管和资产负债表约束，都可能使抛补后的回报偏离简单平价基准。此类偏离表明，考虑市场摩擦后，两条融资路径在经济上并非完全相同。

\begin{tcolorbox}[
  enhanced,
  breakable,
  colback=white,
  colframe=FinanceBlue,
  boxrule=0.7pt,
  sharp corners,
  borderline west={2.2pt}{0pt}{FinanceBlue},
  title=\textbf{财务展望：远期利率并非预测结果},
  coltitle=white,
  colbacktitle=FinanceBlue,
  fonttitle=\bfseries,
  before skip=1.2em,
  after skip=1.2em
]

投资者购买外币资产，并提前出售该资产的到期收益，从而锁定了本币回报。如果没有这种套期保值措施，实际获得的回报将取决于到期时的即期汇率。因此，远期汇率实际上是一种约定的未来汇率，而预期的未来即期汇率则反映了投资者的预期以及承担汇率风险所需的补偿。这两者并不一定要一致。

\end{tcolorbox}


\subsection{非抛补利率平价与货币风险溢价}

非抛补利率平价（UIP）比较未对冲汇率风险时不同币种的预期回报。因此，外国存款投资者仍承受未来即期汇率风险。不同于~CIP，UIP~不通过远期契约消除汇率风险，其平价条件涉及预期回报，而非通过契约套期保值锁定的回报。

让$r_d$和$r_f$分别表示在同一投资期限内的国内和外国利率。外国存款的预期本国货币总回报为

\begin{equation*}
(1+r_f)\frac{\mathrm{E}_0[S_1]}{S_0}.
\label{eq:ch10-uip-foreign-return}
\end{equation*}

一个简单的UIP基准测试将这一预期回报与国内总回报等同起来：

\begin{equation}
1+r_d
=
(1+r_f)\frac{\mathrm{E}_0[S_1]}{S_0}.
\label{eq:ch10-uip-exact}
\end{equation}

对于适度的利率和汇率波动，

\begin{equation}
\mathrm{E}_0\left[\frac{S_1-S_0}{S_0}\right]
\approx
r_d-r_f.
\label{eq:ch10-uip-approx}
\end{equation}

本国利率高于外国利率时，UIP~基准意味着预期本币贬值。 

当投资者需要因承担货币风险而获得补偿时，未对冲的外币存款的预期国内货币回报可能会与国内存款的回报有所不同。用$\rho$表示这两种预期回报之间的差值。

\begin{equation}
\rho
=
(1+r_f)\frac{\mathrm{E}_0[S_1]}{S_0}
-
(1+r_d).
\label{eq:ch10-currency-risk-premium}
\end{equation}

因为

\[
\frac{\mathrm{E}_0[S_1]}{S_0}
=
1+
\mathrm{E}_0\left[\frac{S_1-S_0}{S_0}\right],
\]
替换后的结果如下：LaTeX命令、公式、交叉引用以及数字。

\begin{align}
\rho
={}&r_f-r_d
+\mathrm{E}_0\left[\frac{S_1-S_0}{S_0}\right]
+r_f\mathrm{E}_0\left[\frac{S_1-S_0}{S_0}\right].
\end{align}

对于较为温和的利率和汇率变化情况，可以将产品项视为二阶项而忽略不计。由此得到的一阶近似结果为：

\begin{equation}
\rho
\approx
r_f-r_d
+
\mathrm{E}_0\left[\frac{S_1-S_0}{S_0}\right].
\label{eq:ch10-currency-risk-premium-approx}
\end{equation}

重新排序后得到如下结果：LaTeX命令、公式、交叉引用以及数字。

\begin{equation}
\mathrm{E}_0\left[\frac{S_1-S_0}{S_0}\right]
\approx
r_d-r_f+\rho.
\label{eq:ch10-uip-with-risk-premium}
\end{equation}

当$\rho=0$时，两个头寸的预期回报相等，标准UIP基准也就得到了恢复。当$\rho$不为零时，预期汇率变化也反映了承担货币风险所带来的补偿。


\section[货币风险暴露与跨境风险分配]{货币风险暴露与跨境风险分配}

当现金流、资产或负债以与其价值衡量货币不同的货币计价时，就会出现货币风险。汇率变动会导致价值在合同头寸之间重新分配，而套期保值和金融决策则决定了这种风险如何在企业、投资者和金融中介机构之间进行分配。

\subsection{计价币种与契约现金流}

货币单位指的是合同规定的支付单位。以美元计价的债券承诺支付美元款项，而以人民币计价的贷款则需要用人民币进行偿还。货币单位是否会导致企业面临汇率风险，这取决于企业的计价货币以及其收入、资产和其他负债所使用的货币结构。

对于拥有外币负债的公司来说，以本国货币计算的偿还成本会随着即期汇率的变化而变化。本国货币的贬值会增加履行该义务所需的本国货币金额。如果公司的收入主要以本国货币计价，那么外币债务就会形成货币错配。外币收入或资产可以部分抵消这种风险，并提供一种天然的套期保值手段。

合同计价并不能完全涵盖所有汇率风险。企业可以通过选择其结算货币或套期保值合同现金流来减少直接交易风险，同时仍可能因收入、投入成本、竞争环境以及海外业务价值的变化而面临风险。因此，企业的整体货币风险取决于汇率变动对其运营和财务现金流的影响。


\subsection{交易风险暴露}

当以合同形式约定的外币现金流需要在未来日期进行结算时，就会出现交易风险。此时，该现金流的本国货币价值会受到汇率的影响而变得不稳定。这种情况通常出现在应收账款、应付账款以及外币债务偿还过程中。

对方程~\eqref{eq:ch10-export-receivable-value}关于$S_T$求导，同时保持外币应收款项不变，可以得到

\begin{equation}
\frac{\partial V_D,T}{\partial S_T}=V_{F,T}.
\label{eq:ch10-transaction-exposure-sensitivity}
\end{equation}

$S_T$的上升会增加外币应收款的本国货币价值，以及外币应付款的本国货币成本。因此，其经济效应取决于合同中的现金流方向。

由于交易风险通常与特定的金额和结算日期相关联，因此当交易的时间和金额足够可预测时，往往可以与远期合约或其他套期保值合同相结合。

\begin{example}{出口应收账款的交易风险暴露}
{transaction-exposure-exporter}

这位中国出口商将在90天后收到约500万美元的款项。如果结算时的汇率是每美元7.00人民币，那么这笔款项将转换为3500万人民币。但如果汇率达到每美元7.40人民币，那么这笔款项将转换为3700万人民币。

每美元增加0.40人民币的因素，使得应收款项以人民币计算的价值增加了约200万人民币：

\[
\Delta V_{D,T}
=
V_{F,T}\Delta S_T
=
5\times 0.40
=
\text{RMB }2.0\text{ 百万}.
\]

示例展示了固定外币应收款项对汇率的敏感性：较高的$S_T$会使其在本国货币中的价值增加，而较低的$S_T$则会降低这一价值。

\end{example}


\subsection{折算风险暴露}

当需要将以外币表示的财务报表或资产负债表数据转换为本公司的报告货币时，就会涉及到翻译问题。这一点对于拥有海外业务的跨国企业来说尤为重要。

假设一家母公司以本国货币进行报告，而一家外国子公司则使用另一种货币进行账目记录。即使子公司与母公司之间没有相应的现金流变化，汇率的变化仍可能会改变子公司资产、负债、收入和支出的本国货币价值。

因此，折算风险暴露与交易风险暴露不同。交易风险暴露来自将实际结算的契约外币现金流，折算风险暴露则来自将会计金额换算为报告币种。折算损失本身并不产生即时现金支付。

翻译效果仍然可以影响所报告的权益、收益、杠杆率以及其他财务指标。这些指标的重要性取决于公司的资产负债表结构、报告框架，以及与报告中的会计价值相关的任何合同或监管规定。

会计处理方式决定了翻译影响在财务报表中的体现方式。不同的会计框架有着不同的处理规定，因此在进行风险分析时，必须区分实际的外币头寸与其会计记录之间的差异。

\subsection{经济风险暴露}

经济风险暴露有时也称经营风险暴露，指未来经营现金流现值对汇率变化的敏感性。它超出已签约外币应收款和应付款的范围。

即使一家公司的销售额和合同以本国货币计价，它仍然可能面临经济风险。汇率的波动会改变国内产品和国外产品的相对价格、进口投入的成本，以及公司的竞争地位。这些变化可能会影响销售量、利润率以及未来的运营现金流。

经济风险暴露的方向和大小取决于企业收入与成本结构。本币贬值可能提高外币收入价值，但进口投入、外币债务或竞争者反应可能抵消部分影响。因此，并非总能从现有契约的计价币种推断经济风险暴露。


\begin{table}[htbp]
\centering
\caption{三类货币风险暴露}
\label{tab:ch10-currency-exposure-types}

\begin{tabularx}{\textwidth}{@{}p{3.1cm}X X@{}}
\toprule
\textbf{曝光度}
& \textbf{保留LaTeX命令、公式、交叉引用和数字。}
& \textbf{说明性效果} \\
\midrule

交易风险
& 约定的外币现金流
& 结算前的本国货币价值变动 \\

翻译的曝光度&将外币交易转换为报告货币单位&报告的会计数值会随着翻译速率的变化而变化。\\

经济影响
& 汇率对价格、成本、需求和竞争的影响
& 未来现金流现值的变动 \\

\bottomrule
\end{tabularx}
\end{table}

经济风险也会随着企业调整采购、生产、定价和融资方式而发生变化。基于当前运营安排的中短期风险可能会与企业在应对汇率变化后的长期风险有所不同。交易风险、换算风险以及经济风险的主要来源和影响因素在表~\ref{tab:ch10-currency-exposure-types}中有所总结。
这三种风险形式可能会相互重叠。一个国外的业务活动可能会产生换算风险和经济风险，而特定的外币合同则可能导致企业内部的交易风险。


\subsection{套期保值逻辑与汇率风险分配}

套期保值改变头寸对汇率变动的敏感性。其目的通常是抵消已识别的风险暴露，而非预测汇率方向。

考虑一个出口商，其以外币计价的应收款项为$V_{F,T}$，到期日为$T$。如果出口商以汇率$F_{0,T}$将应收款项的一部分$h$进行远期交易，那么相对于即期汇率的结算收益或损失为

\begin{equation}
hV_{F,T}\times(F_{0,T}-S_T).
\label{eq:ch10-export-forward-payoff}
\end{equation}

将预期支付与剩余未对冲应收款的转换结合起来，就可以得到对冲头寸的本国货币价值：

\begin{equation}
V_{D,T}^{H}
=
V_{F,T}\left[(1-h)S_T+hF_{0,T}\right].
\label{eq:ch10-partial-forward-hedge}
\end{equation}

当 $h=0$ 时，应收款项完全暴露于未来的即期汇率之下。而当 $h=1$，且远期合约的金额和到期日与应收款项相匹配时……

\begin{equation}
V_{D,T}^{H}=F_{0,T}V_{F,T}.
\label{eq:ch10-full-forward-hedge}
\end{equation}

因此，套期保值可以用预先确定的汇率来替代部分或全部的不确定的未来转换率。这降低了指定现金流的汇率敏感性，而与基础交易相关的其他风险仍然存在。

\begin{example}{完全与部分远期套期保值}{forward-hedge-example}

这家公司将在六个月内获得约400万美元的收益，而六个月的远期汇率是每美元约7.25人民币。进行全面的套期保值后，其风险将得到锁定。

$$
4\times 7.25)=29.0
$$
百万人民币。

如果该公司选择对冲60\%的应收款项，那么以本地货币表示的价值就会变为……

$$
4\times\left[0.4S_T+0.6\times 7.25\right].
$$

在未来的汇率为每美元7.00人民币的情况下，该收据的价值为2800万人民币；而在汇率为每美元7.50人民币的情况下，其价值则达到2940万人民币。进行部分套期保值虽然可以降低对未来汇率的敏感性，但并不能完全消除这种风险。

\end{example}

同样的逻辑也适用于其他套期保值工具。期货提供标准化的合约，期权会改变汇率结果的分布，而货币互换则可以改变不同货币之间的现金流。企业还可以通过调整收入、成本、资产和负债的货币来创建天然的套期保值措施。在每种情况下，相关的问题都是：套期保值支付如何抵消潜在的暴露风险。

套期保值比率未必等于一。如果基础现金流金额或时点不确定，当实际风险暴露不同于套期保值金额时，完全套期保值可能使企业产生新的货币头寸。衍生品套期保值还可能在收到基础现金流之前，产生抵押品或流动性要求。

通过商业或融资合同本身也可以分配汇率风险。以本国货币开具发票的出口商，其交易风险更多地转移给了外国买家；而以未来收入的货币进行借贷则可以将部分风险转移到负债方。因此，合同的计价方式、运营选择以及金融套期保值为参与者分担汇率风险提供了不同的途径。

\begin{tcolorbox}[
enhanced,
breakable,
colback=white,
colframe=FinanceBlue,
boxrule=0.7pt,
sharp corners,
borderline west={2.2pt}{0pt}{FinanceBlue},
title=\textbf{财务展望：应针对套期保值进行干预，而非对预测进行调整。},
coltitle=white,
colbacktitle=FinanceBlue,
fonttitle=\bfseries,
before skip=1.2em,
after skip=1.2em
]

套期保值头寸的金额、币种和时点与基础风险暴露相对应时，其效果较好。其目的在于改变现金流结果的分布，而非形成对未来汇率的判断。

\end{tcolorbox}


%========================================================
% Summary
%========================================================

\TailModule{summary}{本章小结}

外汇市场将不同货币表示的请求权连接在一起。交易、投资和跨境融资在现金流和计价货币不一致时会产生货币转换的需求。交易商和电子平台在去中心化的市场中重新分配这些流动，而交易和结算安排则影响着流动性和执行效率。

汇率的解释始于报价惯例。在本章中，汇率指的是本国货币与外国货币之间的比率，因此汇率上升意味着本国货币贬值。实际汇率包含了相对价格水平，而有效汇率则反映了与多个贸易伙伴之间的汇率变动情况。不同国家的汇率制度在应对市场压力的方式上有所不同，有的国家通过汇率调整来应对，而有的国家则采取政策手段来干预。

非抛补利率平价在特定的假设条件下提供了基准。一价定律和购买力平价将商品价格与汇率联系起来，这种联系是通过商品市场套利实现的。抛补利率平价则通过抛补回报的比较，将即期汇率、远期汇率和利率联系起来；而非抛补利率平价则将利息差与预期汇率变化联系起来，这些预期汇率变化是基于对预期回报和风险溢价的前提假设而得出的。因此，远期汇率不必等于未来预期的即期汇率。

货币风险敞口源于合同、会计事项、经营现金流以及融资活动。交易风险敞口与以外币计价的现金流相关；转换风险敞口涉及将海外业务转换为报告货币；经济风险敞口则与未来经营现金流对汇率变动的敏感性有关。通过金融对冲和天然对冲措施，可以在现金流的金额、时间以及货币方面与潜在风险敞口保持足够接近的情况下，减少或重新分配这些风险敞口。


%========================================================
% Key Terms
%========================================================

\TailModule{keyterms}{关键术语}

外汇市场（Foreign exchange market）；即期汇率（spot exchange rate）；直接标价（direct quote）；买卖价差（bid--ask spread）；交叉汇率（cross rate）；三角套利（triangular arbitrage）；市场流动性（market liquidity）；名义汇率（nominal exchange rate）；实际汇率（real exchange rate）；名义有效汇率（nominal effective exchange rate）；实际有效汇率（real effective exchange rate）；固定汇率制度（fixed exchange-rate regime）；浮动汇率制度（floating exchange-rate regime）；管理汇率制度（managed exchange-rate regime）；外汇干预（foreign exchange intervention）；一价定律（law of one price）；购买力平价（purchasing power parity）；抛补利率平价（covered interest parity）；远期汇率（forward exchange rate）；非抛补利率平价（uncovered interest parity）；货币风险溢价（currency risk premium）；计价币种（currency denomination）；交易风险暴露（transaction exposure）；折算风险暴露（translation exposure）；经济风险暴露（economic exposure）；远期套期保值（forward hedge）；套期保值比率（hedge ratio）。

%========================================================
% Financial AI Lab
%========================================================

\TailModule{aiprompt}{金融~AI~实验室}

\textbf{实验：对外汇报价进行审计、计算平价，并进行套期保值}

请一个AI系统使用即期汇率、国内和国外的利率以及一年期的远期利率来分析一年的跨货币投资。在进行任何计算之前，要求它先说明汇率的报价方式。

将分析过程分为四个阶段来进行审核。

\begin{enumerate}

\item \textbf{请检查引号的排列方式。} 请确定分子和分母中的货币种类，并确保对升值和贬值的描述与引号中的表述一致。

\item \textbf{确认抛补利率平价是否正确。} 将两个利率设定在相同的到期日和复利基础上。计算由抛补利率平价所暗示的远期汇率，并将其与提供的远期汇率进行比较。在将任何差异视为套利机会之前，需要考虑交易成本、融资和抵押品条件以及其他相关的市场摩擦因素。

\item \textbf{将定价与预测分离开来。} 让人工智能将远期汇率解释为未来即期汇率的预测。检查它是否能够区分覆盖性定价关系与对未来汇率和货币风险溢价的预期。

\item \textbf{对套期保值进行审计。} 给人工智能一个外币应收款项，要求它在没有套期保值的条件下、部分远期套期保值的条件下以及完全远期套期保值的条件下计算现金流。验证远期头寸的方向，并在本国货币升值和贬值的情况下测试结果。

\end{enumerate}


%========================================================
% Discussion Questions
%========================================================

\TailModule{questions}{讨论题}

\begin{enumerate}

\item 为什么外国投资在其本国货币中能够带来正收益，而对于国内投资者来说，在进行货币兑换后却可能带来负收益呢？

\item 为什么拥有浮动汇率的国家仍然可以干预外汇市场？这种干预与维持固定汇率有什么区别？

\item 为什么抛补利率平价提供的套利基准比非抛补利率平价更强？

\item 该公司的出口合同全部以本国货币计价。那么，为什么这家公司仍然会受到汇率波动的影响呢？

\end{enumerate}


%========================================================
% Applied Problems
%========================================================

\TailModule{appprobs}{应用题}

\begin{enumerate}

\item \textbf{汇率、增值率和交叉汇率。}

当前汇率是人民币约7.24元兑换1美元，而美元则兑换欧元约为1.08欧元。

\begin{enumerate}

\item 计算美元/人民币的 reciprocal 报价。

\item 如果人民币对美元的汇率升至7.38，那么根据本章所使用的报价规则，人民币是升值了还是贬值了呢？

\item 请使用原始的报价来计算人民币对欧元的隐含汇率。

\item 解释如果直接报价的人民币/欧元汇率与隐含的交叉汇率相差超过交易成本及其他执行障碍时，可能会产生什么样的交易压力。

\end{enumerate}

\item \textbf{名义汇率和实际汇率。}

在本章所使用的报价制度下，本国货币在一年内贬值了6\%。国内通货膨胀率为4\%，而国外通货膨胀率则降至2\%。

\begin{enumerate}

\item 利用方程~(\ref{eq:ch10-real-exchange-rate-change})中的近似公式，估算实际汇率的变动百分比。

\item 如果国内通货膨胀率为9\%，请重新计算实际的汇率变化，其他数值保持不变。

\item 解释为什么名义折旧并不一定意味着实际折旧。

\end{enumerate}

\item \textbf{抛补利率平价。}

当前汇率为1单位外币对应7.20单位本国货币。一年期本国利率为4.5\%，而相应的国外利率为2.0\%。

\begin{enumerate}

\item 计算由抛补利率平价所暗示的一年远期汇率。

\item 将向前期保费转换为现汇汇率的百分比形式。

\item 假设观察到的远期利率为7.30。在得出存在可交易套利机会的结论之前，还需要哪些额外的信息呢？

\item 请解释为什么CIP所暗示的远期汇率并不必然代表未来的即期汇率。

\end{enumerate}

\item \textbf{未覆盖的回报与货币风险溢价。}

国内一年的利率为5\%，而国外的利率则低至3\%。预计本国货币在一年内将贬值4\%。

\begin{enumerate}

\item 利用方程~\ref{eq:ch10-currency-risk-premium}中的一阶近似公式，计算未对冲的外币存款相对于本国存款的预期超额回报。

\item 在这种假设下，零风险溢价UIP基准测试是否成立呢？

\item 与零风险溢价UIP相对应的预期汇率变化是多少呢？

\item 解释随时间变化的货币风险溢价如何影响利息差异的解读。

\end{enumerate}

\item \textbf{货币风险与套期保值。}

这家国内出口商将在六个月内收到约800万美元。六个月的远期汇率为每美元7.18人民币，该公司选择75\%的套期保值比例。

\begin{enumerate}

\item 如果未来的即期汇率为每美元6.95人民币，请计算对冲后收据的本国货币价值。

\item 如果未来的汇率为每美元7.45人民币，请重新计算数值。

\item 使用相同的两种汇率数据，比较有套期保值措施下的收入范围与没有套期保值措施下的收入范围。

\item 计算在完全套期保值下锁定的本国货币价值。

\item 即使在完整的货币套期保值下，仍存在一个风险。

\end{enumerate}

\end{enumerate}


%========================================================
% Data Application
%========================================================

\TailModule{dataapp}{数据应用}

\textbf{人民币的双边汇率、名义有效汇率和实际有效汇率}

从权威的中国来源获取每月的人民币对美元的汇率，以及从具有 documented methodology 的权威来源获取中国的名义和实际有效汇率指数。

请使用至少五年的数据样本。在适当的情况下，将这一系列数据转换为统一的起始值100。同时绘制出人民币对美元的双边汇率、名义有效汇率指数以及实际有效汇率指数。明确双边汇率的报价原则，并说明每个指数上升的方向。

请指出至少两个时期，在这期间双边汇率和有效指数出现了不同的变化。解释其他贸易伙伴货币的变化如何导致这种差异。然后比较名义有效指数和实际有效指数，并将它们之间的差异与相对价格或成本的变化联系起来。

请仔细解释这些LaTeX命令、公式、交叉引用以及数字的含义。这些元素在描述不同货币之间的汇率变动时起着重要作用，但它们并不能完全反映一个国家的国际竞争力。请至少列举出两个可能影响整体经济分析的因素。

请报告样本周期、频率、指数定义、数据来源机构，以及应用于数据的任何数据处理方法。


%========================================================
% Further Reading
%========================================================

\TailModule{furtherread}{延伸阅读}

\begin{itemize}

\item Dornbusch, R"udiger (1976),
``Expectations and Exchange Rate Dynamics（预期与汇率动态）,''
\emph{Journal of Political Economy}, 84(6), 1161--1176.

\item Meese, Richard A., and Kenneth Rogoff (1983),
``Empirical Exchange Rate Models of the Seventies: Do They Fit Out of Sample?（20世纪70年代的实证汇率模型：样本外表现如何？）''
\emph{Journal of International Economics}, 14(1--2), 3--24.

\item Fama, Eugene F. (1984),
``Forward and Spot Exchange Rates（远期与即期汇率）,''
\emph{Journal of Monetary Economics}, 14(3), 319--338.

\item Obstfeld, Maurice, and Kenneth Rogoff (1995),
``The Mirage of Fixed Exchange Rates（固定汇率的幻象）,''
\emph{Journal of Economic Perspectives}, 9(4), 73--96.

\item Du, Wenxin, Alexander Tepper, and Adrien Verdelhan (2018),
``Deviations from Covered Interest Rate Parity（抛补利率平价的偏离）,''
\emph{The Journal of Finance}, 73(3), 915--957.

\end{itemize}


\ChapterTail{summary,keyterms,aiprompt,questions,appprobs,dataapp,furtherread}

\chapter{衍生品市场与或有请求权}
\label{chap:Derivatives_Markets_and_Contingent_Claims}

\begin{openingcase}{契约如何在不改变基础业务的情况下改变风险？}

这家制造商计划在三个月内购买1000吨铜。按照当前的价格，这笔采购预计将花费约8000万元人民币。该公司并非试图投机买卖铜价。他们担心的是，如果铜价上涨，生产成本也会相应增加，而在此之前，他们无法及时调整对客户的售价。

该公司可以签订一份远期合约，约定三个月后铜的购买价格。或者，该公司可以选择购买一个看涨期权，该期权可以在铜价大幅上涨时限制成本，同时如果铜价下跌，公司仍能从中受益。无论采取哪种方式，公司对铜的实际需求都不会发生变化。不同的是，公司未来现金流的状态依存模式会有所改变。

如何在不转移基础资产的情况下，让合同重新分配未来不确定情况下的财务后果呢？

\end{openingcase}



\begin{tcolorbox}[
  enhanced,
  colback=FinanceBlue!6,
  colframe=FinanceBlue,
  boxrule=0.6pt,
  sharp corners,
  left=8pt,
  right=8pt,
  top=7pt,
  bottom=7pt,
  before skip=1em,
  after skip=1.6em,
  title=\textbf{核心问题},
  coltitle=white,
  colbacktitle=FinanceBlue,
  fonttitle=\bfseries,
  attach boxed title to top left={xshift=6pt,yshift=-2mm},
  boxed title style={
    boxrule=0pt,
    sharp corners,
    left=5pt,
    right=5pt,
    top=2pt,
    bottom=2pt
  }
]
\begin{enumerate}[
  leftmargin=*,
  label=\textcolor{FinanceBlue}{\textbf{\arabic*.}},
  itemsep=0.55em,
  topsep=0.2em
]


\item 衍生品是如何利用状态依存请求权来拆分和重新分配金融风险敞口的呢？

\item 无套利和复制是如何将衍生品价格与其他交易中的请求权价格联系起来的？

\item 交易所交易、双边合同、抵押品、轧差和中央清算如何影响交易对手风险以及衍生品市场中风险的分配？


\end{enumerate}
\end{tcolorbox}

\section[衍生品存在的原因：或有请求权与风险分配]{衍生品存在的原因：或有请求权与风险分配}

衍生品是支付取决于基础资产、利率、指数、事件或其他指定条件的金融契约。其经济作用不限于投机。衍生品使参与者能够拆分日常经营和金融头寸中捆绑的风险暴露，并将其中选定部分转移给其他市场参与者。

\subsection{状态依存支付}

或有请求权是支付取决于未来状态的金融请求权。设~$T$~时点的可能状态为~$s=1,\ldots,S$，请求权在状态~$s$~下的支付记为~$X_T(s)$。

一个状态代表了与合同支付相关的未来条件。这些条件可能包括基础资产的价格、利率、汇率、商品价格，或者特定信用事件的发生。因此，一项或有请求权描述了合同支付在不同未来状态下的变化方式。

远期契约体现线性的状态依存支付。它约定在未来日期~$T$，按契约订立时确定的交割价格~$K$~买卖基础资产。负有购买义务的一方持有多头，负有交付义务的一方持有空头。

在到期时，长期多头头寸每单位基础资产的支付为

\begin{equation}
X_T^{\mathrm{long}}
=
S_T-K.
\label{eq:ch11-forward-long-payoff}
\end{equation}

短期多头头寸的支付情况则相反：

\begin{equation}
X_T^{\mathrm{short}}
=
K-S_T
=
-X_T^{\mathrm{long}}.
\label{eq:ch11-forward-short-payoff}
\end{equation}

其中，$S_T$~为基础资产到期市场价格。不计交易成本和交易对手损失时，两种头寸的支付等额、方向相反。

期权同样形成状态依存请求权，但持有人拥有行权权利，而非无条件交易义务。其非线性支付结构将在后续章节讨论。

同样的状态依存观点也适用于债务、权益、保险契约和担保。这些请求权在未来的不同状态下分配支付金额，这反映了合同优先权的差异、风险分配的差异，以及支付条件的差异。

\subsection{风险暴露分解与风险拆分}

财务和运营方面的头寸往往同时涉及多种风险来源，这些风险可能会相互作用。例如，一家进口企业的现金流可能受到商品和汇率的影响，而这两个变量之间的共同变动可能会放大或部分抵消该企业本国货币投入成本的变化。因此，一项头寸的风险不仅取决于其各个单独的风险因素，还取决于这些因素之间的相互作用。

风险分解是将整体风险位置分解为经济上不同的来源，并明确每个来源如何影响该位置的现金流或价值。这样就能区分出那些可以保留的风险来源，以及那些可能需要减少或修改的风险来源。

一旦确定了不同的风险敞口，就可以使用衍生品来分别对其进行调整。与特定价格、利率、指数或事件相关的衍生品会创造出一种支付机制，这种支付机制能够响应那些风险来源，而无需改变相关的业务活动或资产状况。通过这种方式，可以将嵌入在更广泛头寸中的风险敞口分离出来，并单独进行调整。

为了说明风险分离，假设在日期$T$时，一个现有头寸的价值为$V_T$，并且面临着多个风险因素。与其中一个因素相关的衍生品支付金额为$X_T$。那么，组合头寸的支付情况就是

\begin{equation}
V_T^{D}=V_T+X_T.
\label{eq:ch11-combined-derivative-position}
\end{equation}

由于~$X_T$~对选定风险因素作出反应，原头寸的这一风险成分可以独立于其他成分调整。如果衍生品支付抵消与该因素相关的部分现金流变化，组合头寸对它的敏感性下降；如果支付强化原有风险暴露，敏感性则可能上升。

套期保值是该机制的一种应用方式。通过采取具有相反支付金额的衍生品头寸，参与者可以在保持原始头寸的其他组成部分基本不变的情况下，减少特定风险敞口。

\subsection{风险分担与市场完备性}

当市场参与者能够交易具有不同状态依存支付结构的请求权时，风险分担变得更加灵活。寻求更稳定现金流的公司可能会将部分风险转移给愿意承担该风险的另一方，以换取补偿，或者因为该风险适合该参与者的其他头寸。衍生品通过创建与特定未来状态相关的合同支付结构，促进了这种风险转移。

风险分担的程度取决于可以从交易中的请求权中构建出的支付范围。考虑一个具有两种未来状态的简单经济体系，分别用$u$和$d$表示。假设有两种交易的证券，它们的支付向量各不相同。

\begin{equation}
\mathbf{x}_1=
\begin{bmatrix}
x_{1u}\\
x_{1d}                                                      
\end{bmatrix},
\qquad
\mathbf{x}_2=
\begin{bmatrix}
x_{2u}\\
x_{2d}
\end{bmatrix}.
\label{eq:ch11-two-security-payoffs}
\end{equation}

持有$a$和$b$两种证券的投资组合，会产生如下状态依存的支付向量：

\begin{equation}
\mathbf{x}_p
=
a\mathbf{x}_1+b\mathbf{x}_2
=
\begin{bmatrix}
ax_{1u}+bx_{2u}\\
ax_{1d}+bx_{2d}
\end{bmatrix},
\label{eq:ch11-spanning-portfolio}
\end{equation}

$a$和$b$的正值表示对相应证券的长期持有，而负值则表示短期抛售。

如果两个支付向量线性独立，它们便张成二维状态支付空间。因此，适当的组合持仓能够产生两种状态下的任意支付向量。所有相关状态依存支付都能由可得交易请求权构造的市场，称为\emph{完备市场}。

如果现有的请求权无法产生所有相关的状态依存支付，那么市场就是不完整的。因此，一些期望的现金流模式就无法精确实现，参与者在交易现有工具后可能仍然承担剩余风险。

在实际市场中，交易成本、交易限制、不可交易的风险、模型的不确定性以及有限的合同类型可能会阻碍某些期望的状态依存支付目标的实现。

衍生品通过引入与现有证券尚未充分代表的风险相联系的请求权，扩展可得的状态依存支付范围，增加参与者之间可转移、可分担的风险种类。衍生品并不创造额外经济资源，而是在不同状态和参与者之间重新分配未来现金流与风险。

\subsection{金融创新与新型请求权}

在不完全市场中，一些状态依存下的期望支付无法从现有的请求权中构建出来。金融创新可以通过创造与特定价格、利率、指数、事件或到期日相关的新衍生品来部分解决这一问题。这些新的请求权可以扩大可转让或对冲的风险范围。

新的请求权可以通过改变参考基准或支付规则来设计。合同可以与新的商品指数、波动性指标或信用事件相关联。其支付结构也可以通过设定上限、下限、障碍特征或路径依赖条件来进行重新设计。通过改变合同的关联因素或支付方式，创新可以改变市场参与者可获得的状态依存现金流。

然而，更紧密的匹配并不能完全消除风险。通过更紧密地跟踪基础资产，定制化的请求权确实可以降低基础风险，但这种做法可能会减少交易的活跃性，或者使得对特定交易对手的依赖程度增加。而标准化的请求权则可以提供更高的流动性，同时与管理中的风险之间的匹配度相对较低。因此，金融创新不仅改变了可以转移的风险类型，也改变了剩余风险的形式。

\section{无套利与复制}

衍生品市场依赖两种相关的机制来保持定价的一致性。无套利机制根据状态依存的支付关系来限制证券的相对价格，而复制机制则利用交易过的证券投资组合来构建等效的支付关系。这两种机制共同为评估衍生品合约提供了基础。

\subsection{套利机会与价格一致性}

一价定律要求，状态依存现金流相同的衍生品头寸具有相同现价。有效比较需要匹配基础参照、名义本金、期限、支付日期、结算币种及其他决定支付的契约条款。这些特征不同可能导致现金流不同，并不必然表明定价不一致。

套利机会可以有两种形式。第一种初始净投资为零，未来所有状态下支付非负，且至少一个状态严格为正。第二种产生正的初始现金流入，却不产生任何未来净支付义务。两种形式均使交易者不承担抵销性未来损失而获得收益。

\paragraph{第一类：从支付优势出发的套利。}

当两种衍生品请求权具有相同的当前价格时，就会出现第一种套利情况。其中一种衍生品在每一个未来状态下提供的价值至少与另一种相同，并且在至少一个状态下其价值严格高于另一种。假设

\begin{equation}
P_A=P_B,
\qquad
X_A(s)\geq X_B(s), \quad \forall s,
\label{eq:ch11-payoff-dominance}
\end{equation}
在至少一种可能的状态下，使用严格的不等式。投资者可以买入$A$，同时卖出$B$。这种策略不需要任何初始净投资：

\begin{equation}
\Pi_0=P_B-P_A=0.
\end{equation}

其未来的净支付满足要求

\begin{equation}
\Pi_T(s)=X_A(s)-X_B(s)\geq 0,
\qquad \forall s,
\end{equation}
在至少一个状态下具有严格正的支付。这种策略能够在不需要初始投资的情况下产生非负的未来支付，从而构成了一种套利机会。

在无套利条件下，如果能够建立相应交易头寸，一项请求权在所有相关状态下支付至少与另一项相等，并在至少一种状态下严格更多，则其现价必须严格更高。

\paragraph{第二种类型：等价支付请求权之间的套利。}

考虑两个交易的衍生品头寸，$A$以及$B$，在每个相关的未来状态下，现金流与当前状态相同。设$P_A$以及那些使用LaTeX命令、公式表达、交叉引用和数字的格式。$P_B$用LaTeX命令、公式、交叉引用和数字来表示它们的当前价格。$X_A(s)$以及那些使用LaTeX命令、公式表达、交叉引用和数字的格式。$X_B(s)$表示它们在状态中的各自支付$s$. 他们的支付等价性
需要

\begin{equation}
X_A(s)=X_B(s),
\qquad \forall s.
\label{eq:ch11-equivalent-derivative-payoffs}
\end{equation}

如果 $P_A < P_B$，那么投资者可以购买资产 $A$，同时做空资产 $B$。初始的净现金流入为……

\begin{equation}
\Pi_0=P_B-P_A>0,
\end{equation}

虽然未来的净支付是

\begin{equation}
\Pi_T(s)=X_A(s)-X_B(s)=0,
\qquad \forall s.
\end{equation}

该策略能够产生正的初始现金流入，同时不会在未来产生净支付义务。如果$P_A>P_B$，那么反向操作这两种头寸就能获得相应的机会。在无套利条件下，相应的衍生品头寸必须具有相同的当前价格。

\begin{example}{支付占优形成的套利}
{ch11-payoff-dominance-arbitrage}

考虑两个具有相同到期日和结算货币的衍生品请求权。它们在两种可能状态下的当前价格和支付情况如下：

\begin{center}
\begin{tabular}{lcc}
\toprule
\textbf{位置}
& \textbf{状态1}
& \textbf{保留LaTeX命令、公式、交叉引用和数字。  
状态2} \\
\midrule
请求权 A（价格 9） & 12 & 20 \\
请求权 B（价格 9） & 10 & 20 \\
\midrule
长文本用LaTeX命令表示，公式用公式编辑器显示，交叉引用则保持不变。数字部分也照原样保留。 & 2 & 0 \\
\bottomrule
\end{tabular}
\end{center}

买入~A、卖空~B~不需要初始净投资，所得头寸在状态~1~产生正支付，在状态~2~支付为零。因此，相同现价不符合无套利条件。在无套利市场中，请求权~A~的现价必须高于请求权~B。

\end{example}

在实际市场中，理论上的套利机会可能无法完全实现。交易成本、买卖价差、卖空限制、融资要求和抵押品需求，以及交易对手风险，都会影响到套利策略的可行性和净收益。 

在没有任何单一证券的支付能够等同于其他支付的情况下，也可以对证券投资组合的支付一致性进行评估。


\subsection{复制与合成头寸}

复制指构造状态依存支付与目标金融请求权相匹配的投资组合。重现另一请求权支付的组合称为\emph{合成头寸}。精确复制时，两种头寸未来现金流完全相同，因此在无套利条件下必须具有相同现值。

考虑一个拥有$J$种交易证券的市场，以及$S$种可能的未来状态。令$X_j(s)$表示在状态$s$下证券$j$的支付额，令$P_j$表示其当前价格。持有每种证券$\theta_j$单位的投资组合会产生相应的支付额。

\begin{equation}
X_p(s)
=
\sum_{j=1}^{J}\theta_j X_j(s),
\qquad s=1,\ldots,S.
\label{eq:ch11-replicating-portfolio-payoff}
\end{equation}

假设一个目标请求权在状态$s$下的支付为$X(s)$。进行复制操作要求投资组合满足

\begin{equation}
\sum_{j=1}^{J}\theta_j X_j(s)
=
X(s),
\qquad s=1,\ldots,S.
\label{eq:ch11-replication-conditions}
\end{equation}

在矩阵表示法中，这个条件可以表示为

\begin{equation}
\mathbf{X}\boldsymbol{\theta}
=
\mathbf{x},
\label{eq:ch11-replication-matrix}
\end{equation}
其中，$\mathbf{X}$ 是交易证券的支付矩阵；$\boldsymbol{\theta}$ 是投资组合持有量的向量；$\mathbf{x}$ 则是目标请求权的支付向量。

当前复制投资组合的成本为

\begin{equation}
P_p
=
\sum_{j=1}^{J}\theta_j P_j.
\label{eq:ch11-replicating-portfolio-cost}
\end{equation}

如果投资组合能够完全复制目标请求权的合同现金流，那么无套利要求该请求权的当前价格$P_X$等于复制成本：

\begin{equation}
P_X=P_p
=
\sum_{j=1}^{J}\theta_j P_j.
\label{eq:ch11-replication-pricing}
\end{equation}

如果目标请求权价格高于复制成本，投资者可以买入复制组合、卖出该请求权；低于复制成本时，则建立相反头寸。两种情况下，未来支付在每一状态相互抵销，投资者获得当前价格差。

\begin{example}{复制或有保障请求权}
{ch11-contingent-protection-replication}

考虑一个拥有三种可能未来状态的经济体：扩张、放缓和严重衰退。市场上交易着三种证券，它们的当前价格分别为人民币约44.20元、人民币约65.00元和人民币约48.00元。这三种证券在三种状态下的每单位支付情况如下所示。

\begin{center}
\begin{tabular}{@{}lccc@{}}
\toprule
\textbf{状态}
& \textbf{证券 1}
& \textbf{证券 2}
& \makecell{\textbf{无风险}\\\textbf{证券}} \\
\midrule
扩展内容       & 100 & 20  & 50 \\
减速现象        & 50  & 50  & 50 \\
严重衰退 & 20  & 100 & 50 \\
\bottomrule
\end{tabular}
\end{center}

假设在扩张期间，商业中断保护请求权的支付为人民币~0；在放缓期间，支付为人民币~50；而在严重衰退期间，支付为人民币~100。其支付向量为

\[
\mathbf{x}
=
\begin{bmatrix}
0\\
50\\
100
\end{bmatrix}.
\]

三种交易证券的投资组合可以再现这种支付情况。让$\theta_1$、$\theta_2$和$\theta_3$分别表示各自的持有量。复制条件为

\begin{align*}
100\theta_1+20\theta_2+50\theta_3&=0,\\
50\theta_1+50\theta_2+50\theta_3&=50,\\
20\theta_1+100\theta_2+50\theta_3&=100.
\end{align*}

解决这些方程可以得到结果。

\[
\theta_1=-0.625,\qquad
\theta_2=0.625,\qquad
\theta_3=1.
\]

因此，该复制投资组合由以下资产组成：0.625单位的证券1的空头头寸，0.625单位的证券2的空头头寸，以及1单位的无风险证券的多头头寸。

其当前成本为

\begin{align*}
P_p
&=-0.625(44.20)+0.625(65)+48\\
&=61.
\end{align*}

该投资组合在所有三种状态下都能实现该保护请求书的支付。在规定的无套利条件下，该请求书的当前价值因此为人民币约61元。

\end{example}

复制不同于更广泛意义上的套期保值。套期保值可能减少现有头寸对特定风险因素的敏感性，却不消除所有状态依存差异；精确复制则要求组合在每一相关状态下都匹配目标请求权的支付。

\section{衍生品市场架构与风险基础设施}

衍生品市场架构决定了交易关系是如何形成的，风险敞口是如何被抵消的，以及交易对手风险是如何被重新分配的。标准化和集中化的交易可以集中流动性并改善价格发现，而抵押品、轧差和清算则限制了双边风险敞口，并重新组织了市场参与者之间损失的分流。

\subsection{交易场所：场内市场与场外市场}

衍生品可以在有组织的交易所进行交易，也可以通过场外交易（OTC）的方式交易。在交易所交易的衍生品通常采用标准化的合约条款，包括合约规模、到期日、结算条款以及其他规格。标准化使得许多参与者能够使用具有共同条款的合约进行交易，从而使各种头寸能够更方便地互换。

场外衍生品交易是由交易对手之间协商达成的，通常遵循标准化的法律文件，同时允许根据经济条件进行定制。因此，名义金额、到期日、参考变量和结算条款都可以进行调整，以更紧密地匹配特定的风险敞口。然而，更大的定制化可能会使得合约的互换性降低，从而增加了头寸转移的难度。

这一区别并非绝对。有些场外衍生品条款高度标准化，并可能进行中央清算；电子交易平台也缩小了传统上场内执行与交易商执行的差异。因此，交易场所、契约设计与清算安排是相关但不同的市场组织维度。

主要的衍生品交易所提供了标准化合约如何在正式交易场所中组织的具体例子。表~\ref{tab:ch11-major-derivative-exchanges}展示了来自中国和主要国际市场的精选案例。

\begin{table}[H]
\centering
\caption{主要场内衍生品市场示例}
\label{tab:ch11-major-derivative-exchanges}

\begin{tabularx}{\textwidth}{p{3.8cm} p{3.2cm} X}
\toprule
\textbf{交换} &
\textbf{市场焦点} &
\textbf{代表性衍生品} \\
\midrule

中国金融期货交易所 (CFFEX) & 衍生品 & 权益指数期货与期权；政府债券期货 \\

上海期货交易所 (SHFE) &  
商品衍生品 &  
金属、能源及其他商品期货和期权 \\

大连商品交易所 (DCE) & 
衍生品 & 
农业、工业及化工期货与期权 \\

郑州商品交易所（ZCE） &
衍生品市场 &
农业、能源及工业期货与期权 \\

广州期货交易所 (GFEX) &  
商品及新兴衍生品 &  
与特定行业和新兴市场相关的合约 \\

上海和深圳证券交易所 & 
证券衍生品 & 
交易所交易ETF及其他批准的期权 \\

\addlinespace
CME小组 &
金融和商品衍生品 &
利率衍生品、权益指数衍生品、外汇衍生品和商品衍生品 \\

Eurex &
金融衍生品 &
权益指数、利率及其他衍生品 \\

ICE &
金融和商品衍生品 &
能源、利率、权益及其他衍生品 \\

\bottomrule
\end{tabularx}
\end{table}

\subsection{流动性与价格发现}

衍生品市场的设计会影响交易活动和信息的集中程度。标准化合约可以将具有相似风险敞口的参与者聚集在同一工具上，这可能会提高交易的流动性，并使市场价格更加容易观察。而定制化的合约则能够更紧密地匹配各个个体的风险敞口，但交易往往分散在更多不同的头寸上。

液体市场也有助于价格的确定。当参与程度较高且交易频繁发生时，报价和交易价格就能吸收新的信息，为相关工具和基础市场提供参考价值。这种信息功能取决于市场的深度、参与度和交易条件；仅有标准化是不够的，它无法确保市场具有活跃性和信息性。


\subsection{交易对手风险暴露、抵押品与轧差}

衍生品的市场价值会随着市场状况的变化而发生变化。当双边合约对某一方具有正面价值时，该方就对交易对手方拥有请求权，如果交易对手方未能履行义务，该方可能会遭受损失。因此，交易对手方的风险敞口会随着合约的市场价值而变化，并且会随时间推移而发展。

交易对手信用风险不同于衍生品市场风险。市场风险来自基础价格、利率或其他风险因素变化，交易对手信用风险则来自契约义务可能无法履行。当市场变化增加风险暴露、同时交易对手信用质量恶化时，两种风险可能相互作用。

保证金要求提供了随衍生品价值变化控制交易对手风险暴露的机制。两类保证金作用不同：初始保证金是覆盖违约头寸平仓或替代所需期间潜在未来风险暴露的抵押品；变动保证金反映契约当前市场价值变化，并随价值变化在交易对手之间转移。

在交易型期货市场中，合约在每个交易日结束时根据市场情况进行结算。每日的盈亏是根据结算价格的变动来计算的，并计入或从参与者的保证金账户中扣除。对于涉及$n$份合约且合约乘数為$m$的长期期货头寸，每日的保证金变化为

\begin{equation}
VM_t=nm(F_t-F_{t-1}),
\label{eq:ch11-daily-variation-margin}
\end{equation}

其中，$F_t$和$F_{t-1}$分别表示当前和之前的结算价格。持有短期头寸的人需要相反的变动保证金。

每日结算在损益发生时转移损益，而非让市场价值变化积累为无担保契约风险暴露。初始保证金则为参与者违约后平仓期间的潜在损失提供额外保护。

抵押品与轧差是降低双边无担保风险暴露的两项重要机制。抵押品随契约价值变化，为义务提供资产支持。具有法律可执行性的轧差，允许合并考虑同一交易对手之间的正负头寸，而非逐笔分别处理。轧差减少可能需要抵押品支持的风险暴露，抵押品则减少剩余义务中的无担保部分。

\subsection{中央清算与交易对手风险再分配}

抵押品和轧差降低双边关系中的交易对手风险暴露，中央清算则改变这些关系的结构。合格衍生品被接受清算后，中央对手方（CCP）介入原交易双方之间，各清算成员就已清算义务面对~CCP。原本双边存在的风险暴露因此在清算体系中重新组织。

\begin{figure}[htbp]

\centering

\begin{tikzpicture}[
    >=Latex,
    dealer/.style={
        draw,
        rounded corners,
        minimum width=20mm,
        minimum height=8mm,
        align=center,
        font=\small
    },
    ccp/.style={
        draw=FinanceBlue,
        fill=FinanceBlue!8,
        rounded corners,
        minimum width=24mm,
        minimum height=12mm,
        align=center,
        font=\small
    },
    bilateral/.style={
        Latex-Latex,
        line width=0.7pt,
        draw=FinanceBlue
    },
    title/.style={
        font=\small\itshape
    }
]

% Left panel: bilateral structure
\node[title] at (2.5,3.5) {双边市场结构};

\node[dealer] (A1) at (0.8,2.4) {交易商A};
\node[dealer] (B1) at (4.2,2.4) {交易商B};
\node[dealer] (C1) at (0.8,0.4) {交易商 C};
\node[dealer] (D1) at (4.2,0.4) {交易商 D};

\draw[bilateral] (A1) -- (B1);
\draw[bilateral] (A1) -- (C1);
\draw[bilateral] (A1) -- (D1);
\draw[bilateral] (B1) -- (C1);
\draw[bilateral] (B1) -- (D1);
\draw[bilateral] (C1) -- (D1);

\node[align=center, font=\scriptsize]
at (2.5,-0.55)
{多次双边曝光};

% Arrow between structures
\draw[->, line width=0.9pt, gray!70]
(5.25,1.4) -- (6.25,1.4);

% Right panel: centrally cleared structure
\node[title] at (9.0,3.5) {中央统一整理的结构};

\node[dealer] (A2) at (7.3,2.5) {交易商A};
\node[dealer] (B2) at (10.7,2.5) {交易商B};
\node[dealer] (C2) at (7.3,0.3) {交易商 C};
\node[dealer] (D2) at (10.7,0.3) {交易商 D};

\node[ccp] (CCP) at (9.0,1.4) {中心\\交易对手方};

\draw[line width=0.8pt, FinanceBlue] (A2) -- (CCP);
\draw[line width=0.8pt, FinanceBlue] (B2) -- (CCP);
\draw[line width=0.8pt, FinanceBlue] (C2) -- (CCP);
\draw[line width=0.8pt, FinanceBlue] (D2) -- (CCP);

\node[align=center, font=\scriptsize]
at (9.0,-0.55)
{通过CCP重新组织了各部分的布局。};

\end{tikzpicture}

\caption{双边与中央清算下的交易对手关系}

\label{fig:ch11-bilateral-central-clearing}

\end{figure}

中央清算还可以支持跨清算头寸的多边轧差。那些原本需要分别处理的两岸关系中的义务，可以在清算系统中进行抵消，从而减少某些总风险敞口。

因此，中央清算重新分配交易对手风险，而非消除风险。交易对手关系围绕~CCP~组织，而不是分散在许多双边关系中，使~CCP~成为已清算市场的核心节点。

%========================================================
% Summary
%========================================================

\TailModule{summary}{本章小结}

衍生品市场为交易合同提供了制度框架，这些合同的现金流取决于基础资产、利率、指数或特定事件。这些合同产生了或有请求权，使得支付可以在未来的不同状态下进行分配。通过将特定的风险暴露与更广泛的商业和财务头寸分离开来，衍生品有助于市场参与者之间的风险转移和分担。可获得的或有请求权的范围取决于可交易的金融工具以及市场的完整性。

或有请求权的价格通过交易头寸之间的无套利关系相互关联。支付的主导性对相对价格施加了限制，而具有相同状态依存支付的请求权必须具有相同的当前价格。当一项请求权可以通过一组交易证券的投资组合来复制时，其价格必须等于该投资组合的成本。因此，交易工具的可用性决定了哪些请求权可以精确复制，而对于那些无法复制的请求权，则存在一系列符合无套利原则的支付价格范围。

衍生品市场架构决定了交易、结算以及交易对手风险管理的模式。交易所交易市场和场外交易市场在标准化程度及合同灵活性上存在差异。抵押品、交易对手风险管理等因素用于管理交易对手风险，而清算机构则通过中央对手方来重新安排双边交易的风险。这些安排虽然能够重新分配风险，但并不能完全消除风险。

%========================================================
% Key Terms
%========================================================

\TailModule{keyterms}{关键术语}
或有请求权（Contingent claims）；状态依存支付（state-contingent payoffs）；风险暴露分解（exposure decomposition）；风险拆分（risk unbundling）；风险分担（risk sharing）；市场完备性（market completeness）；金融创新（financial innovation）；一价定律（law of one price）；套利（arbitrage）；支付占优（payoff dominance）；复制（replication）；合成头寸（synthetic positions）；场内衍生品（exchange-traded derivatives）；场外衍生品（over-the-counter (OTC) derivatives）；流动性（liquidity）；价格发现（price discovery）；交易对手风险暴露（counterparty exposure）；保证金要求（margin requirements）；抵押品与轧差（collateral and netting）；中央清算与中央对手方（central clearing and central counterparty (CCP)）。

%========================================================
% Financial AI Lab
%========================================================

\TailModule{aiprompt}{金融~AI~实验室}

%========================================================
% Financial AI Lab
%========================================================

\TailModule{aiprompt}{金融~AI~实验室}

%========================================================
% Financial AI Lab
%========================================================

\TailModule{aiprompt}{金融~AI~实验室}

\textbf{实验：在衍生品市场中审计交易对手风险与风险管理}

请一个AI系统来评估一家公司OTC衍生品投资组合中的交易对手风险敞口。假设有一家公司在同一日期拥有以下头寸，所有金额均以人民币百万为单位计算。

\begin{center}
\begin{tabular}{lcc}
\toprule
\textbf{对方银行}
& \textbf{合同/协议}
& \textbf{市场价值} \\
\midrule
交易商A & 合同1 & $+6.0$ \\
交易商A & 合同2 & $-4.0$ \\
交易商B & 合同3 & $+3.0$ \\
\bottomrule
\end{tabular}
\end{center}

与交易商A签订的 두 계약은 동일한 법적으로 구속력 있는 종료轧差 협정에 따릅니다. 거래商A는 적용 가능한抵押品折扣率에 따라 회사에 약 120만 RMB의 변동 마진을 지불했습니다. 반면, 거래商B와의 계약은 별도의轧差 협정에 따르며, 아직 어떠한抵押品도 받지 않았습니다.

要求人工智能系统计算该公司当前未受担保的对手方风险敞口，并评估不同的风险管理安排对该风险敞口的影响。

将答案分为四个阶段进行审核。

\begin{enumerate}

\item \textbf{计算交易对手的净风险敞口。}

确定在应用具有法律约束力的轧差后，对每个交易商的净当前敞口。考虑所收到的抵押品，并解释为什么与不同交易商持有的头寸不能自动相互抵消。

\item \textbf{区分市场风险与交易对手信用风险。}

解释基础价格的变化如何影响合约的市场价值，以及公司对交易对手的潜在风险。阐明为什么正的置换价值会使得公司面临交易对手无法履行义务的风险。

\item \textbf{评估保证金和抵押品安排。}

请保留LaTeX命令、公式、交叉引用和数字的格式不变。

区分初始保证金与变动保证金的区别。
解释每日市价跟踪机制和变动保证金机制如何影响当前的持仓状况，同时说明初始保证金如何能够在清算期间保护投资者免受未来潜在风险的影响。指出即使进行了抵押担保，仍可能存在的风险。

\item \textbf{评估中央清算的作用。}

解释引入中央交易对手如何改变原有的双边关系。
讨论多边轧差、保证金要求、
违约风险以及交易对手风险在清算系统中的集中情况。
区分交易对手风险的重新分配与市场风险消除的区别。

\end{enumerate}

最终报告应验证抵押品的计算结果，明确与轧差和抵押品相关的法律和运营假设，并区分风险降低与风险转移以及风险集中。

%========================================================
% Discussion Questions
%========================================================

\TailModule{questions}{讨论题}

\begin{enumerate}

\item 衍生品如何在不改变公司基础业务活动的情况下，对选定的金融风险进行调节？

\item 市场完整性如何影响复制或有请求权的能力？在无法完全复制的情况下，可能会存在哪些定价不确定性？

\item 为什么定制化的OTC合约能够在增加抵押品、保证金要求以及双边轧差的重要性同时，使得对基础资产的风险敞口更加匹配呢？

\item 中央清算是如何重新安排交易对手关系的？为什么这会使中央交易对手的风险管理和韧性变得重要？

\end{enumerate}

%========================================================
% Applied Problems
%========================================================

\TailModule{appprobs}{应用题}

\begin{enumerate}

\item \textbf{状态依存的支付和远期套期保值。}

一位商品生产商预计在日期$T$出售一单位该商品。在该日期的现货价格可能为70、100或140。该生产商签订了一份现金结算的短期合约，合约的交付价格为100，到期日亦然。

\begin{enumerate}

\item 计算每个状态下短期远期头寸的支付情况。

\item 计算生产者在每个状态下的物理销售和未来头寸所带来的现金流总和。

\item 解释远期合约如何改变生产者的现金流在不同州的分配情况。

\end{enumerate}


\item \textbf{暴露分解与风险分离。}

该公司在日期$T$时将购买500个单位的进口商品。以外币计价的商品价格表示为$P_T$，而单位外币在本国货币中的价格则表示为$S_T$。该公司使用本国货币进行购买时的成本为……

\[
C_T=500P_TS_T.
\]

\begin{enumerate}

\item 确定影响企业购买成本的两大市场风险因素。

\item 解释每个风险因素的变化如何影响公司的本国货币现金流出。

\item 假设该公司使用的衍生品仅与商品价格$P_T$挂钩。哪些风险敞口会发生变化，哪些则保持不变？

\end{enumerate}


\item \textbf{或有请求权的复制。}

在双状态市场中交易了两份证券。它们当前的价格和支付情况如下：

\begin{center}
\begin{tabular}{lccc}
\toprule
\textbf{证券}
& \textbf{当前价格}
& \textbf{上状态}
& \textbf{关闭状态} \\
\midrule
证券 A & 8 & 12 & 6 \\
证券 B & 9 & 4 & 10 \\
\bottomrule
\end{tabular}
\end{center}

在“上状态”下，或有请求权可获得20的赔偿；在“下状态”下，可获得2的赔偿。

\begin{enumerate}

\item 确定需要持有的证券A和B的数量，以履行该或有请求权。

\item 请标明哪些部分是长句，哪些部分是短句。同时保留所有LaTeX命令、公式表达、交叉引用以及数字的内容。

\item 计算当前投资组合的成本。

\item 解释为什么在无套利条件下，请求权的当前价值必须等于复制成本。

\end{enumerate}


\item \textbf{套利机会与价格一致性。}

考虑两对衍生品请求权。

对于第一对请求权，请求权A和B在状态1时的支付均为15，在状态2时的支付均为8。它们当前的价值分别为13和10。

对于第二对请求权，C和D的当前价格都是11。它们的支付情况如下：

\begin{center}
\begin{tabular}{lcc}
\toprule
\textbf{请求权}
& \textbf{状态1}
& \textbf{保留LaTeX命令、公式、交叉引用和数字。  
状态2} \\
\midrule
请求权 C & 14 & 20 \\
请求权 D & 12 & 20 \\
\bottomrule
\end{tabular}
\end{center}

\begin{enumerate}

\item 对于第一对情况，确定套利策略，并计算其初始现金流以及未来的支付金额。

\item 对于第二对情况，确定套利策略，并计算其初始现金流以及未来的支付金额。

\item 为什么预期回报较高的投资不一定构成套利机会呢？

\end{enumerate}


\end{enumerate}
%========================================================
% Further Reading
%========================================================

\TailModule{furtherread}{延伸阅读}

\begin{itemize}

\item Duffie, Darrell, and Chi-Fu Huang (1985),
``Implementing Arrow-Debreu Equilibria by Continuous Trading of Few Long-Lived Securities（通过少数长期证券的连续交易实现阿罗—德布鲁均衡）,''
\emph{Econometrica}, 53(6), 1337--1356.

\end{itemize}

\ChapterTail{summary,keyterms,aiprompt,questions,appprobs,furtherread}

\part{金融机构}
\chapter{商业银行：信用中介与流动性创造}
\label{chap:Commercial_Banks_Credit_Intermediation_and_Liquidity_Creation}

\begin{openingcase}{为什么溶剂银行仍然可能面临现金短缺的问题？}

假设一家银行的资产为人民币1000亿。其大部分资产都是向家庭和企业的贷款，此外还有几年后才会到期债券的持有。在负债方面，该银行的大部分资金来自客户可以随时提取或短时间内即可取款的存款。该银行还有能够吸收损失的权益资本。

假设银行的借款人继续还款，而银行仍然能够维持偿付，其资产价值超过负债。即便如此，突然出现的提款潮仍可能导致立即的资金需求。快速出售长期资产可能需要降价出售，而用市场资金替代存款则可能成本较高或无法获得。

同样的资产负债表结构使得银行能够在提供流动性存款的同时，也能进行非流动贷款的融资。然而，当提取资金的需求突然增加时，这种结构也会带来风险。

\end{openingcase}


\begin{tcolorbox}[
  enhanced,
  colback=FinanceBlue!6,
  colframe=FinanceBlue,
  boxrule=0.6pt,
  sharp corners,
  left=8pt,
  right=8pt,
  top=7pt,
  bottom=7pt,
  before skip=1em,
  after skip=1.6em,
  title=\textbf{核心问题},
  coltitle=white,
  colbacktitle=FinanceBlue,
  fonttitle=\bfseries,
  attach boxed title to top left={xshift=6pt,yshift=-2mm},
  boxed title style={
    boxrule=0pt,
    sharp corners,
    left=5pt,
    right=5pt,
    top=2pt,
    bottom=2pt
  }
]
\begin{enumerate}[
  leftmargin=*,
  label=\textcolor{FinanceBlue}{\textbf{\arabic*.}},
  itemsep=0.55em,
  topsep=0.2em
]


\item 商业银行如何利用其资产负债表将存款人和其他提供者的资金转化为贷款和其他金融资产呢？

\item 为什么银行能够通过信息的生成以及合同的设计来创造价值？这包括对风险的筛选、监控，以及关系型贷款等方面。

\item 期限转换、流动性创造、杠杆效应以及资金结构如何在为银行创造价值的同时，也使其面临信用、利率和流动性风险？

\end{enumerate}
\end{tcolorbox}

\section{银行资产负债表与资产负债表中介}

商业银行的核心是一种资产负债表中介：通过存款和其他负债筹资，取得贷款及其他金融资产，并持有权益资本吸收损失。

\subsection{银行资产负债表：资产、融资与交易}

银行的资产负债表显示了其所持有的资产以及这些资产的融资方式。商业银行通常持有贷款、证券、储备金以及其他流动性资产。这些资产是通过存款、批发融资以及其他债务以及权益资本来融资的。

资产负债表的经济含义，体现在各类头寸的构成和契约特征中。贷款使银行承受借款人经营表现的影响。证券可提供收益、流动性和抵押品，准备金及其他流动资产支持支付结算和短期流动性需求。融资端的存款与批发借款，在期限、稳定性、重新定价和提款条件上不同。权益在债权人请求权受损之前吸收损失。

\begin{table}[H]

\centering

\caption{简化的商业银行资产负债表}

\label{tab:ch12-simplified-bank-balance-sheet}

\begin{tabularx}{0.82\textwidth}{@{}X X@{}}

\toprule

\textbf{资产} & \textbf{负债与权益} \\

\midrule

向家庭和企业的贷款 & 交易和储蓄存款 \\

政府债券及其他证券 & 定期存款和定期存单 \\

储备及其他流动性资产 & 批发和市场融资 \\

& 权益资本 \\

\bottomrule

\end{tabularx}

\end{table}

商业银行通过持有对借款人和其他发行者的请求权来进行中介业务，同时向储户和其他资金提供者发行自己的负债。双方的资产和负债在到期时间、流动性、风险和重新定价特性方面可能存在差异。双方资产的管理是共同进行的：贷款决策会影响银行资产的规模和构成，而存款流动、批发融资和权益则会影响这些资产的融资方式以及银行的可用流动性和损失吸收能力。

银行交易直接体现这些相互作用。银行发放贷款，并将款项记入借款人在本行的存款账户时，同时记录一笔新增贷款资产和相应存款负债。

\begin{example}{从贷款发放到结算}{loan-origination-balance-sheet}

在发放新贷款之前，银行拥有以下简化的资产负债表，单位以百万人民币计：

\begin{center}

\begin{tabular}{
@{\hspace{0.8em}}
p{0.25\textwidth} r
@{\hspace{2.0em}}
p{0.24\textwidth} r
@{\hspace{2.5em}}
}

\toprule

\multicolumn{2}{c}{\textbf{资产}} &
\multicolumn{2}{c}{\textbf{负债与权益}} \\

\cmidrule(lr){1-2}\cmidrule(lr){3-4}

贷款 & 70 & 存款 & 82 \\

证券 & 15 & 批发融资 & 8 \\

储备金 & 15 & 权益 & 10 \\

\midrule

\textbf{总计} & \textbf{100} &
\textbf{总计} & \textbf{100} \\

\bottomrule

\end{tabular}

\end{center}

银行提供了新的人民币约500万元的贷款，并将资金存入借款人的存款账户中。在申请过程中，

\[
\Delta L=+5,
\qquad
\Delta D=+5.
\]

总资产和总负债加上权益均增加至约1.05亿元人民币。权益在初始时保持不变。

假设借款人随后将人民币~500~万元存款转入另一家银行的账户。如果通过准备金结算，且没有抵销性流入，贷款发放银行的存款负债和准备金均减少~500~万元，贷款仍保留在资产端。

\end{example}

贷款发放与支付结算是不同的资产负债表事件。发放贷款时创造存款，本身并不提供此后资金离开银行所需的结算资产。银行遭遇存款净流出时，可能需要动用流动资产、获得替代融资，或调整其他头寸。

\subsection{银行资本、杠杆与资产负债表约束}

在这一资产负债表结构中，权益资本是满足负债后对银行资产的剩余请求权。资产价值下降先减少权益，之后才影响债权人请求权。如果损失足以耗尽资本缓冲，债权人就可能承受损失。

一种简单的资产负债表杠杆指标是资产权益比~$A/E$。杠杆越高，给定权益支持的资产规模越大，权益对资产价值变化也越敏感。负债不变时，资产损失反映为权益减少，因此资产较小比例的下降，也可能造成权益大得多的比例降幅。

\begin{example}{利用与损失吸收}{bank-leverage-loss}

该银行的资产为人民币约1000亿元，负债为人民币约920亿元，权益为人民币约80亿元。其资产与权益的比率为

\[
100\div 8=12.5.
\]

假设资产价值下降2\%，即约20亿元人民币，而负债保持不变。权益则从约80亿元人民币下降到约60亿元人民币，减少了25\%。

这个例子说明了杠杆效应如何使得权益对资产价值的变化更加敏感，而总资产则相对不那么敏感。

\end{example}

在银行资产负债表扩张的过程中，有几个限制因素需要考虑。额外的贷款或其他资产增长需要资金支持，同时也会增加资本和流动性的需求。审慎的监管要求还会进一步限制杠杆率、风险加权暴露、流动性状况以及风险集中度。究竟哪个限制因素最为重要，则取决于银行的资产负债表结构以及市场状况。

这些限制因素影响着资产负债表的规模和构成。扩大贷款规模会增加信贷风险，可能需要额外的资金和资本支持；而持有更多流动性资产虽然能增强流动性能力，但可能会降低预期回报。过度依赖短期批发融资可以提供灵活的融资方式，但同时也会增加再融资的风险。

\begin{tcolorbox}[
  enhanced,
  breakable,
  colback=white,
  colframe=FinanceBlue,
  boxrule=0.7pt,
  sharp corners,
  borderline west={2.2pt}{0pt}{FinanceBlue},
  title=\textbf{财务展望：资本与流动性发挥着不同的作用},
  coltitle=white,
  colbacktitle=FinanceBlue,
  fonttitle=\bfseries,
  before skip=1.2em,
  after skip=1.2em
]

当银行资产的价值下降时，资本能够吸收这些损失。流动性则确保了在支付和提取发生的时候能够顺利进行。

银行可能权益为正，却没有足够流动资源应付即时资金流出。反之，银行也可能持有大量流动资产，但损失已大到削弱其资本状况。

资本与流动性之间存在关联，但它们分别涉及资产负债表中的不同脆弱性方面。

\end{tcolorbox}



\section{信用评估与贷款决策}

贷款通过一种信贷流程进入商业银行的资产负债表，该流程将借款人的信息转化为贷款决策。商业银行利用与借款人相关的信息以及公开可用的信息来评估信用质量、设定信用限额和贷款条件，并在情况发生变化时重新评估风险敞口。因此，信贷分配是一个持续的过程，而不是在贷款发起时的单一决策。

在银行信贷关系的各个阶段，主要的信息需求以及相应的贷款决策都在表~\ref{tab:ch12-credit-process}中有所总结。

\begin{table}[H]
\centering
\caption{关于银行信贷关系的信息与决策}
\label{tab:ch12-credit-process}
\begin{tabularx}{0.90\textwidth}{@{}p{0.22\textwidth} X X@{}}
\toprule
\textbf{阶段} & \textbf{信息重点} & \textbf{银行决策} \\
\midrule
初步评估
& 偿还能力、财务状况、现金流稳定性、抵押品以及经营状况
& 批准、拒绝或限制拟定的发布方式。 \\
合同设计
& 资金的预期使用方式、还款模式、不确定性因素以及银行所能提供的保护措施。
& 设定金额、到期日、还款条款、定价、抵押品以及契约限制条款 \\
监控与更新
& 支付行为、账户活动情况、最新的财务信息以及业务状况的变化。
& 维护、扩展、强化、重构或终止相关内容的展示。 \\
\bottomrule
\end{tabularx}
\end{table}

\subsection{借款人信息与信用评估}

在信用评估中，会考虑借款人的履行合同义务的能力以及所承担风险的大小。这一评估基于多种信息来源进行，每一种信息来源都提供了关于信用质量不同方面的证据。

财务报表提供了关于盈利能力、杠杆率、流动性和资产负债表结构的信息。现金流记录有助于评估运营收入是否能够支持债务偿还，而还款历史和账户活动则提供了有关还款行为的证据。行业状况、客户集中度、所有权结构以及贷款的目的等因素，也会进一步影响预期的还款能力。如果合同规定的付款未能按约定履行，抵押品则可以提供额外的回收来源。

银行通过财务分析、评分系统、统计模型以及专业判断来整合这些信息。每种方法的重要性因借款人和贷款类型而异。标准化的信息有助于进行更公平的评估，而当还款能力取决于那些仅凭财务比率难以衡量的因素时，关于借款人的具体信息就变得更为重要了。

这些信息随后会被转化为内部的信用评估结果，这种评估可以采用评级、风险等级或其他银行在信贷过程中使用的分类方式来表示。不过，这种评估并不能机械地决定是否给予贷款。它只是为关于信贷额度、合同条款以及后续监控的强度和方式等决策提供参考依据而已。

\begin{example}{从借款人信息到信用评估}{borrower-information-credit-assessment}

有两家公司申请了类似的营运资金支持计划。两家公司的年度销售额均约为2000万元人民币，其杠杆比率也相当相似。

A公司能够产生稳定的运营现金流，拥有足够的债务偿还能力，并且其还款记录一直都很稳定。而B公司则面临更大的现金流波动风险，其履行债务义务的能力较为有限，而且有延迟付款的历史记录。

银行通常会认为公司A更具信用价值，因此会给其更高的内部信用评级，同时违约概率也会较低。在其他条件相同的情况下，公司A将能够获得更高的信贷额度、更低的贷款利差，以及相对较宽松的抵押品要求。

\end{example}



\subsection{贷款条款、信用限额与风险暴露}

信用评估随后会被转化为贷款决策。银行会设定信用额度和其他贷款条款，以确定可以承担的风险规模以及这种风险随时间的变化情况。主要的合同条款及其对银行信用风险的影响在表~\ref{tab:ch12-loan-terms-exposure}中进行了总结。这些条款共同决定了银行对借款人的风险敞口的规模、久期和条件。

\begin{table}[htbp]

\centering

\caption{选定的贷款条款及其对银行风险敞口的影响}

\label{tab:ch12-loan-terms-exposure}

\begin{tabularx}{0.90\textwidth}{@{}p{0.26\textwidth} X@{}}

\toprule

\textbf{贷款期限} &
\textbf{曝光影响} \\

\midrule

贷款金额或信用额度 &
规定了合同中允许使用的最大信用额度 \\

到期日 &
设定了需要偿还或续签的期限。\\

偿还结构 &
规定了存续本金随时间的变化情况 \\

利率结构 &  
决定了是否以及以何种方式重新设定合同利率 \\

抵押品和契约限制条款 &
通过抵押资产和对借款人行为的限制或要求来提供合同保护 \\

绘图与更新条件 &  
规定在何时可以提取额外资金，以及信用额度是否会在当前期限之后继续有效 \\

\bottomrule

\end{tabularx}

\end{table}

偿还和提款条款决定了在授予信贷后，存续的敞口会如何演变。用$B_t$表示第$t$期末的存续本金，$D_t$表示该期间的新提款额，$R_t$表示本金偿还额。一个简单的平衡关系为：

\begin{equation}
B_t=B_{t-1}+D_t-R_t.
\label{eq:ch12-loan-balance-dynamics}
\end{equation}

对于定期贷款，本金通常在借款时一次性支付，之后随着还款的进行而逐渐减少。而循环信贷机制则有所不同：借款人可以在合同约定的信用额度内进行提款、还款，并继续提款。因此，信用额度与当前存续的本金数额是分开的。

这一区别形成两个相关的风险暴露指标。\emph{已融资风险暴露}是已经发放且仍未偿还的金额；\emph{未提用承诺}是承诺额度中尚未使用的部分，借款人按契约进一步提款后，可能转化为已融资风险暴露。因此，即使契约信用限额不变，使用率提高也可能增加银行的已融资风险暴露。

\begin{example}{循环信贷融资与银行风险敞口}{revolving-credit-exposure}

一家公司获得了为期一年的循环信贷额度，额度为人民币约1000万元。年初时，有人民币约400万元尚未偿还。在下一个季度内，该公司又借入了人民币约300万元，并偿还了人民币约100万元。使用公式~\ref{eq:ch12-loan-balance-dynamics}来计算。

\[
B_t=4+3-1=6
\]

百万人民币。

因此，该银行拥有约600万的人民币资金敞口，而根据这项融资安排，还有约400万的人民币资金并未被使用。如果借款人日后能够使用掉这部分未使用的资金，那么资金敞口将会增加，但合同规定的信贷限额则不会发生变化。

\end{example}

贷款条款确定了初始的敞口规模，但一旦信贷关系建立后，借款人的状况可能会发生变化。因此，新的信息可能会导致银行重新评估该敞口，进一步考虑是否需要追加贷款，以及是否应该继续维持这一贷款关系。

\subsection{监督、续贷与信贷调整}

通过监控，银行能够利用还款情况、账户活动、财务信息、抵押品价值、契约限制条款的遵守情况以及业务状况的变化来更新其评估。契约限制条款指的是借款人是否继续满足贷款合同中规定的限制和要求。

这些指标的变动可能会引发信用审查。未能按时支付、经营现金流的持续疲软，或者严重的契约限制条款违反，都可能表明状况恶化；而更良好的支付记录或财务状况则可能支持现有信用的持续。进行此次审查的目的是为了确定银行当前的风险敞度是否与其对借款人的最新评估一致。

续签决策是一个重要的节点。银行在决定是否继续、扩大、减少或终止这一信贷安排时，可以将当前的信息与借款人在整个贷款期限内的表现结合起来考虑。因此，即使没有发生违约情况，通过借贷关系积累的信息仍然会影响后续的信贷决策。

信用调整取决于重新评估的结果以及合同中规定的权利条款。在贷款期限内，银行可以在允许的情况下限制借款人的额外提取行为，或者要求借款人采取特定的保护措施。在贷款到期时，银行可以重新确定信用限额以及其他合同条款。因此，通过定期监测和重新评估，银行能够将对借款人状况变化的影响转化为对银行风险敞口的变化调整。

\begin{example}{从监控到信用调整}{monitoring-credit-adjustment}

某公司拥有高达800万人民币的循环信贷额度。销售额基本保持稳定，但经营现金流入减少，客户应收账款增加，因此信用额度的使用比例从40\%上升到了85\%。 

这些发展本身并不意味着默认情况的发生，但综合来看，它们可能会引发对信贷状况的审查。银行可以重新评估现金流的压力来源，更新对借款人的评估，并重新考虑现有的风险状况。

如果银行认为这种状况只是暂时的，那么银行可以继续维持该贷款安排。但如果情况持续恶化，根据合同条款，银行可能会限制进一步的提款行为，或者要求采取额外的保护措施。在贷款期限到期时，银行可能会降低信用额度，甚至拒绝继续提供全额贷款。

\end{example}

因此，信贷决策不仅会影响银行资产的构成，还会影响到为这些资产提供资金支持所需的资金数额。这一因素直接关系到银行的融资结构和流动性管理问题。

\section{银行融资与流动性创造}

商业银行通常通过较短期限且流动性较高的负债来为长期投资提供资金。创造流动性取决于银行的融资结构以及其应对现金流出需求的能力。

\subsection{存款融资与批发融资}

银行融资的构成非常重要，因为不同的融资来源在提取权、到期时间、重新定价以及集中程度方面存在差异。商业银行主要依赖存款以及各种形式的批发融资，而这些资金来源对利率变化、市场状况以及信心的变化反应各不相同。

存款是商业银行负债，但不同产品和客户的资金行为不同。交易性存款与储蓄存款可能在较短通知期内支取，定期存款则通常提供期限更明确的融资。稳定性还取决于存款人集中度、存款保险覆盖、利率敏感性及对银行风险的判断。零售存款往往更分散，大额企业或机构余额则可能形成更集中的融资风险暴露。

批发融资包括银行间借款、定期存单、担保借款以及其他基于市场的负债。当资产负债表增长或支付流需要额外的资金时，批发融资可以补充存款资金。其有用性不仅取决于成本，还取决于期限和续期条件：短期融资需要更频繁的再融资，而当市场状况恶化时，这类融资可能会变得不那么可用。

\begin{table}[H]

\centering

\caption{资金来源及选定的流动性特征}

\label{tab:ch12-bank-funding-characteristics}

\begin{tabularx}{0.92\textwidth}{@{}p{0.23\textwidth} X X@{}}

\toprule

\textbf{资金来源} &
\textbf{合同特征} &
\textbf{对银行融资的影响} \\

\midrule

交易和储蓄存款&通常可以在短时间内提供给客户使用，具体详情请参阅产品条款说明。&稳定性取决于客户结构、集中度、利率以及市场信心等因素。\\

定期存款 &  
资金在指定的到期日或通知期限内保持不变 &  
在到期前，资金的供应较为可预测；而随着到期日的临近，就需要做出续存决策了 \\

定期存款 &  
可交易或不可交易的银行负债，具有特定的到期日 &  
融资成本和滚动投资取决于市场条件和投资者需求 \\

银行间借款 &
从其他金融机构获得的短期借款 &
提供灵活的融资方式，但可能需要频繁进行再融资操作 \\

抵押品融资市场融资 &
以抵押品为担保的借款 &
资金供应取决于可用的抵押品和融资条件 \\

\bottomrule

\end{tabularx}

\end{table}

融资结构与流动性创造紧密相关。银行可以用为持有人提供较高流动性的负债，融资相对缺乏流动性的资产。维持这种结构，需要足够的流动性能力，以应付提款、到期负债及其他到期现金流出。

\subsection{流动性能力与资产负债表调整}

许多银行的负债要求在银行资产到期之前进行结算。银行的流动性能力取决于在相关时间范围内能够转化为结算资金的资源。这些资源可以包括现金和储备金、可出售资产的收益，以及额外的借款能力。

一个简单的表示方式如下：保留LaTeX命令、公式、交叉引用以及数字的内容。

\begin{equation}
LC_t=C_t+M_t+F_t^{\mathrm{avail}},
\label{eq:ch12-liquidity-capacity}
\end{equation}

其中，$C_t$表示立即可用的现金和储备金，$M_t$表示在一段时间内可以从可销售资产中合理获得的现金，而$F_t^{\mathrm{avail}}$则表示可以合理获得的额外资金。为了避免重复计算，当抵押品或借款能力依赖于相同的资产时，应调整这些金额。因此，流动性能力取决于时间范围以及必须筹集资金的状况。

从某项资产中获得的现金金额可能与其报告的市场价值有所不同。交易成本、结算时间、价格变动以及抵押品使用的限制都会影响到资产销售所得的资金。同样，借款能力也取决于抵押品的可用性和价值、交易对手的借款意愿以及当前的融资条件。

当出现现金流出时，银行可以调整其资产负债表中的各项指标。它可以使用储备金、出售可交易的证券，或者通过抵押符合条件的资产来获得有担保的借款，同时还可以获得其他替代性融资，或者减少新资产的增长。这些应对措施通过不同的渠道来满足流动性需求，并会在之后产生不同的资产负债表状况。

\begin{example}{遇到存款流出情况}{meeting-deposit-outflow}

一家银行拥有约60亿元人民币的储备金，约80亿元人民币的可用证券尚未作为抵押品，还有约40亿元人民币的担保借款能力，这些能力由其他符合条件的资产来支持。该银行出现了约100亿元人民币的净存款外流。

一种可能的调整是使用60亿元人民币的储备金，并获得40亿元人民币的担保借款。另一种方法是使用储备金出售40亿元人民币的可交易证券。这两种方案都能提供所需的流动性，但它们对资产负债表的影响却有所不同。

第一种应对方式会增加抵押品所对应的负债，同时也会加重资产负担。第二种应对方式则会减少银行持有的可交易证券。因此，流动性管理不仅涉及满足当前的资金流出需求，还涉及到之后仍然可用的流动性资源。

\end{example}

流动性能力决定了银行能够多快地应对当前的现金流出需求。如果流出金额较大或持续时间较长，那么银行就需要进行更大的资产负债表调整。

\subsection{融资流出与流动性调整}

当存款从银行流出，或者到期需要偿还的负债在没有新资金补充的情况下被偿还时，就会出现资金流出。这种资金流出会减少银行的资金供应，从而产生现金需求。除非有资产端的现金流入或新的资金来源来抵消这种需求。

用 $O_t$ 表示在较短时间内发生的资金流出，$I_t$ 表示来自到期或出售资产的现金收入，而 $F_t$ 则表示在同一时期获得的新的资金。现金和储备金的变化可以表示为：

\begin{equation}
\Delta C_t=I_t+F_t-O_t.
\label{eq:ch12-liquidity-adjustment}
\end{equation}

该关系图展示了资金流出是如何影响银行的流动性状况的。当资产端的现金流入和新融资不足以弥补资金流出时，现金和储备金就会减少。而当替代资源充足时，银行就可以在不耗尽现有现金储备的情况下应对资金流出问题。

所需的调整不仅取决于资金流出金额，还取决于其时机和集中程度。对于大型批发负债的不再续期来说，可能需要立即进行替代，而如果将相同金额的提取分散到多个存款人身上，则情况会有所不同。相关的问题是，在可用的时间范围内，需要替换多少资金。

预期回报

因此，资金流出会改变银行的流动性水平以及资产负债表的结构。更为显著或持续性的调整还会影响到银行的融资成本、资产持有情况，以及银行风险的其他方面。

\section{银行综合风险管理}

银行风险是通过共同的资产负债表相互影响的。信贷质量、利率或融资方面的冲击可能会影响到企业的收益、资本和流动性，而由此产生的资产负债表调整则可能会带来更多的风险。

\subsection{贷款损失、收益与资本}

贷款投资组合的恶化可能会降低银行的收益，并通过留存收益削弱其资本状况。这一过程从确认信用损失开始，最终影响到当前的收入和权益。

贷款损失拨备费用是在适用会计框架下确认信用损失的一项费用。计提拨备减少当期收益，相应贷款损失准备减少贷款净账面金额。此后若用已有准备核销贷款，贷款总额和准备同时减少；损失已通过较早的拨备费用影响收益。

简化的银行收入报表可以这样编写：

\begin{equation}
\textit{NI}=\textit{NII}+\textit{NIR}-\textit{OE}-P_L-T,
\label{eq:ch12-bank-net-income}
\end{equation}

其中，$\textit{NI}$表示净收入，$\textit{NII}$表示净利息收入，$\textit{NIR}$表示非利息收入，$\textit{OE}$表示运营费用，$P_L$表示贷款损失准备费用，而$T$则表示税费。该方程揭示了准备费用对银行收益的直接影响。

如果股利和其他分配用$D$表示，而该期间发行的新权益用$\textit{NE}$表示，那么账面权益的变动大致如下：

\begin{equation}
E_1=E_0+\textit{NI}-D+\textit{NE},
\label{eq:ch12-equity-evolution}
\end{equation}
除了对权益的其他直接调整之外，这两个关系展示了传导机制：

\[
\begin{array}{c}
\text{\emph{损益表}} \\[2pt]
\text{贷款恶化}
\rightarrow
\text{相关费用条款}
\rightarrow
\text{较低的净收入}
\\[5pt]
\Downarrow
\\[-1pt]
\text{\emph{资产负债表}} \\[2pt]
\text{降低留存收益}
\rightarrow
\text{较弱的权益状况}
\end{array}
\]

\begin{example}{资源配置、收益以及资本}{provisioning-earnings-capital}

一家银行年初拥有6亿人民币的权益。在扣除准备金和税费后，其实际收益为1.2亿人民币。该银行记录的贷款损失准备金额为7000万元人民币，同时支付了2000万元人民币的股利。忽略税费和其他权益调整因素。

净收入为5000万人民币，年末的权益为

\[
600+50-20=630
\]

百万人民币。

如果在相同假设下，准备金费用上升至1.3亿元人民币，那么净利润将变为负1,000万元人民币，年末权益将降至5.7亿元人民币。因此，更大的已确认贷款损失会减少收益，而在其他条件不变的情况下，银行的权益将不足以吸收后续的损失。

\end{example}

\subsection{重新定价与资产负债错配}

利率变化通过收入渠道和估值渠道对银行产生影响。这种影响取决于资产和负债的利率何时重新设定，以及它们的经济价值对市场收益率变化的敏感程度。

让 $\textit{RSA}$ 和 $\textit{RSL}$ 分别表示在给定时期内需要重新估价的利率敏感型资产和负债的金额。对于某个特定的资产负债表状况，净利息收入的大致变化为

\begin{equation}
\Delta \textit{NII}\approx \textit{RSA}\cdot\Delta r_A
-RSL\cdot\Delta r_L,
\end{equation}
其中，$\Delta r_A$和$\Delta r_L$分别表示相应资产和负债的变动幅度。如果这些变动幅度大致相同，即$\Delta r_A\approx\Delta r_L\approx\Delta r$，那么表达式就可以简化为

\begin{equation}
\Delta \textit{NII}\approx (\textit{RSA}-\textit{RSL})\Delta r
=G\Delta r.
\label{eq:ch12-net-interest-income-gap}
\end{equation}

重新定价缺口~$G$~衡量利率敏感性资产与负债之差。正缺口意味着利率上升通常提高净利息收入，负缺口则相反。

这一近似刻意作了简化。存款利率对市场利率的调整可能不到一比一，贷款可能仅在约定日期重设利率，或参照不同基准，客户行为也可能改变未偿余额。因此，重新定价缺口只是短期收入敏感性的初步近似，并非利率风险暴露的完整度量。

经济价值分析考察错配的另一维度。即使近期现金流不变，固定利率资产与负债对市场收益率变化的敏感性也可能不同。设~$D_A$~和~$D_L$~分别表示资产和负债的有效久期，$A$~和~$L$~表示其当前经济价值。久期缺口为

\begin{equation}
D_{GAP}=D_A-\frac{L}{A}D_L.
\label{eq:ch12-duration-gap}
\end{equation}

有效久期衡量的是价值对收益率微小变化的近似百分比敏感性。对于市场收益率的微小平行变化，

\begin{equation}
\Delta E\approx -A\,D_{GAP}\,\Delta y,
\label{eq:ch12-duration-gap-equity}
\end{equation}

其中，$\Delta E$~表示由此产生的权益经济价值变化。因此，在上述近似条件下，正久期缺口意味着收益率上升可能降低权益经济价值。

\begin{tcolorbox}[
  enhanced,
  breakable,
  colback=white,
  colframe=FinanceBlue,
  boxrule=0.7pt,
  sharp corners,
  borderline west={2.2pt}{0pt}{FinanceBlue},
  title=\textbf{财务展望：收入效应与经济价值效应可能存在差异},
  coltitle=white,
  colbacktitle=FinanceBlue,
  fonttitle=\bfseries,
  before skip=1.2em,
  after skip=1.2em
]

同样的利率变动可能会以不同的方式影响银行的收入和经济价值。例如，利率上升可能会增加短期的净利息收入，但同时会降低权益的经济价值。

\end{tcolorbox}

\subsection{风险通过资产负债表的相互作用}

银行风险是相互影响的，因为信用损失、利率波动以及资金供应方面的冲击都会对同样的资产负债表产生影响。

某一领域的冲击可能改变另一领域的风险暴露。贷款损失可能削弱收益与资本；利率变化可能同时改变资产价值与融资成本；融资流出可能要求出售资产或借入替代资金。这些应对又可能影响抵押品可得性、已实现损益、融资成本，以及银行吸收进一步冲击的能力。

这些相互作用的强度取决于资产和负债的构成、融资期限、抵押品的可用性、套期保值措施，以及银行的初始资本和流动性状况。

\begin{figure}[H]
\centering

\begin{tikzpicture}[
  >=Latex,
  box/.style={
    draw,
    rounded corners,
    align=center,
    text width=8cm,
    minimum height=0.95cm,
    inner sep=5pt
  },
  main/.style={
    draw=FinanceBlue,
    fill=FinanceBlue!8,
    rounded corners,
    align=center,
    text width=9.6cm,
    minimum height=1.0cm,
    inner sep=5pt
  },
  flow/.style={
    -Latex,
    line width=0.8pt,
    FinanceBlue
  }
]

\node[box] (shock) at (0,3.0)
{信用、利率或融资冲击};

\node[main] (bs) at (0,1.3)
{资产负债表对现金流和价值的影響};

\node[box] (capacity) at (0,-0.4)
{收益、流动性与资本};

\node[box] (adjustment) at (0,-2.1)
{资产负债表反应};

\draw[flow] (shock) -- (bs);
\draw[flow] (bs) -- (capacity);
\draw[flow] (capacity) -- (adjustment);

\draw[flow]
(adjustment.east)
.. controls +(2.0,0) and +(2.0,0) ..
(bs.east);

\end{tikzpicture}

\caption{商业银行内的风险传递与资产负债表调整}
\label{fig:ch12-integrated-risk-interactions}
\end{figure}
资产负债表调整由此可能改变银行随后面临的信用、利率和融资风险暴露。

\subsection[资本应对与资产负债表调整]
{资本应对与资产负债表调整}

信用损失会降低账面权益，而资产快速增长如果权益增长不成正比，则可能会降低权益与资产的比例。银行可以通过保留收益或发行新股权来加强资本状况，从而减少资产基础，或者结合这两种措施来应对这一问题。一种简单、未加权的资本衡量方法就是……

\begin{equation}
k=\frac{E}{A},
\end{equation}
其中，$E$表示账面权益，而$A$则表示总资产。  
在持有资产不变的情况下，权益的增加会提高这一比率；而在保持权益不变的情况下，资产的减少会产生同样的效果。

通过保留收益、减少分配或发行新股份，可以增强权益。资产调整可以包括减缓新的贷款发放，让资产在无需替换的情况下到期，或者出售贷款和证券。这些选择在时机和经济成本上有所不同，同时也会改变未来的收益和资产负债表的构成。

假设一家银行的初始权益为$E_0$，资产为$A_0$，在确认资产损失$L_s$时，负债保持不变，那么……

\begin{equation}
E_1=E_0-L_s,
\label{eq:ch12-capital-after-loss}
\end{equation}
而破产后的资产则……

\begin{equation}
A_1=A_0-L_s.
\end{equation}

如果没有新的权益被筹集，而银行寻求一个目标比率$k^{*}$，那么与这一目标相一致的资产水平是

\begin{equation}
A^{*}=\frac{E_1}{k^{*}}.
\label{eq:ch12-assets-consistent-capital-ratio}
\end{equation}

当 $A_1 > A^{*}$ 时，差值 $A_1 - A^{*}$ 表示在这些简化假设条件下所需的资产减少量。

\begin{example}{恢复资本比率的两条途径}{capital-response-paths}

一家银行的资产为人民币约1000亿元，权益为人民币约80亿元。该公司确认了约20亿元人民币的资产损失，因此目前资产为人民币约980亿元，权益为人民币约60亿元。该公司的权益与资产比率从8\%下降到了约6.12\%。

假设管理层希望恢复8\%的比例。

一种方法是发行新的权益。如果募集资金仍然留在资产负债表上，那么发行约20亿元人民币的权益将使得资产总额达到约1000亿元人民币，而权益则达到约80亿元人民币：

\[
\frac{8}{100}=8\%.
\]

第二种方法是不增加新的权益，而是减少资产。在权益仍为约60亿元人民币的情况下，资产水平与8\%的比例相符合。

\[
\frac{6}{0.08}=75
\]

十亿人民币。

因此，该银行需要将资产减少约230亿元人民币，以弥补亏损后的资产规模。这两种解决方案虽然保持了相同的比率，但资产负债表中的各项指标却有所不同。

\end{example}

因此，资本结构调整不仅会影响资本比率，还会影响到银行的未来收益、流动性、融资需求以及风险暴露情况。

%========================================================
% Summary
%========================================================

\TailModule{summary}{本章小结}

商业银行是一种资产负债表中介机构。它持有贷款、证券、储备金以及其他资产，同时通过存款、批发负债和权益来为这些资产提供资金。贷款的发放可以同时产生贷款资产和存款负债；而后续的支付则可以将存款和储备金在银行之间转移。在债权人的请求权受损之前，权益会先吸收资产损失；而且，杠杆效应会增加权益对资产价值变化的敏感性。

银行信贷决策是随着时间而变化的。借款人的相关信息有助于评估信用状况，而贷款条款和额度则决定了银行的风险规模和结构。在贷款发放后，持续的监控能够确保获得最新的信息；而当借款人情况发生变化时，贷款的续期或合同调整也会对银行的风险产生影响。

银行融资将存款与具有不同到期日、集中程度、重新定价条件和提取条件的批发负债相结合。当相对流动性较高的负债用于资助流动性较低或期限较长的资产时，就会产生流动性创造。因此，流动性管理取决于可用于满足现金流出需求的资源，以及当融资条件发生变化时所使用的资产负债表调整措施。

银行风险通过共同的资产负债表相互关联。贷款恶化可能会减少收益和资本，利率波动可能会影响净利息收入和经济价值，而资金流出则可能会消耗流动性资源或需要新的融资方式。对一种压力的应对可能会改变对另一种压力的风险承受能力。权益可以通过留存收益或新的权益来重建，而资产调整则可能会改变资产负债表的规模和构成，从而影响未来的收益、流动性和风险。

%========================================================
% Mind Map
%========================================================

\TailModule{mindmap}{思维导图}

\begin{figure}[H]
\centering
\begin{tikzpicture}[
  >=Latex,
  center/.style={
    draw=FinanceBlue,
    fill=FinanceBlue!8,
    rounded corners,
    align=center,
    text width=3.2cm,
    minimum height=1.0cm
  },
  box/.style={
    draw,
    rounded corners,
    align=center,
    text width=3.3cm,
    minimum height=0.95cm
  },
  flow/.style={-Latex, line width=0.8pt, FinanceBlue}
]

\node[center] (bank) at (0,0) {商业银行\\资产负债表};

\node[box] (balance) at (0,3.0)
{资产负债表中介、\\资本，以及杠杆率};

\node[box] (credit) at (-5.0,0.9)
{引用评估，\\贷款条款，以及监控措施。};

\node[box] (funding) at (5.0,0.9)
{存款、资金供应以及\\流动性管理};

\node[box] (risk) at (0,-3.0)
{风险相互作用\\资产负债表调整};

\draw[flow] (balance) -- (bank);
\draw[flow] (credit) -- (bank);
\draw[flow] (funding) -- (bank);
\draw[flow] (bank) -- (risk);

\end{tikzpicture}
\caption*{商业银行通过统一的资产负债表来管理信贷决策、资金供应、流动性和资本问题。}
\end{figure}

%========================================================
% Key Terms
%========================================================

\TailModule{keyterms}{关键术语}

商业银行（Commercial bank）；资产负债表中介（balance-sheet intermediation）；存款（deposits）；批发融资（wholesale funding）；权益资本（equity capital）；杠杆（leverage）；信用评估（credit assessment）；信用限额（credit limit）；已融资风险暴露（funded exposure）；未提用承诺（undrawn commitment）；贷款监督（loan monitoring）；贷款续期（loan renewal）；融资期限（funding maturity）；流动性创造（liquidity creation）；流动性能力（liquidity capacity）；融资流出（funding outflows）；净利息收入（net interest income）；重新定价缺口（repricing gap）；资产负债错配（asset--liability mismatch）；久期缺口（duration gap）；权益经济价值（economic value of equity）；资本比率（capital ratio）；资产负债表调整（balance-sheet adjustment）。

%========================================================
% Financial AI Lab
%========================================================

\TailModule{aiprompt}{金融~AI~实验室}

\textbf{实验室：对综合银行资产负债表进行分析与审计}

请让人工智能系统分析一家包含贷款、证券、储备金、存款、批发融资和权益的简化版商业银行。提供有关借款人信用审查、存款外流以及市场利率变化的信息。

将分析过程分为四个阶段进行实施。

\begin{enumerate}

\item \textbf{检查资产负债表恒等式。} 在贷款发放、支付结算、资金流出以及确认的损失发生后，确认资产等于负债加上权益。对于每一变化，确定相应的资产负债表条目。

\item \textbf{请评估信用决策的结果。}  
引入较为薄弱的借款人现金流信息，并判断是否应保持现有的信用额度。在评估还款能力、合同条款以及风险状况时，不应仅仅接受没有依据的违约概率。

\item \textbf{将流动性与资本分离开来。}  
实施存款流出政策，并确定在指定时间内能够提供现金的资源。需要注意的是，分析应区分银行的流动性能力、银行资产的经济价值以及银行的资本状况。

\item \textbf{追踪风险互动。}  
实施利率调整以及贷款损失准备金政策。分析这些调整对净利息收入、经济价值、收益和资本的影响，并确认任何拟议的资产负债表调整在内部是协调一致的。

\end{enumerate}

%========================================================
% Discussion Questions
%========================================================

\TailModule{questions}{讨论题}

\begin{enumerate}

\item 为什么商业银行要充当资产负债表中介的角色，而不是直接将储户现有的存款转移给借款人呢？

\item 如何根据借款人的具体情况来决定银行愿意提供的信贷额度以及相关的合同条款呢？

\item 为什么一个已承诺的信用额度会与银行当前的可用资金额度有所不同呢？

\item 资金集中度、债务期限以及利率敏感性如何影响着银行的流动性管理呢？

\item 如何在同一利率上升的情况下，增加短期的净利息收入，同时降低权益的经济价值？

\item 为什么为了增加流动性而进行的资产出售，仍然会通过已实现的利润或损失来影响企业的收益或资本状况呢？

\end{enumerate}

%========================================================
% Applied Problems
%========================================================

\TailModule{appprobs}{应用题}

\begin{enumerate}

\item \textbf{贷款发起与结算。}

一家银行最初拥有约8000万的人民币贷款、1500万的证券以及500万的储备金。其资金来源包括8200万的人民币存款、800万的批发融资以及1000万的权益。

\begin{enumerate}

\item 银行提供了新的400万人民币贷款，并将资金存入借款人的存款账户。请展示这一操作对资产负债表的直接影响。

\item 借款人将400万人民币的存款转存至另一家银行。如果结算导致储备减少400万人民币，请展示相应的资产负债表。

\item 如果银行希望将储备金恢复到300万人民币的水平，请提出两种能够产生所需现金的资产负债表调整方案。

\end{enumerate}

\item \textbf{信用额度与风险敞口。}

一家公司的流动资金为1200万人民币。初始资金为500万人民币。在本季度内，该公司提取了400万人民币，同时偿还了200万人民币的债务。

\begin{enumerate}
\item 计算最终的融资规模。
\item 计算在信用额度下未使用的金额。
\item 假设运营现金流减弱，而应收账款上升。请解释银行可以在不需要立即偿还全部存续余额的情况下，采取两种方式来调整风险敞口。
\item 解释为什么提高贷款利率并不是应对信用评估结果不佳的唯一办法。
\end{enumerate}

\item \textbf{资金流出与流动性能力。}

一家银行拥有70亿人民币的储备金，以及90亿人民币的证券，这些证券可以在一天内转换为85亿人民币的现金。此外，该银行还有30亿人民币的可动用抵押借款能力，并得到了其他抵押品的支持。然而，该银行出现了140亿人民币的净存款外流。

\begin{enumerate}
\item 根据方程~(\ref{eq:ch12-liquidity-capacity})计算流动性能力。
\item 所宣称的产能是否足以满足输出需求呢？
\item 如果银行使用所有储备金以及通过证券销售获得的约70亿元人民币现金，那么在上述假设下，还剩下多少流动资金呢？
\item 如果证券投资组合中的部分证券已经被用作抵押品，那么答案将会如何改变？
\end{enumerate}


\item \textbf{重新定价与经济价值风险分析。}


在接下来的一年里，这家银行拥有约850亿元人民币的利率敏感资产，以及约700亿元人民币的利率敏感负债。

\begin{enumerate}
\item 计算重新定价差距。
\item 如果使用方程~(\ref{eq:ch12-net-interest-income-gap})来计算，那么当相关利率上升80个基点时，年度净利息收入的变动幅度是多少呢？
\item 同一家银行的资产为人民币1000亿，负债为人民币920亿，资产久期为4.2，负债久期为1.8。计算其久期差距。
\item 估算权益经济价值的一阶变化，以应对收益增加50个基点的情况。
\item 解释为什么收入和经济价值指标在同一时期并不需要保持相同的增长方向。
\end{enumerate}

\item \textbf{资本响应与资产负债表调整。}

一家银行的资产为人民币约1200亿元，权益为人民币约96亿元。该银行确认了人民币约24亿元的亏损。

\begin{enumerate}
\item 计算亏损发生前后的权益与资产比率，假设亏损导致资产和权益减少的幅度相同。
\item 如果银行希望实现8\%的权益与资产比例，并且不增加任何新的权益，那么请计算出符合该比例的资产规模。
\item 计算亏损后的资产负债表所显示的资产减少情况。
\item 提出两种缓解资产负债表收缩的方法，并分别说明每种方法的局限性。
\end{enumerate}

\end{enumerate}

%========================================================
% Data Application
%========================================================

\TailModule{dataapp}{数据应用}

\textbf{中国的商业银行资产负债表、资金状况及资本结构}

参考以下翻译：

使用权威的中国来源，如中国人民银行、国家金融监督管理局，或上市商业银行发布的财务报表。选择一家银行或一个一致的银行业系列数据，收集至少五年的主要资产、存款、其他融资、单独报告的流动资产以及权益方面的数据。

构建能够回答三个问题的比率：银行的资产是如何分配的？资产负债表是如何融资的？有多少权益能够支持资产基础？有用的衡量指标可能包括贷款与总资产的比率、存款与总负债的比率、流动资产与总资产的比率，以及权益与总资产的比率。在整个样本期间使用一致的会计定义，并区分会计权益与监管资本指标。

找出一个时期，在这段时间内，资产增长、存款增长、资金构成或权益发生了显著变化。通过资产负债表的关系来分析这种变化，而不是孤立地解释某个比率。如果贷款的增长速度超过了存款的增长速度，那么需要确定是哪些负债或权益为这一差异提供了资金支持。如果权益的增长速度低于资产的增长速度，那么就需要计算由此导致的简单杠杆变化。

如果两家银行拥有可比的数据，那么请比较它们的资产组合、存款结构、对非存款资金的依赖程度以及流动资产状况。探讨这些差异如何影响其资产负债表在资金流出或利率变化情况下的反应。不要仅依据报告的流动资产来衡量一家银行的流动性能力。

请列出该机构的信息、样本周期、会计定义、相关文件，以及分析中使用的所有处理步骤。

%========================================================
% Case Analysis
%========================================================

\TailModule{caseanalysis}{案例分析}

\textbf{案例：硅谷银行——个体银行资产负债表分析}

2023年3月10日，加利福尼亚州金融保护与创新部门关闭了Silicon Valley Bank，并由联邦存款保险公司担任接管人。联邦储备系统随后的审查发现，该银行存在资产负债表快速增长、业务模式集中、严重依赖未投保的存款，以及流动性和利率风险管理方面的不足。

随着2022年市场利率的上升，SVB持有的长期固定利率证券的市场价值有所下降。与此同时，风险投资活动的减弱导致存款流入减少，而一些由科技或风险投资支持的客户则提取了存款以用于融资。2023年3月8日，SVB金融集团宣布进行资产负债表重组，其中包括出售一批可供出售的证券投资组合、确认税后亏损，以及计划进行资本筹集。

从银行资产负债表的角度来分析这一情况。需要解决四个问题：证券投资组合如何应对较高的市场利率？存款结构如何影响资金的稳定性？哪些资产或资金来源能够在所需的时间范围内提供流动性？实际发生的损失对银行的资本状况产生了什么影响？而计划中的资本筹集措施又该如何应对？

接下来，我们来探讨一下这些维度之间的相互作用。例如，为了增加流动性而进行的资产出售，可能会将市场价值的未实现下跌转化为实际损失，从而降低收益和资本水平。因此，这一事件表明，利率风险、资金流出、流动性调整以及资本等因素是如何通过相同的资产负债表相互影响的。

这里的分析仍然停留在单个银行的层面。关于系统范围内的传染效应、紧急流动性支持以及审慎监管方面的内容，将在后续章节中进行探讨。

\emph{主要来源：联邦储备系统理事会，  
《对硅谷银行的监管与监管评估》，2023年4月；存款保险局，  
2023年硅谷银行倒闭信息，  
SVB金融集团，SEC文件，2023年3月。}

%========================================================
% Further Reading
%========================================================

\TailModule{furtherread}{延伸阅读}

\begin{itemize}

\item Berger, Allen N., and Gregory F. Udell (1995),
``Relationship Lending and Lines of Credit in Small Firm Finance（小企业融资中的关系型贷款与信贷额度）,''
\emph{The Journal of Business}, 68(3), 351--381.

\item Kashyap, Anil K., Raghuram Rajan, and Jeremy C. Stein (2002),
``Banks as Liquidity Providers: An Explanation for the Coexistence of Lending and Deposit-Taking（作为流动性提供者的银行：对放贷与吸收存款并存的解释）,''
\emph{The Journal of Finance}, 57(1), 33--73.

\item Drechsler, Itamar, Alexi Savov, and Philipp Schnabl (2021),
``Banking on Deposits: Maturity Transformation without Interest Rate Risk（依托存款的银行经营：无利率风险的期限转换）,''
\emph{The Journal of Finance}, 76(3), 1091--1143.

\end{itemize}

\ChapterTail{summary,mindmap,keyterms,aiprompt,questions,appprobs,dataapp,caseanalysis,furtherread}

\chapter{投资银行：资本筹集与资本市场中介}
\label{chap:Investment_Banks_Capital_Raising_and_Capital_Market_Intermediation}

\begin{openingcase}{直接融资为什么需要投资银行？}

这家发展迅速的公司计划通过向公众投资者出售新股来筹集约60亿元人民币的资金。该公司的财务报表已经过审计，相关信息可以在招股书中找到。投资者最终将持有该公司发行的股票。

然而，该公司并不会独自与投资者进行接触。它会聘请一家投资银行，并支付相应的费用以获得其服务。这家投资银行会深度参与项目的筹备、定价以及发行过程，并且可能会承担与本次发行相关的一部分风险。

\end{openingcase}

\begin{tcolorbox}[
  enhanced,
  colback=FinanceBlue!6,
  colframe=FinanceBlue,
  boxrule=0.6pt,
  sharp corners,
  left=8pt,
  right=8pt,
  top=7pt,
  bottom=7pt,
  before skip=1em,
  after skip=1.6em,
  title=\textbf{核心问题},
  coltitle=white,
  colbacktitle=FinanceBlue,
  fonttitle=\bfseries,
  attach boxed title to top left={xshift=6pt,yshift=-2mm},
  boxed title style={
    boxrule=0pt,
    sharp corners,
    left=5pt,
    right=5pt,
    top=2pt,
    bottom=2pt
  }
]
\begin{enumerate}[
  leftmargin=*,
  label=\textcolor{FinanceBlue}{\textbf{\arabic*.}},
  itemsep=0.55em,
  topsep=0.2em
]

\item 为什么在证券已经建立了发行人与投资者之间的直接请求权的情况下，发行人仍然使用投资银行呢？

\item 投资银行如何利用信息、声誉、估值、分销网络以及资产负债表能力来支持证券的发行、交易和战略交易呢？

\end{enumerate}
\end{tcolorbox}

\section{资本市场中介与认证}

投资银行帮助发行人获取资本市场资金，而不会取代投资者持有的证券。他们的角色与商业银行的资产负债表中介业务不同，后者通常通过存款和其他银行负债来为借款人的请求权提供资金支持。

\subsection{发行人、投资者与资本市场准入}

发行方进入资本市场时有着不同的融资需求。成长中的公司可以通过发行权益来为投资提供资金，而无需承担固定的偿还义务；而公司则可以通过发行债券来获得短期融资、重新融资现有债务或分散其资金来源。政府和金融机构也根据不同的融资目标和合同结构来发行证券。

进入资本市场需要将融资需求转化为投资者能够评估和购买的证券。发行人必须提供相关信息，明确请求权的条款，吸引潜在投资者，并确定能够吸引足够投资者需求的条款。法律、信息披露和交易成本也决定了发行的规模和设计在经济上的重要性。

投资银行帮助连接市场的这两端。他们就发行的结构和时机提供建议，组织筹备和分配过程，并收集有关投资者需求的信息。这种协调可以减少搜索和交易成本，同时提高发行人与投资者之间的信息流通。图~\ref{fig:ch13-direct-finance-intermediation}总结了这些关系。

\begin{figure}[H]

\centering

\begin{tikzpicture}[
  >=Latex,
  box/.style={
    draw,
    rounded corners,
    align=center,
    minimum width=1.8cm,
    minimum height=0.9cm,
    inner sep=5pt
  },
  service/.style={Latex-Latex, line width=0.85pt, FinanceBlue},
  claim/.style={Latex-Latex, dashed, line width=0.75pt, draw=gray!75}
]

\node[box] (issuer) at (-4.6,0) {发行方};
\node[box] (ib) at (0.5,0) {投资银行};
\node[box] (investors) at (5.2,0) {投资者};

\draw[service] (issuer) --
  node[midway, above=2pt, font=\small,
  text=FinanceBlue, align=center, text width=2.8cm]
  {咨询与交易事项\\协调处理。}
  (ib);

% Investment bank to investors: distribution
\draw[->, line width=0.85pt, FinanceBlue]
  ([yshift=0.14cm]ib.east) --
  node[midway, above=3pt, font=\small,
       text=FinanceBlue, align=center]
  {证券\\分配}
  ([yshift=0.14cm]investors.west);

% Investors to investment bank: demand information
\draw[->, dashed, line width=0.75pt, FinanceBlue]
  ([yshift=-0.14cm]investors.west) --
  node[midway, below=1pt, font=\small,
       align=center]
  { 保留LaTeX命令、公式、交叉引用和数字。\\信息}
  ([yshift=-0.14cm]ib.east);

% Funds: Issuer -> Investors
\coordinate (funds-start) at
  ($(investors.south west)!0.25!(investors.south east)$);

\coordinate (funds-end) at
  ($(issuer.south west)!0.75!(issuer.south east)$);

\coordinate (funds-level) at
  ($(funds-start)!0.5!(funds-end)+(0,-0.75cm)$);

\draw[->, dashed, line width=0.75pt, FinanceBlue]
  (funds-start)
  -- ++(0,-0.75cm)
  -| (funds-end);

\node[below=2pt, font=\small] at (funds-level) {资金};

% Funds: Investors -> Issuer
\coordinate (funds-start) at
  ($(investors.south west)!0.75!(investors.south east)$);

\coordinate (funds-end) at
  ($(issuer.south west)!0.25!(issuer.south east)$);

\coordinate (funds-level) at
  ($(funds-start)!0.5!(funds-end)+(0,-1.5cm)$);

\draw[->, dashed, line width=0.75pt, FinanceBlue]
  (funds-end)
  -- ++(0,-1.5cm)
  -| (funds-start);

\node[below=2pt, font=\small] at (funds-level) {证券};
  
\end{tikzpicture}

\caption{投资银行对资本市场准入的支持}

\label{fig:ch13-direct-finance-intermediation}

\end{figure}


\subsection[尽职调查、信息认证与声誉]{尽职调查、信息认证与声誉}

证券的发行依赖于对发行人、证券以及交易风险的可靠披露。尽职调查通过审查可能影响披露、估值或执行的信息来支持这一过程。尽职调查可能包括对财务报表、商业计划、重要合同、资本结构、所有权、诉讼以及特定交易风险等方面的审查。法律顾问、审计师以及其他专家通常会与投资银行一起参与这一过程。

声誉会影响这项调查背后的激励因素。发行人可能只是偶尔进入资本市场，而投资银行则通常会参与多次交易。因此，如果筛选机制薄弱或信息披露不可靠，就会降低投资者的信心，从而使投资银行在未来的交易中面临更大的困难或更高的成本。

这种声誉风险使投资银行的参与可能发挥认证作用。投资者若无法以合理成本重复同样的调查，就可能认为其参与具有信息价值，因为中介未来业务部分取决于维持市场信誉。


\begin{tcolorbox}[
  enhanced,
  breakable,
  colback=white,
  colframe=FinanceBlue,
  boxrule=0.7pt,
  sharp corners,
  borderline west={2.2pt}{0pt}{FinanceBlue},
  title=\textbf{财务展望：认证并不能保证成功},
  coltitle=white,
  colbacktitle=FinanceBlue,
  fonttitle=\bfseries,
  before skip=1.2em,
  after skip=1.2em
]

认证能够缓解信息摩擦，却不能消除投资风险。投资银行可以调查发行人、审查披露，并通过参与发行使自身声誉承担风险。这些行为可能使其参与本身向投资者传递信息。

但投资银行并不保证发行人未来现金流或证券市场价值。认证影响信息的可信度，无法消除基础请求权的不确定性。

\end{tcolorbox}

\subsection{证券设计与发行人—投资者匹配}

证券设计将发行人的融资需求转化为投资者可以评估和持有的请求权。相关的选择可能包括债务与权益、期限、清偿顺位、支付结构、抵押品、契约限制条款、货币种类以及发行规模等。

这些条款在发行人和投资者之间分配风险和控制权。较长的到期时间可以减少近期再融资的压力，同时增加投资者对利率或信用变化的承受力。抵押品可以加强债权人的保护，同时限制发行人对质押资产的使用。权益则避免了固定的债务偿还，但会改变所有权和请求权。因此，证券的设计涉及权衡取舍，而不是单一的首选结构。

合适的设计还取决于潜在的投资者群体。机构投资者的投资期限、流动性需求、货币风险敞口、监管限制以及允许的资产类别各不相同。即使某项证券在财务上表现良好，但如果其条款不符合潜在买家的要求，那么这项证券可能很难被出售。

投资银行在促成这笔交易的过程中发挥着重要作用。他们利用与投资者的关系、可比的证券、当前市场状况以及预售反馈来评估需求。随着这些信息的不断积累，发行方可以调整债券的规模、期限、定价条款或其他相关因素。因此，证券设计和投资者匹配的过程是同步进行的，直到交易达成能够满足发行方融资需求的同时又能吸引足够多投资者的需求为止。

\section{承销与证券分销}

承销将计划中的证券发行转化为实际的市场交易行为。投资银行负责组织发行过程，与投资者协商证券的定价，并可能承担部分分配风险。

\subsection{承销承诺与承销团安排}

承销协议决定发行人与投资银行如何分担分销风险。在包销安排下，承销商从发行人或出售股份的股东处购买证券，再向投资者出售，因此承担无法按拟定价格完成配售的风险。在尽力销售的代销安排下，投资银行承诺推销证券，但不作相同的购买承诺。

对于具有购买承诺的承销发行，设~$P_o$~为公开发行价格，$P_u$~为承销商向发行人或出售股份的股东支付的价格。每单位证券的承销总价差为

\begin{equation}
g=P_o-P_u,
\label{eq:ch13-underwriting-spread}
\end{equation}
对于$N$种证券而言，按照规定的发行条款，总的毛差价为

\begin{equation}
G=N(P_o-P_u).
\label{eq:ch13-total-underwriting-spread}
\end{equation}

价差补偿承销、分销及相关交易服务，并不等于利润。承销商承担经营和融资成本，投资者需求减弱还可能降低尚未配售证券的实现价值。

大规模发行常组建承销团，结合不同机构的分销能力，并分担承销风险暴露。牵头行协调交易，其他成员承担约定的承销与分销职责。承销团可扩大投资者覆盖，但也需要协调职责、收益分配、客户覆盖与配售。

\begin{example}{一笔包销发行}{firm-commitment-underwriting}

一家公司通过承销方式发行了2000万股新股。公开发行的价格为每股约25元人民币，而承销商则从发行人处以每股约24.25元人民币的价格购得了这些股票。

总差价为

\[
25-24.25=\text{RMB }0.75
\]

每股的价格，而总毛利润差距为

\[
20{,}000{,}000\times0.75=\text{RMB }15\text{ 百万}.
\]

在支付其他费用之前，发行人会先获得收益。

\[
20{,}000{,}000\times24.25=\text{RMB }485\text{ 百万}.
\]

如果在做出承销承诺后，投资者需求减弱，那么部分股票可能需要以低于发行价格的价格出售，或者暂时持有。因此，合同中的价差并不能消除承销商在分配条件方面的风险。

\end{example}


\subsection{首次公开发行与发行定价}

首次公开发行（IPO）是指公司将股票首次向公众投资者出售。投资银行负责协助组织此次发行，协调信息披露和营销工作，评估投资者的需求，并支持确定发行价格和分配股票。

在IPO中，可能会包含普通股、优先股，或者两者都有。普通股是由公司新发行的股份，为发行人提供新的资金。而优先股则由现有股东出售，从而将所有权转移给购买者，同时不会为公司筹集新的资金。因此，发行股票的构成会影响到公司的融资效果以及销售所得资金的分配情况。

在IPO定价时存在困难，因为该股票的公开交易价格尚未确定。估值过程可能会基于财务预测、同类公司的状况、以往的交易情况以及贴现现金流分析来进行。同时，投资者的反馈也能提供有关需求的额外信息。这些输入数据通常用于确定一个估值范围，而不是一个单一的可观测数值。

在存在这种不确定性的情况下，发行价格需要慎重制定。较高的价格虽然能增加每股的收益，但可能会降低需求；而较低的价格则有利于股份的分销，但如果股份后的价格上升，那么初始投资者可能会获得更大的收益。因此，定价决策必须综合考虑发行公司的资金状况以及成功发行所需的条件。

IPO~抑价通常指从发行价到首日市场价格的正初始回报，常用首日收盘价计量。其幅度因企业、时期和市场而异，可能反映信息摩擦、不确定性、投资者参与激励或发行流程特征。

\begin{tcolorbox}[
  enhanced,
  breakable,
  colback=white,
  colframe=FinanceBlue,
  boxrule=0.7pt,
  sharp corners,
  borderline west={2.2pt}{0pt}{FinanceBlue},
  title=\textbf{财务展望：盘后定价不足是一种结果现象},
  coltitle=white,
  colbacktitle=FinanceBlue,
  fonttitle=\bfseries,
  before skip=1.2em,
  after skip=1.2em
]

在二级市场交易开始之后，才会出现首日上涨的情况。相比之下，发行价格是在对投资者需求以及股票交易价格不确定的情况下预先确定的。

因此，抑价本身并不表明发行人或承销商能够事先知道后来的市场价格。这是一项事后结果，可能同时反映定价决策，以及交易开始时得到解决的不确定性。

\end{tcolorbox}

\subsection{债券发行与定价}

承销工作主要围绕确定债券向投资者发行的条件展开。市场收益率、发行人的信用状况、到期日、清偿顺位、流动性以及合同条款都影响着投资者所需的回报水平。

因此，定价始于对市场的参考点。政府的票息曲线、二级市场的信用利差、近期发行情况以及可比的借款人提供了所需票息的基准。承销商和发行人利用这些基准数据，结合投资者需求的迹象，来确定合适的票息和发行价格。

对于固定利率的债券，票面利率不必等于市场要求的收益率。当票息与投资者要求的收益率不同时，债券可以以高于或低于面值的价格发行。实际上，票息和发行价格通常会一起确定，这样债券就可以与类似的证券进行直接比较，并且可以以合理的条件进行分配。

\begin{example}{从市场基准到最终发行条款}{bond-issue-pricing}

有一家公司计划发行为期五年的债券。该政府债券的收益率为2.8\%，而同类企业债券的利率则比政府债券高出约120个基点。

与投资者的初步讨论表明，对债券的需求在110个基点范围内较为有限，而在125个基点附近则存在较强的需求。考虑到这些迹象，发行人和承销商最终决定将发行利率定为125个基点，这意味着债券的发行收益率大约为

\[
2.8\%+1.25\%=4.05\%.
\]

票息和发行价格可以围绕这个要求的收益率来设定。示例展示了市场基准和投资者反馈如何共同影响新债券发行的最终条款。

\end{example}

投资者的需求可能因期限、信用质量、货币、流动性以及其他特征而有所不同。因此，来自银行、保险公司、养老基金、资产管理公司以及其他投资者的反馈可能会影响债券的最终期限、发行规模、定价以及资金分配。

债务发行的条件可能会继续调整，直到发行方的融资目标和投资者的必要报酬率足够接近，使得该交易能够顺利进行。

\begin{figure}[H]

\centering

\begin{tikzpicture}[
  >=Latex,
  box/.style={
    draw,
    rounded corners,
    align=center,
    minimum width=3.1cm,
    minimum height=0.9cm,
    inner sep=5pt
  },
  flow/.style={-Latex, line width=0.85pt, FinanceBlue}
]

\node[box] (bench) at (-4.6,1.25)
{市场基准数据\\和信贷条件};

\node[box] (demand) at (-4.6,-0.75)
{投资者需求\\和偏好};

\node[box] (yield) at (0,0.25)
{所需收益率};

\node[box] (terms) at (4.0,0.25)
{票息\\发行价格};

% Two inputs converge into a common vertical line
\coordinate (mergeTop) at (-2.0,1.25);
\coordinate (mergeBottom) at (-2.0,-0.75);
\coordinate (mergeCenter) at (-2.0,0.25);

\draw[line width=0.85pt, FinanceBlue] (bench.east) -- (mergeTop);
\draw[line width=0.85pt, FinanceBlue] (demand.east) -- (mergeBottom);

% Shared connection and output
\draw[line width=0.85pt, FinanceBlue]
  (mergeTop) -- (mergeBottom);

\draw[flow] (mergeCenter) -- (yield.west);
\draw[flow] (yield.east) -- (terms.west);
\end{tikzpicture}

\caption{新发行债券的定价}

\label{fig:ch13-bond-issue-pricing}

\end{figure}

\subsection{簿记建档、定价与投资者配售}

簿记建档是承销商在最终定价和配售前收集投资者需求意向的过程。投资者可表明不同价格或收益率下愿意购买的数量，使承销商评估需求的强度及分布。

让 $q_i(p)$ 表示投资者 $i$ 在价格 $p$ 下的指示需求。总体指示需求为……

\begin{equation}
Q(p)=\sum_{i=1}^{n} q_i(p).
\label{eq:ch13-bookbuilding-demand}
\end{equation}

订单簿提供的信息比单纯的总需求更为丰富。承销商可以观察到订单的规模和价格敏感性、参与投资者的构成，以及随着拟议条款的修改而变化的需求情况。这些信息有助于评估发行的价格以及该需求的不确定性程度。

因此，最终定价是依据订单簿来进行的，而不是遵循纯粹的机械规则。对发行人更为有利的定价可能会提高每证券的收益，但在交易完成之前，需求的变化空间就会变小。在确定最终条款时，发行人和承销商还可能考虑执行风险、投资者参与度以及预期的二级市场状况。

当意向需求与发行数量不同时，配售决定证券如何分配。有些发行主要按比例配售，另一些则允许更大裁量。配售可以考虑订单规模、价格敏感性、投资者特征，以及簿记建档中传递的信息。裁量可能改善发行配售，但若受与交易无关的关系或目标影响，也会产生利益冲突。

\begin{example}{债券簿记建档中的价格敏感性}{bond-bookbuilding-sensitivity}

一家公司计划发行约10亿元人民币的五年期债券。在簿记建档过程中，投资者提交了总计约16亿元人民币的订单，预期的收益率为3.0\%。不过，其中约5亿元人民币的订单是以至少3.2\%的收益率为条件的。

如果最终定价定为3.0\%，那么这些对价格敏感的订单可能就不会被执行，这样一来，与最终定价条款相一致的订单就只剩下约11亿元人民币的价值了。不过，随着市场情况的变化，这些订单也可能在分配之前被重新调整。

因此，承销商必须评估提交订单的总量，以及它们的价格敏感性和可靠性。设定较高的收益率可能会吸引更多的需求，但这也会增加发行人的借款成本。因此，簿记建档涉及到平衡投资者的需求与发行人面临的融资成本之间的关系。
\end{example}

\section[做市与二级市场中介]{做市与二级市场中介}

在二级市场上，交易商通过报价、执行客户订单以及当买卖订单不一致时临时持有证券来支持交易活动。这些交易活动利用资产负债表的能力来吸收短期的库存不平衡问题，并促进可交易证券的交易。

\subsection{交易商库存、资产负债表与流动性供给}

交易商可以在找到其他愿意购买该证券的交易商之前，从客户那里购买该证券。这笔交易会进入交易商的库存中，这样客户就可以进行交易，而无需等待相反的订单。

令 $I_t$ 表示周期 $t$ 结束时交易商的库存，$Q_t^{B}$ 表示从客户处采购的商品数量，$Q_t^{S}$ 则表示向客户销售的商品数量。一个简单的库存关系可以表示为：

\begin{equation}
I_t=I_{t-1}+Q_t^{B}-Q_t^{S}.
\label{eq:ch13-dealer-inventory}
\end{equation}

持有库存会使交易商面临价格波动的风险，同时还需要资金支持和资产负债表能力。如果交易商从客户那里购买债券，而在债券出售之前市场收益率上升，那么库存的价值就会下降。因此，资本需求、融资成本、套期保值成本以及内部风险限制等因素都会影响到交易商接受客户订单的意愿。

\begin{example}{客户交易中的库存风险}{dealer-inventory-risk}

该交易商购买了一份市场价值达5000万元人民币的债券头寸，并计划当天晚些时候将其转售。如果在出售该头寸之前，市场价值下降0.4\%，那么其实际损失将约为……

\[
50{,}000{,}000\times 0.004=\text{RMB }200{,}000.
\]

\end{example}

因此，交易商的流动性供应会随着市场状况的变化而不同。当波动性上升时，融资成本会增加，或者风险限制会变得更加严格，交易商可能会持有较少的库存或给出的报价不那么积极。即使证券的合同条款没有发生变化，流动性也可能减弱。

承担库存的作用，使交易商中介不同于简单订单撮合。交易平台可以不持有头寸而连接买卖双方，交易商则通过暂时使用自身资产负债表提供即时成交服务。

\subsection[买卖价差、信息与交易成本]
{买卖价差、信息与交易成本}

交易商通常会报出一个买入价，表示它愿意购买；同时也会报出一个卖出价，表示它愿意出售。所报的价差就是……

\begin{equation}
s=P_{\mathrm{ask}}-P_{\mathrm{bid}}.
\label{eq:ch13-dealer-spread}
\end{equation}

这种价差反映了多种交易成本和风险来源。订单处理和市场准入涉及直接成本，而库存头寸则使交易商面临价格和资金风险。逆向选择增加了信息因素，因为当交易商与那些对证券价值有更好了解的投资人进行交易时，就会引入更多的信息因素。

这些证券的大小因证券和交易的不同而有所差异。对于流动性较高的证券来说，小额交易可能不需要大量的库存投入；而对于流动性较低的证券而言，大额交易则更难以进行对冲处理，从而让交易商面临更大的价格风险。

\begin{tcolorbox}[
  enhanced,
  breakable,
  colback=white,
  colframe=FinanceBlue,
  boxrule=0.7pt,
  sharp corners,
  borderline west={2.2pt}{0pt}{FinanceBlue},
  title=\textbf{财务展望：所报价差并非全部交易成本的体现},
  coltitle=white,
  colbacktitle=FinanceBlue,
  fonttitle=\bfseries,
  before skip=1.2em,
  after skip=1.2em
]

报价买卖价差描述给定数量下可得的价格。如果最优报价提供的数量有限，较大交易可能需要按较不利的价格执行。

因此，交易成本既可能包括报价价差，也可能包括执行所需数量时的额外价格让步。交易规模相对于市场可得流动性较大时，这一区别更加重要。

\end{tcolorbox}

信息不对称还会进一步影响交易商的报价。如果接到的订单中包含更多关于未来价值的信息，那么针对这些订单进行交易可能会给交易商带来损失。更宽的价差、更小的报价数量或者不那么激进的定价可以降低这种风险，不过这也会增加其他投资者即时执行交易的成本。

因此，买卖价差反映了在不确定性情况下提供即时交易服务的成本。这些价差往往取决于证券的流动性、交易的规模、交易商的风险敞口，以及交易发生的信息环境。

\section{顾问服务与战略交易}

投资银行还为企业提供关于那些涉及所有权、控制权或资本结构变化的交易的建议。他们的职责是帮助评估各种选择方案，协商交易条款，在需要时安排融资，并协调那些其价值依赖于企业特定信息和合同设计的交易的执行过程。

\subsection{并购顾问}

并购活动涉及对企业资产和现金流的控制权的转移或重组。投资银行可以通过评估交易方案、估算价值、识别交易对手方，以及支持谈判和执行过程，为收购方、出售方、目标企业或专门委员会提供咨询建议。

交易决策不仅仅涉及估值问题。买家可能会期望通过合并两家公司来获得运营或财务上的好处，但能否实现这些好处仍然存在不确定性。卖方在考虑价格的同时，还需要考虑支付方式、融资条件、监管批准情况以及交易能够成功完成的概率等因素。因此，估值为谈判提供了依据，而并非是决定交易价格的唯一依据。

一个简单的框架将独立价值与从组合中获得的增量价值区分开来。令 $V_{acquirer}$ 和 $V_{target}$ 分别表示收购方和目标公司的独立价值，$S$ 表示将它们组合后预期的总收益，$C$ 则表示交易和整合成本的现值。组合后的价值就是……

\begin{equation}
V_{merger}=V_{acquirer}+V_{target}+S-C.
\label{eq:ch13-combined-transaction-value}
\end{equation}

如果收购方以总对价$P_{target}$收购目标企业，那么在这一简化框架下，收购方现有股东预期的净现值将为……

\begin{equation}
\textit{NPV}_{acquirer}=S-C-(P_{target}-V_{target}).
\label{eq:ch13-acquirer-transaction-npv}
\end{equation}

术语“$P_{target}-V_{target}$”指的是相对于目标资产预估独立价值的收购溢价。此类交易可能会带来正的合并收益，而这些收益中的大部分将通过协商确定的溢价分配给目标股东。

\begin{example}{协同效应与协商溢价}{ma-synergy-premium}

该目标项目的预计独立价值约为10亿元人民币。买家期望获得约2.5亿元人民币的总收益，而整合成本现值约为5000万元人民币。因此，预期净增量价值为2亿元人民币。

如果约定的购买价格为11.5亿元人民币，那么相对于目标资产预估独立价值的溢价约为1.5亿元人民币。基于这些假设，买方预期的交易净现值如下：

\[
200-150=\text{RMB }50\text{ 百万}.
\]

\end{example}

顾问帮助客户评估这些权衡因素，并将估值结果转化为具体的交易条款。咨询服务可以提高决策的质量和执行效率，但它并不能确保该交易最终能够创造价值。

\subsection{重组、资本重整与困境交易}

企业可能会出于战略原因或应对财务压力而调整其资本结构。重组会改变金融请求权的组合、期限或所有权结构，而不一定改变企业的运营资产。当支付义务变得难以维持时，重组可能需要对债务条款、优先级、回收权或控制权进行更实质性的调整。

投资银行可以为公司、债权人或投资者提供关于这些替代方案的评估建议。在重组过程中，分析可能会关注杠杆率、债务结构、融资灵活性以及权益和控制的分配。在财务困境的情况下，更多的注意力会集中在企业价值、预期回收情况、短期融资需求以及各索赔群体之间的价值分配上。可能的交易包括债务或权益的发行、债务期限延长、交换提议、新的融资方式以及资产出售。

一种负债管理交易体现了这一机制。面临近期到期义务的公司可能会提供期限较长的证券作为抵押品来偿还现有债务。延长到期时间可以降低即时的再融资压力，但债权人会对比新提供的证券与原有证券的价值、优先级、抵押品以及预期的回收金额。因此，这些条款必须能够为足够多的持有人提供足够的价值，以吸引他们的参与。

回收行为也可以重新分配价值以及控制权。高级债权人、次级债权人和股东在回收方式、风险以及对公司未来控制权的看法上可能有所不同。顾问可以协助进行估值、交易设计以及谈判，而每个债权人的可执行权利则仍然受合同和适用法律的约束。

\subsection{估值、谈判与交易执行}

估值提供了一个比较交易方案的工具，但它通常依赖于假设，而不是得出一个唯一确定的价值。贴现现金流分析依赖于预测和贴现率，而同类公司分析和先例交易分析方法则取决于所选基准指标的相关性。当预期的协同效应、融资能力或控制优势不同时，战略价值也可能因购买方的不同而有所差异。因此，投资银行通常会使用估值范围和敏感性分析来展示估计价值如何随着关键假设的变化而变化。

谈判过程会将这些评估结果转化为具体的交易条款。价格只是其中一个因素而已。支付方式可以是现金、股票、债务，或者多种方式的组合；而绩效补偿条款、条件性支付、退出费、融资条件以及结算条款则会影响风险和价值的分配。因此，虽然这些交易的价格看起来相似，但实际上它们给参与各方带来的经济负担却可能大相径庭。

执行环节决定了是否能在约定的条款下完成所协商的交易。顾问们负责协调融资、文件编制、时间表安排以及信息流通，这些工作涉及法律、会计、监管等多个领域。融资条件的变化、监管审查、股东批准、尽职调查的结果或市场状况都可能影响交易的完成概率和条款。

\begin{tcolorbox}[
  enhanced,
  breakable,
  colback=white,
  colframe=FinanceBlue,
  boxrule=0.7pt,
  sharp corners,
  borderline west={2.2pt}{0pt}{FinanceBlue},
  title=\textbf{财务展望：价值和价格并非等同的概念。},
  coltitle=white,
  colbacktitle=FinanceBlue,
  fonttitle=\bfseries,
  before skip=1.2em,
  after skip=1.2em
]

估值是指在既定假设条件下，对某项交易的经济价值进行估算。协商达成的价格不仅反映了这些估值结果，还考虑了谈判条件、竞争报价、融资条款、合同保护以及执行风险等因素。

因此，定价范围可以为谈判提供指导，但并不能直接确定最终的交易价格。

\end{tcolorbox}

%========================================================
% Summary
%========================================================

\TailModule{summary}{本章小结}

投资银行通过信息生产、认证、证券设计、分销与执行，支持资本市场交易。其作用不同于传统资产负债表中介，因为投资者通常仍持有融资实体直接发行的请求权。

承销有助于将证券从发行人转移给投资者。公司的承诺使承销商面临配售风险，同时承销还涉及到股份分配的能力和风险，而簿记建档则提供了定价和分配的信息。

在二级市场中，交易商利用库存和资产负债表能力来提供即时性。买卖价差反映了交易成本、库存风险、信息风险以及当前的市场状况。

投资银行还为并购、资本重组和重组等业务提供咨询。在这些活动中，它们的经济作用主要体现在信息收集、声誉建设、估值分析、交易专长以及合理运用资产负债表能力等方面。

%========================================================
% Key Terms
%========================================================

\TailModule{keyterms}{关键术语}

投资银行（Investment bank）；资本市场中介（capital-market intermediation）；认证（certification）；尽职调查（due diligence）；证券设计（security design）；承销（underwriting）；代销（best efforts）；承销价差（underwriting spread）；承销团（underwriting syndicate）；首次公开发行（initial public offering）；公司新发行股份（primary shares）；现有股东出售股份（secondary shares）；债券承销（debt underwriting）；簿记建档（bookbuilding）；认购意向（indication of interest）；投资者配售（investor allocation）；做市（market making）；交易商库存（dealer inventory）；买卖价差（bid--ask spread）；战略顾问（strategic advisory）；合并（merger）；收购（acquisition）；收购溢价（acquisition premium）；资本重整（recapitalization）；重组（restructuring）；负债管理（liability management）；交易执行（transaction execution）。

%========================================================
% Financial AI Lab
%========================================================

\TailModule{aiprompt}{金融~AI~实验室}

\textbf{实验室：对人工智能融资建议的审核}

请选择一个最近已筹集到权益或债务资本的上市公司。请一个AI系统担任投资银行顾问，并推荐一种发行结构。该建议应明确融资目标、证券类型、发行规模、可能的投资者群体、相关的估值或收益率基准，以及主要的执行风险。

将建议分为四个阶段进行审核。

\begin{enumerate}

\item \textbf{确认融资需求以及现有的资本结构。}

利用财务报表和正式披露信息来确认公司的债务、权益、流动性状况，以及资金的使用方式。将经过验证的事实与人工智能所提出的假设区分开来。

\item \textbf{评估证券设计。}

评估所提出的证券条款是否符合发行人的融资需求以及潜在的投资者群体。对于债务融资，需要考虑到期日、还款结构、清偿顺位以及相关的抵押品。对于权益融资，则需要考虑发行规模、所有权效应以及发行形式。确定一种替代的融资结构，并解释其中的主要权衡。

\item \textbf{审计定价证据。}

对于权益发行，将拟定的估值与可观察到的市场价格或合适的可比公司数据进行比较。对于债券发行，将拟定的收益率或息差与适当的基准曲线以及可比发行人进行比较。确保日期、货币、到期日以及单位都是一致的。

\item \textbf{与承销承诺分开考虑的建议。}

确定推荐意见中的哪些部分涉及信息、估值和交易建议，而哪些部分可能会让投资银行面临承销风险。解释在 firm commitment 模式下与 placement 或 best-efforts 安排相比，分配风险是如何变化的。

\end{enumerate}

%========================================================
% Discussion Questions
%========================================================

\TailModule{questions}{讨论题}
\begin{enumerate}

\item 为什么在证券发行中使用投资银行并不能直接将融资转化为间接融资呢？

\item 投资银行声誉如何能够减少信息摩擦？为什么认证并不能消除投资者判断的重要性呢？

\item 一家承诺型承销商承担什么样的分销风险？而一般的代理公司通常并不承担这种风险。

\item 为什么在BookBuild交易中，最终确定的报价可能会与某些投资者报价的最高价格有所不同呢？

\item 交易商如何支持流动性供应，以及交易商承担哪些资产负债表风险？

\item 在收购过程中，为什么目标公司的股东能够获得巨大的收益，而收购方现有股东的资产价值却几乎为零呢？

\end{enumerate}

%========================================================
% Applied Problems
%========================================================

\TailModule{appprobs}{应用题}

\begin{enumerate}

\item \textbf{承销经济学。}

一家公司在公开募股中发行了3000万股股票，公开价格约为每股18元人民币。承销团以每股17.55元人民币的价格从发行人手中购得了这些股票。

\begin{enumerate}
\item 计算每股的承销总差额。
\item 计算总毛差幅。
\item 计算发行人在支付其他费用之前所收到的收益。
\item 解释为什么总利差与承销团队的利润并不相同。
\end{enumerate}

\item \textbf{主要和次要份额。}

此次公开募股中，公司发行了8000万股优先股，同时现有股东还出售了2000万股普通股。本次发行的价格约为每股12元人民币。

\begin{enumerate}
\item 计算与优先股相关的总收入。
\item 计算与次级股份相关的总收入。
\item 在扣除费用和开支之后，公司还能获得多少资金呢？
\item 解释这两个组成部分如何不同地影响所有权和融资问题。
\end{enumerate}

\item \textbf{簿记建档与需求。}

发行方计划出售5000万股股票。预计的需求量为：以人民币20元价格交易的股票为9000万股，以人民币21元价格交易的股票为6500万股，以人民币22元价格交易的股票为4200万股。

\begin{enumerate}
\item 在哪些定价条件下，需求会超过供应的数量呢？
\item 为什么仅这些信息还不足以确定一个合适的报价呢？
\item 确定订单簿中承销商除了总体数量之外还可以审查的两种定性特征。
\item 解释更高的发行价格如何同时影响发行公司的收益以及执行风险。
\end{enumerate}

\item \textbf{交易商库存。}

该交易商在一天的开始时持有市值约为1000万人民币的债券。它从客户那里购买了市值约为3500万人民币的债券，并在当天出售了价值2800万人民币的债券。

\begin{enumerate}

\item 计算最终的库存数量，以所述交易价值为准。

\item 如果最终库存在出售前价值下降了0.6\%，那么请计算出相应的市价损失。

\item 请列出另外两项可能影响提供流动性的经济因素。

\item 解释更严格的内部风险限制如何即使债券的合同现金流保持不变，也能影响客户的执行。

\end{enumerate}

\item \textbf{收购价格与谈判事宜。}

该目标项目的预估独立价值为20亿元人民币左右。投标方估计的 총协同效应为4.2亿元人民币，而整合和交易成本现值总计为1.2亿元人民币。

\begin{enumerate}
\item 计算预计的净增量价值。
\item 如果投标者提供的报价为22.4亿元人民币，那么相对于单独项目的价值而言，这一报价所包含的溢价是多少呢？
\item 使用公式~(\ref{eq:ch13-acquirer-transaction-npv})，在给定假设条件下计算投标方预期的交易净现值。
\item 如果整合成本上升到2.2亿元人民币，请重新计算净现值。
\item 解释为什么协商出的价格可能与基于这些假设的估值结果有所不同。
\end{enumerate}

\item \textbf{债券发行与市场反馈。}

有一家公司计划发行为期五年的债券。该政府债券的收益率为3.0\%。相比之下，同类企业债券的收益率比政府债券高出约110个基点。初步的投资者反馈显示，在100个基点的价差下，市场需求较弱；而在120个基点的价差下，市场需求则较为旺盛。

\begin{enumerate}

\item 计算在110个基点利差下所预期的基准收益率。

\item 计算在120个基点利差下所暗示的发行收益率。

\item 请解释为什么发行方即使接受更高的利差也是可以接受的，尽管这会增加借款成本。

\item 在定价之前，除了基准收益率之外，还有两个因素可能导致最终还款条款发生变化。

\end{enumerate}

\end{enumerate}

%========================================================
% Case Analysis
%========================================================

\TailModule{caseanalysis}{案例分析}

\textbf{公开发行}

阿里巴巴集团控股有限公司于~2014~年~9~月在美国首次公开发行，定价为每份美国存托股（ADS）68~美元。基础发行包括~320,106,100~份~ADS：其中阿里巴巴发行~123,076,931~份新增~ADS，现有股东出售~197,029,169~份~ADS。按发行价格计算，基础发行总所得约为~217.7~亿美元。

六家投资银行担任联合簿记管理人。招股说明书披露，每份~ADS~的承销折扣和佣金为~0.816~美元。阿里巴巴及部分出售股份的股东还授予承销商选择权，最多可额外购买~48,015,900~份~ADS。该选择权随后被全部行使，使总发行所得增至约~250.3~亿美元。

运用该交易考察投资银行的多项功能。首先，区分阿里巴巴出售的新增~ADS~与现有股东出售的存量~ADS。哪部分为公司筹集新资本？哪部分将所有权从原股东转移给新投资者？

其次，计算承销折扣占~68~美元发行价的百分比。这项报酬反映哪些服务和风险？为什么不应将其解释为纯承销利润？

第三，考察联合簿记管理人的作用。承销团如何扩大投资者覆盖、在簿记建档中汇集信息，并分担分销风险暴露？发行人希望提高发行价格，与承销商关注执行风险和投资者参与之间，可能产生什么张力？

最后，解释额外购买~ADS~的选择权，如何在基础发行规模之外提供灵活性并支持分销过程。分析应聚焦承销与分销，而非此后股价表现。

\emph{主要来源：阿里巴巴集团控股有限公司，于2014年9月向美国证券交易委员会提交的最终招股说明书；阿里巴巴集团关于IPO定价以及承销商选择权的公告，2014年9月。}

%========================================================
% Further Reading
%========================================================

\TailModule{furtherread}{延伸阅读}

\begin{itemize}

\item Beatty, Randolph P., and Jay R. Ritter (1986),
``Investment Banking, Reputation, and the Underpricing of Initial Public Offerings（投资银行、声誉与首次公开发行抑价）,''
\emph{Journal of Financial Economics}, 15(1--2), 213--232.

\item Benveniste, Lawrence M., and Paul A. Spindt (1989),
``How Investment Bankers Determine the Offer Price and Allocation of New Issues（投资银行家如何决定新发行证券的发行价格与配售）,''
\emph{Journal of Financial Economics}, 24(2), 343--361.

\item Chemmanur, Thomas J., and Paolo Fulghieri (1994),
``Investment Bank Reputation, Information Production, and Financial Intermediation（投资银行声誉、信息生产与金融中介）,''
\emph{The Journal of Finance}, 49(1), 57--79.

\item Ellis, Katrina, Roni Michaely, and Maureen O'Hara (2000),
``When the Underwriter Is the Market Maker: An Examination of Trading in the IPO Aftermarket（当承销商兼任做市商：对IPO发行后市场交易的考察）,''
\emph{The Journal of Finance}, 55(3), 1039--1074.

\item Servaes, Henri, and Marc Zenner (1996),
``The Role of Investment Banks in Acquisitions（投资银行在收购中的作用）,''
\emph{The Review of Financial Studies}, 9(3), 787--815.

\end{itemize}

\ChapterTail{summary,keyterms,aiprompt,questions,appprobs,caseanalysis,furtherread}

\chapter{资产管理与机构投资}
\label{chap:Asset_Management_and_Institutional_Investment}

\begin{openingcase}{管理人为谁的利益服务？}

一个养老基金代表数千名受益人进行退休储蓄的投资。这些受益者最终承担投资带来的收益和损失，但他们不会单独选择证券或自行交易投资组合。相反，该基金将大部分责任委托给专业的资产管理者来负责。

有一位管理者发现了一种投资机会，这种投资方式需要大幅偏离基准水平，且可能需要数年时间才能获得回报。如果这项投资成功，那么收益主要将归那些享受养老金福利的人所有。然而，如果这项投资在短期内表现不佳，那么养老基金可能会削减或终止对该管理者的委托责任，这样一来，该管理者的费用收入也会相应减少。因此，这位管理者可能会犹豫是否要采取符合受益人长期利益的决策，因为这样的决策会带来短期内的风险。

\end{openingcase}


\begin{tcolorbox}[
  enhanced,
  colback=FinanceBlue!6,
  colframe=FinanceBlue,
  boxrule=0.6pt,
  sharp corners,
  left=8pt,
  right=8pt,
  top=7pt,
  bottom=7pt,
  before skip=1em,
  after skip=1.6em,
  title=\textbf{核心问题},
  coltitle=white,
  colbacktitle=FinanceBlue,
  fonttitle=\bfseries,
  attach boxed title to top left={xshift=6pt,yshift=-2mm},
  boxed title style={
    boxrule=0pt,
    sharp corners,
    left=5pt,
    right=5pt,
    top=2pt,
    bottom=2pt
  }
]
\begin{enumerate}[
  leftmargin=*,
  label=\textcolor{FinanceBlue}{\textbf{\arabic*.}},
  itemsep=0.55em,
  topsep=0.2em
]


\item 为什么投资者会委托进行投资组合决策？而指令要求、基准指标、费用以及信托责任又如何影响着投资者与资产管理者之间的关系呢？

\item 不同的资产管理工具和机构投资者在流动性、控制权、投资期限以及风险分配方面有何差异？

\item 资产管理者如何生成信息？他们如何影响市场以及企业治理？同时，他们如何应对被动投资和技术驱动的投资呢？


\end{enumerate}
\end{tcolorbox}

\section{委托资产管理的经济逻辑}

资产管理将资本所有权与许多投资决策分离开来。投资者承担投资组合的盈亏，而专业管理者则有权在规定的职责范围内选择、交易、监控证券，有时甚至有权投票决定证券的处置方式。

\subsection{直接投资与委托投资}

个人投资者可以直接持有和管理金融资产，或者将投资组合的决策委托给专业的基金经理。直接投资让投资者能够掌控证券的选择、交易以及投资组合的监控，但这需要时间、信息以及投资方面的专业知识。而在委托投资的情况下，投资者仍然对投资组合承担经济责任，而基金经理则有权在规定的授权范围内做出具体的投资决策。具体的法律安排取决于资产是通过集合平台、单独账户还是其他形式来持有的。

当投资组合管理需要专业知识、规模或基础设施时，委托给专业人士来管理可以创造价值。对于个人投资者来说，这些资源往往难以获得。管理者可以将研究、交易系统、风险测量、报告和管理等成本分摊到多个客户身上。通过汇集投资，还可以以比单个投资者更低的交易和监控成本来实现多元化投资。

这一结构不同于吸收存款的中介。商业银行主要通过在自身资产负债表上发行存款和其他负债，融资其资产组合。资产管理人则通常管理客户资产或集合基金，而非主要用自身负债融资这些资产。因此，组合损益通常由投资者承担，但具体仍取决于载体法律结构，以及管理人提供的契约承诺或资本。

一个简单的委托投资关系可以表示为：

\begin{equation}
W_T=W_0(1+r_{gross})-F,
\label{eq:ch14-investor-terminal-wealth}
\end{equation}
其中，$W_0$ 是投资者的初始资本，$r_{gross}$ 是期间投资组合的毛回报率，$F$ 则代表总费用和其他费用，以货币金额表示。与直接投资相比，委托投资改变了做出投资组合决策的人选，并引入了管理成本，而投资者通常仍然需要承担投资组合的业绩风险。

\subsection{资产管理契约、费用与激励}

委托资产管理是由定义决策权、投资限制和补偿条款的合同来管理的。投资指令可能会规定适用的资产、基准指标、风险限额、杠杆率、投资期限以及报告要求。在集合投资工具中，基金条款还可能规定流动性及提款的相关事宜。所有这些规定共同构成了管理层的决策空间以及评估绩效的依据。

基准提供了一个参考投资组合或回报指标，可以通过该指标来评估投资组合的表现。如果$r_P$表示投资组合的回报，$r_B$则表示基准回报，那么主动回报就是……

\begin{equation}
r_{active}=r_P-r_B.
\label{eq:ch14-active-return}
\end{equation}

基准测试能够更清晰地显示相对性能和投资组合的偏差，但这也可能影响投资决策。如果管理者始终以基准作为衡量标准，可能会限制那些可能导致显著相对表现不佳的较大偏差。

补偿通常包括基于资产的管理费用。如果年度管理费率为$f$，而$A$则代表相关的费用基数，那么管理费用可以表示为：

\begin{equation}
F_M=fA.
\label{eq:ch14-management-fee}
\end{equation}

一些安排还包括基于绩效的补偿机制。简化的绩效奖励制度如下：

\begin{equation}
F_P=\alpha\max\{r_P-h,0\}A,
\label{eq:ch14-performance-fee}
\end{equation}
其中，$\alpha$表示绩效费用率，$h$表示门槛值或参考收益水平，$A$则表示费用的基础金额。更详细的合同条款可能会包括最高限额、损失承担条款，或者基准调整相关的内容。

费用设计影响管理人激励。按资产规模收费，使收入与管理资产规模相联系，提高吸引和留住投资者资金的价值。业绩报酬使报酬对投资结果更敏感，但在某些条件下，不对称参与收益与损失也可能鼓励更高风险承担。最终激励取决于整个契约，以及管理人是否与客户共同投入资本。

\begin{example}{管理费与投资者净回报}{management-fee-net-return}

该基金年初拥有1亿元人民币的资金，其投资组合的毛回报率为8\%。假设管理费的支出为初始资产的1\%，且没有其他费用或现金流。

管理费为

\[
F_M=0.01\times100=1
\]

投资组合的收益使基金的资产在扣除费用前增加到1.08亿元人民币。扣除管理费后，净资产为1.07亿元人民币，对应的简单净回报率为7\%。

总回报与净回报之间的每百分点差异，实际上反映了按照既定惯例所计算出的管理费用。

\end{example}

\subsection[委托代理问题与受托责任]{委托代理问题与受托责任}

委托资产管理会形成一种本金与代理人的关系。投资者将投资权限委托给管理者，但管理者的目标并不总是完全符合投资者的利益。补偿、职业考量、资产积累以及商业关系都可能影响管理者行使权力的方式。

机构之间的利益冲突可以有多种表现形式。当补偿机制不对称时，管理者可能会承担额外的风险；为了降低职业发展的风险，他们可能会坚持遵循某个基准标准；或者出于某种考虑，他们可能会进行超出投资合理范围的交易行为。这些行为并非委托授权过程中固有的现象；其重要性取决于合同设计、组织结构以及监管环境等因素。

受托责任约束委托权力的使用。其法律内容因司法辖区和机构形式而异，但基本经济原则是：管理人应充分考虑客户利益，行使受托权力，并适当管理利益冲突。披露、内部控制、独立监督和对利益冲突交易的限制，均可支持这一目标。

投资者还可以通过监控和退出来约束基金经理的行为。机构投资者可以协商投资指令，评估基金经理的表现，并更换基金经理；而散户投资者则更多地依赖标准化的信息披露、基金治理机制以及赎回或转换基金的能力。当未来的投资指令和资产流动依赖于投资者的信心时，声誉因素就成为了一个额外的制约因素。

\begin{tcolorbox}[
  enhanced,
  breakable,
  colback=white,
  colframe=FinanceBlue,
  boxrule=0.7pt,
  sharp corners,
  borderline west={2.2pt}{0pt}{FinanceBlue},
  title=\textbf{财务展望：诚信责任并不能消除代理问题},
  coltitle=white,
  colbacktitle=FinanceBlue,
  fonttitle=\bfseries,
  before skip=1.2em,
  after skip=1.2em
]

受托责任约束管理人如何使用委托权力，却不会使投资者与管理人的激励完全一致。投资结果不确定，管理行为未必完全可观察，利益冲突也可能持续存在。因此，即使存在受托义务，契约、监督、治理和声誉仍然重要。

\end{tcolorbox}

\section{集合投资载体}

一种常见的组织委托投资的方式是通过集体投资机构。这种机构不是单独管理每个投资者的投资组合，而是汇集许多投资者的资金，并将这些资源共同用于特定的目标。

\subsection{开放式共同基金}

开放式共同基金发行投资者可向基金申购或赎回的份额。申购创造新份额，赎回注销现有份额，因此已发行在外份额数量随投资者资金流动而变化。

每股资产净值是多少？

\begin{equation}
\textit{NAV}_t=
\frac{A_t-L_t}{N_t},
\label{eq:ch14-nav}
\end{equation}

其中，$A_t$~为基金资产市值，$L_t$~为基金负债，$N_t$~为已发行在外份额数量。单位净值将基金组合净价值换算为每份价值。

投资者流动可能会影响投资组合的管理。认购资金可以用于投资更多的资产，而赎回则要求基金向离开的投资者提供现金。基金可以满足因持有现金而产生的赎回需求，同时还可以用于弥补认购资金、允许借款或出售资产。当赎回金额相对于可用流动性而言较大时，即使基础投资的状况没有变化，基金经理也可能需要出售资产。

开放式资金池可以提供多样性和交易管理的便利性，但同时也使得投资者的流动性与投资组合的流动性之间产生了联系。因此，满足赎回需求的成本部分取决于基础资产的可销售性。

\begin{example}{按单位净值申购共同基金}{mutual-fund-subscription-nav}

该共同基金的资产为人民币约5亿元，负债为人民币约500万元，存续股份数量为4500万股。其每股净资产值为

\[
\textit{NAV}=\frac{500-5}{45}=11.00.
\]

一位投资者以110万元人民币的出资参与了该项目的投资。扣除相关费用后，投资者实际获得的收益为……

\[
\frac{1.1\text{ 百万}}{11.00}=100{,}000
\]

新发行的基金份额使基金资产增加至5.011亿元人民币，存续份额达到4510万份。按照既定的假设，每股净资产价值仍保持在11元人民币左右。

\end{example}

\subsection[交易所交易基金与创设赎回机制]{交易型基金与创建——\\赎回机制}

交易所交易基金（ETF）持有基础资产组合，其份额在二级市场交易。多数投资者买卖~ETF~份额时，不直接与基金交易。但授权参与人可通过与基金交易，大批量创设或赎回份额。

创设时，授权参与人交付指定的一篮子证券、现金或两者组合，获得新发行的~ETF~份额；赎回时，将~ETF~份额交回基金，换取相应资产篮子或现金结算。这一机制使~ETF~份额供给能够根据其自身及基础组合的市场条件调整。

让 $P_t^{ETF}$ 表示 ETF 股票的市场价格，而 $\textit{NAV}_t$ 则表示其每股的资产净值。一个简单的溢价或折价衡量方法就是……

\begin{equation}
\pi_t=
\frac{P_t^{ETF}-\textit{NAV}_t}{\textit{NAV}_t}.
\label{eq:ch14-etf-premium-discount}
\end{equation}

ETF~以足够大的溢价交易时，授权参与人可创设新份额并在二级市场出售，促使溢价缩小；存在足够大的折价时，赎回可发挥相反作用。这一机制的有效性取决于交易成本、融资条件、税收、交易时间和基础资产流动性。

\begin{example}{ETF溢价交易时的份额创设}{etf-creation-premium}

假设一只ETF的市值为每股10.00人民币，但在二级市场上交易价格为10.10人民币。忽略交易成本和其他摩擦因素，授权参与者可以以相当于每股10.00人民币的价格购买基础资产篮子，将篮子交付给基金，然后获得新创建的ETF股份。

如果这些股票以10.10人民币的价格出售，那么价格差异就会创造出套利机会。这种机会会增加ETF股票的供应，从而有助于降低溢价率。

实际上，只有当价格差异足够大以覆盖交易成本、融资成本、税收以及其他摩擦时，创建才变得有吸引力。

\end{example}

不同于普通开放式共同基金，ETF~不要求每次投资者买卖份额时基金自身都进行交易。因此，二级市场可以吸收大量投资者交易，而不直接改变基础组合。但创设赎回机制仍为交易压力在~ETF~份额与基础资产之间传递提供渠道。

\subsection{基金流动性、单位净值与赎回风险}

请求权投资组合的市场流动性与其底层资产的流动性虽然相关，但又是不同的。即使某些投资组合资产的交易频率较低，基金仍可能提供频繁的申购和赎回操作。当投资者的流动性与可用现金和市场流动性相比较为有限时，这种安排更容易维持。

大量赎回可能要求出售资产。组合资产相对缺乏流动性时，出售可能产生交易成本、价格让步或市场冲击。如果赎回投资者获得的价格没有充分反映这些成本，部分成本可能由留在基金中的投资者承担。在这种条件下，对进一步资金流出的预期可能形成提前赎回激励。

基金的设计方式会影响这些成本的分配。现金缓冲和流动性储备可以减少对立即出售资产的需求，而赎回费、变动定价、通知期限以及其他流动性管理工具则可以将更多成本转移给那些产生这些成本的投资者。这些工具的运用取决于基金的结构、资产的流动性状况以及相关的规章制度。

流动性不足的资产也会使证券的净值计算变得复杂。当市场价格无法获取或获取频率较低时，估值可能会更多地依赖于模型、交易商的报价、可比证券或其他估算方法。虽然证券的净值仍然是衡量投资者请求权的一个有用指标，但当可观测的市场价格提供的信息较少时，估值的不确定性可能会增加。

\begin{figure}[H]
\centering
\begin{tikzpicture}[
  >=Latex,
  box/.style={
    draw=FinanceBlue,
    rounded corners=2pt,
    align=center,
    text width=3.0cm,
    minimum height=0.9cm,
    inner sep=5pt
  },
  flow/.style={-Latex, line width=0.8pt, FinanceBlue}
]

\node[box] (redeem) at (0,0) {投资者赎回资金};
\node[box] (cash) at (4.3,0) {需要现金。};
\node[box] (sales) at (8.6,0) {投资组合资产销售};

\node[box] (costs) at (8.6,-2.0) {交易成本与市场影响};
\node[box] (remain) at (4.3,-2.0) {由其余投资者部分承担的成本};
\node[box] (incentive) at (0,-2.0) {可能的激励措施，以提前实现兑换目标。};

\draw[flow] (redeem) -- (cash);
\draw[flow] (cash) -- (sales);
\draw[flow] (sales) -- (costs);
\draw[flow] (costs) -- (remain);
\draw[flow] (remain) -- (incentive);

\end{tikzpicture}
\caption{开放式基金的赎回压力与流动性成本}
\label{fig:ch14-redemption-liquidity}
\end{figure}

\section{机构投资者与另类资产管理}

除了集体投资工具之外，委托投资还通过各种具有不同负债结构、投资期限、流动性需求以及治理安排的机构来组织实施。这些特征影响着它们能够持有的资产、能够承担的风险，以及它们与金融市场和企业发行人的互动方式。

\subsection{养老基金与负债驱动型投资}

养老基金投资资产的目的是为了履行未来的福利义务。因此，它们的投资策略取决于这些支付的时机、通货膨胀的影响程度以及支付金额的不确定性，同时也需要考虑预期的资产回报情况。

设预期的养老金支付在日期$t=1,\ldots,T$时为$B_t$。负债的现值简化表示为

\begin{equation}
L_0=\sum_{t=1}^{T}\frac{B_t}{(1+y_t)^t},
\label{eq:ch14-pension-liability-pv}
\end{equation}
其中，$y_t$ 表示应用于第 $t$ 天支付的贴现率。衡量资金状况的一个简单方法是……

\begin{equation}
FR=\frac{A_0}{L_0},
\label{eq:ch14-pension-funding-ratio}
\end{equation}
其中，$A_0$ 表示养老金资产的市场价值。

负债驱动型投资（LDI）相对于养老金负债的价值和敏感性管理资产。例如，长期债券的久期若更接近负债久期，可减少利率变化对资金充足状况的影响。基金也可能持有其他资产，以寻求更高预期回报。具体配置取决于资金充足状况、给付设计、缴费政策、监管约束，以及对资金充足程度波动的容忍度。

长期负债还可以使养老基金在短期福利支付和抵押品需求仍然可控的情况下，拥有持有流动性较低的资产更大的能力。虽然较长的投资期限并不能消除流动性风险；但当有足够的流动性来满足短期债务需求时，它可以扩大可行的投资范围。

\begin{tcolorbox}[
  enhanced,
  breakable,
  colback=white,
  colframe=FinanceBlue,
  boxrule=0.7pt,
  sharp corners,
  borderline west={2.2pt}{0pt}{FinanceBlue},
  title=\textbf{财务前景：LDI专注于资金状况，而不仅仅是资产状况。},
  coltitle=white,
  colbacktitle=FinanceBlue,
  fonttitle=\bfseries,
  before skip=1.2em,
  after skip=1.2em
]

负债驱动型投资相对于养老金负债评估资产风险。资产组合单独看可能表现良好，但若负债价值上升更多，资金充足状况仍会恶化。因此，关键在于资产与负债价值如何共同变化。

\end{tcolorbox}

\subsection{主权财富基金、捐赠基金与长期资本}

主权财富基金、捐赠基金、基金会以及类似的机构通常会在较长的时期内进行投资，但它们的资本来源和支出义务却存在显著差异。主权财富基金可能通过商品收入、财政盈余、外汇储备或代际储蓄来融资，而捐赠基金则通常用于支持机构的经常性支出，并力求在一段时间内保持购买力的稳定。

这些差异会影响流动性需求和资产配置。较长的投资期限可以增加持有流动性较低的资产的能力，比如私募股权、基础设施和房地产等，前提是短期现金需求仍然可控。因此，仅仅拥有较长的投资期限是不够的；该机构还必须具备维持长期头寸的流动性和管理能力。



\begin{tcolorbox}[
  enhanced,
  breakable,
  colback=white,
  colframe=FinanceBlue,
  boxrule=0.7pt,
  sharp corners,
  borderline west={2.2pt}{0pt}{FinanceBlue},
  title=\textbf{财务前景：长远的规划并不能消除流动性方面的限制问题},
  coltitle=white,
  colbacktitle=FinanceBlue,
  fonttitle=\bfseries,
  before skip=1.2em,
  after skip=1.2em
]

较长的投资期限可以增加机构持有流动性较低资产的能力，但仅凭期限是不够的。支出需求、可能的提款、抵押品要求以及治理限制仍然可能需要现金支持。因此，长期投资的能力取决于投资期限和流动性能力。

\end{tcolorbox}

支出规则和机构目标进一步影响着投资组合的选择。拥有相对稳定年度收益的捐赠基金可能会更重视长期实际回报，而具有稳定目标的主权基金在不利经济条件下 withdrawals 可能会增加流动性资产的持有量。

长期资本投资还可以降低在市场出现短期压力时需要进行强制出售的需要。拥有足够流动性的机构能够在价格出现短暂下跌时继续持有头寸，而这样做仍然符合它们的职责要求以及风险承受能力。

\subsection{私募股权与风险投资}

私募股权与风险投资通常通过封闭式基金组织投资，有限合伙人承诺投入资本，普通合伙人作出投资决策。投资者一般无法随时赎回权益，使基金有更大能力长期持有非流动性投资。

私人市场投资的一个核心特点是集中所有权和实时监控。收购基金可能会获得实质性的控制权，从而能够影响公司的融资、治理、管理和运营决策。风险资本投资者通常会分阶段为年轻企业提供资金支持，并可能通过参与董事会和管理层来进一步干预公司事务。在信息有限的情况下，这种合作方式尤为重要，因为企业的价值很大程度上取决于其未来的发展情况。

管理人报酬通常由管理费与挂钩投资收益的业绩分成构成，后者常称附带收益。这使普通合伙人直接参与投资结果，但仍可能存在激励冲突。基金条款、分配规则和管理人投入的资本，都会影响普通合伙人与有限合伙人如何分享收益、损失与激励。

当投资组合中的公司并非持续被交易时，估值就变得更加困难。基金的价值可能依赖于模型、可比交易或其他在可观察市场交易中得出的估值方法。因此，报告的估值可能会比公开市场的价格调整得更慢，并且存在更大的估计不确定性。

\subsection{对冲基金与另类投资策略}

对冲基金通常遵循相对灵活的投资策略，可能会使用卖空、衍生品、杠杆、相对价值头寸、事件驱动策略或集中性持仓。这些工具使得基金经理能够表达正面和负面的观点，对冲选定的风险，并构建在传统长期持有投资组合中难以实现的头寸。

这些策略可能会引入资金风险、利润率风险、流动性风险以及模型风险，此外还有基础头寸的市场风险。因此，这些风险不仅取决于所持有的资产，还取决于如何为这些头寸提供融资、进行组合操作以及采取何种抵押措施。

投资者流动性条款通常旨在反映策略的流动性状况。锁定期、限制措施等规定可能会减少因提款而需要进行强制资产出售的情况，但同时也会限制投资者快速退出市场的能力。

权益可以放大资金权益的变化。如果一个基金的资产$A$部分由权益$E$来融资，那么一个简单的资产负债表杠杆指标就是$A/E$。

这一比率并未涵盖所有形式的经济杠杆效应。衍生品、空头头寸以及其他风险敞口可能会产生一些在资产负债表资产中未能完全反映的敏感性问题。因此，风险分析还需要考虑投资组合的风险敞口、抵押品要求以及融资条件。

\section[作为信息与治理中介的资产管理人]{作为信息与治理中介的资产管理人}

一旦投资决策被委托出去，资产管理者就可以影响不仅仅是投资组合的分配。通过研究、交易、投票和参与等方式，他们可以影响市场价格所反映的信息，以及他们所持有证券的公司治理。

\subsection{信息生产、研究与证券选择}

活跃的资产管理团队会收集和使用信息，以确定那些预期回报相对于价格和风险来说具有吸引力的证券或投资组合。研究可能会参考财务报表、行业和宏观经济状况、公司信息、定量模型以及其他允许的数据来源。其价值取决于是否能够在价格完全确定之前获取相关信息并将其纳入投资决策中。

信息的生产与实施成本很高。研究需要分析师、数据和技术支持，而投资观点则必须通过交易转化为投资组合中的具体仓位。随着其他市场参与者获取并利用类似的信息，这种信息优势的价值可能会下降。

如果$r_{gross}$表示在管理费和实施成本之前的投资组合回报，$f$和$c$分别表示相应的成本率，那么简化的净回报就是

\begin{equation}
r_{net}=r_{gross}-f-c.
\label{eq:ch14-net-portfolio-return}
\end{equation}

只有当信息、证券选择以及执行方面的贡献足以抵消这些成本时，主动管理才能为投资者带来价值。

规模可以以不同的方式影响这一过程。研究、数据和技术可能具有规模经济效应，而大型投资组合在资源有限的情况下进行部署时可能会更加困难，或者在进行交易时可能会影响价格。因此，所管理的额外资产的价值取决于其运营市场的策略以及市场的流动性。

\subsection{机构持股与公司治理}

机构投资者可以通过积极的监督以及调整自己的投资布局来影响公司治理结构。这些渠道的有效性取决于企业的持股规模、投票权分配、股权集中度，以及投资管理的相关规定。

管理指的是代表客户或受益人积极监督被投资公司的行为。根据客户的投资需求，资产管理者可以行使与客户投资组合中的股份相关的投票权，包括对董事选举、高管薪酬、合并以及其他股东决议的投票权。他们还可以与董事会和管理层就公司战略、业绩和治理问题进行沟通。通过协调这些活动，资产管理者可以在股东监督方面实现规模经济，并有助于解决与分散式股份所有权相关的集体行动问题。

机构持股也可能引发治理问题。大型资产管理公司在许多公司中拥有巨大的投票影响力，而与被投资公司的业务关系或客户利益之间的冲突则可能引发利益冲突。同时，对竞争企业的共同持股也引起了关注，因为机构持股的重叠可能会影响企业之间的竞争动机。这些问题的严重性仍有争议，不同市场以及不同的所有权结构下，实际情况也有所不同。

\begin{tcolorbox}[
  enhanced,
  breakable,
  colback=white,
  colframe=FinanceBlue,
  boxrule=0.7pt,
  sharp corners,
  borderline west={2.2pt}{0pt}{FinanceBlue},
  title=\textbf{财务视角：被动投资并不排除主动管理},
  coltitle=white,
  colbacktitle=FinanceBlue,
  fonttitle=\bfseries,
  before skip=1.2em,
  after skip=1.2em
]

指数基金遵循基于基准的投资策略，限制证券选择及组合权重的裁量，但并不妨碍管理人行使投票权、与被投资公司沟通或监督公司治理。仅因治理担忧卖出指数成分证券，可能产生跟踪误差，因此投票和沟通可成为被动机构投资者发挥影响的重要渠道。

\end{tcolorbox}

\subsection{机构交易与市场流动性}

机构投资者交易证券是为了实施投资组合决策、重新平衡投资组合、适应投资者流动以及管理风险。他们的订单在其他交易的一方时可以提供流动性，而在寻求即时执行时则可能需要流动性。

大型订单需要在执行速度和市场影响之间做出权衡。较快的执行速度可以降低价格波动带来的风险，但此时订单可能无法完全执行，因此需要更高的价格让步。相反，较慢的执行速度虽然可以减少瞬间的市场影响，但却会增加在执行过程中面临价格变动的风险。

实施不足（IS）衡量的是与投资决策时现行价格相对的执行成本。对于完全执行的购买行为，令$P_0$表示决策价格，$\bar{P}$表示平均执行价格，$Q$表示购买的数量。忽略佣金和其他交易成本。

\begin{equation}
IS=Q(\bar{P}-P_0).
\label{eq:ch14-implementation-shortfall}
\end{equation}

正值表示平均执行价格高于决策价格，这意味着相对于该基准值而言，存在执行成本。而在出售的情况下，符号则相反。

机构交易可能会在以下情况下影响市场流动性：当大额或相关性较高的订单超过当前价格下的可成交深度时。由此产生的价格变动取决于订单规模、执行策略、市场深度以及其他参与者是否愿意采取相反的行动。

\section{资产管理的转型}

新的投资结构和技术正在改变资产管理的运作方式。它们可以降低投资组合构建和执行的成本，同时将一些决策从自由裁量权转移到规则和模型之上。

\subsection{被动投资、指数化与基于规则的配置}

被动投资的目标是复制特定基准指标的业绩表现，而不是通过自主选择证券来创造超额回报。指数基金和被动ETF是常见的实现方式，不过类似的策略也可以用于独立管理的账户和机构投资组合中。

基于规则的配置，用预先设定的标准决定组合风险暴露或资产权重。指数化属于规则投资，但并非所有规则策略都是被动策略。一些策略会根据特定市场条件、投资信号或配置规则，系统地调整权重或风险暴露。

一个简单的例子是一个固定权重的投资组合，其中60\%的资金用于股票，40\%用于债券。随着市场价格的变化，实际权重可能会偏离这些目标。一个预先确定的重新平衡规则要求管理者定期或在偏差超过指定阈值时恢复原来的权重。其他基于规则的策略可以根据预定义的波动性目标、估值指标或动量信号来调整投资组合的持仓比例。这些方法减少了对主观决策的依赖，不过在选择指标、阈值和重新平衡规则方面仍然需要投资方面的判断。

被动策略可以在降低研究和投资组合管理成本的同时，为市场或因素提供透明的投资方向。被动投资组合的回报与其基准回报之间的差距被称为跟踪差异。该差异是根据方程~\eqref{eq:ch14-active-return}中的主动回报定义来衡量的。负的跟踪差异表明相对于基准的业绩表现不佳。这可能是由于管理费用、税收、交易成本、现金持有量、抽样以及重新平衡过程中的摩擦所导致的，而证券借贷收入则可能抵消部分这些影响。

被动投资限制了持续证券选择的自由裁量权，但并未消除投资选择的权利。证券提供商通过关于合格证券、纳入标准以及权重分配方法的决策来确定基准构成。资产管理公司则负责选择如何复制这一基准、执行交易、管理现金持有量以及控制跟踪偏差。因此，投资自由裁量权实际上受到基准方法的限制，并且集中在投资组合的运作上。

\subsection{智能投顾与算法投资组合管理}

Robo-advisory利用数字系统将投资者信息转化为投资组合建议，并自动化部分实施和再平衡过程。关于投资期限、风险承受能力、收入需求以及其他限制条件的信息会被转化为资产配置以及一系列投资工具。

自动化可以在这些任务标准化的情况下，降低构建和维护多样化投资组合的成本。它还可以使推荐意见更好地反映更广泛的投资者特征。最终建议的质量取决于收集到的信息、用于解释这些信息的方法，以及系统中内置的约束条件。

算法投资组合管理将这些方法的应用范围扩展到了零售建议之外。资产管理者可以在研究、投资组合构建和交易过程中使用优化算法、执行算法、机器学习技术以及庞大的数据集。这些工具可以提升信息处理能力，并自动化部分投资管理工作。不过，投资目标、风险限制以及法律责任仍然制约着这些工具的广泛应用。

\subsection{模型风险、自动化与人工监督}

基于模型的资产管理将数据、假设和决策规则转化为投资组合的选择。当模型中包含的信息或关系无法充分反映投资问题时，或者当模型的目标与投资者的实际目标不一致时，就会产生模型风险。

自动化可以提升基于模型的决策的实施速度和规模。但它也可能放大错误或常见的反应。如果许多投资组合对相同的信号做出相同的反应，那么在市场流动性有限的情况下，交易压力可能会集中。这种影响取决于投资组合的规模、市场深度以及所使用的策略的多样性。

人工智能能够扩展信息处理的能力，并自动化研究、预测、分类以及通信等过程的部分工作。这些应用或许能够提高效率，但其可靠性仍取决于数据的质量、模型的设计以及输出结果的验证。

人类监督仍然非常重要，因为投资目标和可接受的风险风险是无法通过自动化系统来确定的。监督工作包括明确投资目标、设定风险界限、审查模型假设和输出结果、识别异常情况，以及决定在何种情况下需要人工干预以纠正自动化的规则设置错误。

\begin{table}[H]

\centering

\caption{资产管理自动化：潜在收益与仍存约束}

\label{tab:ch14-automation-asset-management}

\begin{tabularx}{0.94\textwidth}{@{}p{0.30\textwidth} X X@{}}

\toprule

\textbf{函数} & \textbf{自动化带来的潜在收益} & \textbf{剩余的约束条件} \\

\midrule

客户端评估  
& 处理成本更低，一致性更高  
& 客户端信息不完整或不准确  
\\

研究
& 大规模数据集的处理效率提升
& 数据质量与模型可靠性 \\

投资组合构建
& 分配和再平衡规则的一致应用
& 目标明确与估算误差 \\

交易方式  
& 更快且更系统的执行流程  
& 市场影响、流动性以及不断变化的市场状况  
\\

\bottomrule

\end{tabularx}

\end{table}

%========================================================
% Summary
%========================================================

\TailModule{summary}{本章小结}

资产管理将投资组合决策从投资者手中委托给专业管理者。这种委托可以发挥规模经济和信息生产、交易、分散投资和管理的专业优势，但同时也带来了代理问题。投资指令、基准指标、费用、受托责任以及治理安排有助于使管理者的决策与投资者的利益保持一致。通常情况下，投资者仍然需要承担投资组合的盈亏风险。

集体投资工具在发行、交易和赎回请求权方面有所不同。传统的开放式共同基金以资产净值为基础来发行和赎回股份，而ETF则会在交易所进行交易，并将二级市场的交易与创建—赎回机制相结合。因此，基金份额的流动性不仅取决于交易安排，还取决于基础资产的流动性以及交易和赎回成本的分配。

机构投资者在负债、投资期限、流动性需求以及治理结构方面存在差异。养老基金会根据未来的福利义务来管理资产；主权财富基金和捐赠基金通常具有更长的投资期限；私募股权和风险资本基金则结合了承诺的资金与积极的投资决策；而套期保值基金则采用更为灵活的策略，包括使用杠杆和做空交易。

资产管理者也会产生信息，行使所有权，并参与证券市场。被动式、基于规则的和自动化的策略改变了这些功能的执行方式，而不是完全消除它们。在投资决策变得越来越系统化且依赖技术的情况下，模型设计、市场流动性、激励机制以及人工监管仍然非常重要。

%========================================================
% Key Terms
%========================================================

\TailModule{keyterms}{关键术语}

资产管理（Asset management）；委托投资（delegated investment）；投资授权（investment mandate）；委托代理问题（principal-agent problem）；受托责任（fiduciary duty）；管理费（management fee）；基准（benchmark）；主动回报（active return）；尽责管理（stewardship）；共同基金（mutual fund）；单位净值（net asset value）；交易所交易基金（exchange-traded fund）；赎回风险（redemption risk）；负债驱动型投资（liability-driven investing）；资金充足率（funding ratio）；私募股权（private equity）；风险投资（venture capital）；对冲基金（hedge fund）；机构持股（institutional ownership）；被动投资（passive investing）；指数基金（index fund）；基于规则的配置（rules-based allocation）；跟踪偏离（tracking difference）；执行缺口（implementation shortfall）；智能投顾（robo-advisory）；模型风险（model risk）。

%========================================================
% Financial AI Lab
%========================================================

\TailModule{aiprompt}{金融~AI~实验室}

\textbf{实验室：审核由人工智能生成的基金推荐方案}

请一个AI系统来比较两款投资于同一大类资产类别的投资基金，并推荐其中一款。要求该系统解释该基金的架构、基准指标、费用、流动性、投资组合构成，以及相关的历史跟踪差异和主要投资风险。

将建议分为四个阶段进行审核。

\begin{enumerate}

\item \textbf{请检查资金结构。}

请判断每辆汽车是否属于开放式共同基金、ETF、封闭式基金、私人基金，还是其他类型的投资结构。同时，请确认投资者是如何参与和退出投资的，以及交易是否是在市场价格进行的，还是基于净资产值进行的。

\item \textbf{请查看相关基准测试和成本信息。}

请使用官方的资金披露文件来确认所声称的基准值、管理费用以及其他相关费用。如果使用一致的周期和回报率来计算性能比较，那么计算结果将更加准确。

\item \textbf{检查流动性机制。}

将投资者可获得的流动性与基础资产的流动性进行对比。找出在何种情况下，赎回操作、市场影响或ETF的溢价与折价可能会变得重要。

\item \textbf{请检查推荐逻辑是否正确。}

确定该推荐是基于投资者特征还是基于未提供的未来结果。具体来说，需要明确关于投资期限、风险承受能力、流动性需求、税务状况或预期回报的假设。

\end{enumerate}

最终报告应当明确区分可验证的基金特征与预测、假设以及需要更多信息的投资者个人判断。在报告中应保留所有LaTeX命令、公式表达、交叉引用以及数字内容。

%========================================================
% Discussion Questions
%========================================================

\TailModule{questions}{讨论题}

\begin{enumerate}

\item 为什么即使投资者和基金经理都更倾向于较高的投资组合回报，委托资产管理仍然会产生本金-代理人问题呢？

\item ETF的创设-赎回机制与开放式共同基金的赎回机制有何不同？为什么这种差异对投资组合的交易具有重要影响？

\item 为什么一个拥有长期负债的养老基金会持有与那些投资者可以每天赎回的共同基金不同的投资组合呢？

\item 即使被动资产管理者不依赖自主的证券选择，他们如何能够影响公司治理？

\item 哪些资产管理功能可以通过自动化来更高效地执行？而在投资管理的哪些方面，仍然需要特别注意激励机制、市场流动性、治理结构以及人类判断的作用呢？

\end{enumerate}

%========================================================
% Applied Problems
%========================================================

\TailModule{appprobs}{应用题}

\begin{enumerate}

\item \textbf{费用与委托执行相关事项。}

该基金年初拥有约2亿元人民币的资金，并实现了9\%的毛回报。管理费为初始资产的1.2\%。绩效费相当于毛回报超过6\%部分的15\%，该部分是基于初始资产计算的，且在支付管理费之前进行计算的。忽略其他现金流和成本。

\begin{enumerate}

\item 计算管理费用。

\item 计算绩效费用。

\item 计算扣除各项费用后的最终资产价值。

\item 计算投资者简化的净回报。

\end{enumerate}

\item \textbf{共同基金资产净值与赎回事宜。}

该开放式基金的资产为人民币8.4亿元，负债为人民币1200万元，存续股份数量为7200万股。

\begin{enumerate}

\item 计算每股的净资产值。

\item 一位投资者以每股净资产的价格赎回了60万股股票。计算所需的现金数额，同时忽略相关费用和市场影响。

\item 假设该基金为了满足赎回需求而产生了约60万元人民币的交易成本，这些成本由剩余的资金承担。请计算出最终的每股净资产值，并说明谁应承担这些成本。

\item 请选出一种能够改变此类赎回成本分配方式的基金设计工具。

\end{enumerate}

\item \textbf{ETF溢价与折价。}

该ETF的每股市值约为4.00人民币，实际交易价格为4.08人民币。

\begin{enumerate}

\item 计算相对于净资产值的百分比溢价。

\item 描述在溢价足够覆盖交易和融资成本的情况下，潜在的创建交易的过程。

\item 解释为什么在基础证券流动性不足或市场关闭的情况下，溢价或折扣并不需要立即消失。

\end{enumerate}

\item \textbf{养老金资金筹措与负债敏感性。}

养老基金的资产为人民币950亿元，预计负债的现值为人民币1000亿元。

\begin{enumerate}

\item 计算资金比例。

\item 如果负债额上升至1060亿元人民币，而资产保持不变，那么需要计算新的资金比率。

\item 解释长期贴现率下降如何增加长期负债的现值。

\end{enumerate}

\item \textbf{机构交易与实施缺口。}

有一家基金会决定在参考价格为20元人民币时购买500,000股股票。平均执行价格约为20.12元人民币。

\begin{enumerate}

\item 使用LaTeX命令、公式、交叉引用以及数字来表示计算结果。  
根据方程~\ref{eq:ch14-implementation-shortfall}来计算实施过程中的不足，同时忽略相关费用的影响。

\item 将短缺金额表示为订单参考价值的百分比。

\item 解释快速完成订单与降低市场影响之间的权衡关系。

\end{enumerate}

\end{enumerate}

%========================================================
% Data Application
%========================================================

\TailModule{dataapp}{数据应用}

\textbf{在中国，公开募集的资金、ETF以及被动投资策略}

使用中国证券投资基金业协会、上海证券交易所、深圳证券交易所、基金招募说明书和定期基金报告等官方或原始来源的数据。选择少量可在共同期间比较的主动管理基金，以及指数基金或~ETF。

首先，需要提供有关基金规模、认购或赎回情况的信息，以及管理费用的相关数据。对于ETF来说，还需要收集市场价格以及选定日期的净资产值或预估净资产值，然后利用公式~\ref{eq:ch14-etf-premium-discount}来计算溢价或折扣。同时，要明确说明市场价格与净资产值指标在时间上的对应关系。

其次，将各被动基金的回报与它们所设定的基准指标进行长期比较，并计算其跟踪差异。利用基金披露信息来分析费用、现金持有量、抽样方法、重新平衡操作等因素是否有助于解释这些差异。不过，不要将过去的跟踪表现视为未来表现的预测指标。

第三，检查所选基金如何行使所有权。在信息披露可用的前提下，可以使用投票记录、定期报告或其他与管理相关的信息来评估被动投资组合的构建是否伴随着积极的投票或参与行为。

请列出各基金的名称、份额类别、样本周期、基准定义、费用规定、数据来源，以及针对收益或净资产值的任何调整措施。务必区分实际的数据与关于投资者需求、公司治理或市场影响的推测性描述。

%========================================================
% Further Reading
%========================================================

\TailModule{furtherread}{延伸阅读}

\begin{itemize}

\item Chevalier, Judith, and Glenn Ellison (1997),

``Risk Taking by Mutual Funds as a Response to Incentives（共同基金对激励的风险承担反应）,''

\emph{Journal of Political Economy}, 105(6), 1167--1200.

\item Sirri, Erik R., and Peter Tufano (1998),

``Costly Search and Mutual Fund Flows（有成本的搜寻与共同基金资金流动）,''

\emph{The Journal of Finance}, 53(5), 1589--1622.

\item Berk, Jonathan B., and Richard C. Green (2004),

``Mutual Fund Flows and Performance in Rational Markets（理性市场中的共同基金资金流动与业绩）,''

\emph{Journal of Political Economy}, 112(6), 1269--1295.

\item Appel, Ian R., Todd A. Gormley, and Donald B. Keim (2016),

``Passive Investors, Not Passive Owners（被动投资者，并非被动所有者）,''

\emph{Journal of Financial Economics}, 121(1), 111--141.

\item Ben-David, Itzhak, Francesco Franzoni, and Rabih Moussawi (2018),

``Do ETFs Increase Volatility?（ETF是否增加波动？）''

\emph{The Journal of Finance}, 73(6), 2471--2535.

\end{itemize}

\ChapterTail{summary,keyterms,aiprompt,questions,appprobs,dataapp,furtherread}

\chapter{保险与风险分担机构}
\label{chap:Insurance_and_Risk_Sharing_Institutions}


\begin{openingcase}{小额保险费与大额承诺}

一个家庭拥有一套价值约200万人民币的房子。虽然在特定年份发生严重火灾的可能性不大，但如果真的发生了火灾，造成的经济损失可能会非常惨重。为了避免承担这种风险，这个家庭购买了保险，并支付每年仅相当于保险公司最终可能需要支付的保险费的一小部分。

这笔交易看起来是不对称的。家庭今天支付的是一笔相对较小且确定的金额，而保险公司则承担潜在的巨大潜在负债。然而，保险公司同时会签订成千上万这样的合同。

\end{openingcase}

\begin{tcolorbox}[
  enhanced,
  colback=FinanceBlue!6,
  colframe=FinanceBlue,
  boxrule=0.6pt,
  sharp corners,
  left=8pt,
  right=8pt,
  top=7pt,
  bottom=7pt,
  before skip=1em,
  after skip=1.6em,
  title=\textbf{核心问题},
  coltitle=white,
  colbacktitle=FinanceBlue,
  fonttitle=\bfseries,
  attach boxed title to top left={xshift=6pt,yshift=-2mm},
  boxed title style={
    boxrule=0pt,
    sharp corners,
    left=5pt,
    right=5pt,
    top=2pt,
    bottom=2pt
  }
]
\begin{enumerate}[
  leftmargin=*,
  label=\textcolor{FinanceBlue}{\textbf{\arabic*.}},
  itemsep=0.55em,
  topsep=0.2em
]


\item 为什么保险能够改善风险分担，以及何时通过集中管理使得整体损失更加可预测？

\item 保险契约如何在解决逆向选择和道德风险的同时，还能保持投保人的积极性？

\item 保险公司如何为保险契约定价、评估和管理其负债，如何管理资产和资本，以及如何通过再保险和资本市场来转移风险？

\end{enumerate}
\end{tcolorbox}

\section{保险、风险汇集与可保性}

保险公司通过合同来转移不确定事件的财务后果，用已知的保险费换取对特定损失的保障。通过结合许多这样的合同，保险公司可以汇集风险，使总损失更加可预测。而这种安排的可行性取决于潜在风险是否可以被保险所涵盖。

\subsection{个体风险与保险需求}

个人可能会面临一种低概率但相对较大的损失，相对于初始财富而言。设$W$表示初始财富，假设随机损失$\tilde{L}$以概率$p$取值$L$，而以概率$1-p$则取值为零。如果没有保险，最终的财富状况……

\begin{equation}
W_T=
\begin{cases}
W-L, & \text{以一定的概率 } p,\\
W, & \text{以一定的概率 } 1-p.
\end{cases}
\label{eq:ch15-uninsured-wealth}
\end{equation}

预期的经济损失为

\begin{equation}
\mathbb{E}[\tilde{L}]=pL.
\label{eq:ch15-expected-loss-simple}
\end{equation}

仅有预期损失，无法概括投保人在不同状态下的风险暴露。保险重新分配各状态下的财富：无论是否出险，投保人均支付保险费；损失发生时则获得赔偿。效用函数对财富呈凹性时，减少状态间财富差异可能提高预期效用。因此，风险厌恶的投保人可能愿意支付超过预期损失的保险费，以换取对重大不利结果的保障。

假设保险费为$P$时，能够获得全额覆盖。忽略排除条款和索赔的障碍后，与这种损失相关的最终财富就成为了

\begin{equation}
W_T^{I}=W-P.
\label{eq:ch15-fully-insured-wealth}
\end{equation}

完全保险以更稳定的财富状况替代状态依存财富。即使保险费超过预期损失，只要减少财富风险的价值超过保费附加，保险仍可能有价值。

保险还可以减少为了应对低频但损失严重的事件而需要持有的预防性流动资产的需求。当在发生损失后借款或清算资产变得成本高昂或极其困难时，这种效应就更加明显了。

\begin{example}{相对于家庭财富的重大损失}{household-insurance-demand}

一个家庭初始财富为100万人民币，全年面临2\%的概率会遭受40万人民币的损失。如果不进行保险处理，那么如果发生损失，最终财富将仅为60万人民币；而如果没有发生损失，最终财富则仍为100万人民币。

预期损失是

$$
0.02\times 400{,}000=8{,}000.
$$

在保险费为$P$的情况下，无论处于哪种状态，最终的财富都约为人民币$1{,}000{,}000-P$。因此，一个厌恶风险的家庭可能会接受高于人民币8,000的保险费，只要减少面临重大不利状况的风险所带来的价值足以覆盖保险费的增加。

\end{example}


\subsection{风险汇集与大数定律}

当保险公司将许多风险汇集在一起时，风险汇集的效果会更好。作为参考，假设有$n$个投保人，他们具有独立的损失指标$X_i$，其中……

$$
X_i=
\begin{cases}
1, & \text{如果投保人 } i \text{ 感受到了损失},\\
0, & \text{否则，请保留LaTeX命令、公式、交叉引用以及数字的内容。},
\end{cases}
$$
且每次索赔发生的概率均为$p$。索赔的数量为

\begin{equation}
N=\sum_{i=1}^{n}X_i.
\label{eq:ch15-number-of-claims}
\end{equation}

其期望值为

\begin{equation}
\mathbb{E}[N]=np,
\label{eq:ch15-expected-claims}
\end{equation}
此外，在独立性方面，

\begin{equation}
\operatorname{Var}(N)=np(1-p).
\label{eq:ch15-claims-variance}
\end{equation}

观察到的索赔频率定义为

\begin{equation}
\hat p=\frac{N}{n}
=\frac{1}{n}\sum_{i=1}^{n}X_i.
\end{equation}

在独立分布下，其方差为

\begin{equation}
\operatorname{Var}(\hat p)
=\frac{1}{n^2}\operatorname{Var}(N)=\frac{p(1-p)}{n}.
\label{eq:ch15-frequency-variance}
\end{equation}

随着风险池扩大，索赔频率的方差按投保人数量的倒数下降。总索赔数量仍不确定，但平均每张保单的赔付经验变得更可预测。这构成风险汇集的统计基础，使保险公司能够更可靠地估计组合损失。

可预测性的提高取决于各种损失之间的关联性。当损失之间相关性较弱时，集中处理会更加有效；而当出现共同冲击时，可能会同时引发许多索赔，从而削弱分散风险的效果。

\begin{example}{汇集独立索赔}{pooling-independent-claims}

假设有10,000份类似的保单，每份保单的保额均为20,000人民币，且每年的索赔概率为1\%。

索赔的预期数量是

$$
10{,}000\times 0.01=100.
$$

索赔的标准差为

$$
\sqrt{10{,}000\times 0.01\times 0.99}
\approx 9.95.
$$

预期的索赔频率为1\%，而其标准差为

$$
\frac{9.95}{10{,}000}
\approx 0.001,
$$

或者大约0.10个百分点。

任何一位投保人的损失仍然不确定，而整个群体的索赔频率则变得更加稳定。

\end{example}


\subsection{可保性与风险汇集的限度}

更大保险池的多样化效益取决于各次损失之间的依赖性。为了说明这一点，设$X_i$表示第$i$份保单的损失。假设各次损失的方差均为$\sigma^2$，且每对损失都具有共同的相关性$\rho$。那么平均损失的方差就是

\begin{equation}
\operatorname{Var}(\bar{X})=
\frac{\sigma^2}{n}
+
\frac{n-1}{n}\rho\sigma^2.
\label{eq:ch15-correlated-pooling}
\end{equation}

当~$\rho=0$~时，随着风险池扩大，平均损失方差趋于零。当~$\rho>0$~时，第一部分随风险池规模下降，但相关性成分保留。因此，风险汇集可以减少个体特有波动，却不能消除共同冲击的风险暴露。

这种区分有助于解释为什么有些风险比其他风险更容易被承保。孤立的家庭火灾和许多汽车事故在保单的具体规定上存在很大差异，通常可以将这些事故合并起来处理。而地震、地区性洪水、流行病以及系统性网络事件则可能会同时导致许多投保人遭受损失，这就使得总风险更加难以分散。

因此，可保性不只取决于预期损失。事件与承保损失须足够明确且可核实，损失频率和程度须能够较可靠地估计，潜在总赔付也须与保险费及可用风险承担能力相匹配。信息和激励问题还可能限制可持续条款下能够提供的保障。

\begin{tcolorbox}[
enhanced,
breakable,
colback=white,
colframe=FinanceBlue,
boxrule=0.7pt,
sharp corners,
borderline west={2.2pt}{0pt}{FinanceBlue},
title=\textbf{财务展望：当资源整合仍不足以解决问题时},
coltitle=white,
colbacktitle=FinanceBlue,
fonttitle=\bfseries,
before skip=1.2em,
after skip=1.2em
]

当大部分损失风险由共同因素引起时，保险公司必须决定保留多少风险，以及将多少风险转移出去。地理上的多样化可以降低对局部事件的暴露风险，而再保险和与保险相关的证券可以将部分风险转移到其他资产负债表上。资本可以为那些仍由保险公司承担的损失提供缓冲。

这些机制并不能消除潜在的风险。它们只是将风险的财务后果重新分配给了不同的承担者，并决定了保险公司能够承受多大的集中性损失。

\end{tcolorbox}


\section{保险契约与信息问题}

保险契约规定了哪些损失是涵盖在内的，哪些损失由被保险人承担，以及索赔是如何被核实和支付的。这些条款还有助于管理在保险开始之前和之后出现的各种信息和激励问题。


\subsection{保险保障、契约支付与索赔}

保险单规定保险事件、保障期间、给付或赔偿、除外责任、免赔额、限额及索赔条件。这些条款共同决定何时触发支付，以及保险公司支付多少。

对于涵盖实际损失$L$的赔偿保单，令$d$表示免赔额，$C$表示保险公司的最大支付金额。简化的支付函数如下：

\begin{equation}
I(L)=\min\left\{\max(L-d,0), C\right\}.
\label{eq:ch15-insurance-payoff}
\end{equation}

被保险人所保留的损失份额为

\begin{equation}
R(L)=L-I(L).
\label{eq:ch15-retained-loss}
\end{equation}

因此，契约重新分配财务损失，而非消除损失。免赔额使部分损失由被保险人保留，赔付限额则限制保险公司的最大风险暴露。

这种支付结构对于赔偿保险来说最为直接。人寿保险通常在特定事件发生时提供一定的利益，而不是对实际财产损失的赔偿。健康保险则可能结合赔偿、合同利益、医疗服务提供者安排以及费用分摊等方式。因此，支付的形态取决于被保险事件的性质，以及可以观察和验证的内容。

索赔管理负责执行合同中的支付条款。保险公司会确定某项索赔是否属于保险范围，核实相关的损失或事件，按照保单条款进行赔付，并作出相应的支付。这些核实和结算程序也会影响到与合同相关的信息和激励措施。

\begin{example}{免赔额与保障限额}{deductible-and-limit}

该保单的免赔额为人民币~20,000，保险公司的最高赔偿金额为人民币~300,000。

如果承保损失为人民币80,000元，保险公司将予以赔偿

$$
I(80{,}000)=80{,}000-20{,}000=60{,}000.
$$

如果承保损失达到500,000元人民币，保险公司将予以赔偿

$$
I(500{,}000)=300{,}000,
$$

导致被保险人损失人民币约20万。

免赔额决定了保险公司需要承担多少较小的损失，而赔偿限额则限制了保险公司对较大损失的承受能力。

\end{example}


\subsection{逆向选择、核保与风险分类}

潜在投保人掌握自身风险信息，而保险公司无法完全观察时，可能产生逆向选择。如果不同风险类型面对相同条款，高风险者可能更愿购买或维持保障，低风险者则可能减少参与，导致承保风险池比定价时预计的更高风险。

假设有两种风险类型。类型$L$的预期损失为$\ell_L$，类型$H$的预期损失为$\ell_H$，

\begin{equation}
\ell_H>\ell_L.
\label{eq}
\end{equation}

基于两类风险共同确定的混合保险费，相对于~$L$~类的预期损失通常偏高，相对于~$H$~类则偏低。如果参与结构向~$H$~类转移，承保风险池的平均预期损失就会上升，对保险费形成上行压力，或要求调整契约条款。

核保通过收集和评估信息来解决这一信息问题，从而决定是否以及以何种条件为某项风险进行保险。根据产品和法律环境的不同，保险公司可能会考虑损失历史、财产特性、健康信息、商业活动、地理位置等因素，以及其他被允许的指标。

风险分类利用这些信息将具有相似预期索赔特征的保险风险进行分组。更详细的分类可以提高保险费与合同条款与预期损失之间的匹配度，不过其有效性取决于基础信息的质量、相关性以及允许使用的程度。

核保与风险分类通过改善承保风险评估及定价所依赖的信息，帮助保险公司应对逆向选择。相关评估主要关注损失的预期发生频率与程度，必要时还包括不同风险暴露之间的依赖关系。

\subsection{道德风险、风险分担与激励设计}

保障开始后，保险改变行为时可能产生道德风险。保障减轻被保险人承担的财务后果，可能削弱其预防损失、维护投保资产、控制损失程度或控制承保支出的激励，尤其在这些行动有成本且无法完全观察时。

保险契约可以让保险公司与投保人分担损失，而非完全转移损失。免赔额为~$d$~时，被保险人承担承保损失的初始部分。在比例分担安排下，被保险人还承担超过免赔额部分的一定比例损失。若保险公司承担免赔后损失的比例为~$\alpha$，则简化支付为


\begin{equation}
I(L)
=\alpha\max(L-d,0),
\qquad 0\leq\alpha\leq 1,
\label{eq:ch15-coinsurance-payoff}
\end{equation}
遵守任何覆盖限制规定。保留LaTeX命令、公式、交叉引用和数字。

风险分担使得被保险人面临更大的损失风险，这反而能够增强预防和成本控制的动力。更全面的保险覆盖则能提供更强的抵御不利状况的保护。因此，合同的设计需要平衡风险分担与激励效应之间的关系。

监控和索赔核实可以补充这种风险分担的形式。保险公司可能会要求采取预防措施，检查被保财产，核实损失和文件资料，使用经过批准的供应商网络，或者以特定的行为作为保险覆盖的条件。这样的机制可以提高激励效果并降低索赔成本，但其有效性取决于监控的可见性和监控成本。

\begin{tcolorbox}[
enhanced,
breakable,
colback=white,
colframe=FinanceBlue,
boxrule=0.7pt,
sharp corners,
borderline west={2.2pt}{0pt}{FinanceBlue},
title=\textbf{财务前景：报道的广度可能会削弱激励措施的效果。},
coltitle=white,
colbacktitle=FinanceBlue,
fonttitle=\bfseries,
before skip=1.2em,
after skip=1.2em
]

完全转移风险能提供最强的损失保护，但也可能降低被保险人采取行动的财务后果。保留部分损失可以维持激励措施，但代价是让投保人面临更大的不良后果风险。

因此，风险分担的适当程度取决于风险的性质、行为的可观察性、监控的成本，以及投保人承担剩余损失的能力。

\end{tcolorbox}


\section{保险定价与负债管理}

保险定价决定了为应对不确定的未来索赔及相关成本所需的保险费。这些索赔会产生负债，这些负债可能会持续到收取保险费之后，因此保险公司需要随着时间来管理储备金、资产、流动性、资本以及风险转移问题。


\subsection{预期损失、保险费与风险附加}

保险定价始于未来索赔支付的预期价值。用$L$表示与一份保单相关的保险损失的实际现值。预期损失是

\begin{equation}
EL=\mathbb{E}[L].
\label{eq:ch15-expected-loss}
\end{equation}

简化的保险费可以表示为

\begin{equation}
P=EL+E+R,
\label{eq:ch15-premium-decomposition}
\end{equation}

其中，$E$表示预期的获取和运营费用，而$R$则代表了不确定性的风险负荷，以及为了应对未被预期索赔所涵盖的风险而需要的资本。

预期损失既反映了索赔的频率，也反映了索赔的严重程度。在一个简单的设定中，该时期内最多只发生一次索赔。如果以概率$p$发生索赔，那么在已知有索赔发生的条件下，预期损失为$\mathbb{E}[S\mid \text{claim}]$。

\begin{equation}
EL=
p,\mathbb{E}[S\mid \text{索赔}].
\label{eq:ch15-frequency-severity}
\end{equation}

风险附加反映承担预期损失周围不确定性的成本。索赔波动较大、相关性较高、难以分散，或存在显著模型不确定性时，其重要性通常更大。再保险、监管、竞争、投资条件和契约设计，也可能影响最终收取的保险费。

因此，定价是面向未来的。历史损失数据可以提供信息，但当索赔频率、严重性、投入成本、法律条件、灾难风险或保单条款发生变化时，保险费可能需要进行调整。

\begin{example}{从预期损失到保险费}{expected-loss-to-premium}

保险公司估计每年的索赔概率为4\%，并且在发生索赔的情况下，预期索赔损失为约50,000人民币。每保单的预期损失为

$$
EL=0.04(50{,}000)=2{,}000.
$$

假设预期的收购和运营费用为500元人民币，而保险公司适用的风险系数为300。那么所确定的保险费为

$$
P=2{,}000+500+300=2{,}800.
$$

预期损失与所标示的保险费之间的差额，反映了在这种简化定价方式下所需的费用和风险负荷。

\end{example}


\subsection{保险负债估值与准备金}

定价将未来索赔的估算转化为当前的保险费条款。一旦保险合约签订，那些未来的索赔就变成了负债，其金额和支付时间必须在合同的整个期限内进行衡量。

让 $C_t$ 表示未来日期 $t=1,\ldots,T$ 时不确定的索赔现金流。作为一种简化的估值基准，预期索赔支付的现值可以表示为

\begin{equation}
V_L=\sum_{t=1}^{T}
\frac{\mathbb{E}[C_t]}{(1+r_t)^t},
\label{eq:ch15-insurance-liability-pv}
\end{equation}

其中，$r_t$表示与到期日和估值基准相适应的贴现率。该表达式说明了预期现金流和贴现的作用；更全面的负债衡量标准还可能反映费用、合同特性、不确定性以及其他相关估值框架所需的调整因素。

根据适用的会计或监管框架，储备金被表示为当前的负债。根据合同的不同，这些储备金可能反映那些已经发生但尚未支付或未报告的索赔，以及与长期保险契约相关的未来收益。例如，人寿和年金负债可能会受到死亡率、寿命、失效行为、费用、担保以及贴现率的影响。

因此，储备金并非仅仅是为未来索赔而预留的现金。它是一项由保险公司资产和资本支持的负债。索赔情况、假设或贴现率的变化可能会改变所计算的负债以及保险公司的财务状况。

预期支付的时间也很重要。与短期负债相比，长期负债通常对贴现率的变化更为敏感，因此将保险负债的估值与更广泛的资金时间价值框架联系起来是必要的。



\subsection{资产负债管理与偿付能力}

估值和预留将未来的保险义务转化为具有特定时间、敏感性和不确定性的负债。然后，保险公司管理资产以抵消这些负债，并保持资本充足，以应对资产负债表两边的逆向情况。因此，资产与负债管理着重于资产与负债现金流、流动性、久期以及风险的匹配。

设~$A$~表示资产经济价值，$L$~表示保险及其他负债的经济价值。经济盈余的一种简化度量为

\begin{equation}
S=A-L.
\label{eq:ch15-economic-surplus}
\end{equation}

资产负债表任一侧变化，都可能影响经济盈余。利率变化是重要来源之一。利率下降通常提高长期负债价值，具有类似久期特征的资产则可抵消部分风险暴露。

简化的久期-缺口衡量方法为

\begin{equation}
DG=D_A-\frac{L}{A}D_L,
\label{eq:ch15-insurer-duration-gap}
\end{equation}

其中，$D_A$~和~$D_L$~分别表示资产和负债的有效久期。久期缺口为经济盈余对利率变化的敏感性提供一阶度量。当保险负债包含保证、退保选择权、失效行为或其他非线性特征时，实际风险暴露可能更复杂。

流动性为资产-负债管理增加了另一个维度。财产保险公司可能在灾难发生后面临集中性的索赔支付，而人寿保险公司则可能会遇到资金退保或福利支付的问题。即使保险公司拥有足够的长期资产价值，但如果现金流出发生在资产能够以合理的成本转化为资金之前，那么该公司仍可能面临短期压力。

偿付能力涉及保险公司吸收资产和负债不利变化的能力。资本为应对意外索赔、资产损失、模型误差以及其他不利情况提供损失吸收能力。不同司法辖区的监管资本要求各不相同，但其基本经济作用是在实际结果不如预期负债和资产价值所反映的情况时提供缓冲。


\begin{figure}[H]
\centering
\begin{tikzpicture}[
> =Latex,
box/.style={
 draw,
 rounded corners,
 align=center,
 text width=5.0cm,
 minimum height=0.95cm,
 inner sep=4pt
 },
 main/.style={
 draw=FinanceBlue,
 fill=FinanceBlue!8,
 rounded corners,
 align=center,
 text width=5.0cm,
 minimum height=1.0cm,
 inner sep=4pt
 },
 flow/.style={-Latex, line width=0.8pt, FinanceBlue},
 link/.style={<->, line width=0.8pt, FinanceBlue}
 ]
\node[box] (assets) at (-3.4,1.6)
{资产\\流动性、久期、信用质量};

\node[box] (liab) at (3.4,1.6)
{负债\\时间、金额、敏感性};

\node[main] (alm) at (0,-0.5)
{资产--负债 \\ 管理};

\node[box] (capital) at (0,-2.2)
{保留LaTeX命令、公式、交叉引用和数字。  
大写\\损失吸收能力};

\draw[link] (assets) -- (liab);
\draw[flow] (assets) -- (alm);
\draw[flow] (liab) -- (alm);
\draw[flow] (capital) -- (alm);

\end{tikzpicture}

\caption{资产负债管理与损失吸收能力}
\label{fig}
\end{figure}


\section{保险机构与风险转移}

保险公司通过保险契约来承担风险，但它们并不需要将所有风险都保留在自己的资产负债表上。根据风险的性质和集中度，保险公司可以选择保留部分风险，将其中的一部分转移给再保险公司，或者将某些风险转移给资本市场投资者。

\subsection{原保险公司与主要险种}

在保险体系中，主要保险公司直接从家庭、企业以及其他投保人那里承担风险，涵盖人寿、健康以及财产保险等主要领域。

人寿保险提供与死亡或生存相关的福利。某些合同覆盖较长的期限，可能包含储蓄成分、担保或投保人的选择，这使得负债容易受到死亡率、寿命长短、退保行为和利率的影响。

健康保险涵盖医疗费用或特定的健康相关福利。索赔可能频繁发生，而使用情况可能会因保险覆盖而有所变化，因此在许多健康保险安排中，信息管理和激励措施显得非常重要。

财产保险保单涵盖财产损失、负债、机动车保险、商业中断险以及其他非人寿风险。索赔结算的时间跨度可以从较短的时段到多年不等，而灾难性事件可能导致各保单之间产生巨大的相关损失。

这些差异导致不同的负债状况，从而影响保险费设定、储备期限、流动性需求、资产与负债管理以及再保险的需求。



\begin{table}[H]

\centering

\caption{主要险种与负债特征}

\label{tab}

\begin{tabularx}{0.92\textwidth}{@{}p{0.20\textwidth} X X@{}}

\toprule

\textbf{行号} & \textbf{典型的保险风险敞口} & \textbf{负债特征} \\

\midrule

人寿保险
& 死亡率、生存率和长期家庭保障
& 通常具有较长的久期；可能包含担保和投保人的选择选项 \\

健康保险
& 医疗费用及与健康相关的福利
& 索赔可能频繁发生；使用和信息方面的问题可能会很严重 \\

属性与伤亡情况&财产损失、负债、机动车损害以及业务中断&索赔久期因行而异；灾难性和负债风险可能导致大额或延迟的支付\\

\bottomrule

\end{tabularx}

\end{table}


\subsection{再保险与保险风险再分配}


上述负债特征差异，会产生不同的风险转移需求。再保险使原保险公司支付再保险费，将部分承保风险暴露转移给再保险公司，从而在机构间重新分配承保损失。


再保险可以限制对个别重大损失的承受能力，减少灾难事件带来的损失，稳定承保结果，并提高保险公司开展额外业务的能力。保护的形式取决于保险费与损失如何在主保险公司和再保险公司之间分配。

在比例再保险中，保险费与损失是根据约定的比例来分配的。如果再保险公司的份额为$\alpha$，那么分担的损失就是

\begin{equation}
C(L)=\alpha L.
\label{eq:ch15-proportional-reinsurance-payoff}
\end{equation}

在损失超赔再保险中，再保险公司负责支付超出保留金$K$的赔付金额，但最高支付额有限，为$M$：

\begin{equation}
C(L)=\min\left\{\max(L-K,0),,M\right\}.
\label{eq:ch15-excess-of-loss}
\end{equation}

因此，主要保险公司将承担高达$K$的损失，以及超过$K+M$的任何金额。合同改变了保险公司的损失承担方式，尤其是在不利的情况下。

\begin{example}{超额损失再保险}{excess-loss-reinsurance}

该保险公司购买了价值约1亿元人民币的灾难性保险再保险，并且提供的最大赔付金额为约2亿元人民币。

如果涉及到的灾难性损失达到8000万人民币，那么再保险商将不支付任何赔偿。

如果损失达到1.8亿元人民币，再保险商将进行赔付。

[
180-100=80
]

百万人民币。

如果损失达到4亿元人民币，再保险公司将支付最多2亿元人民币的赔偿。主要保险公司将获得最初的1亿元人民币赔偿，而超过再保险层部分的1亿元人民币则由再保险公司承担。

\end{example}

通过接受来自多个主要保险公司的再保险风险，再保险公司可以将不同公司、地区以及业务领域的风险进行整合。这种风险分散的效果取决于这些风险之间的相互依赖性，而高度相关的事件可能会导致大量风险集中在再保险层面。

再保险也形成对再保险公司本身的风险暴露。因此，保障价值不仅取决于契约条款，也取决于承保损失发生时，再保险公司履行义务的能力。


\subsection{保险连接证券与巨灾债券}

再保险是将保险风险转移给其他保险报表，而证券-linked证券则可以将这种转移扩展到资本市场投资者。灾难债券就是最典型的例子。

典型巨灾债券结构中，投资者本金存入抵押安排，投资者因承担指定事件风险获得补偿。风险期间内若未发生契约触发事件，本金到期返还；若触发条件满足，则依契约释放部分或全部本金，减轻发起人对承保事件的风险暴露。

设债券的本金为$F$。忽略票息支付，简化的终端本金支付为

\begin{equation}
X_T=
\begin{cases}
F, & \text{如果没有触发任何事件发生的话。},\\
(1-\theta)F, & \text{如果触发条件导致本金的减少 } \theta,
\end{cases}
\label{eq:ch15-cat-bond-payoff}
\end{equation}

其中，$0\leq\theta\leq1$ 的值由合约的触发条件和损失结构来决定。

触发因素可能基于赞助方的实际损失、行业损失、物理事件参数、模拟出的损失，或者这些指标的组合。触发因素的选择会影响基础风险的分配，以及结算的透明度和速度。

灾难债券拓宽了用于吸收集中保险损失的资本来源。当灾难风险较大或需要额外的承担风险能力时，灾难债券可以补充再保险的作用。投资者因此能够因承担具有不同驱动因素的风险而获得补偿，这种风险与传统的金融市场风险有所不同。

这种转移并非无成本之举。结构设计费用、不确定性建模、基础风险、法律安排以及所需的投资者回报都会影响该证券的经济价值。因此，保险-linked证券应该是对传统再保险的补充，而不是简单地替代它。

\begin{tcolorbox}[
enhanced,
breakable,
colback=white,
colframe=FinanceBlue,
boxrule=0.7pt,
sharp corners,
borderline west={2.2pt}{0pt}{FinanceBlue},
title=\textbf{财务展望：将风险转移扩展到保险资产负债表之外},
coltitle=white,
colbacktitle=FinanceBlue,
fonttitle=\bfseries,
before skip=1.2em,
after skip=1.2em
]

保险风险可以通过一系列资产负债表来传递：

$$
\text{投保人}
\rightarrow
\text{主要保险公司}
\rightarrow
\text{再保险人}
\rightarrow
\text{资本市场投资者}.
$$

每一步都会重新分配因特定损失而产生的财务后果，但同时并不会消除潜在的不确定性。

\end{tcolorbox}


%========================================================
% Summary
%========================================================

\TailModule{summary}{本章小结}

保险公司通过状态依存合同来转移不确定的损失。投保人支付一定的保险费以获得保护，以应对特定的损失，而保险公司则汇集这些风险，并保持足够的财务能力来处理索赔。

当个体损失之间的相关性足够弱时，风险汇集可以使整体损失变得更加可预测。然而，相关性和共同的冲击会限制这种好处，从而增加对资本和风险转移的需求。

保险契约在风险分担与激励措施之间实现平衡。核保和风险分类可以应对逆向选择问题，而免赔额、比例分担制、监控以及索赔核实则可以缓解道德风险。

定价需要考虑预期损失、费用以及承担不确定性的成本。储备措施用于衡量未来的义务，资产与负债管理将资产与负债进行协调，而资本则提供吸收损失的能力。

风险不必留在主要保险公司手中。再保险可以将风险转移给不同的保险公司，而与保险相关的证券则可以将某些风险转移给资本市场投资者。


%========================================================
% Key Terms
%========================================================

\TailModule{keyterms}{关键术语}

保险（Insurance）；风险汇集（risk pooling）；大数定律（law of large numbers）；可保性（insurability）；相关损失（correlated losses）；逆向选择（adverse selection）；核保（underwriting）；道德风险（moral hazard）；免赔额与比例分担（deductibles and coinsurance）；预期损失（expected loss）；风险附加（risk loading）；保险负债（insurance liability）；准备金（reserves）；资产负债管理（asset--liability management）；偿付能力（solvency）；再保险（reinsurance）；巨灾债券（catastrophe bonds）；保险连接证券（insurance-linked securities）。

%========================================================
% Financial AI Lab
%========================================================

\TailModule{aiprompt}{金融~AI~实验室}

\textbf{实验室：对保险定价及风险转移分析的审计分析}

请一个AI系统来分析一家简化的财产保险公司。提供保单数量、预期索赔频率、预期索赔严重程度、运营费用、拟定的保险费或风险负荷、免赔额，以及简化的总损失再保险层。

将分析过程分为四个阶段进行实施。

\begin{enumerate}

\item \textbf{请确认预期损失和保险费。}

确保索赔频率和条件性严重性能够一致地结合使用。区分每份保单的预期损失与总体预期损失，并检查各项费用以及任何风险负荷是如何影响所确定的保险费的。

\item \textbf{检查合同设计。}

研究免赔额如何改变保险公司的赔付金额以及投保人的自留损失。确定所提供的信息是否足以计算其对索赔预期的影响；仅使用平均损失可能无法准确反映扣除免赔额后的预期赔付金额。将任何关于道德风险的索赔与实际提供的假设区分开来。

\item \textbf{追踪再保险层。}

在附件点下方、覆盖层内以及其极限之上，应用几种 aggregate-loss 结果。重新计算每笔情况下的再保险公司支付金额和主要保险公司的保留损失。

\item \textbf{测试相关损耗场景。}

将独立索赔替换为一种灾难性状况，这种状况会同时导致许多保单的损失。解释为什么更大的风险池可以减少特异性差异，而不会消除对共同冲击的暴露风险。

\end{enumerate}

最终报告应区分算术结果以及关于损失分布、投保人行为、风险之间的依赖性以及合同解释方面的假设。


%========================================================
% Discussion Questions
%========================================================

\TailModule{questions}{讨论题}

\begin{enumerate}

\item 为什么一个风险厌恶的家庭愿意支付高于预期损失的保险费呢？

\item 大数法则如何支持风险汇集，为什么正相关会限制这种效应？

\item 免赔额和共同承担费用如何影响风险分担以及投保人的激励措施？

\item 为什么当投保人有关于自身风险的私人信息时，汇集的保险费会产生逆向选择压力呢？

\item 保险公司储备与保险公司资本之间的经济区别是什么？

\item 再保险公司和灾难债券是如何将保险风险从主要保险公司转移出去的？

\end{enumerate}


%========================================================
% Applied Problems
%========================================================

\TailModule{appprobs}{应用题}


\begin{enumerate}

\item \textbf{预期损失与保险需求。}

一个家庭面临每年有3\%的概率会损失约20万人民币。

\begin{enumerate}

\item 计算预期的年度损失。

\item 将预期损失与每年7,000人民币的保险费进行比较。

\item 解释为什么风险厌恶型家庭仍然可能倾向于选择这种保险契约。

\item 请指出除了风险厌恶之外，还有哪些因素可能会影响家庭的决策。

\end{enumerate}

\item \textbf{合并独立的索赔。}

一家保险公司负责处理25,000种类似的风险。每种风险每年的索赔概率均为2\%，且各种索赔的结果是相互独立的。

\begin{enumerate}

\item 计算索赔的预期数量。

\item 计算索赔数量的标准差。

\item 计算索赔频率的标准差。

\item 解释索赔结果之间的正相关性如何会影响汇集资金所带来的分散化效益。

\end{enumerate}

\item \textbf{免赔额和支付限额。}

该保单的免赔额为人民币~10,000元，保险公司的最高赔偿金额为人民币~150,000元。

\begin{enumerate}

\item 计算保险公司对8,000元人民币、60,000元人民币和250,000元人民币损失所承担的赔偿金额。

\item 计算每个案例中被保险人所承担的损失金额。

\item 解释免赔额如何影响投保人在首次损失中仅获得部分赔偿的激励措施。

\end{enumerate}

\item \textbf{保险费分解。}

保险公司估计每保单的预期损失为人民币约1800元，预期费用约为人民币450元，风险负荷为人民币约250元。

\begin{enumerate}

\item 根据方程~\ref{eq:ch15-premium-decomposition}计算所要求的保险费。

\item 计算所示保险费与预期损失之间的差额。

\item 将这笔保险费表示为所指示保险费的百分比。

\item 解释为什么保险费与预期损失之间的区别不应被解释为会计利润。

\end{enumerate}

\item \textbf{负债评估与利率。}

保险公司预计在三年后支付约2000万元人民币。采用4\%的恒定年贴现率进行计算。

\begin{enumerate}

\item 计算预期付款的现值。

\item 使用3\%的贴现率重新计算现值。

\item 请说明这一变革的方向。

\item 将这一结果与长期保险公司资产-负债管理问题联系起来。

\end{enumerate}

\item \textbf{超额损失再保险。}

主要保险公司承担总灾难损失中的首笔人民币约5000万元，并购买再保险，以获得额外高达人民币约1.5亿元的保障。

\begin{enumerate}

\item 计算再保险商对总损失的责任赔偿金额，具体为：3000万人民币、1亿人民币、1.8亿人民币，以及2.6亿人民币。

\item 计算每起案件中的主要保险公司所承担的损失。

\item 确定再保险层的连接点和耗尽点。

\item 解释合同是如何改变主要保险公司所保留的损失分配的上限的。

\end{enumerate}

\end{enumerate}


%========================================================
% Case Analysis
%========================================================

\TailModule{caseanalysis}{案例分析}

\textbf{案例：2021年河南洪水与保险风险分担}

2021年7月，河南省遭受了严重的洪水袭击，短时间内产生了大量的索赔申请。截至7月28日，该省的保险业已经收到了约412,300份索赔报告，初步估计的保险损失约为人民币98.04亿元。当时，已有约43,400份索赔得到了解决，赔偿金额总计人民币6.09亿元。

索赔涉及多条保险条款。其中，机动车保险相关的索赔约为226,400件，造成的损失估计为人民币约6.412亿元；而家庭财产和农业保险方面的索赔也相当可观。在同一地区和同一事件下出现的大量损失，为研究当多个风险同时受到影响时保险是如何运作的提供了有用的分析框架。

从四个角度来分析这一案例。首先，解释地理相关性如何影响风险汇集的多样化效益。其次，探讨保单限额、免责条款、免赔额以及保险价值如何导致实际损失与索赔损失之间的差异。第三，解释为什么在索赔的金额和时机仍然不确定时，保险公司可能需要流动性。第四，讨论再保险如何重新分配由主要保险公司保留的部分损失。

上述数据仅为初步结果，不应被视为最终的行业损失或索赔金额。

\emph{原始资料来源：2021年7月29日河南省防汛新闻发布会；2021年7月30日发布的官方信息，援引了河南省银行和保险监管局的声明。}


%========================================================
% Further Reading
%========================================================

\TailModule{furtherread}{延伸阅读}

\begin{itemize}

\item Arrow, Kenneth J. (1963),
``Uncertainty and the Welfare Economics of Medical Care（不确定性与医疗保健的福利经济学）,''
\emph{American Economic Review}, 53(5), 941--973.

\item Pauly, Mark V. (1968),
``The Economics of Moral Hazard: Comment（道德风险经济学：评论）,''
\emph{American Economic Review}, 58(3), 531--537.

\item Rothschild, Michael, and Joseph E. Stiglitz (1976),
``Equilibrium in Competitive Insurance Markets: An Essay on the Economics of Imperfect Information（竞争性保险市场中的均衡：不完全信息经济学研究）,''
\emph{The Quarterly Journal of Economics}, 90(4), 629--649.

\item Froot, Kenneth A. (2001),
``The Market for Catastrophe Risk: A Clinical Examination（巨灾风险市场：案例考察）,''
\emph{Journal of Financial Economics}, 60(2--3), 529--571.

\end{itemize}


\ChapterTail{summary,keyterms,aiprompt,questions,appprobs,caseanalysis,furtherread}

\end{document}
